% a06_BALtorus_c: part 2 of 3 of the appendix app:a06_BALtorus (input from sec/a06_BALtorus.tex): Steps 3 and 4 of the proof of thm:a06:T1.
% Step 3 continues in the part files sec/a06_BALtorus_d1.tex ... sec/a06_BALtorus_d7b.tex (groups (B)-(D)), then
% sec/a06_BALtorus_ct1.tex ... ct4.tex (Step 3(E): counterterm sizes, vacuum pieces) and ct5.tex ... ct8.tex (Step 3(F): the remaining Ba2 estimates), input below.

\paragraph{Step 3: the unprinted steps.} The statement $[\mathrm B\text{-base}]_5$ of \cref{thm:a06:BALA}(iii) is the list of the estimates which \cite{Ba1,Ba2,Ba3} assert without a printed argument, in four groups: (A) the Green's function estimates of \cite{Ba4} and \cite[Sect.~2]{Ba1} (the transfer of \cite[Cor.~2.3]{Ba4} to the propagators with mixed block pins, the domain-difference clause \cite[(1.20)]{Ba4}, the lower bound \cite[(1.15)]{Ba4}, the $j$-uniform analyticity neighbourhood of \cite[Lemma~2.4]{Ba4}, \cite[(1.8)]{Ba4} on a general region, \cite[Prop.~2.2]{Ba1}); (B) the steps of the lower bound (\cite[Props.~3.1--3.2, (3.20)--(3.21), (3.34)--(3.36), (3.66)--(3.69)]{Ba1}); (C) the steps of the upper bound (\cite[(2.11)--(2.14), Props.~2.1--2.2, 2.6, Lemma~2.7, (2.78), Sect.~3A, the $k$-uniformity of the per-step constants]{Ba2}); (D) \cite[Prop.~1]{Ba3} with its twenty sketched steps, the counterterm moves of \cite[p.~422]{Ba3} and the split \cite[(1.26)--(1.31)]{Ba3}. Of these, group (A) is proved in full: \cite[Prop.~2.3]{Ba1} (restated on \cite[p.~574]{Ba4}) is an implication of the printed statements \cite[(1.15), (1.16), (1.18), (1.20), (5.5), Cor.~2.3]{Ba4} read at source (with two harmless print findings: the equation numbers of \cite[Sect.~3]{Ba4} are shifted by four against the journal numbering, and the bracket of \cite[(3.9)]{Ba4} has the wrong sign, which the second resolvent identity corrects, only norms of the bracket being used), and the $j$-uniformity of \cite[Lemma~2.4]{Ba4} is the following theorem.

\begin{theorem}[{\cite[Lemma~2.4]{Ba4} with constants independent of the scale}]\label{thm:a06:A5}
Let $d\ge1$, $L\ge3$ odd, $0<a_-\le a_+$. There are $\delta_*>0$, $c_0$ and, for each $\alpha\in[0,1)$, $c_1(\alpha)$, depending only on $d,a_-,a_+,\alpha$, such that \cite[(2.35)--(2.36)]{Ba4} hold for every $j\ge0$, every $m_j^2\ge0$, every $a_j\in[a_-,a_+]$ and every rectangular parallelepiped (or periodic box) that is a union of unit blocks, with $\delta_0=\delta_*/2$ in the $\ell^1$ distance; \cite[(2.37)]{Ba4} holds with constants depending on $L$ but not on $j$.
\end{theorem}

\begin{proof}
The operator is $-\Delta^\xi+m^2+aQ_j^*Q_j$ on $\xi\Z^d$, $\xi=L^{-j}$, with $Q_j$ the average over the unit blocks. Its Fourier symbol on the unit-lattice momenta $p'\in[-\pi,\pi]^d$ is $E(p')=\Delta^\xi(p')(1+aS(p'))$ with $S(p')=\sum_{l}w(p'+l)/\Delta^\xi(p'+l)$, the sum over the $N^d$ aliases $l$ of $p'$ ($N=L^j$) and $w(p)=\prod_\mu g(p_\mu)^2$, $g(t)=\xi\sin(t/2)/\sin(\xi t/2)$ \cite[(2.44)--(2.48)]{Ba4}. On the strip $|\Im p'|_\infty\le\delta$: each alias term is bounded by elementary estimates on $|g|$ and on $|\Delta^\xi|$ from below, uniformly in $j$; the denominator $E$ is bounded below on the strip by a constant times $|\Delta^\xi(p')|$ at $\delta\le\delta_*$, since the alias sum $S$ is positive on the real section and its imaginary part is controlled by $\delta$ times a $j$-free constant; so the kernel $(G_jQ_j^*)(x,y)$ is the Fourier integral of a function analytic on the strip with $j$-free bounds, and shifting the contour by $\delta_*$ gives the exponential decay \cite[(2.35)]{Ba4} with $c_0$, $\delta_*$ free of $j$; the first-difference part of \cite[(2.35)]{Ba4} is the same integral with one factor $\hat\partial_\mu(p)=(e^{i\xi p_\mu}-1)/\xi$, and the H\"older quotient \cite[(2.36)]{Ba4} carries in addition $(e^{ip\cdot v'}-1)/|v'|^{\alpha}$, $v'=x-x'$; both are bounded on the strip uniformly in $\xi$. Neumann faces and periodic boxes are reduced to the full lattice by reflection and periodisation, which preserve the block structure since the box is a union of blocks. \cite[(2.37)]{Ba4} is the transfer to the mixed-pin operator and inherits the $L$-dependence of that operator's constants.
\end{proof}

Of groups (B)--(D), the following are proved from the printed text or by exact computation: \cite[(3.20)--(3.21)]{Ba1}, \cite[(3.66)--(3.69)]{Ba1}; \cite[(2.11)--(2.13)]{Ba2} given $|{\dm}|\le c\eps^{-1}$ (supplied by \cref{thm:a06e:CT1}); \cite[Prop.~2.2]{Ba2}; \cite[Lemma~2.7]{Ba2} given the path estimate of \cite[p.~560]{Ba2} at scale $k$; the replacement $A''\to A'$ of \cite[pp.~561--562]{Ba2} with error $O(\eps^\kappa)|\Lambda_7|$ for every $\kappa$; \cite[(2.78)]{Ba2}; \cite[(3.29)]{Ba2} through \cite[Prop.~3.1$'$]{Ba4}; \cite[p.~438]{Ba3}; and the interaction-term half of \cite[Prop.~2.6]{Ba2} at a nonzero background, stated as \BALint\ (\cref{hyp:k01:BALint}) and proved through classes complete at every background. The remaining pieces are proved in Steps~3(B)--3(F) below: the perturbative form of the effective action \cite[(3.34)--(3.36)]{Ba1} (Step~3(B), an algebraic identity between formal series); the counterterm moves of \cite[pp.~411--422]{Ba3}, i.e.\ the vacuum part of the effective action measured against the localised counterterms and its behaviour under the domain, localisation and counterterm changes of a step (Steps~3(D$_1$)--3(D$_2$)); the large-field chain of \cite[Sect.~3A]{Ba2} with the $k$-uniformity of the unqualified $O(1)$ lines of \cite{Ba2} and the $\eps_0$-freedom of the per-step constants (Steps~3(C$_1$)--3(C$_2$)); and the sketched steps of \cite[Prop.~1]{Ba3} itself, whose regrouping step is \RENorb\ (\cref{hyp:k01:RENorb}) and whose renormalisation cancellation is the induction of \cite[pp.~444--445]{Ba3} (Steps~3(D$_3$)--3(D$_4$)); the sizes of the mass counterterm, $|\dm|\le c_m\eps^{-1}$, $\dm_{K,k}(L^k\eps)^2=O(L^k\eps)$ and the finite part $\dm_{\mathrm{fin}}$ in the form in which it is consumed (Step~3(E), \cref{thm:a06e:CT1,thm:a06e:DM,thm:a06e:T0}), together with the vacuum counterterm pieces $E_{K,k}$, \cite[(2.14)]{Ba2} for the constant the chain carries, and \cite[Props.~3.1--3.2]{Ba1}, \cite[p.~579]{Ba2} given Proposition~1 (\cref{thm:a06e:VK,cor:a06e:C1b,prop:a06e:B12}); and the remaining asserted estimates of \cite[Sects.~2--3]{Ba2}: \cite[p.~568, (2.47)]{Ba2} (\cref{prop:a06f:C10}), the quadratic-form step \cite[(2.108)--(2.109)]{Ba2} (\cref{lem:a06f:C53,lem:a06f:C54}), the field-dependent part of \cite[Prop.~2.1, (2.57)]{Ba2} (\cref{thm:a06f:C2}), the passage from \cite[(3.5)]{Ba2} at $k=0$ to \cite[(3.21)]{Ba2} with the payment of \cite[Sect.~3C]{Ba2} (\cref{lem:a06f:C142,lem:a06f:C144}), the region sum \cite[(3.42)]{Ba2} at $L=5$ with the counterterm moves paid in it (\cref{thm:a06f:R7,prop:a06f:D2C}), and the $g_k$-localisation of $Q^{(k)}$ \cite[p.~579]{Ba2} (\cref{prop:a06f:H}) (Step~3(F)). Every constant in these steps is explicit and free of $\eps$, the step and the torus, and free of $\eps_0$ except through the caps on $\eps_0$ and the final-step constants that $[\mathrm B\text{-base}]_5$ allows (the bound on $\dm_{\mathrm{fin}}$ at fixed $\eps_0$, \cref{cor:a06e:T0p}, enters only through the $\eps_0$-free form of \cref{thm:a06e:T0}); so $[\mathrm B\text{-base}]_5$ holds and \cref{thm:a06:BALA}(iii) applies.

% a06_BALtorus_d1: part of Step 3 of the proof of thm:a06:T1 (input from sec/a06_BALtorus_c.tex).
% Group (B): the perturbative form of the effective action, Ba1 (3.34)-(3.36).

\subsubsection*{Notation for Steps 3(B)--3(D)}
Throughout the rest of Step 3, $d=3$, $L=5$, $k$ is the step, $t_k\defeq L^k\eps\le\eps_0\le1$ is Ba{\l}aban's scale parameter (written $t$ when $k$ is fixed), $\eta\defeq L^{-k}$, and all scale-$k$ operators act on the unit lattice $T^{(k)}_1$ of scale $k$ with the fine lattice $\eta\Z^3$ inside each unit cell. The rescaled couplings are $e_k\defeq e\,t_k^{1/2}$, $\lambda_k\defeq\lambda t_k$; $p(t)\defeq b_0(1+\ln t^{-1})^{\bar p}$ and $r(t)\defeq R(1+\ln t^{-1})^{\varrho}$, $\varrho>1$, are the field bound and the layer width of \cite[(2.7), (2.16)--(2.17)]{Ba2}; $\Lambda^{(k)}_i$, $i=-1,\dots,8$, are the regions of \cite[(2.8)]{Ba2} at scale $k$ (unions of large blocks, $\Lambda^{(k)c}_{i+1}$ the union of the large blocks within distance $r(t_k)$ of $\Lambda^{(k)c}_i$), and a primed region $\Lambda^{(k-1)'}_i$ is the scale-$(k-1)$ region read on $T^{(k)}_1$. $M$ is the side of Ba{\l}aban's large blocks (\cite[p.~607]{Ba1}: the block partition with $L$ replaced by $ML$; written $M_b$ in \cref{thm:a06:BALA}), not the scalar mass of \cref{thm:a06:T1}, which is $m$ or $M_1$ here. The expansion order is $\bar n$; $a_k$ is the block-spin coefficient of \cite[(2.13)]{Ba1}, with $a_{k+1}=aa_k/(aL^{-2}+a_k)$, $a_1=a$, so that $a_*\defeq a(1-L^{-2})\le a_k\le a$ for all $k$. $Q_k(A)$ is the covariant $k$-step block average, $P_k=Q_k^*Q_k$, $G_k(\Omega,A)=(-\Delta^{\eta,N}_{A,\Omega}+m_k^2+a_kP_k(A))^{-1}$ the scalar propagator on a union $\Omega$ of blocks with Neumann conditions, $\Delta^{(k)}(\Omega,A)=a_k-a_k^2Q_kG_k(\Omega,A)Q_k^*$ the block-spin operator, $C^{(k)}_\Lambda=((\Delta^{(k)}+aL^{-2}P)|_\Lambda)^{-1}$ the fluctuation covariance, as in \cite[Sect.~2]{Ba1} and \cite[Sect.~1]{Ba4}. For a kernel or external field, $\norm{\cdot}_{1,\alpha_0}$ is the norm \cite[(1.32)]{Ba3} (sup, one covariant derivative, its H\"older quotient of exponent $\alpha_0$), in each variable for kernels of two variables. ``Proposition~1'' is \cite[Prop.~1, (1.33)]{Ba3}: for a localised renormalised class $G_{\mathrm{ren}}$ with localisation $\Lambda=\{\square(v)\}$, $|E(G_{\mathrm{ren}},\Lambda,\Phi)|\le O(1)e_k^{d_v}\lambda_k^{d_s}e^{-\delta_0d(\Lambda)}\prod\norm{h\Phi}_{1,\alpha_0}$, $d(\Lambda)$ the length of a shortest tree through the cubes of $\Lambda$; its proof at $L=5$ is Step~3(D) below.

\subsubsection*{Step 3(B): the perturbative form of the effective action, \cite[(3.34)--(3.36)]{Ba1}}

\cite[(3.36)]{Ba1} asserts, ``by close inspection'', that after $k$ steps the polynomial part of the effective action is $S^{(k)}=\frac12\langle A,\Delta^{(k)}A\rangle+$ the truncation to order $\bar n$ of the formal series $\mathfrak E_k$ of \cite[(3.35)]{Ba1}, built with the contour convention \cite[(3.34)]{Ba1}. This is an algebraic identity between formal power series, proved here. The rescalings \cite[(1.22)]{Ba1} are changes of variables with field-independent Jacobians and change nothing below except additive constants.

\begin{definition}[grading, truncation, re-weighting]\label{def:a06d:B1}
\textup{(a)} In \cite[(3.35)]{Ba1} replace $e'$ by a formal variable $s$ wherever it occurs (on every fluctuation field, and in $\dm(e',\lambda')$, $E(e',\lambda')$), and read \cite[(3.34)]{Ba1} with $j$ from $0$ (which is \cite[(1.3)]{Ba3}). Write $\mathfrak E_k(s,\lambda';A,\phi)\defeq E_k(s,\lambda',eA^{(k),\eps},\phi)$, a formal power series in $(s,\lambda')$ whose coefficients are smooth functions of the block fields; the degree of a monomial is the power of $s$ plus the power of $\lambda'$. \textup{(b)} $T_{\bar n}$ truncates to degree $\le\bar n$ and ``$\mathrm{ev}$'' evaluates at $(s,\lambda')=(e,\lambda)$; in this notation \cite[(3.36)]{Ba1} reads $S^{(k)}=\frac12\langle A,\Delta^{(k)}A\rangle+\mathrm{ev}\,T_{\bar n}\mathfrak E_k$. \textup{(c)} In step $k+1$, $A=A'+aL^{-2}C^{(k)}Q^*B$ \cite[(3.41)]{Ba1} and $A^{(k),\eps}=A'^{(k),\eps}+A^{(k+1),\eps}$ \cite[(3.42)]{Ba1}; for a coefficient function $f$ of the background $eA^{(k),\eps}$, $\iota_kf$ is the Taylor series of $\tau\mapsto f(\tau A'^{(k),\eps}+eA^{(k+1),\eps})$ at $\tau=0$ with $\tau$ then replaced by $s$, so that each power of the new fluctuation carries one power of $s$ (Ba{\l}aban's order count: every $A'$- or $\tilde A$-leg in \cite[(1.8)--(1.15)]{Ba3} carries one factor $e(L^k\eps)$).
\end{definition}

\begin{lemma}[truncation commutes with a Gaussian integral]\label{lem:a06d:B2}
Let $d\mu\propto e^{-F_0}dX$ be a Gaussian law on a finite-dimensional space ($F_0$ a positive quadratic form), and let $W$ be a formal series of degree $\ge1$ whose coefficients are polynomials in $X$. Then $-\log\langle e^{-W}\rangle_\mu$ is a well-defined formal series; its degree-$n$ part is a finite combination of Gaussian moments of the parts $W_1,\dots,W_n$ of $W$ of degrees $1,\dots,n$; hence $T_{\bar n}(-\log\langle e^{-W}\rangle_\mu)=T_{\bar n}(-\log\langle e^{-T_{\bar n}W}\rangle_\mu)$.
\end{lemma}

\begin{proof}
$\langle e^{-W}\rangle=\sum_{m\ge0}(-1)^m\langle W^m\rangle/m!$; the degree-$n$ part of $W^m$ vanishes for $m>n$ and for $m\le n$ involves only $W_1,\dots,W_{n-m+1}$. So $\langle e^{-W}\rangle=1+u$ with $u$ of degree $\ge1$, and the degree-$n$ part of $-\log(1+u)=\sum_{m\ge1}(-1)^mu^m/m$ involves only the parts of $u$ of degree $\le n$. Each coefficient is a Gaussian moment of a polynomial, hence finite.
\end{proof}

This is the cumulant expansion: Ba{\l}aban's ``terms of the order $\le\bar n$'' of an integrated action are exactly $T_{\bar n}$ of this series.

\begin{lemma}[composition of the scalar kernels]\label{lem:a06d:B3}
Let $Q\colon\R^n\to\R^m$ be linear with $QQ^*=c\,\mathbf 1_m$ ($c>0$), and $\alpha,\beta>0$. Then for all $\psi\in\R^m$, $\chi\in\R^n$,
\[
\int_{\R^n}d\phi\,\exp\big(-\tfrac12\alpha|\psi-Q\phi|^2-\tfrac12\beta|\phi-\chi|^2\big)=(2\pi/\beta)^{n/2}(1+\alpha c/\beta)^{-m/2}\exp\big(-\tfrac12\gamma|\psi-Q\chi|^2\big),\qquad\gamma\defeq\frac{\alpha\beta}{\beta+\alpha c}.
\]
\end{lemma}

\begin{proof}
Write $\phi=\chi+\beta^{-1/2}\xi$ with $\xi$ standard Gaussian; then $Q\phi=Q\chi+Z$ with $Z\sim N(0,(c/\beta)\mathbf 1_m)$, and for $Z\sim N(0,\sigma^2\mathbf 1_m)$, $\E\exp(-\frac12\alpha|v+Z|^2)=(1+\alpha\sigma^2)^{-m/2}\exp(-\frac12\frac{\alpha}{1+\alpha\sigma^2}|v|^2)$. Take $v=Q\chi-\psi$, $\sigma^2=c/\beta$.
\end{proof}

The covariant one-step average $(Q(\beta)\phi)(z)=L^{-d}\sum_{y\in B(z)}U(\beta(\Gamma_{z,y}))\phi(y)$ has $Q(\beta)Q(\beta)^*=L^{-d}\mathbf 1$ for every background $\beta$ (the $U$ are orthogonal, the blocks disjoint). With $\alpha=aL^{d-2}$ \cite[(3.38)]{Ba1}, $\beta=a_k$ and $c=L^{-d}$ the lemma gives $\gamma=L^{d-2}aa_k/(a_k+aL^{-2})=L^{d-2}a_{k+1}$, the recursion \cite[(2.13)]{Ba1}; the factor $L^{d-2}$ is removed by the rescaling of $\psi$, and the prefactor depends neither on the background nor on $\psi,\chi$. At a fixed external field this is \cite[(2.10)--(2.16)]{Ba1}; the next two lemmas extend it to level-dependent fields.

\begin{lemma}[composite averages]\label{lem:a06d:B4}
Let $\beta$ be a fine field and $\Phi$ assign a phase to each composite contour $\Gamma^{(k)}_{y,x}$. Then $(Q(\beta)Q_k(\Phi)\phi')(z)=(Q_{k+1}(\Phi')\phi')(z)$ with $\Phi'(\Gamma^{(k+1)}_{z,x})\defeq\beta(\Gamma_{z,x_k})+\Phi(\Gamma^{(k)}_{x_k,x})$.
\end{lemma}

\begin{proof}
$U(a)U(b)=U(a+b)$ since all $U$ commute; $B(z)\circ B^k(y)=B^{k+1}(z)$ and $\Gamma^{(k+1)}_{z,x}=\Gamma_{z,x_k}\cup\Gamma^{(k)}_{x_k,x}$ by \cite[(2.2)]{Ba1}; the weights $L^{-d}\eta^d$ are those of the $(k+1)$-average.
\end{proof}

\begin{lemma}[the contour convention is the iterate, and $j$ starts at $0$]\label{lem:a06d:B5}
For $k\ge1$ put $\Phi_k(s)(\Gamma^{(k)}_{y,x})\defeq eA^{(k),\eps}(\Gamma^{(k)}_{y,x})+s\sum_{j=0}^{k-1}A'^{(j),\eps}(\Gamma^{(j+1)}_{x_{j+1},x})$, which is \cite[(3.34)]{Ba1} with $j$ from $0$ and the fluctuations weighted by $s$. \textup{(a)} With $\beta\defeq\iota_k(eA^{(k),\eps})=sA'^{(k),\eps}+eA^{(k+1),\eps}$ on the first-step contour $\Gamma_{z,x_k}$ prescribed by \cite[(3.37)]{Ba1}, $\beta(\Gamma_{z,x_k})+(\iota_k\Phi_k(s))(\Gamma^{(k)}_{x_k,x})=\Phi_{k+1}(s)(\Gamma^{(k+1)}_{z,x})$. \textup{(b)} At $k=0$ the first kernel carries $U(eA(\Gamma_{z,x}))$ with the full fine field $A=A'_0+A^{(1),\eps}$ \cite[(3.15), p.~615]{Ba1}, and $\iota_0(eA(\Gamma_{z,x}))=sA'_0(\Gamma^{(1)}_{x_1,x})+eA^{(1),\eps}(\Gamma^{(1)}_{z,x})=\Phi_1(s)(\Gamma^{(1)}_{z,x})$: its $s$-term is the $j=0$ term. \textup{(c)} Hence the printed range $j=1,\dots,k-1$ of \cite[(3.34)]{Ba1} is not what the procedure produces: at $k=1$ it would drop $sA'_0(\Gamma^{(1)})$ and with it every contour vertex \cite[(1.14)]{Ba3} with $A'_0$-legs, which the first step does produce \cite[(3.15)--(3.16)]{Ba1}; with $j=0,\dots,k-1$ the convention is the iterate.
\end{lemma}

\begin{proof}
(a) By \cite[(2.2)]{Ba1}, $\Gamma^{(k+1)}_{z,x}=\Gamma_{z,x_k}\cup\Gamma^{(k)}_{x_k,x}$ and $\Gamma^{(k+1)}_{x_{k+1},x}=\Gamma^{(k+1)}_{z,x}$ since $x_{k+1}=z$; by \cite[(2.3)]{Ba1} line integrals add over unions of contours, so $\beta(\Gamma_{z,x_k})+sA'^{(k)}(\Gamma^{(k)}_{x_k,x})+eA^{(k+1)}(\Gamma^{(k)}_{x_k,x})=sA'^{(k)}(\Gamma^{(k+1)}_{x_{k+1},x})+eA^{(k+1)}(\Gamma^{(k+1)}_{z,x})$, and the terms with $j<k$ are unchanged because the points $x_j$, $j\le k$, of $x$ are the same at both levels. (b) and (c) are read off the definitions.
\end{proof}

\begin{lemma}[{the Gaussian chain on a domain; \cite[(3.41), (3.42)]{Ba1} and \cite[(1.27)]{Ba3}}]\label{lem:a06d:B6}
Let $H>0$ be the covariant Neumann operator $-\Delta^\eta_{B,\Omega}+m^2$ on a union $\Omega$ of blocks at a fixed background $B$, or the vector operator at $B=0$ componentwise; let $R_k$ be the $k$-step composite average onto $T^{(k)}_1\cap\Omega$, $Q$ the one-step average with $QQ^*=c\,\mathbf 1$, and $\beta_k,\alpha>0$. Consider the Gaussian density of $(\phi',\phi,\psi)$ proportional to $\exp(-\frac12\langle\phi',H\phi'\rangle-\frac12\beta_k|\phi-R_k\phi'|^2-\frac12\alpha|\psi-Q\phi|^2)$, and put $G_k\defeq(H+\beta_kR_k^*R_k)^{-1}$, $\Delta^{(k)}\defeq\beta_k-\beta_k^2R_kG_kR_k^*$, $C^{(k)}\defeq(\Delta^{(k)}+\alpha Q^*Q)^{-1}$, $\gamma\defeq\alpha\beta_k/(\beta_k+\alpha c)$, $G_{k+1}\defeq(H+\gamma(QR_k)^*(QR_k))^{-1}$. Then: \textup{(a)} given $\phi$, $\phi'\sim N(\beta_kG_kR_k^*\phi,G_k)$, and the marginal density of $\phi$ is $\propto\exp(-\frac12\langle\phi,\Delta^{(k)}\phi\rangle)$; \textup{(b)} given $\psi$, $\phi\sim N(\alpha C^{(k)}Q^*\psi,C^{(k)})$, and $\E[\phi'\mid\psi]=\beta_kG_kR_k^*\alpha C^{(k)}Q^*\psi=\gamma G_{k+1}(QR_k)^*\psi$; \textup{(c)} $\Cov(\phi'\mid\psi)=G_{k+1}=G_k+\beta_k^2G_kR_k^*C^{(k)}R_kG_k$.
\end{lemma}

(a) is the definition of $A^{(k)}$, $\phi^{(k)}$ and $\Delta^{(k)}$ \cite[(3.29), (2.21)]{Ba1}; (b) is \cite[(3.41)]{Ba1} (the conditional fluctuation $A'=A-\E[A\mid B]$ is independent of $B$ with covariance $C^{(k)}$) together with the second equality of \cite[(3.42)]{Ba1}, $a_kG_kQ_k^*\,aL^{-2}C^{(k)}Q^*B=a_{k+1}L^{-2}G^\eta_{k+1}Q^*_{k+1}B$; (c) is, on the torus at $B=0$, one step of \cite[(1.27)]{Ba3}, and on $\Omega$ the same identity with Neumann lines.

\begin{proof}
(a): complete the square in $\phi'$. (b): by (a) the $(\phi,\psi)$-density is $\propto\exp(-\frac12\langle\phi,\Delta^{(k)}\phi\rangle-\frac12\alpha|\psi-Q\phi|^2)$; complete the square in $\phi$. The full density factorises as $f(\phi',\phi)g(\phi,\psi)$, so $\phi'$ and $\psi$ are conditionally independent given $\phi$ and $\E[\phi'\mid\psi]=\E[\E[\phi'\mid\phi]\mid\psi]=\beta_kG_kR_k^*\E[\phi\mid\psi]$. Integrating $\phi$ first (\cref{lem:a06d:B3} with $\chi=R_k\phi'$), the $(\phi',\psi)$-density is $\propto\exp(-\frac12\langle\phi',H\phi'\rangle-\frac12\gamma|\psi-QR_k\phi'|^2)$, whence $\E[\phi'\mid\psi]=\gamma G_{k+1}(QR_k)^*\psi$ and $\Cov(\phi'\mid\psi)=G_{k+1}$. (c): the law of total covariance, $\Cov(\phi'\mid\psi)=\E[\Cov(\phi'\mid\phi,\psi)\mid\psi]+\Cov(\E[\phi'\mid\phi,\psi]\mid\psi)=G_k+\beta_k^2G_kR_k^*C^{(k)}R_kG_k$.
\end{proof}

\begin{lemma}[diagonal truncation]\label{lem:a06d:B7}
Let $H(s,\tau,\lambda')$ be a formal series and $h(s,\lambda')\defeq H(s,s,\lambda')$. Then $T_{\bar n}h$ at $(s,\lambda')=(e,\lambda)$ equals $T^{\mathrm{tot}}_{\bar n}H$ at $(s,\tau,\lambda')=(e,e,\lambda)$, where $T^{\mathrm{tot}}_{\bar n}$ truncates the total degree.
\end{lemma}

\begin{proof}
The coefficient of $s^m\lambda'^\beta$ in $h$ is $\sum_{\alpha+\gamma=m}H_{\alpha\gamma\beta}$.
\end{proof}

\begin{theorem}[{\cite[(3.36)]{Ba1} at every $k$}]\label{thm:a06d:B8}
Let $S^{(0)}\defeq S^\eps$ \cite[(1.20)]{Ba3}, and for $k\ge0$ let $S^{(k+1)}_{\mathrm{pol}}$ be the polynomial part, of order $\le\bar n$, that the step $k+1$ of \cite[Sect.~3]{Ba1} produces from $S^{(k)}_{\mathrm{pol}}$: the kernels of \cite[(3.37)]{Ba1} (gauge $T_{a,L}$, scalar $T_{a,L,A^{(k),\eps}}$), the translation \cite[(3.41)]{Ba1}, the expansion of all coefficients in the new fluctuation \cite[(3.44)]{Ba1}, the Gaussian integral and the cumulant expansion keeping the terms of order $\le\bar n$ (the restrictions $\chi$ and the terms of order $>\bar n$ are split off; they are the error terms of \cite[(3.26)]{Ba1} and are not part of this statement). Then for every $k\ge0$,
\[
S^{(k)}_{\mathrm{pol}}=\tfrac12\langle A,\Delta^{(k)}A\rangle+\mathrm{ev}\,T_{\bar n}\mathfrak E_k+c_k,
\]
where $\mathfrak E_k$ is \cite[(3.35)]{Ba1} with \cite[(3.34)]{Ba1} read from $j=0$, and $c_k$ is a constant independent of the fields and of the couplings (a normalisation, entering $E_0$ only). With the printed range $j\ge1$ the identity is false at $k=1$ (\cref{lem:a06d:B5}(c)).
\end{theorem}

\begin{proof}
Induction on $k$. \emph{Base $k=0$.} $\mathfrak E_0$ contains no fluctuation: $\mathfrak E_0=\frac12\langle\phi,(-\Delta_{eA}+m^2)\phi\rangle+\sum\eps^d(\lambda'|\phi|^4+\frac12\dm(s,\lambda')|\phi|^2)+E(s,\lambda')+\mathrm{const}$, and $\frac12\langle A,\Delta^{(0)}A\rangle=\frac12\langle A,(-\Delta+\lmu^2)A\rangle$; $T_{\bar n}$ acts as the identity, since $\dm$ has orders $2\le\alpha+2\beta\le4\le\bar n$ \cite[p.~417]{Ba3} and $E$ has orders $\le\bar n$. So $S^{(0)}=\frac12\langle A,\Delta^{(0)}A\rangle+\mathrm{ev}\,T_{\bar n}\mathfrak E_0$. At $k=0$ the inner kernel is absent ($Q_0=1$): the gauge part of the step is \cref{lem:a06d:B6}(b) with $H=\Delta^{(0)}$, and the scalar part is \cite[(3.35)]{Ba1} at $k=1$ by definition, with $a_1=a$ and the $j=0$ contour term of \cref{lem:a06d:B5}(b).

\emph{Step $k\to k+1$.} Define
\[
\mathrm{Step}_kF(B,\psi)\defeq-\log\Big[N_k\int dA\int d\phi\,\exp\big(-\tfrac12aL^{d-2}|B-QA|^2-\tfrac12aL^{d-2}|\psi-Q(\iota_keA^{(k),\eps})\phi|^2\big)e^{-\iota_kF(A,\phi)}\Big],
\]
followed by the rescaling \cite[(1.22)]{Ba1}.

(i) \emph{Exact identity} for $F=\frac12\langle A,\Delta^{(k)}A\rangle+\mathfrak E_k$, degree by degree (Fubini for Gaussian moments of polynomials). Gauge part: by \cref{lem:a06d:B6}(b),
$\int dA\exp(-\frac12aL^{d-2}|B-QA|^2-\frac12\langle A,\Delta^{(k)}A\rangle)f(A^{(k),\eps})=N\exp(-\frac12\langle B,\Delta^{(k+1)}B\rangle)\int d\mu_{C^{(k)}}(A'_k)f(A'^{(k),\eps}+A^{(k+1),\eps})$, and $\iota_k$ puts the factor $s$ on $A'^{(k),\eps}$. Scalar part: write $e^{-\mathfrak E_k}$ as the nested integral \cite[(3.35)]{Ba1}; the $\phi$-integral is $\int d\phi\exp(-\frac12\alpha|\psi-Q(\beta)\phi|^2-\frac12a_k|\phi-Q_k(\Phi_k)\phi'|^2)=N'\exp(-\frac12a_{k+1}|\psi-Q(\beta)Q_k(\Phi_k)\phi'|^2)$ by \cref{lem:a06d:B3}, and $Q(\beta)Q_k(\iota_k\Phi_k)=Q_{k+1}(\Phi_{k+1})$ by \cref{lem:a06d:B4,lem:a06d:B5}(a). Covariant Laplacian: on every bond $b$, $(\iota_k\Phi_k)_b=sA'^{(k)}_b+eA^{(k+1)}_b+s\sum_{j<k}A'^{(j)}_b=(\Phi_{k+1})_b$. The $\lambda'$-terms and $\dm(s,\lambda')$, $E(s,\lambda')$ are unchanged. Hence $\mathrm{Step}_k(\frac12\langle A,\Delta^{(k)}A\rangle+\mathfrak E_k)=\frac12\langle B,\Delta^{(k+1)}B\rangle+\mathfrak E_{k+1}+c'_k$.

(ii) \emph{Truncation.} The degree-$0$ part of the integrand of $\mathrm{Step}_k$ is Gaussian in $X=(A'_k,\phi)$: its gauge part is given by (i), its scalar part is $\frac12\langle\phi,\Delta^{(k)}(eA^{(k+1),\eps})\phi\rangle+\frac12\alpha|\psi-Q(eA^{(k+1),\eps})\phi|^2$. Since $\iota_k$ only raises degrees, $T_{\bar n}\iota_kT_{\bar n}=T_{\bar n}\iota_k$, and \cref{lem:a06d:B2} gives $T_{\bar n}\mathrm{Step}_k(T_{\bar n}F)=T_{\bar n}\mathrm{Step}_k(F)$.

(iii) \emph{Evaluation.} Ba{\l}aban's step expands in the new fluctuation at the true coupling $e$ and keeps total order $\le\bar n$; by \cref{lem:a06d:B7} with $\tau$ the weight of the new fluctuation ($\iota_k$ is the diagonal $\tau=s$), his truncation is $\mathrm{ev}\,T_{\bar n}$ of the diagonal. Hence $S^{(k+1)}_{\mathrm{pol}}=\mathrm{ev}\,T_{\bar n}\mathrm{Step}_k(S^{(k)}_{\mathrm{pol}})=\frac12\langle B,\Delta^{(k+1)}B\rangle+\mathrm{ev}\,T_{\bar n}\mathfrak E_{k+1}+c_{k+1}$.
\end{proof}

The theorem is the polynomial-part identity that ``close inspection'' asserts, an implication of the printed definitions \cite[(2.1)--(2.3), (3.29), (3.33), (3.35), (3.37)]{Ba1}, \cite[(1.2)]{Ba4} and of \cite[(3.34)]{Ba1} read with $j$ from $0$, i.e.\ of the defining formula \cite[(1.3)]{Ba3}. It does not cover the analytic content of a step (the restrictions, the order-$>\bar n$ remainders, the error terms of \cite[(3.26)]{Ba1}): that is \cite[Props.~3.1--3.2]{Ba1}, which follow from Proposition~1 (Step~3(D)).

% a06_BALtorus_d2: part of Step 3 of the proof of thm:a06:T1 (input from sec/a06_BALtorus_c.tex).
% Group (D), first half: the counterterm moves of Ba3 pp. 411-422, the vacuum functionals of a datum and their exact algebra.

\subsubsection*{Step 3(D$_1$): the counterterm moves of \cite[pp.~411--422]{Ba3}, exact algebra}

Two statements of \cite{Ba3} on the vacuum part of the effective action carry no printed argument: that the field-independent part generated after $k$ steps on the domains of the step, measured against the localised counterterms $E_k(\cdot,\Omega_1)$ of \cite[(1.30)]{Ba3}, is $O(t_k)$ per unit block (``obvious consequences of Proposition~1'', \cite[pp.~421--422]{Ba3}), and that the domain and localisation changes of a step, with the counterterm move of \cite[p.~422]{Ba3}, cost only collar terms (``even simpler and we omit them''). We write them as the statements $(\mathrm{VAC}^\flat)_k$ and $(\mathrm{DOM})_k$ (\cref{thm:a06d:V12,thm:a06d:V14}) and prove them from Proposition~1. This part sets up the objects and their exact identities; the bounds are in Step~3(D$_2$).

\begin{definition}[datum at scale $k$]\label{def:a06d:V1}
A datum is $\mathfrak D=(\Omega,\Omega_1,g,c)$: $\Omega$ a union of big blocks of $T^{(k)}_1$ (the domain hypothesis of \cite[Props.~2.1--2.3]{Ba1}, \cite[p.~411]{Ba3}), on whose $\eta$-points the fine scalar field lives with the Neumann covariant Laplacian; $\Omega_1\subset\Omega$ a union of unit blocks; $g$ smooth with $\supp g\subset\Omega_1$, $g=1$ where $\dist(\cdot,\partial\Omega_1)\ge M$ and $\norm g_{C^2}\le C_g$ (\cite[p.~412]{Ba3} has $C_g=1$; a fixed $C_g$ is allowed, as for the partitions of unity of \cite[p.~420]{Ba3}); $c(x)$ the coefficient of the vertices \cite[(1.7)]{Ba3}, namely \cite[(1.29)]{Ba3} built with $g$ at the scale of $\mathfrak D$, $c=\dm_{\mathrm{fin}}t_k^2+\dm_{K,k}t_k^2+\dm_k(se(t_k)g,\lambda'\lambda(t_k),\Omega_1,x)$. The depth is $\rho(\mathfrak D)\defeq\dist(\Omega_1,\Omega^c)$ in scale-$k$ units. Ba{\l}aban's step-$k$ datum is $\mathfrak D_k\defeq(B^k(\Lambda^{(k-1)'}_2),B^k(\Lambda^{(k-1)'}_7),g_k,c_k)$; torus data are $\mathfrak D^T_{W,g}\defeq(T,W,g,c_{W,g})$, and the plain torus datum $\mathfrak D^T$ has $W=T$, $g=1$, $c=\dm$.
\end{definition}

\begin{definition}[the vacuum functionals of a datum]\label{def:a06d:V2}
$\Xi_k[\mathfrak D](A,\phi)$ is the nested integral inside the logarithm of \cite[(1.4)]{Ba3}, with $(\Omega,\Omega_1,g_k)$ and the coefficient replaced by those of $\mathfrak D$, $e'$ replaced by $s$, and the counterterm factor $e^{-E_1}$ omitted (the vector covariances are the torus $C^{(j)}$ of \cite[(1.4)]{Ba3}). $\mathrm{Gen}_k[\mathfrak D]\defeq T_{[1,\bar n]}\log[\Xi_k[\mathfrak D](0,0)/\Xi_k[\mathfrak D](0,0)|_{s=\lambda'=0}]$ (the generated vacuum constants); $F_k[\mathfrak D](A,\phi)\defeq T_{[0,\bar n]}\log[\Xi_k[\mathfrak D](A,\phi)/\Xi_k[\mathfrak D](0,0)]$ (the field-dependent part, including its degree-$0$ Gaussian part); $\mathrm{Ct}_k[\mathfrak D]\defeq E_k(se(t_k)g,\lambda'\lambda(t_k),\Omega_1)$ of \cite[(1.30)]{Ba3}, under \cref{conv:a06d:LOC}; $\mathcal R_k[\mathfrak D]\defeq\mathrm{Gen}_k[\mathfrak D]-\mathrm{Ct}_k[\mathfrak D]$.
\end{definition}

By \cite[(1.4)]{Ba3}, $E_k(\cdot,\Omega,0,0)=E_1(e',g_k,\lambda',\Omega_1)-\log\Xi_k[\mathfrak D_k](0,0)$, and by \cite[(1.30)]{Ba3}, $E_1(e',g_k,\lambda',\Omega_1)=\Efin+E_{K,k}+\mathrm{Ct}_k[\mathfrak D_k]$; so the field-independent part of the polynomial $\mathcal P^{(k)}(\Omega_1,\cdot)$ of \cite[(1.5)]{Ba3}, in the logarithmic convention, is $\mathcal R_k[\mathfrak D_k]-\Efin-E_{K,k}$, with $\Efin+E_{K,k}$ a global front factor.

\begin{convention}[how the localised counterterms are built]\label{conv:a06d:LOC}
\textup{(a)} Every vertex \cite[(1.7)]{Ba3} carries the coefficient $c$ of \cite[(1.29)]{Ba3} as a number, not opened, both in the expansion of \cite[(1.4)]{Ba3} and in the graphs of $\mathrm{Ct}_k$. \textup{(b)} $\mathrm{Ct}_k[\mathfrak D]$ is the sum of the bare vacuum graphs (vertices \cite[(1.6), (1.7), (1.8), (1.10)]{Ba3} with $n'=0$, $\tilde B=0$) of orders $1\le\alpha+\beta\le\bar n$, with every line of index $<k$: scalar lines $G_k(T,0)$ (the rescaled $G^\eps_k(0)$ of \cite[(1.25)--(1.28)]{Ba3}) and vector lines $G^{\mathrm{vec}}_k(T)$, the vertices localised exactly as the generated ones of \cite[(1.4)]{Ba3}: \cite[(1.6)--(1.7)]{Ba3} summed over $\Omega_1$, $g$ on each vector leg. \textup{(c)} In $\dm_k(se(t_k)g,\dots)$ the factor $g$ sits on the vector legs of the coefficient's graphs, so a divergent self-energy subgraph of a generated graph and its \cite[(1.7)]{Ba3} partner (\cite[p.~420]{Ba3}: ``the same graph $G$ but with both external legs localized in $x$'') carry $g$ on corresponding legs.
\end{convention}

(b) and (c) are printed: \cite[(1.31), p.~419]{Ba3} displays the localised second-order photon counterterm with torus lines and $g_k$ on both vector legs, and \cite[p.~419]{Ba3} reads ``localize the vertices (1.6) and (1.7) restricting summations to the set $\Omega_1$. The other vertices are localized by functions $g_k$ multiplying each leg of the vector field''. (a) is forced: if the \cite[(1.7)]{Ba3} vertices of $\mathrm{Ct}$ carried the full torus $\dm$, $\mathrm{Ct}$ would contain on $\{g<1\}$ self-energy subgraphs with legs $gg'$ paired with an unlocalised coefficient, an uncancelled $O(e(t)^2\eta^{-1})$ per site; if the coefficients were opened into lines of index $<k$, $\mathrm{Gen}-\mathrm{Ct}$ would contain $(\dm_{\mathrm{fin}}t^2+\dm_{K,k}t^2)G_k(x,x)=O(t\eta^{-1})$ per site. On the torus datum, $\Efin+E_{K,k}+\mathrm{Ct}_k[\mathfrak D^T]=E_1$. (Ba{\l}aban remarks on \cite[p.~418]{Ba3} that it ``is sufficient'' to keep the orders $2\le\alpha+2\beta\le6$ in $E_1$; under that restricted $E_1$ the convergent bare graphs of higher order are extra standalone renormalised families in $\mathcal R_k$, bounded by \cref{lem:a06d:V9}(a) below, and nothing else changes.)

\begin{definition}[graphs, families, localisations]\label{def:a06d:V4}
Expanding $\log\Xi_k[\mathfrak D](0,0)$ in $(s,\lambda')$ gives connected vacuum graphs \cite[pp.~414--415]{Ba3}, and $\mathrm{Gen}_k=\mathrm{Gen}^{\mathrm{bare}}_k+\mathfrak B_k$: bare graphs (all vertices from \cite[(1.6), (1.7), (1.8), (1.10)]{Ba3}, from the expansion of $-\Delta^\eta_{sgA'}$ and of $V$) and contour graphs (at least one product \cite[(1.18)]{Ba3} of the form $(1.13)\cdot(1.14)$ or $(1.14)\cdot(1.14)$, from the expansion of $a_k|Q_k(sgA')\phi'|^2-a_k|Q_k(0)\phi'|^2$; \cite[(1.12)]{Ba3} does not occur since the block field is $0$). Scalar lines are $G_k(\Omega,0)$ (at $\phi=0$ the Gaussian part is $-\frac12\langle\phi',(-\Delta^\eta+m^2t_k^2+a_kP_k(0))\phi'\rangle$ on $\Omega$); vector lines are $G^{\mathrm{vec}}_k(T)$, the covariance of $A'$ \cite[p.~414]{Ba3}. Localisation follows \cite[p.~420]{Ba3}, the families are the $G_{\mathrm{ren}}$ of \cite[p.~420]{Ba3}, $X(\Lambda)$ is the sum of the terms of $X$ with localisation $\Lambda=\{\square(v)\}$, and $d(\Lambda)$ the tree length of \cite[(1.33)]{Ba3}.
\end{definition}

\begin{lemma}[the generated bare part and the counterterm differ only in their lines]\label{lem:a06d:V5}
Under \cref{conv:a06d:LOC}, $\mathrm{Gen}^{\mathrm{bare}}_k[\mathfrak D]$ and $\mathrm{Ct}_k[\mathfrak D]$ are sums over the same graphs, with the same localisations and vertex coefficients, with lines $(G_k(\Omega,0),G^{\mathrm{vec}}_k(T))$ and $(G_k(T,0),G^{\mathrm{vec}}_k(T))$ respectively. Hence $\mathcal R_k[\mathfrak D]=\mathfrak B_k[\mathfrak D]+\Delta_k[\mathfrak D]$ with
\[
\Delta_k[\mathfrak D]\defeq\sum_G\sum_\ell G\big[G_k(T,0)\text{ on the scalar lines before }\ell;\ \delta_k\text{ on }\ell;\ G_k(\Omega,0)\text{ after }\ell\big],\qquad\delta_k\defeq G_k(\Omega,0)-G_k(T,0)=-\delta G_k(T,\Omega,0),
\]
the scalar lines of each $G$ ordered so that the lines of each closed scalar loop are consecutive.
\end{lemma}

\begin{proof}
At zero block fields the bare vertex functions of \cite[(1.4)]{Ba3} come from the Taylor expansion of $U(sgA'_b)=\exp(qs\eta g(b)A'_b)$ inside $-\Delta^\eta_{sgA'}$ and from $V$: they are \cite[(1.8), (1.10)]{Ba3} with $n'=0$, $\tilde B=0$, together with \cite[(1.6), (1.7)]{Ba3}; the bare vacuum graphs of $S^\eps$ \cite[p.~416]{Ba3} come from the same expansion of the lattice covariant Laplacian, rescaled, and \cref{conv:a06d:LOC}(a), (b) give them the same localisation and coefficients. The generated scalar lines are $G_k(\Omega,0)$; the counterterm lines have index $<k$, and by \cite[(1.27)]{Ba3} the index-$<k$ part of $G^\eps_K(0)$ is $G^\eps_k(0)$, which rescales to $G_k(T,0)$; the vector lines agree. Multilinearity in the scalar lines, $\prod_\ell a_\ell-\prod_\ell b_\ell=\sum_\ell(\prod_{\ell'<\ell}b_{\ell'})(a_\ell-b_\ell)(\prod_{\ell'>\ell}a_{\ell'})$, gives $\Delta_k$. $S^\eps$ has no block kernel, so no vertex \cite[(1.13)--(1.15)]{Ba3} occurs in $\mathrm{Ct}_k$, and the contour graphs, which have no counterterm, stay in $\mathfrak B_k$.
\end{proof}

The single terms of $\Delta_k$ are never bounded one by one: the vacuum graph made of one \cite[(1.7)]{Ba3} vertex and one scalar loop gives $c(x)\delta_k(x,x)$ with $c\sim e(t)^2\eta^{-1}$, cancelled only together with the graph carrying the self-energy $S$ with $\delta_k$ on the closing line, $[S-c]\delta_k=[S-S(T)]\delta_k+[S(T)-c]\delta_k$. So $\Delta_k$, and every difference below, is used only through renormalised families whose members carry the telescoping order induced from $G$ on $G/S$ (\cref{lem:a06d:K2}); with the ordering of \cref{lem:a06d:V5} every closed scalar loop not containing $\ell$ carries uniform full lines, all $G_k(T,0)$ or all $G_k(\Omega,0)$, and no closed loop is ever decomposed by fixed scales or orderings. In the measure convention of \cite[(2.1)]{Ba3} a lone tadpole has $D=-1$ and Proposition~1 never bounds it alone.

\begin{lemma}[exact vacuum staging on a fixed datum]\label{lem:a06d:V6}
Fix a datum $\mathfrak D$ at scale $k$ and let $\mathfrak D\!\uparrow$ be the same datum read at scale $k+1$ (same fine domains, vertex functions and coefficient $c$). Let $d\mu_k$ be the Gaussian law of the step-$k$ fluctuations at zero $(k+1)$-block fields with the composing covariances $C^{(k)}$ (vector, torus) and $C^{(k)}(\Omega)\defeq(\Delta^{(k)}(\Omega,0)+\alpha Q^*Q)^{-1}$ (scalar, \cref{lem:a06d:B6} with $H$ the Neumann operator on $\Omega$, $B=0$). For a datum with localisation $g$ let $\iota^g_kf$ be the Taylor series of $\tau\mapsto f(\tau gA'^{(k),\eps}+eA^{(k+1),\eps})$ at $\tau=0$ with $\tau\to s$ ($\iota^g_k=\iota_k$ on the torus; \cref{lem:a06d:B4,lem:a06d:B5}(a) hold verbatim for it, line integrals of $gA'$ being additive). Put
\[
\mathrm{out}_k[\mathfrak D]\defeq T_{[1,\bar n]}\log\int d\mu_k\exp\big(\iota^g_kF_k[\mathfrak D]-(\iota^g_kF_k)_0+\mathcal K_k\big),\quad(\iota^g_kF_k)_0\defeq F_k|_{s=\lambda'=0,A=0}=-\tfrac12\langle\phi,\Delta^{(k)}(\Omega,0)\phi\rangle,\quad\mathcal K_k\defeq-\tfrac12\alpha\big(|Q(sgA'^{(k)})\phi|^2-|Q(0)\phi|^2\big),
\]
the vacuum constant generated by step $k+1$. Then: \textup{(a)} $\mathrm{Gen}_{k+1}[\mathfrak D\!\uparrow]=\mathrm{Gen}_k[\mathfrak D]+\mathrm{out}_k[\mathfrak D]$; \textup{(b)} the lines of $\mathrm{Gen}_{k+1}[\mathfrak D\!\uparrow]$ are $G_{k+1}(\Omega,0)=G_k(\Omega,0)+\beta_k^2G_k(\Omega,0)R_k^*C^{(k)}(\Omega)R_kG_k(\Omega,0)$ and $G^{\mathrm{vec}}_{k+1}=G^{\mathrm{vec}}_k+(\text{lifted }C^{(k)})$, so $\mathrm{out}_k[\mathfrak D]$ is the part of $\mathrm{Gen}_{k+1}[\mathfrak D\!\uparrow]$ with at least one line of index $k$; \textup{(c)} with $\mathcal V_k[\mathfrak D]\defeq\mathrm{Ct}_{k+1}[\mathfrak D\!\uparrow]-\mathrm{Ct}_k[\mathfrak D]$ (the bare vacuum graphs with all lines of index $\le k$ and at least one of index $k$, localised and coefficiented as in $\mathfrak D$), exactly
\[
\mathcal R_{k+1}[\mathfrak D\!\uparrow]=\mathcal R_k[\mathfrak D]+\mathrm{out}_k[\mathfrak D]-\mathcal V_k[\mathfrak D].
\]
\end{lemma}

\begin{proof}
Fubini on the nested integral: the $(k+1)$-step composite kernel on $\Omega$ factors into the step-$(k+1)$ kernels and the $k$-step composite (\cref{lem:a06d:B3,lem:a06d:B4,lem:a06d:B5}(a), all valid on $\Omega$), the gauge part being \cref{lem:a06d:B6}(b). So $\Xi_{k+1}[\mathfrak D\!\uparrow](0,0)=N\int d\mu_k\exp(\iota^g_kF_k[\mathfrak D]-(\iota^g_kF_k)_0+\mathcal K_k)\,\Xi_k[\mathfrak D](0,0)$ up to terms of degree $>\bar n$: only the zero-background degree-$0$ part $-\frac12\langle\phi,\Delta^{(k)}(\Omega,0)\phi\rangle$ and the term $\frac12\alpha|Q(0)\phi|^2$ combine with the step kernels into the Gaussian weight of $d\mu_k$, and every other term of $\iota^g_kF_k$, and $\mathcal K_k$, has degree $\ge1$ and is kept. (The degree-$0$ part of $F_k$ depends on the background $A$ through the scalar fluctuation determinants, Ba{\l}aban's $Z^{(j)}$ factors \cite[(2.52), (2.92), (2.105)]{Ba2}, and through $-\frac12\langle\phi,\Delta^{(k)}(A)\phi\rangle$; $\iota^g_k$ turns that dependence into terms of degree $\ge1$, which belong to $\mathrm{out}_k$, as on \cite[p.~576]{Ba2}: ``in the density $\varrho'^{(k)}$ only the factors $Z^{(j)}$ depend on this field''. $\mathcal K_k$ is the fluctuation dependence of the new scalar kernel at $\psi=0$; it gives the \cite[(1.12)]{Ba3}-type contour vertices on the new outer segment $\Gamma_{z,x_k}$, and without it (a) is false.) Take logarithms; constants pass through, and \cref{lem:a06d:B2} truncates: this is (a). (b) is \cref{lem:a06d:B6}(c), with $H$ the Neumann operator on $\Omega$ for the scalar and on the torus for the vector; expanding multilinearly in the lines identifies $\mathrm{out}_k$ with the terms having at least one index-$k$ line. (c): $\mathcal R_{k+1}-\mathcal R_k=(\mathrm{Gen}_{k+1}-\mathrm{Gen}_k)-(\mathrm{Ct}_{k+1}-\mathrm{Ct}_k)$, the coefficient $c$ being the same in $\mathrm{Gen}$ and $\mathrm{Ct}$ and held fixed between $k$ and $k+1$.
\end{proof}

On the torus datum $\Delta=0$ and (c) reads $\mathrm{out}_k=\mathcal V_k+\mathfrak B_{k+1}-\mathfrak B_k$; on a general datum (c) holds with the $\Delta$ terms, and a term of $\mathcal R_k$ is never bounded in absolute value except where a datum changes.

\begin{lemma}[comparison of data]\label{lem:a06d:V7}
Let $\mathfrak D,\mathfrak D'$ be data at the same scale. \textup{(a)} (Propagator domain.) If $\mathfrak D'$ is $\mathfrak D$ with $\Omega$ replaced by an admissible $\Omega'\subset\Omega$ and the same vertex datum, then $\mathcal R_k[\mathfrak D]-\mathcal R_k[\mathfrak D']$ is a sum of families of the two forms of \cref{lem:a06d:V5} (contour or $\Delta$), each with one designated line $\delta G_k(\Omega,\Omega',0)$ (in Ba{\l}aban's list, \cite[p.~414]{Ba3}); likewise $F_k[\mathfrak D]-F_k[\mathfrak D']$ is a sum of families with one designated $\delta G_k(\Omega,\Omega',\cdot)$ line or one designated $\delta$ in the external-field map $a_kG_k(\cdot)Q_k^*$ \cite[p.~415]{Ba3}. \textup{(b)} (Vertex datum.) If $\mathfrak D'$ is $\mathfrak D$ with $(\Omega_1,g,c)$ replaced by $(\Omega'_1,g',c')$, then $\mathcal R_k[\mathfrak D]-\mathcal R_k[\mathfrak D']$ (likewise $\mathrm{out}_k$, $F_k$) is a sum of families of the forms of \cref{lem:a06d:V5} with one designated vertex of one of three kinds: (b1) a factor $(g-g')$ on one vector leg, with the same change on the corresponding leg of the \cite[(1.7)]{Ba3} partner's coefficient graph (\cref{conv:a06d:LOC}(c)); (b2) a \cite[(1.6)]{Ba3} or \cite[(1.7)]{Ba3} vertex in $\Omega_1\triangle\Omega'_1$; (b3) a \cite[(1.7)]{Ba3} vertex carrying a change of the $g$-independent finite part of $c$ ($\dm_{\mathrm{fin}}$, $\dm_{K,\cdot}$, or a coefficient graph with an index-$k$ line, all regular; $|\text{change}|\le C_Tt_k(1+\ln t_k^{-1})^5$ per site by \cref{thm:a06e:T0}, and $\dm_{K,k}(L^k\eps)^2=O(L^k\eps)$ by \cref{thm:a06e:DM}, as asserted on \cite[p.~418]{Ba3}). The designated factor is supported in $Z(\mathfrak D,\mathfrak D')\defeq(\Omega_1\triangle\Omega'_1)\cup\supp(g-g')\cup\supp(c-c')$ up to exponentially decaying tails of the regular coefficient graphs. \textup{(c)} (Step covariances.) For two choices $C,C'$ of the step-$k$ covariances, $\mathrm{out}_k[\mathfrak D;C]-\mathrm{out}_k[\mathfrak D;C']$ is a sum of step-$k$ cumulant families with one designated contraction by $\delta C\defeq C-C'$, lifted by $\beta_kG_kR_k^*$ (\cref{lem:a06d:B6}(b)).
\end{lemma}

\begin{proof}
Telescope multilinearly: in the lines for (a), in the vertex multipliers for (b), in the contraction kernels for (c). In (b) the telescoping is symmetric, $g'=g+(g'-g)$ and $\mathbf 1_{\Omega'_1}=\mathbf 1_{\Omega_1}+(\mathbf 1_{\Omega'_1}-\mathbf 1_{\Omega_1})$ on all legs and sites at once, applied identically to $\mathrm{Gen}$, $\mathrm{Ct}$ and the coefficient graphs; every term then carries at least one factor $(g'-g)$ or a site difference, the members of a renormalised pair carry the same set of changed legs and sites (\cref{conv:a06d:LOC}(c)), and the number of terms is $\le2^{\text{legs}+\text{sites}}\le C(\bar n)$. (A lone \cite[(1.7)]{Ba3} vertex is divergent and the \cite[(1.6)]{Ba3} tadpole needs its partner, so the kinds are not bounded vertex by vertex.) By \cref{conv:a06d:LOC} $\mathrm{Gen}$ and $\mathrm{Ct}$ carry the same vertices, so each difference term again has the form $\mathfrak B+\Delta$ of \cref{lem:a06d:V5}: no bare vacuum graph with an uncancelled overall divergence survives; by (c) of the convention the change on a self-energy subgraph and on its partner are made on corresponding legs, so renormalised families stay renormalised. A coefficient change splits into the $g$-localised part of \cite[(1.29)]{Ba3}, $\dm_k(\cdot g)$, which changes with $g$ and is paired with the generated subgraphs (kind (b1)), and the $g$-independent finite parts, which change only with the scale index ($\dm_{K,k}\to\dm_{K,k+1}$), the moved piece being an index-$k$ slice, a regular graph (kind (b3)). Here $(g'-g)$ is itself a multiplier of the class ($C^2$-norm $\le2C_g$) supported in $Z$.
\end{proof}

% a06_BALtorus_d2b: continuation of a06_BALtorus_d2 (input from sec/a06_BALtorus_c.tex): grouped families and the geometry of the regions.

\begin{lemma}[grouped families]\label{lem:a06d:K2}
Let $X$ be any of $\Delta_k[\mathfrak D]$, $\mathcal R_k[\mathfrak D]-\mathcal R_k[\mathfrak D']$, $\mathrm{out}_k[\mathfrak D;C]-\mathrm{out}_k[\mathfrak D;C']$, written with the line telescoping of \cref{lem:a06d:V5} (order induced on $G/S$ from $G$), the symmetric expansion of \cref{lem:a06d:V7}(b) for vertex multipliers and sites, and the multilinear expansion in contraction kernels. Then $X$ is a finite sum, of at most $C(\bar n)$ terms, of renormalised families in the sense of \cite[p.~420]{Ba3} (each divergent two-scalar-leg subgraph $S$ accompanied by $G/S$ with the \cite[(1.7)]{Ba3} vertex and the same changed lines, legs and sites outside $S$), each with exactly the designated object of \cref{lem:a06d:V9}(b)--(d).
\end{lemma}

\begin{proof}
A term with $\delta_k$ on a line of $S$ has $S$ convergent ($\delta_k$ is regular on the diagonal); a term with $\delta_k$ off $S$ pairs with the $G/S$ term with $\delta_k$ on the corresponding line, and $S-c_S=[S-S(T)]+[S(T)-c_S]$, both convergent ($c_S$ uses torus lines, \cref{conv:a06d:LOC}(b), \cite[(1.31)]{Ba3}). The vertex expansion is symmetric, so the subsets of changed legs and sites of $S$ and of its partner coincide. Contraction kernels are lifted index-$k$ kernels (regular) or their deep differences. Polarisation subgraphs keep bubble and seagull with the same multipliers on corresponding legs, and the local cancellation survives $C^2$ multipliers, the subtraction $\sum_{x'}\Pi(x,x')[m(x')-m(x)-\nabla m(x)\cdot(x'-x)]$ being UV-finite in $d=3$. The grouping is exactly \cite[p.~414]{Ba3}'s line list $\{G(\Omega),G(\Omega_2),\delta G\}$.
\end{proof}

\begin{lemma}[geometry]\label{lem:a06d:V8}
\textup{(a)} $\dist(\Lambda_{i+1},\Lambda_i^c)>r$ for every $i$ (the definition \cite[(2.8)]{Ba2}, with $r=r(t_k)$ at scale $k$). \textup{(b)} At scale $k$, with $c_d$ a dimensional constant: $\dist(B^k(\Lambda^{(k)}_7),B^k(\Lambda^{(k)}_5)^c)>2r(t_k)-c_d$, $\dist(B^k(\Lambda^{(k)}_7),B^k(\Lambda^{(k)}_2)^c)>5r(t_k)-c_d$, and $\dist(B^k(\Lambda^{(k)}_7),B^k(\Lambda^{(k-1)'}_7)^c)>7r(t_k)-c_d$, because $\Lambda^{(k)}_7\subset\Lambda^{(k)}_0\subset\Lambda^{(k-1)'}_7$ \cite[p.~570]{Ba2}. \textup{(c)} $\rho(\mathfrak D_k)\ge(5/L)r(t_{k-1})-c_d\ge r(t_k)-c_d$ at $L=5$ (the primed sets are the scale-$(k-1)$ sets read on $T^{(k)}_1$, distances divided by $L$, inclusions kept). \textup{(d)} If $7r(t_k)-c_d\ge M$ then $g_k\equiv1$ on $B^k(\Lambda^{(k)}_7)$, and the coefficient \cite[(1.29)]{Ba3} there differs from the full $\dm$ only by tails $e^{-\delta(7r-M)}$. \textup{(e)} Put $\mathrm{Col}_k\defeq\Lambda^{(k-1)'}_7\cap\Lambda^{(k)c}_7$ and let $D<7r(t_k)-c_d$; every scale-$k$ block $b\subset\Lambda^{(k)}_7$ with $\dist(b,\Lambda^{(k)c}_7)\le D$ lies within $D+1$ of a block of $\mathrm{Col}_k$, hence for $\delta>0$, $\sum_{b\subset\Lambda^{(k)}_7}e^{-\delta\dist(b,\Lambda^{(k)c}_7)}\le C(\delta)|\mathrm{Col}_k|_k+e^{-\delta(7r(t_k)-c_d)}N_k$ with $C(\delta)\defeq\sum_{D\ge0}(2D+3)^3e^{-\delta D}$ and $N_k\defeq|\Lambda^{(k)}_7|_k$. \textup{(f)} Any zone $Z$ of blocks within distance $D_Z<7r(t_k)-c_d$ of $\partial B^k(\Lambda^{(k)}_7)$, on either side, has $|Z|_k\le(1+(2D_Z+3)^3)|\mathrm{Col}_k|_k$.
\end{lemma}

\begin{proof}
(a): by \cite[(2.8)]{Ba2}, $\Lambda^c_{i+1}$ is the union of the large blocks at distance $\le r$ from $\Lambda^c_i$. (b): iterate (a) along a shortest path; each crossing of a shell costs at least $r$ up to the block size. (c), (d): from (b) and the definition of $g_k$. (e): the nearest point of $\Lambda^{(k)c}_7$ to $b$ is within $D$, and since $(\Lambda^{(k-1)'}_7)^c$ is farther than $7r-c_d>D$ that point lies in $\mathrm{Col}_k$; the map $b\mapsto$ (a collar block within $D+1$) has fibres of size $\le(2D+3)^3$; sum over $D$. (f): the blocks of $Z$ outside $\Lambda^{(k)}_7$ lie in $\mathrm{Col}_k$ (at most $|\mathrm{Col}_k|_k$ of them), those inside are covered by (e).
\end{proof}

% a06_BALtorus_d3: part of Step 3 of the proof of thm:a06:T1 (input from sec/a06_BALtorus_c.tex).
% Group (D), second half: the bounds (VAC-flat)_k and (DOM)_k, the collar clause of Ba2 (2.57), and the bad localisations of Ba3 p. 422.

\subsubsection*{Step 3(D$_2$): the counterterm moves of \cite[pp.~411--422]{Ba3}, bounds}

\paragraph{Inputs.} (I1) Proposition~1 in the scope Ba{\l}aban states for it: (I1a) the expressions of the expansion \cite[pp.~414--415]{Ba3} (vertices \cite[(1.6)--(1.15)]{Ba3} and products \cite[(1.18)]{Ba3}; propagators $G_k(\Omega,\tilde B)$, $G_k(\Omega_2,\tilde B)$ and the vector $G_k$ on any admissible $\Omega\supset\Omega_2$, the torus included; vacuum expressions included, since vacuum contour graphs are among ``an arbitrary expression in the expansion'', \cite[p.~414]{Ba3}, with $D(G)\ge2$); (I1b) $\delta G_k(\Omega,\Omega_2,\tilde B)$ and \cite[(1.16)]{Ba3} treated as external fields \cite[pp.~414, 419]{Ba3}; (I1c) external fields that are functions of many variables, in particular ``expressions obtained by an integration and a cumulant expansion in $k$th step, where some external legs are contracted by propagators $C^{(k)}_{\Lambda^{(k)}_5}$'' \cite[p.~419]{Ba3}; (I1d) the localisation of the vector legs by smooth partitions of unity \cite[p.~420]{Ba3}, i.e.\ smooth multipliers on single legs from a fixed $C^2$-bounded class, with the $O(1)$ uniform in them, in the domains and in $\tilde B$ \cite[p.~421]{Ba3}. (I2) Norms of the external fields: (I2a) $\delta G$ at depth, from \cite[Thm., (1.10)--(1.12)]{Ba4} in its restricted form made two-sided (\cref{lem:a06A:twosided}): for cubes $\square,\square'$ at depths $\rho_\square,\rho_{\square'}\ge R_0$ from $\Omega'^c$, $\norm{h_\square\delta G_k(\Omega,\Omega')h_{\square'}}_{1,\alpha_0}\le Ce^{-(\delta_0/4)(\rho_\square+\rho_{\square'}+\dist(\square,\square'))}$; (I2b) block-level covariance differences ($C_\Lambda$ against $C^{(k)}$, \cref{lem:a06A:CLdiff}; domain differences of $\Delta^{(k)}$ and $C_\Lambda$, \cref{lem:a06A:domdiff}), lifted by $a_kG_kQ_k^*$ whose $C^{1,\alpha_0}$ regularity is \cite[Thm., (1.9)--(1.10)]{Ba4} and \cite[Prop.~2.2, (2.58)]{Ba2}, vector versions at $B=0$ componentwise; (I2c) the lifted index-$k$ kernels $\beta_kG_kR_k^*C^{(k)}R_kG_k\beta_k$, with the same inputs and the decay of $C^{(k)}$ (\cref{lem:a06A:blockinv}(c)). (I3) The size of the finite mass vertices: in scale-$k$ units the non-renormalising coefficient $c_k=t_k^2(\dm_{\mathrm{fin}}+\dm_{K,k})$ of the \cite[(1.7)]{Ba3} vertex obeys $|c_k|\le C_Tt_k(1+\ln t_k^{-1})^5$ with $C_T$ free of $\eps_0$ (\cref{thm:a06e:T0}; asserted on \cite[p.~418]{Ba3} as $\dm_{K,k}t_k^2=O(t_k)$, $\dm_{\mathrm{fin}}=O(1)$, the first of which is \cref{thm:a06e:DM}), entering only through Proposition~1's order count (\cite[p.~421]{Ba3}: ``the vertices (1.7) give also a positive power of $L^k\eps$'') and in kind (b3) of \cref{lem:a06d:V7}; every bound of Steps~3(D$_1$)--3(D$_2$) stated with a factor $t_k$ from such a vertex holds with $t_k(1+\ln t_k^{-1})^5$ in its place, which its consumers absorb (\cref{lem:a06e:T1}). Geometry: \cite[(2.7)--(2.8)]{Ba2} and the bound $e^{-\delta c_*R(1+\log t^{-1})^\varrho}\le C_\kappa t^\kappa$ for every $\kappa$, uniformly in $t\in(0,1]$ ($\varrho>1$; \cref{lem:a06A:polylog}).

\begin{lemma}[gluing two families along contractions]\label{lem:a06d:V10}
Let $F_1,F_2$ be families obeying \cite[(1.33)]{Ba3} with multi-variable external fields on their legs, $|E(F_i,\Phi)|\le P_ie^{-\delta_0d_i}\norm{h\Phi}_{1,\alpha_0}$, and let $K_1,\dots,K_r$ be two-variable kernels joining a leg of $F_1$ (cube $\square_j$) to a leg of $F_2$ (cube $\square'_j$), with $\norm{h_{\square_j}K_jh_{\square'_j}}_{(1,\alpha_0)\otimes(1,\alpha_0)}\le\kappa_je^{-\delta'\dist(\square_j,\square'_j)}$ (the $C^{1,\alpha_0}$-norm in each variable uniformly in the other, with the mixed derivative and H\"older terms). Then $|E(F_1\circ_KF_2,\Phi_{\mathrm{rest}})|\le P_1P_2\prod_j\kappa_j\,e^{-\delta_0(d_1+d_2)}e^{-\delta'\sum_j\dist(\square_j,\square'_j)}\norm{h\Phi_{\mathrm{rest}}}$.
\end{lemma}

\begin{proof}
$E(F_1\circ_KF_2)=E(F_1,\Psi\otimes\Phi_1)$ with $\Psi(x_1,\dots,x_r)\defeq E(F_2,K_1(x_1,\cdot)\otimes\dots\otimes K_r(x_r,\cdot)\otimes\Phi_2)$. $\Psi$ is linear in each $K_j(x_j,\cdot)$; its sup, derivative and H\"older quotient in $x_j$ are $E(F_2,\cdot)$ applied to $K_j(x_j,\cdot)$, $\partial_{x_j}K_j(x_j,\cdot)$ and the H\"older quotients of $K_j$ in $x_j$, each bounded by $P_2e^{-\delta_0d_2}$ times the mixed norm. So $\norm{h\Psi}_{1,\alpha_0}\le P_2e^{-\delta_0d_2}\prod_j\kappa_je^{-\delta'\dist}$; apply \cite[(1.33)]{Ba3} to $F_1$. For $m\ge3$ families (a cumulant $\kappa_m$) induct, joining $F_m$ to the group already joined.
\end{proof}

Proposition~1's many-variable norm (\cite[p.~420]{Ba3}: \cite[(1.32)]{Ba3} ``extends in a natural way'') is $C^{1,\alpha_0}$ per variable; mixed derivatives enter only through the joins of \cref{lem:a06d:V10}, where the lifted kernels and their deep differences factorise with one derivative on each lifting factor; $\delta G$ lines never serve as joining kernels, and self-contractions inside one family are many-variable external fields of that family (I1c).

\begin{lemma}[localisation sums]\label{lem:a06d:V11}
Let $N\le\bar N$ be the number of indices and $\delta>0$. Then $\sum_{\Lambda:\square_{\iota_0}=\square_0}e^{-\delta d(\Lambda)}\le S_N(\delta)\defeq N^{N-2}(e^\delta\sigma(\delta))^{N-1}$, uniformly in the torus, with $\sigma(\delta)\defeq\sup_z\sum_\square e^{-\delta\dist(\square,z)}$. Consequently, with $S\defeq\max_{N\le\bar N}S_N(\delta)$: for a set $Z$ of blocks, $\sum_{\Lambda:\text{designated cube in }Z}Ct_ke^{-\delta d(\Lambda)}\le C\bar NSt_k|Z|_k$, and $\sum_{\Lambda:\text{anchor in the vertex domain}}Ct_ke^{-\delta\rho}e^{-\delta d(\Lambda)}\le C\bar NSt_ke^{-\delta\rho}N_k$.
\end{lemma}

\begin{proof}
$e^{-\delta d}\le\sum_{\text{trees}}e^{-\delta L_T}$; there are $N^{N-2}$ trees (Cayley); sum the cubes from the leaves inward, each edge contributing $\le e^\delta\sigma(\delta)$.
\end{proof}

\begin{lemma}[bounds per family and per localisation, from Proposition~1]\label{lem:a06d:V9}
There are $C_P<\infty$ and $\delta>0$, depending only on $(d,N,\norm q,a,L,M,\bar n,\alpha_0,R,C_g)$ and on the constants of (I1)--(I3), not on $\eps$, $k$, the torus or the data, such that for every localisation $\Lambda$: \textup{(a)} every grouped contour family ($\mathfrak B$-type, \cref{lem:a06d:K2}) on any datum obeys $|\cdot(\Lambda)|\le C_Pt_ke^{-\delta d(\Lambda)}$; \textup{(b)} a grouped $\Delta$-type family, or any grouped family with a designated $\delta G$ line or a designated $\delta$ in the external-field map, whose designated object sits at or joins cubes at depth $\ge\rho\ge R_0+2$ from the boundary where the two propagators differ, obeys $|\cdot(\Lambda)|\le C_Pt_ke^{-\delta\rho}e^{-\delta d(\Lambda)}$; \textup{(c)} every grouped step-$k$ cumulant family (a term of $\mathrm{out}_k$ on any datum) obeys $|\cdot(\Lambda)|\le C_Pt_ke^{-\delta d(\Lambda)}$, and with a designated $\delta C$ contraction at depth $\ge\rho$ it obeys $|\cdot(\Lambda)|\le C_Pt_ke^{-\delta\rho}e^{-\delta d(\Lambda)}$; \textup{(d)} the families of (a)--(c) with one designated vertex of kind (b1), (b2) or (b3) of \cref{lem:a06d:V7} obey the same bounds, and their $\Lambda$ contains a cube that meets the zone $Z$.
\end{lemma}

\begin{proof}
\emph{Scope.} Each family lies in the scope of (I1): (a) is (I1a) without external field; (b) is (I1b) with the designated $\delta G$ as a two-variable external field; (c) is (I1c), with \cref{lem:a06d:V10} when a cumulant joins several families; (d) is (I1d): for (b1) the designated leg carries the multiplier $(g'-g)$, in the class and supported in $Z$; for (b2) the scalar vertices are localised by unit cubes anyway \cite[p.~420]{Ba3}; for (b3) the vertex is a finite mass vertex with coefficient $\le Ct_k$ (I3), in Proposition~1's class inside the families of \cref{lem:a06d:V7} (its loop line is $\delta_k$, a lifted index-$k$ line, or not a pure full-line loop), never alone. (A lone $\mathrm{out}$-term is not bounded either: a figure-eight with one full and one index-$k$ line is $\sim\lambda t\eta^{-1}$; the pairs in $\mathrm{out}_k$ are inherited from the renormalised vacuum part of $F_k$.) \emph{Order.} Every vacuum family has order $o\defeq d_v/2+d_s\ge1$: all $A'$-legs are contracted, so $d_v$ is even, and if $d_v=0$ the family contains a \cite[(1.6)]{Ba3} vertex or a \cite[(1.7)]{Ba3} vertex of positive order; so $e(t_k)^{d_v}\lambda(t_k)^{d_s}=e^{d_v}\lambda^{d_s}t_k^o\le Ct_k$ for $t_k\le1$; a step-$k$ cumulant family is a join of families of order $\ge1$. \emph{External fields.} In (b) and (c), (I2a)--(I2c) bound the designated objects by $Ce^{-(\delta'/4)(\rho_\square+\rho_{\square'}+\dist(\square,\square'))}$ with $\rho_\square,\rho_{\square'}\ge\rho$, so $e^{-\delta'(\rho_\square+\rho_{\square'})/4}\le e^{-\delta'\rho/2}$ and $e^{-\delta'\dist(\square,\square')/4}$ supplies the tree link between the two cubes; the other contractions in (c) are lifted index-$k$ kernels with $\norm{h\cdot h'}\le Ce^{-\delta'\dist}$ (I2c). \emph{Assembly.} \cite[(1.33)]{Ba3} gives $O(1)Ct_ke^{-\delta_0d}\prod(\text{norms})$; the tree length over all cubes is at most the sum of the constituent tree lengths and the link lengths. Take $\delta\defeq\frac12\min(\delta_0,\delta'/4)$ and $C_P\defeq(\text{number of families of order}\le\bar n)\cdot O(1)\cdot C$.
\end{proof}

\begin{theorem}[$(\mathrm{VAC}^\flat)_k$, $0\le k\le K$]\label{thm:a06d:V12}
Let $\mathfrak D$ be a datum at scale $k$ with $\rho(\mathfrak D)\ge R_0+3$. Then for every localisation $\Lambda$, $|\mathcal R_k[\mathfrak D](\Lambda)|\le C_Vt_ke^{-\delta d(\Lambda)}$ with $C_V\defeq2N_{\mathrm{fam}}C_P$, $N_{\mathrm{fam}}$ the number of families of order $\le\bar n$; consequently $|\mathcal R_k[\mathfrak D]|\le C_V\bar NSt_k|\Omega_1|_k$. At $\mathfrak D=\mathfrak D_K$ this gives $|\mathcal R_K|\le C_V\bar NSt_K|T|_K=C_V\bar NSt_K^{-2}|T|\le C_V\bar NSL^2\eps_0^{-2}|T|$, since $t_K\in(\eps_0/L,\eps_0]$; at $\mathfrak D=\mathfrak D_k$ the hypothesis holds, $\rho(\mathfrak D_k)\ge r(t_k)-c_d\ge R-c_d\ge R_0+3$ (\cref{lem:a06d:V8}(c) at $L=5$), given $R\ge R_0+3+c_d$ (admissible: $R$ is chosen after $M$, \cite[(2.7)]{Ba2}).
\end{theorem}

\begin{proof}
$\mathcal R_k=\mathfrak B_k+\Delta_k$ (\cref{lem:a06d:V5}). $\mathfrak B_k$ is bounded by \cref{lem:a06d:V9}(a). $\Delta_k$ is bounded by \cref{lem:a06d:V9}(b) through its grouped families (\cref{lem:a06d:K2}): the designated $\delta_k=G_k(\Omega,0)-G_k(T,0)$ sits between vertices of $\Omega_1$ plus one block, at depth $\ge\rho(\mathfrak D)-1\ge R_0+2$, as (b) requires. Sum over localisations with \cref{lem:a06d:V11}.
\end{proof}

\begin{corollary}[{the vacuum clause of \cite[Prop.~2.1, (2.57)]{Ba2} at zero background}]\label{cor:a06d:V13}
Let $\mathfrak D'\defeq(\Omega^{(k)},B^k(\Lambda^{(k)}_7),g',c')$ with $g'$ admissible for $B^k(\Lambda^{(k)}_7)$ and $c'$ its \cite[(1.29)]{Ba3}. Then $\mathcal R_k[\mathfrak D_k]-\mathcal R_k[\mathfrak D']$ is a sum of families with a designated vertex in $Z_{\mathrm{col}}\defeq B^k(\mathrm{Col}_k)\cup\{g'<1\}$, and $|\mathcal R_k[\mathfrak D_k]-\mathcal R_k[\mathfrak D']|\le Ct_k|\mathrm{Col}_k|_k$ with $C\defeq C_P\bar NSC(\delta)(1+(2M+5)^3)$.
\end{corollary}

At $A^{(k)}=0$, $\phi=0$ the quartic term of \cite[(2.57)]{Ba2} vanishes and $\Efin$, $E_{K,k}$ are unchanged, so this is the $c$-number clause of \cite[(2.57)]{Ba2}, with $O(t_k)\le O(t_k^{\kappa_0})$ per collar block, under \cref{conv:a06d:LOC} (the mass counterterm enters as a vertex coefficient, and the replacement of propagators acts on lines only).

\begin{proof}
Apply \cref{lem:a06d:V7}(b): the vertex domains differ by $\Omega_1\triangle\Omega'_1=B^k(\mathrm{Col}_k)$; on $B^k(\Lambda^{(k)}_7)$, $g_k=1$ (\cref{lem:a06d:V8}(d)), so $\supp(g_k-g')\subset B^k(\mathrm{Col}_k)\cup\{g'<1\}$ and $c_k-c'$ is supported there up to tails $e^{-\delta(7r-M)}$. Apply \cref{lem:a06d:V9}(d), then \cref{lem:a06d:V8}(f) with $D_Z=M+1$, then \cref{lem:a06d:V11}; the coefficient tails are summed by $\sum_xe^{-\delta\dist(x,Z)}\le C(\delta)|Z|$.
\end{proof}

% a06_BALtorus_d3b: continuation of a06_BALtorus_d3 (input from sec/a06_BALtorus_c.tex): the step k+1 as Ba2 treats it, (DOM)_k, and the bad localisations of Ba3 p. 422.

\begin{lemma}[{the new fluctuation at step $k+1$, as \cite[pp.~575--581]{Ba2} treats it}]\label{lem:a06d:K3}
At step $k+1$, \cite{Ba2} writes $\theta_{k+1}A^{(k)}=\theta_{k+1}A'^{(k)}+\tilde B^{(k)}+\theta_{k+1}B^{(k+1)}$ \cite[(2.88)--(2.89)]{Ba2}, $\theta_{k+1}$ a smooth cutoff of $B^k(\Lambda^{(k)}_2)$ changing in a slice of thickness $M$, keeps $(1-\theta_{k+1})A^{(k)}$ in the background $\tilde B$ \cite[p.~579]{Ba2}; expands the $Z^{(j)}$ factors \cite[(2.92)]{Ba2}, $j<k$, on their domains $B^j(\Lambda^{(j)}_2)$ with restrictions $\Lambda^{(j)}_5$, moves them to $B^k(\Lambda^{(k)}_0)$, composes them by \cite[(2.105)]{Ba2} into the polynomial $Q^{(k)}$ \cite[(2.106)]{Ba2}; drops the terms with vertices localised in $B^k(\Lambda^{(k)c}_7)$ and relocalises with a smooth $g^\sharp$ ($\supp g^\sharp\subset\Lambda^{(k)}_7$, $g^\sharp=1$ on $\Lambda^{(k)}_8$) \cite[p.~579]{Ba2}; replaces the lines by $G_k(B^k(\Lambda^{(k)}_2),B^{(k+1)})$ \cite[Prop.~2.6, (2.107)]{Ba2}; and integrates $A',\phi'$ conditionally on $\Lambda^{(k)c}_5$ with $C^{(k)}_{\Lambda^{(k)}_5}$ \cite[(2.115)]{Ba2}. At zero $(k+1)$-fields (so $\tilde B^{(k)}=0=\theta_{k+1}B^{(k+1)}$) write $\mathrm{out}^\theta_k$ for the vacuum generated before the drop, and $\mathrm{out}^\theta_k=\mathrm{out}^\sharp_k+D_k$, where $D_k$ collects the terms with at least one new leg carrying $(1-g^\sharp)\theta_{k+1}$ and $\mathrm{out}^\sharp_k$ is kept in the action; the legs of the new-kernel term $\mathcal K_k$ of \cref{lem:a06d:V6} count as new legs. Then, with $\mathrm{out}_k$ of \cref{lem:a06d:V6} on $\mathfrak D'=(\Omega^{(k)},W,g',c')$, $W\defeq B^k(\Lambda^{(k)}_7)$,
\[
\mathrm{out}^\sharp_k=\mathrm{out}_k[\mathfrak D';C_{\mathrm{act}}]+X_0+X_3,
\]
with $|X_0|\le Ct_ke^{-\delta(5r(t_k)-M-c_d)}N_k$ (the domain moves and the conditioned background $\tilde B=\sum_{l\le k}(1-\theta_{l+1})\theta_lA^{(l)}$ at zero new fields, every slice of which lies in $B^{l-1}(\Lambda^{(l-1)}_2)\setminus B^l(\Lambda^{(l)}_2)\subset B^k(\Lambda^{(k)}_2)^c$, at distance $\ge5r(t_k)-M-c_d$ from $W$), $|X_3|\le Ct_k(2D+5)^3|\mathrm{Col}_k|_k$ with $D\defeq\max(M,D^\sharp)$, $D^\sharp$ the transition width of $g^\sharp$ (families with one new leg carrying $(g^\sharp-g')$, supported within $D$ of $\partial W$), and $|D_k|\le Ct_k(2+(2D+3)^3)|\mathrm{Col}_k|_k$ (the support of $(1-g^\sharp)\theta_{k+1}$ lies in $B^k(\Lambda^{(k)}_2\setminus\Lambda^{(k)}_7)\cup(\theta\text{-slice})\cup\{g^\sharp<1\}$, the first two in $B^k(\mathrm{Col}_k)$ since $\Lambda^{(k)}_2\subset\Lambda^{(k)}_0\subset\Lambda^{(k-1)'}_7$ and $\dist(\Lambda^{(k)}_2,\Lambda^{(k-1)'c}_7)>2r(t_k)>M$). In all three, $t_k(2D+5)^3\le C't_k^{\kappa_0}$, $\kappa_0=1/5$, even for $D\sim r(t_k)$.
\end{lemma}

\begin{proof}
All three are differences of grouped families (\cref{lem:a06d:K2}) with one designated object: a deep domain or background change ($X_0$: the Combes--Thomas bound for the background resolvent difference, (I2a) for domains, (I2b) for the $\Lambda^{(j)}_5$ restrictions; $W$ is at depth $\ge5r(t_k)-c_d$ in $B^k(\Lambda^{(k)}_2)$ and nested inside every $B^j(\Lambda^{(j)}_2)$, $\Lambda^{(j)}_5$), or a smooth multiplier $(g^\sharp-g')$ or $(1-g^\sharp)\theta_{k+1}$ on one new leg ($C^2$-bounded, (I1d)). Bound each by \cref{lem:a06d:V9} and sum with \cref{lem:a06d:V11,lem:a06d:V8}(f); $\theta_k=1=\theta_{k+1}$ on $\supp g^\sharp$, so $g^\sharp\theta_k\theta_{k+1}=g^\sharp$. The last claim is $t(1+\log t^{-1})^{3\varrho}\le Ct^{1/5}$.
\end{proof}

\begin{theorem}[$(\mathrm{DOM})_k$, $0\le k<K$]\label{thm:a06d:V14}
Ba{\l}aban's step $k+1$ (\cite[p.~412]{Ba3}, \cite[pp.~569--582]{Ba2}) consists of: (O1) the collar removal \cite[(2.57)]{Ba2}, down to a vertex datum $(B^k(\Lambda^{(k)}_7),g',c')$; (O2) the translation and the expansion in the new fluctuation; (O3) the integration of the step-$k$ fluctuations with the covariances $C^{(k)}_{\Lambda^{(k)}_5}(B^k(\Lambda^{(k)}_2),0)$ (scalar) and $C^{(k)}_{\Lambda^{(k)}_5}$ (vector) of \cite[p.~419]{Ba3}, and the cumulant expansion; (O4) the counterterm move $k\to k+1$ of \cite[p.~422]{Ba3}, $E_{K,k}=E_{K,k+1}+\mathcal V^{\mathrm{glob}}_k$ and $\dm_{K,k}\to\dm_{K,k+1}$ with the index-$k$ slice moved into $\dm_{k+1}$, followed by localisation; (O5) the rescaling to scale $k+1$ and the new datum $\mathfrak D_{k+1}$. Let $\mathrm{Vac}^{\mathrm{act}}_{k+1}$ be the field-independent part of degree $1..\bar n$ of the resulting polynomial, in the logarithmic convention with the front factor $e^{-\Efin-E_{K,k+1}}$, the removed collar constant $\mathcal R_k[\mathfrak D_k]-\mathcal R_k[\mathfrak D']$ (\cref{cor:a06d:V13}) and the term $D_k$ of \cref{lem:a06d:K3} being dropped and booked by their absolute values. Then
\[
\mathrm{Vac}^{\mathrm{act}}_{k+1}=\mathcal R_{k+1}[\mathfrak D_{k+1}]+\mathcal E'_k+\mathcal O_k,
\]
where $\mathcal E'_k$, $\mathcal O_k$ are finite sums of localised vacuum families with $|\mathcal E'_k|\le Ct_k(2D+5)^3|\mathrm{Col}_k|_k+Ct_ke^{-\delta c_*r(t_k)}N_k\le C't_k^{\kappa_0}|\mathrm{Col}_k|_k+C't_k^{16/5}N_k$ ($D$ as in \cref{lem:a06d:K3}, $\kappa_0=1/5$), $|\mathcal O_k|\le Ct_k|\Lambda^{(k)c}_7|_k$, every family of $\mathcal O_k$ having a designated vertex in $\Lambda^{(k)c}_7$; the constants do not depend on $\eps$, $k$, the torus or the regions. Equivalently, the field-independent output of step $k$, $\mathrm{out}^{\mathrm{act}}_k-\mathcal V^{\mathrm{glob}}_k$, equals $\mathcal R_{k+1}[\mathfrak D_{k+1}]-\mathcal R_k[\mathfrak D']+\mathcal E'_k+\mathcal O_k$.
\end{theorem}

\begin{proof}
Put $W\defeq B^k(\Lambda^{(k)}_7)$, $\Omega^\natural\defeq B^k(\Lambda^{(k)}_2)$, $\mathfrak D'\defeq(\Omega^{(k)},W,g',c')$, $\mathfrak D^\natural\defeq(\Omega^\natural,W,g',c')$; $\Omega^\natural\subset\Omega^{(k)}$ by the nesting $\Lambda^{(k)}_2\subset\Lambda^{(k)}_0\subset\Lambda^{(k-1)'}_7\subset\Lambda^{(k-1)'}_2$ \cite[p.~570]{Ba2}.

1. After (O1) the vacuum part is $\mathcal R_k[\mathfrak D']$, the front factor unchanged; the dropped constant is bounded by \cref{cor:a06d:V13}.

2. (Input domain.) By \cref{lem:a06d:V7}(a), $\mathcal R_k[\mathfrak D']=\mathcal R_k[\mathfrak D^\natural]+X_1$ and $F_k[\mathfrak D']=F_k[\mathfrak D^\natural]+(\delta\text{-families})$, the designated object $\delta G_k(\Omega^{(k)},\Omega^\natural)$ at depth $\ge5r(t_k)-c_d$ (\cref{lem:a06d:V8}(b)).

3. (Step covariances.) $\mathrm{out}^{\mathrm{act}}_k\defeq\mathrm{out}^\sharp_k=\mathrm{out}_k[\mathfrak D';C_{\mathrm{act}}]+X_0+X_3$ (\cref{lem:a06d:K3}), and $\mathrm{out}_k[\mathfrak D';C_{\mathrm{act}}]=\mathrm{out}_k[\mathfrak D^\natural;C_{\mathrm{comp}}]+X_2$ by \cref{lem:a06d:V7}(a) for the input $F_k$ (step~2, then integrated) and \cref{lem:a06d:V7}(c) for $C_{\mathrm{act}}-C_{\mathrm{comp}}$: scalar $C^{(k)}_{\Lambda^{(k)}_5}(\Omega^\natural,0)-C^{(k)}(\Omega^\natural,0)$, vector $C^{(k)}_{\Lambda^{(k)}_5}-C^{(k)}$ (I2b), every designated object at depth $\ge2r(t_k)-c_d$ from $\Lambda^{(k)c}_5$. By \cref{lem:a06d:V9}(b), (c) and \cref{lem:a06d:V11}, $|X_1|+|X_2|\le Ct_ke^{-\delta(2r(t_k)-c_d)}N_k$.

4. (Exact staging.) \cref{lem:a06d:V6}(c) on $\mathfrak D^\natural$: $\mathcal R_k[\mathfrak D^\natural]+\mathrm{out}_k[\mathfrak D^\natural;C_{\mathrm{comp}}]-\mathcal V_k[\mathfrak D^\natural]=\mathcal R_{k+1}[\mathfrak D^\natural\!\uparrow]$.

5. (The counterterm move.) The front factor splits, $e^{-E_{K,k}}=e^{-E_{K,k+1}}e^{-\mathcal V^{\mathrm{glob}}_k}$ with $\mathcal V^{\mathrm{glob}}_k=\mathcal V_k[\mathfrak D^T]$ (largest line index exactly $k$, full coefficient $\dm$). $\mathrm{Ct}$ does not depend on the propagator domain, so $\mathcal V_k[\mathfrak D^\natural]=\mathcal V_k[\mathfrak D^T_{W,g'}]$. \cref{lem:a06d:V6}(c) on the two torus data (where $\Delta=0$) gives $\mathcal V_k[\mathfrak D^T]=\mathrm{out}_k[\mathfrak D^T]-\mathfrak B_{k+1}[\mathfrak D^T\!\uparrow]+\mathfrak B_k[\mathfrak D^T]$ and the same for $\mathfrak D^T_{W,g'}$, hence
\[
Y_k\defeq\mathcal V^{\mathrm{glob}}_k-\mathcal V_k[\mathfrak D^\natural]=\big[\mathrm{out}_k[\mathfrak D^T]-\mathrm{out}_k[\mathfrak D^T_{W,g'}]\big]-\big[\mathfrak B_{k+1}[\mathfrak D^T\!\uparrow]-\mathfrak B_{k+1}[\mathfrak D^T_{W,g'}\!\uparrow]\big]+\big[\mathfrak B_k[\mathfrak D^T]-\mathfrak B_k[\mathfrak D^T_{W,g'}]\big].
\]
By \cref{lem:a06d:V7}(b) each bracket is a sum of families with a designated vertex in $Z^T\defeq W^c\cup\{g'<1\}$ or in the zone where the coefficient changes; so $Y_k$ needs no bound on index-$k$ slices alone, only Proposition~1 for out-type and contour families. Sort by the cube of the designated vertex: families with that cube in $\Lambda^{(k)c}_7$ form $-\mathcal O_k$, families with it in $W$ (then within $M+1$ of $\partial W$) go into $\mathcal E_k$. By \cref{lem:a06d:V9}(a), (c), (d), \cref{lem:a06d:V11} and \cref{lem:a06d:V8}(f): $|\mathcal O_k|\le Ct_k|\Lambda^{(k)c}_7|_k$ and $|Y_k+\mathcal O_k|\le Ct_k(1+(2M+5)^3)|\mathrm{Col}_k|_k$ ($\mathfrak B_{k+1}$ lives at scale $k+1$, which changes the constants by at most $L^3$).

6. (The new datum.) Compare $\mathfrak D_{k+1}$ with $\mathfrak D^\natural\!\uparrow$. The propagator domain changes from $\Omega^\natural$ to $B^{k+1}(\Lambda^{(k)'}_2)$, a deep change (\cref{lem:a06d:V7}(a), depth $\ge5r-c_d$). The vertex datum changes from $(W,g',c')$ to $(B^{k+1}(\Lambda^{(k)'}_7),g_{k+1},c_{k+1})$, including the coefficient move of \cite[(1.28)--(1.29)]{Ba3}, $c'-c_{k+1}=[\dm_{[k]}(1)-\dm_{[k]}(g_{k+1})]+[\dm_k(g')-\dm_k(g_{k+1})]$ with $\dm_{[k]}$ the index-$k$ slice (kinds (b3) and (b1)), all supported within $D_6\defeq M+LM+2$ of $\partial W$ up to exponential tails. By \cref{lem:a06d:V7}(b), \cref{lem:a06d:V9}(d) and \cref{lem:a06d:V8}(f), $|\mathcal R_{k+1}[\mathfrak D^\natural\!\uparrow]-\mathcal R_{k+1}[\mathfrak D_{k+1}]|\le Ct_k(1+(2D_6+3)^3)|\mathrm{Col}_k|_k+Ct_ke^{-\delta(5r-c_d)}N_k$. (If the new legs carry $g^\sharp$, start from the mixed datum, $g'$ on the old legs and $g^\sharp$ on the new; the difference is $X_3$.) The rescaling is exact.

7. (Assembly.) $\mathrm{Vac}^{\mathrm{act}}_{k+1}=\mathcal R_k[\mathfrak D']+\mathrm{out}^{\mathrm{act}}_k-\mathcal V^{\mathrm{glob}}_k=\mathcal R_k[\mathfrak D^\natural]+X_1+\mathrm{out}_k[\mathfrak D^\natural;C_{\mathrm{comp}}]+X_2+X_0+X_3-\mathcal V_k[\mathfrak D^\natural]-Y_k=\mathcal R_{k+1}[\mathfrak D^\natural\!\uparrow]+X_1+X_2+X_0+X_3-Y_k=\mathcal R_{k+1}[\mathfrak D_{k+1}]+\mathcal E'_k+\mathcal O_k$, with $\mathcal E_k\defeq X_1+X_2-(Y_k+\mathcal O_k)+(\mathcal R_{k+1}[\mathfrak D^\natural\!\uparrow]-\mathcal R_{k+1}[\mathfrak D_{k+1}])$ and $\mathcal E'_k\defeq\mathcal E_k+X_0+X_3$.

8. (The exponents.) $t_k\le t_k^{\kappa_0}$, and $e^{-\delta c_*R(1+\log t^{-1})^\varrho}\le C_\kappa t^\kappa$ for every $\kappa$ gives $te^{-\delta c_*r(t)}\le Ct^{16/5}$, likewise for $X_0$'s factor. ($C_\kappa$ is free of $\eps_0$ but can be very large; it is absorbed in the final constants $E_\pm$.)
\end{proof}

No term of $\mathcal R_k$ is bounded in absolute value except its collar part (\cref{cor:a06d:V13}) and differences with a designated object deep inside or in a zone; the restrictions $\chi$, the non-polynomial remainders and the degree-$0$ normalisations are Step~3(C). The field-dependent analogue (\cite[p.~412]{Ba3}: \cite[(1.5)]{Ba3} after $k+1$ steps on the new domains) follows from \cref{lem:a06d:V6,lem:a06d:V7} in the same way, with Proposition~1 under the weaker restrictions on new scalar fields \cite[p.~421]{Ba3}; its interaction-term half at a nonzero background is \BALint\ (\cref{hyp:k01:BALint}).

\paragraph{The bad localisations of \cite[p.~422]{Ba3}.} ``Passing from $k$ to $k+1$ we have to change \dots\ the terms in $\dm_{K,k}(x)(L^k\eps)^2$ and $E_{K,k}$ for which the operators in the decomposition (1.27) have the indices $j\le k$ and at least one of them is equal to $k$ \dots\ then localize them properly, which means that we have to estimate the terms with bad localizations''; the moved line of index $k$ ``is a regular function of both variables and we can treat it as an external field''. Write $S_{(k)}$ for the slice of index $k$ of \cite[(1.27)]{Ba3}, $G_{k+1}=G_k+S_{(k)}$.

\begin{lemma}[the transfer identity]\label{lem:a06d:D21}
On the torus at zero background, $\dm_{K,k}-\dm_{K,k+1}=\dm_{k+1}-\dm_k=:m_k$, and $m_k$ is the sum of the graphs of \cite[(1.23)]{Ba3} whose lines are all $G_k$ or $S_{(k)}$ with at least one $S_{(k)}$ line; the same holds for $E_{K,k}-E_{K,k+1}$.
\end{lemma}

\begin{proof}
By \cite[(1.27)]{Ba3}, $G_{K,k}=S_{(k)}+G_{K,k+1}$ and $G_{k+1}=G_k+S_{(k)}$, for scalar and vector lines alike. Expand each graph of $\dm_K$ multilinearly with every line written $G_k+S_{(k)}+G_{K,k+1}$: $\dm_{K,k}$ collects the terms with at least one line in $\{S_{(k)},G_{K,k+1}\}$, $\dm_{K,k+1}$ those with at least one line in $G_{K,k+1}$, and their difference is the terms with all lines in $\{G_k,S_{(k)}\}$ and at least one $S_{(k)}$, i.e.\ $\dm_{k+1}-\dm_k$. The counterterm vertices inside these graphs are vertices, not lines (\cref{conv:a06d:LOC}(a)), and are not expanded.
\end{proof}

\begin{lemma}[the bad-localisation terms of the moved mass counterterm]\label{lem:a06d:D22}
At step $k+1$ let the moved terms $m_k$, rescaled to the $\eta$-lattice, be localised by the rules of \cref{conv:a06d:LOC} with $\Omega^+_1\defeq B^k(\Lambda^{(k)}_7)$ and $g^+$ ($\supp g^+\subset\Omega^+_1$, $g^+=1$ at distance $\ge LM$ from $\partial\Omega^+_1$ in scale-$k$ units, \cite[p.~412]{Ba3}); assume $R\ge5M+6$, so that $r-1>LM+4$. Put $\mathrm{Bad}\defeq(\Omega^+_1)^c\cup\{z\in\Omega^+_1:\dist(z,\partial\Omega^+_1)<LM+2\}$ and $b_k(x)\defeq m^{\mathrm{glob}}_k(x)-m^{\mathrm{loc}}_k(x)$ for $x\in\Omega^+_1$, the coefficient change of the \cite[(1.7)]{Ba3} vertex caused by the move. Then: \textup{(a)} any localised family of the step-$(k+1)$ expansion containing one \cite[(1.7)]{Ba3} vertex with coefficient $b_k(x)t^2$ at $x$ is a finite sum of Proposition~1 families, each with one $S_{(k)}$ line treated as an external field and at least one vertex localised in $\mathrm{Bad}$, and $|E|\le O(1)e_k^{d_v}\lambda_k^{d_s}\prod(\text{external norms})\,C_Ke^{-(\delta'/2)(\dist(\square(x),\mathrm{Bad})-2)}e^{-(\delta'/2)d(\{\square\})}$, $\delta'\defeq\min(\delta_0,\delta_K)$, $C_K,\delta_K$ the constants of the slice kernel; \textup{(b)} the standalone vertex with both legs external satisfies $\sum_{x\in\Omega^+_1}\eta^3|b_k(x)t^2||\Phi(x)|^2\le Ct^{1/2}p^2\lambda^{-1/2}[|\Lambda^{(k)}_7\setminus\Lambda^{(k)}_8|+e^{-(\delta'/2)(r-LM-6)}|\Lambda^{(k)}_7|]$, and every family with $W\ge2$ containing the $b_k$ vertex obeys the same bound with $t^{1/2-\beta}$; \textup{(c)} in Ba{\l}aban's form: for $x\in B^k(\Lambda^{(k)}_7\setminus\Lambda^{(k)}_8)$ the size is a power of the rescaled couplings ($O(e_k^2+\lambda_k)$ for the coefficient), for $x\in B^k(\Lambda^{(k)}_8)$ it carries in addition $e^{-(\delta'/2)(r-LM-6)}$, the printed $e^{-\delta_1r}$ with $\delta_1=\delta'/4$ once $r\ge2(LM+6)$.
\end{lemma}

\begin{proof}
$m^{\mathrm{loc}}_k$ is the same sum of graphs as $m^{\mathrm{glob}}_k$ with weight $\mathbf 1_{\Omega^+_1}$ on the \cite[(1.6)--(1.7)]{Ba3} vertices and $g^+$ on the vector legs; telescoping the weights gives terms in which at least one internal vertex $z$ carries $1-\mathbf 1_{\Omega^+_1}(z)$ or $1-g^+\ne0$, so $z\in\mathrm{Bad}$. Each graph of $m_k$ has an $S_{(k)}$ line (\cref{lem:a06d:D21}); cut one in a fixed order and regroup into the renormalised families of the cut graph (as in \cref{thm:a06d:S23}, the slice line playing the role of a cut line; the $S_{(k)}$ line belongs to no divergent subgraph of the cut graph); the slice kernel is regular in both variables with $|S_{(k)}|\le C_Ke^{-\delta_K\dist}$ in the two-variable norm, the printed ``regular function of both variables''. Apply Proposition~1 to the resulting families, joining the components of a disconnected cut graph through the slice decay (\cref{lem:a06d:V10}). The tree contains $\square(x)$ and a cube in $\mathrm{Bad}$, so $d(\{\square\})\ge\dist(\square(x),\mathrm{Bad})-2$, and half of the exponent is kept for the localisation sum. The order is $\ge2$, giving $O(e_k^2+\lambda_k)=O(t)$; with $|\Phi|^2\le C^2_\Phi p^2\lambda_k^{-1/2}$ the standalone vertex is $O(t^{1/2}p^2\lambda^{-1/2})$ per cube times the tree factor. Every cube within $LM+3$ of $\partial\Omega^+_1$ lies in $B^k(\Lambda^{(k)}_7\setminus\Lambda^{(k)}_8)$ since $\dist(\Lambda_8,\Lambda^c_7)\ge r-1>LM+4$; for cubes in $B^k(\Lambda^{(k)}_8)$, $\dist(\square(x),\mathrm{Bad})\ge r-1-LM-3$. Count the cubes and use $p(t)^ne^{-\delta r(t)}\le t^\kappa$ for the exponential part.
\end{proof}

\begin{lemma}[the vacuum counterterm moves; ``even simpler'']\label{lem:a06d:D23}
$e_k\defeq[E_{k+1}-E_k]^{\mathrm{glob}}-[E_{k+1}-E_k]^{\mathrm{loc}}$ (localised with $\Omega^+_1$, $g^+$) is a $c$-number with $|e_k|\le Ct|T^{(k)}_1\setminus\Lambda^{(k)}_8|\le Ct(1+(17(r+M))^d)|\Lambda^{(k)c}_0|$; if $\Lambda^{(k)c}_0=\emptyset$ then $\mathrm{Bad}=\emptyset$ and $e_k=0$.
\end{lemma}

\begin{proof}
As in \cref{lem:a06d:D22} without external fields: Proposition~1 for vacuum families after cutting $S_{(k)}$ gives $O(t)$ per localisation (order $\ge2$). Root at a cube in $\mathrm{Bad}$, which lies in $T^{(k)}_1\setminus\Lambda^{(k)}_8$ up to the $LM+2$ layer, and sum. By \cite[(2.8)]{Ba2} every block of $\Lambda_0\setminus\Lambda_8$ is within $8(r+M)$ of $\Lambda^c_0$ (each layer is $r$ plus one large block of side $M$), and at most $(16(r+M)+1)^d\le(17(r+M))^d$ unit blocks lie within that distance of a given block of $\Lambda^c_0$.
\end{proof}

Both moves are bounded by $Ct^{1/2-\beta}(r+M)^d|\Lambda^{(k)c}_0|+C_\kappa t^\kappa|\Lambda^{(k)}_7|$, against which the large-field ledger of \cite[Sect.~3C]{Ba2} pays $(d+1)(1+\ln t^{-1})$ per unit of $|\Lambda^{(k)c}_0|$; that ledger is \cite[(3.42)]{Ba2}, part of the $L=5$ constants of \cref{thm:a06:BALA}.

% a06_BALtorus_d4: part of Step 3 of the proof of thm:a06:T1 (input from sec/a06_BALtorus_c.tex).
% Group (C), first half: the large-field chain of Ba2 Sect. 3A, (3.1)-(3.20), and the k-uniformity of Ba2's unqualified O(1) lines.

\subsubsection*{Step 3(C$_1$): the large-field chain of \cite[Sect.~3A]{Ba2}}

\cite[Sect.~3A]{Ba2} inverts the renormalisation transformations in the large-field region: the inductive hypothesis \cite[(3.5)]{Ba2} at level $k$ is reduced to the same statement at level $k-1$ at the cost of a factor $\exp(O(1)(L^k\eps)^{\kappa_0}|T_\eps|)$, through the Gaussian integral \cite[(3.9)]{Ba2} with the perturbing kernels $H'_k,H''_k$ \cite[(3.8), (3.12)]{Ba2}, the linear terms $F'_k,F''_k$ \cite[(3.10)]{Ba2}, the lower bounds \cite[(3.11)]{Ba2} and the Gaussian comparison \cite[(3.17)--(3.20)]{Ba2}. What is consumed is that each inverse step costs $\exp(C_\kappa t_k^\kappa|T^{(k)}_1|)$ with $C_\kappa$ independent of $k$, $\eps$, $\eps_0$ and the torus. We prove this with explicit constants. Besides the notation of Step~3(B), $M_v\defeq c_Ap(t)/(\lmu t)$ and $M_s\defeq c_sp(t)(\lambda t)^{-1/4}$ are the small-field bounds \cite[(2.16)--(2.17)]{Ba2} at scale $k$ ($c_A\defeq2+2Ld\lmu$, $c_s\defeq2+2Ld$, under the cap $\lambda t\le1$), $M(t)\defeq c_Mt^{-1}p(t)\ge\max(M_v,M_s)$ with $c_M\defeq\max(c_A/\lmu,c_s\lambda^{-1/4})$; scale-$(k-1)$ collar fields rescaled to scale $k$ are bounded by $c_LM_v(t)$ with $c_L\defeq L^{3/2}(1+\ln L)^{\bar p}$; $S_\delta\defeq\sup_x\sum_ze^{-\delta|x-z|}$; the separation of a primed scale-$(k-1)$ layer from the next is $D_k\defeq(r(t)-c_b)/L\ge r(t)/(2L)$ for $R\ge2c_b$, $c_b$ the large-block diameter. Two sup bounds replace every ``for $\eps_0$ small'' by a constant: for $\kappa\ge0$, $\delta>0$, $a>0$, $m\ge0$,
\[
E(\kappa,\delta)\defeq\sup_{0<t\le1}t^{-\kappa}e^{-\delta r(t)}<\infty,\qquad P(a,m)\defeq\sup_{0<t\le1}t^ap(t)^m<\infty,
\]
(with $t=e^{1-u}$, $u\ge1$: $\exp(\kappa(u-1)-\delta Ru^\varrho)$ is bounded since $\varrho>1$, and $u^{m\bar p}e^{-a(u-1)}$ is bounded), depending only on $(\kappa,\delta,R,\varrho)$ resp.\ $(a,m,b_0,\bar p)$. Three kernel facts are used throughout (sup-distance, $d(\cdot)$ a $1$-Lipschitz distance function): (K-a) if $|K(x,x')|\le ce^{-\delta(|x-x'|+d(x)+d(x'))}$ and $d\ge D$ on $Y$, then $|K|\le ce^{-2\delta D}e^{-\delta|x-x'|}$ on $Y\times Y$; (K-b) if $|K_0(x,x')|\le ce^{-\delta|x-x'|}$ for $x\in Y$, $x'\in Z$, $\dist(Y,Z)\ge D$, then $|K_0|\le ce^{-\delta D/2}e^{-\delta|x-x'|/2}$; (K-c) if $|K_1(x,z)|\le c_1e^{-\delta|x-z|}$, $|K_3(z',x')|\le c_3e^{-\delta|z'-x'|}$ and $|K_2(z,z')|\le c_2e^{-\delta(|z-z'|+d(z)+d(z'))}$, then $|(K_1K_2K_3)(x,x')|\le c_1c_2c_3S^2_{\delta/4}e^{-(\delta/4)(|x-x'|+d(x)+d(x'))}$ (with $E\defeq|x-z|+|z-z'|+|z'-x'|+d(z)+d(z')$, the triangle inequality and the Lipschitz property give $|x-x'|+d(x)+d(x')\le2E$, so the exponent $\frac{\delta}{4}(|x-x'|+d(x)+d(x'))+\frac\delta4(|x-z|+|z'-x'|)\le\delta E$; sum over $z,z'$); telescoping a product of $n$ such kernels gives the decay of (K-c) with rate $\delta/4^{n-1}$.

\begin{lemma}[the kernel types of the inverse step]\label{lem:a06d:C21}
At scale $k$ let $Y$ be a set of unit-lattice points and $D\defeq D_k$. Each of the following kernels $K(x,x')$ obeys $|K(x,x')|\le C_He^{-\delta_Hr(t)}e^{-\delta_H|x-x'|}$, with $\delta_H\defeq\min(\delta_7,\delta^S_7/4,\theta_0/4,\theta_0(\gamma_L,L)/6,\theta_W)/(16L)$ (the rates of \cite[Prop.~2.2, (2.27), (2.29)]{Ba1}, of the domain differences of $C_\Lambda$, of the Combes--Thomas bound and of the conditioning-region difference: \cref{prop:a06A:Ba1P22}, \cref{lem:a06A:domdiff}(b), \cref{lem:a06A:CT} and \cref{lem:a06A:CLdiff}) and $C_H$ explicit, all free of $k$, $\eps$, $\eps_0$ and the torus:
(a) \emph{domain change}: $K=O(\Omega_1)-O(\Omega_2)$, $\Omega_1\subset\Omega_2$ unions of big blocks, $O\in\{\Delta^{(k)}(\cdot,\tilde A),C_\Lambda(\cdot,\tilde A)\}$ at a regular background, with $\dist(Y,\Omega^{(k)c}_1)\ge D$;
(a$'$) \emph{conditioning-region change}: $K=\Delta^{(k)}_{\Lambda^{(k-1)}_5}-\Delta^{(k)}=(aL^{-2})^2Q(C-C_{\Lambda_5})Q^*$ (from \cite[(2.46)]{Ba2}, $\Delta^{(k)}_{\Lambda_5}=aL^{-2}-(aL^{-2})^2QC^{(k-1)}_{\Lambda_5}Q^*$), on $Y$ with $\dist(Y,\Omega^{(k-1)}\setminus\Lambda^{(k-1)}_5)\ge D$; this $\Lambda$-change, not covered by (a), is the only change in the photon form of \cite[(2.47), p.~568]{Ba2} not already of types (a)--(d);
(b) \emph{products}: $O_1\cdots O_n-O'_1\cdots O'_n$ ($n\le3$) with factors among $\Delta^{(k)}$, $C_\Lambda$, $a_kG_kQ_k^*$ and adjoints, each pair $O_i,O'_i$ of type (a) or equal;
(c) \emph{cut-off}: $K=\mathbf 1_YK_0\mathbf 1_Z$ with $|K_0(x,x')|\le ce^{-\delta|x-x'|}$ and $\dist(Y,Z)\ge D$;
(d) \emph{background change}: $K=O(\tilde A+\tilde B^{(k)})-O(\tilde A)$ with $|\tilde B^{(k)}|\le C_\kappa t^\kappa$ \cite[(2.90)]{Ba2}, where $e^{-\delta_Hr(t)}$ is replaced by $Ct^{\kappa+1/2}$ for every $\kappa$ (not used in Sect.~3A, which has no background change; used in Step~3(C$_2$)).
\end{lemma}

\begin{proof}
(a): \cite[(2.29)]{Ba1} (\cref{prop:a06A:Ba1P22}) gives $c_7e^{-\delta_7(|x-x'|+d(x)+d(x'))}$ with $d=\dist(\cdot,\Omega^{(k)c}_1)\ge D$ on $Y$, and the domain difference of $C_\Lambda$ the same with $\delta^S_7/4$ (\cref{lem:a06A:domdiff}(b)); then (K-a). (a$'$): the difference $C-C_{\Lambda_5}$ at scale $k-1$ decays at rate $\theta_0(\gamma_L,L)/6$ (\cref{lem:a06A:CLdiff}) in $|y-y'|+\dist(y,\Omega^{(k-1)}\setminus\Lambda^{(k-1)}_5)+\dist(y',\text{same})$, and on $\Lambda^{(k)}_5$ that distance is at least two layers; then the coarse transfer $QKQ^*$ and (K-a). (b): (K-c) and its telescoping, rate $\ge\delta/16$ for $n\le3$. (c): (K-b). In (a)--(c), $e^{-\delta D}\le e^{\delta c_b/L}e^{-\delta r(t)/L}$. (d): the analytic family of block operators in the background $\tilde A+\tau\tilde B^{(k)}$ (Step~3(C$_2$), \cref{lem:a06d:S1}) with $w=\norm qe_kC_\kappa t^\kappa$.
\end{proof}

\begin{lemma}[{\cite[(3.8), (3.10)]{Ba2}}]\label{lem:a06d:C22}
\textup{(i)} $H_k$ of \cite[p.~584]{Ba2} (``we have applied formula (2.108) together with the remark following it''): by \cite[(2.108)]{Ba2} read with the background $\tilde A^{(k)}$, the difference between $\langle\phi,\Delta^{(k)}(B^k(\Lambda^{(k-1)'}_2),\tilde A^{(k)})\phi\rangle$ on $\Lambda^{(k-1)'}_6$ and its three localised pieces is the sum of type-(a) kernels on $(\Lambda^{(k-1)'}_6\cap\Lambda^{(k)c}_3)^2$ and on $(\Lambda^{(k-1)'}_6\cap\Lambda^{(k)c}_3)\times(\Lambda^{(k)}_3\cap\Lambda^{(k)c}_4)$ ($\Omega_1=B^k(\Lambda^{(k-1)'}_2\cap\Lambda^{(k)c}_5)$), a type-(c) kernel on $(\Lambda^{(k-1)'}_6\cap\Lambda^{(k)c}_3)\times\Lambda^{(k)}_4$ (distance $\ge r-c_b$), and a type-(a) kernel on $(\Lambda^{(k)}_3)^2$ ($\Omega_1=B^k(\Lambda^{(k)}_2)\subset\Omega_2=B^k(\Lambda^{(k-1)'}_2)$); no background change occurs, since by \cite[(3.4)]{Ba2} $\tilde A^{(k)}-\tilde A^{(k+1)}=(1-\theta_{k+1})\theta_kA^{(k)}$ vanishes on $B^k(\Lambda^{(k)}_2)$. Hence \cite[(3.8)]{Ba2} holds for $h_k$ with $\delta_1=\delta_0=\delta_H$ and $O(1)=C'_H\defeq4C_H$. \textup{(ii)} $H'_k,H''_k,F'_k,F''_k$ (\cite[p.~585]{Ba2}: ``terms associated with the changes in (2.47) and the discussion following, and the term $H_k$''): by \cite[(2.47), p.~568]{Ba2} the changes are the domain splits of $\Delta^{(k)}_{\Lambda^{(k-1)}_5}(B^{k-1}(\Lambda^{(k-1)}_2),A^{(k)})$ (type (a)), the replacement of $\Delta^{(k)}_{\Lambda^{(k-1)}_5}$ by $\Delta^{(k)}$ on $\Lambda^{(k-1)'}_6$ in the third line of \cite[(2.47)]{Ba2} for the scalar and the photon form (type (a$'$)), the restriction of the third and ninth terms of \cite[(2.46)]{Ba2} to $\Lambda^{(k-1)'}_5\cap\Lambda^{(k-1)'c}_6$ (type (c)), and the Neumann conditions on $\partial B^{k-1}(\Lambda^{(k-1)}_2\cap\Lambda^{(k-1)c}_7)$ in the seventh to ninth terms, whose operators are the products $\Delta C_{\Lambda_5}\Delta$ and $\Delta C_{\Lambda_5}Q^*$ (type (b)); so $H'_k,H''_k$ obey \cite[(3.8)]{Ba2} with a fixed multiple of $C_H$. \textup{(iii)} $F'_k$ (resp.\ $F''_k$) is the sum of the type-(c) and type-(a)/(b) remainders of the linear terms and of $\sum_{x'\notin\Lambda^{(k)}_5}h'(x,x')v(x')$ with $v$ the collar fields, $|v|\le c_LM(t)$; hence $|F'_k(x)|\le C_HS_{\delta_H}e^{-\delta_Hr(t)}c_LM(t)\le C_Fe^{-\delta_Hr(t)/2}$ with $C_F\defeq c_LC_HS_{\delta_H}c_M\sup_tt^{-1}p(t)e^{-\delta_Hr(t)/2}<\infty$, and the same for $F''_k$.
\end{lemma}

\begin{proof}
Each item is the cited type with the cited separations (\cref{lem:a06d:C21}); the separations are those of the layers. Ba{\l}aban defines $H'_k,H''_k,F'_k,F''_k$ by reference (\cite[p.~568]{Ba2}: ``Terms omitted are described more precisely later''); the list of changes is his, and each is of one of the types.
\end{proof}

\begin{lemma}[{\cite[(3.8)]{Ba2} implies \cite[(3.12)]{Ba2} for every $\kappa$}]\label{lem:a06d:C23}
If $|h(x,x')|\le ce^{-\delta_1r(t)}e^{-\delta_0|x-x'|}$ on $Y$, then the operator $H$ on $\ell^2(Y)$ has $\norm H\le cS_{\delta_0}e^{-\delta_1r(t)}\le cS_{\delta_0}E(\kappa,\delta_1)t^\kappa$ for every $\kappa$.
\end{lemma}

\begin{proof}
Schur test, then the definition of $E(\kappa,\delta_1)$.
\end{proof}

\begin{lemma}[{\cite[(3.11)]{Ba2} with $\gamma_1=\frac12$, without field restriction or regularity of $\tilde A$}]\label{lem:a06d:C24}
Let $t\le\tau_1\defeq\min\{1,a_*^{1/2}/\lmu,a_*^{1/2}/m\}$. Then $G'_k\ge\frac12\lmu^2t^2$ and $G''_k\ge\frac12m^2t^2$ on $\ell^2(\Lambda^{(k)}_5)$.
\end{lemma}

\begin{proof}
$G'_k$ (resp.\ $G''_k$) is the quadratic part in $A_k|_{\Lambda_5}$ (resp.\ $\phi_k|_{\Lambda_5}$) of the first four terms of \cite[(3.9)]{Ba2}. The block-spin terms $\frac12\sum_la_{l-k}(L^{l-k})^{d-2}\sum|A_l-Q_{l-k}A_k|^2$ (resp.\ with $Q_{l-k}(\tilde A^{(k+1)})$) contribute nonnegative forms. The forms $\langle\Lambda_5A_k,\Delta^{(k)}_{\Lambda^{(k-1)}_5}\Lambda_5A_k\rangle$ and $\langle\Lambda_5\phi_k,\Delta^{(k)}_{\Lambda^{(k-1)}_5}(B^k(\Lambda^{(k-1)'}_2),\tilde A^{(k)})\Lambda_5\phi_k\rangle$ are compressions of conditional forms: the subscript records the conditional integration of \cite[(2.46)]{Ba2} over the intermediate field on $\Lambda^{(k-1)}_5$ with its collar values as separate variables, so their quadratic part is the partial elimination of that field in the two-level form (block-spin term of the step) $+\langle A',\Delta^{(k-1)}A'\rangle$, whose full elimination is $\Delta^{(k)}$ by the composition identity \cite[(2.40)--(2.41)]{Ba1}. For $H\ge0$, $m_0^2\ge0$, $a>0$ and $\norm Q\le1$, the Schur complement $\Delta\defeq a-a^2Q(H+m_0^2+aQ^*Q)^{-1}Q^*$ satisfies $\langle\psi,\Delta\psi\rangle=\min_\phi[a\norm{\psi-Q\phi}^2+\langle\phi,(H+m_0^2)\phi\rangle]\ge\min_u[a\norm{\psi-u}^2+m_0^2\norm u^2]=\frac{am_0^2}{a+m_0^2}\norm\psi^2$; a partial minimisation over a smaller set is larger, and a compression of an operator $\ge c$ is $\ge c$ on its range. The fundamental operators carry the masses $\lmu^2t^2$ (photon, $-\Delta^\eta+\lmu^2(L^k\eps)^2$) and $m^2t^2$ (scalar, $-\Delta^{\eta,N}_{A^{(k)}}+m^2(L^k\eps)^2$, \cite[p.~585]{Ba2}), and $\norm{Q_k(A)}\le1$ for every real background. So each form is $\ge a_km_0^2t^2/(a_k+m_0^2t^2)\ge\frac12m_0^2t^2$, $m_0\in\{\lmu,m\}$, since $m_0^2t^2\le a_*\le a_k$.
\end{proof}

This answers the ``corollary of Proposition~3.1'' of \cite[p.~586]{Ba2}: no restriction on the fields and no regularity of $\tilde A^{(k)}$ is needed, only $-\Delta^{\eta,N}_{A,\Omega}\ge0$ for every real background.

\begin{lemma}[{\cite[(3.15), (3.16)]{Ba2}}]\label{lem:a06d:C25}
Let a field $v_l$ on the collar $\Lambda^{(l-1)'}_{-1}\cap\Lambda^{(l)c}_5$ at scale $l\ge k+1$ obey $|v_l|\le M_v(t_l)$ (vector) or $M_s(t_l)$ (scalar), and let $v^{(k)}_l\defeq L^{-(l-k)/2}v_l(L^{-(l-k)}\cdot)$ be its image in scale-$k$ units. Then $|A^{(k)}_l|\le L^{-3(l-k)/2}M_v(t)$, $|\phi^{(k)}_l|\le L^{-3(l-k)/4}M_s(t)$ (Ba{\l}aban's $L^{-(l-k)d/2}$, $L^{-(l-k)d/4}$), and the fields $\Phi',\Phi''$ of \cite[(3.13)--(3.14)]{Ba2} satisfy $|\Phi'(x)|\le aM_v(t)/(L^{7/2}-1)$, $|\Phi''(x)|\le aM_s(t)/(L^{11/4}-1)$ on $\Lambda^{(k)}_5$, both $\le aM(t)\le ac_MP(1,1)t^{-2}$.
\end{lemma}

\begin{proof}
$t_l=L^{l-k}t$ and $p(t_l)\le p(t)$, so $L^{-(l-k)/2}c_Ap(t_l)/(\lmu t_l)\le L^{-3(l-k)/2}M_v(t)$, and likewise for the scalar. $|Q^*_{l-k}(\tilde A)\phi_l|=|\phi_l|$ (orthogonal transport) and $a_{l-k}\le a$, so $|\Phi'(x)|\le\sum_{l>k}aL^{-2(l-k)}L^{-3(l-k)/2}M_v(t)$, a geometric series; the same for $\Phi''$. Finally $M(t)\le c_MP(1,1)t^{-2}$.
\end{proof}

\begin{lemma}[the linear coefficient]\label{lem:a06d:C26}
Let $J'_k$ ($J''_k$) be the linear coefficient of $A_k|_{\Lambda_5}$ ($\phi_k|_{\Lambda_5}$) in \cite[(3.9)]{Ba2} without $F'_k$ ($F''_k$). Then $|J'_k(x)|,|J''_k(x)|\le C_JM(t)$ on $\Lambda^{(k)}_5$, hence $\norm J\le C_JM(t)|\Lambda^{(k)}_5|^{1/2}$, with $C_J\defeq a+c_{\mathrm{coll}}+c_Lc_f$: $c_{\mathrm{coll}}$ the $\ell^\infty\to\ell^\infty$ norm of the kernel of the conditional operator $\Delta^{(k)}_{\Lambda^{(k-1)}_5}$ from the collar $\Lambda^{(k-1)'}_5\cap\Lambda^{(k)c}_5$ to $\Lambda^{(k)}_5$ (its decay is that of $aL^{-2}-(aL^{-2})^2QC_{\Lambda_5}Q^*$ with $|C_{\Lambda_5}(y,y')|\le c_2e^{-\theta_0|y-y'|}$), $c_f$ that of the kernels of the third and ninth terms of \cite[(2.46)]{Ba2} (products $\Delta^{(k-1)}C^{(k-1)}_{\Lambda_5}Q^*$ of decaying kernels, $c_f\le aL^{-2}c_7c_2S_{\delta_7}S_{\theta_0}$).
\end{lemma}

\begin{proof}
$\Phi'$ by \cref{lem:a06d:C25}; the collar term by the conditional operator's decay and $|A_k|\le M_v$ on the collar; $f'_k$ is a sum of such kernels applied to scale-$(k-1)$ collar fields bounded by $c_LM(t)$.
\end{proof}

This replaces Ba{\l}aban's $O((L^k\eps)^{-2})|\Lambda^{(k)}_5|^{1/2}$ \cite[p.~587]{Ba2} by $C_Jc_Mt^{-1}p(t)|\Lambda_5|^{1/2}$.

\begin{lemma}[Gaussian comparison]\label{lem:a06d:C27}
Let $V$ be a real Euclidean space of dimension $n$, $G,H$ symmetric on $V$ with $G\ge g>0$ and $\norm H\le h\le g/2$, and $J,F\in V$. Then
\[
\int_Ve^{-\frac12\langle x,(G-H)x\rangle+\langle J+F,x\rangle}dx\le\int_Ve^{-\frac12\langle x,Gx\rangle+\langle J,x\rangle}dx\cdot\exp\Big[\frac{nh}g+\frac{h\norm J^2}{g^2}+\frac{2\norm F\norm J}g+\frac{\norm F^2}g\Big].
\]
\end{lemma}

\begin{proof}
Both integrals are Gaussian, $\int e^{-\frac12\langle x,Mx\rangle+\langle b,x\rangle}=(2\pi)^{n/2}(\det M)^{-1/2}e^{\frac12\langle b,M^{-1}b\rangle}$ for $M>0$ (here $G-H\ge g/2$). Determinant (\cite[(3.18)--(3.20)]{Ba2}): $\det(G-H)=\det G\det(I-X)$, $X\defeq G^{-1/2}HG^{-1/2}$, $\norm X\le h/g\le\frac12$; for an eigenvalue $\lambda$ of $X$, $-\log(1-\lambda)\le2|\lambda|$, so $-\frac12\log\det(I-X)\le\sum|\lambda_i|\le nh/g$. Exponent: $(G-H)^{-1}-G^{-1}=G^{-1}H(G-H)^{-1}$ has norm $\le2h/g^2$ and $\norm{(G-H)^{-1}}\le2/g$, so $\frac12\langle J+F,(G-H)^{-1}(J+F)\rangle-\frac12\langle J,G^{-1}J\rangle\le h\norm J^2/g^2+2\norm F\norm J/g+\norm F^2/g$.
\end{proof}

\begin{corollary}[{\cite[(3.9)]{Ba2} against the unperturbed integral}]\label{cor:a06d:C28}
Let $t\le\tau_1$ and assume $C'_HS_{\delta_H}e^{-\delta_HR}\le\mu_{\mathrm{lo}}^2/4$ and $\delta_HR\ge2$, $\mu_{\mathrm{lo}}^2\defeq\min(\lmu^2,m^2)$ (a condition on $R$, chosen before $\eps_0$). Then for every $\kappa>0$, \cite[(3.9)]{Ba2} is at most \cite[(3.9)]{Ba2} at $H'_k=H''_k=0$, $F'_k=F''_k=0$, times $\exp(C_{\kappa}t^\kappa|\Lambda^{(k)}_5|)$, with
\[
C_\kappa\defeq\sup_{0<t\le1}t^{-\kappa}\Big[\frac{2(d+N)h(t)}{\mu_{\mathrm{lo}}^2t^2}+\frac{4h(t)C_J^2M(t)^2}{\mu_{\mathrm{lo}}^4t^4}+\frac{4(d+N)^{1/2}C_Fe^{-\delta_Hr(t)/2}C_JM(t)}{\mu_{\mathrm{lo}}^2t^2}+\frac{2(d+N)C_F^2e^{-\delta_Hr(t)}}{\mu_{\mathrm{lo}}^2t^2}\Big],\quad h(t)\defeq C'_HS_{\delta_H}e^{-\delta_Hr(t)},
\]
finite since every term is $e^{-\delta_Hr(t)/2}$ times a power of $t^{-1}p(t)$.
\end{corollary}

\begin{proof}
$V\defeq\ell^2(\Lambda^{(k)}_5)^d\oplus\ell^2(\Lambda^{(k)}_5)^N$, $n=(d+N)|\Lambda^{(k)}_5|$; $G\defeq G'\oplus G''\ge g\defeq\frac12\mu_{\mathrm{lo}}^2t^2$ (\cref{lem:a06d:C24}); $H\defeq H'\oplus H''$ with $h=h(t)$ (\cref{lem:a06d:C22,lem:a06d:C23}), and $h\le g/2$ for all $t\le1$ by the cap ($e^{-\delta_Hr(t)}\le e^{-\delta_HR}t^2$ when $\delta_HR\ge2$); $\norm J\le C_JM(t)|\Lambda_5|^{1/2}$ (\cref{lem:a06d:C26}); $\norm F\le(d+N)^{1/2}C_Fe^{-\delta_Hr/2}|\Lambda_5|^{1/2}$ (\cref{lem:a06d:C22}(iii)). \cref{lem:a06d:C27}, with the scalar background treated as an independent external field (\cite[p.~568]{Ba2}), gives the bound.
\end{proof}

\begin{lemma}[{the collar forms, \cite[p.~584]{Ba2}}]\label{lem:a06d:C29}
On the collar $\Lambda^{(k-1)'}_5\cap\Lambda^{(k)c}_5$, where $|A_k|\le M_v$, $|\phi_k|\le M_s$, every form $\frac12\langle v,Kv\rangle$ with $K$ of the types of \cref{lem:a06d:C21} obeys $|\frac12\langle v,Kv\rangle|\le\frac12C'_HS_{\delta_H}e^{-\delta_Hr(t)}M(t)^2|\text{collar}|\le C'_\kappa t^\kappa|\Lambda^{(k-1)'}_5\cap\Lambda^{(k)c}_5|$ with $C'_\kappa\defeq\frac12C'_HS_{\delta_H}c_L^2c_M^2\sup_tt^{-\kappa-2}p(t)^2e^{-\delta_Hr(t)}<\infty$ (Schur test).
\end{lemma}

\begin{proposition}[{the inductive step \cite[(3.5)]{Ba2}$_k\Rightarrow$\cite[(3.5)]{Ba2}$_{k-1}$}]\label{prop:a06d:C211}
Assume \cite[(3.5)]{Ba2} at level $k$. Then it holds at level $k-1$ with the additional factor $\exp(C_{\mathrm{LF},\kappa}t_k^\kappa|T^{(k)}_1|)$, $C_{\mathrm{LF},\kappa}\defeq C_\kappa+3C'_\kappa$: the $T$-transformations are integral operators with nonnegative kernels, so inequalities inside the nested brackets are preserved; the transformations not depending on $A_k,\phi_k|_{\Lambda^{(k)}_5}$ are moved out (Fubini--Tonelli); the forms of \cite[(3.5)]{Ba2} at level $k$ are merged by \cite[(2.108)]{Ba2} into one form plus $H_k$ (\cref{lem:a06d:C22}(i)), the forms evaluated on the collar cost \cref{lem:a06d:C29}, the remaining integral is \cite[(3.9)]{Ba2} (\cref{lem:a06d:C22}(ii)) and \cref{cor:a06d:C28} compares it with the $H=F=0$ integral; the latter is, by the definitions \cite[(2.45)--(2.50)]{Ba2} read backwards, the level-$k$ factor of $\varrho''^{(k)}$ \cite[(2.46)]{Ba2}, and \cite[(2.48)]{Ba2} replaces $\varrho''$ by $\varrho$ (for $k<K$ its interior part is of $H'$-type and sits inside \cref{cor:a06d:C28}, its collar part is a \cref{lem:a06d:C29} form; at $k=K$, where \cite[p.~583]{Ba2} restricts $B,\psi$ everywhere, the whole of \cite[(2.48)]{Ba2} is a sum of \cref{lem:a06d:C29} forms, which is the base case \cite[(3.1)]{Ba2}). The sum over levels of the added factors is $\sum_lC_{\mathrm{LF},\kappa}t_l^\kappa|T^{(l)}_1|=\sum_lC_{\mathrm{LF},\kappa}t_l^{\kappa-d}|T_\eps|$, Ba{\l}aban's $\exp(\sum_lO(1)(L^l\eps)^{\kappa_0}|T_\eps|)$ with $\kappa_0=\kappa-d$ and $O(1)=C_{\mathrm{LF},\kappa}$ independent of $l$; the induction is finite ($K$ steps).
\end{proposition}

Every constant above is explicit and free of $k$, $\eps$, $\eps_0$ and the torus; the caps used are $t\le\tau_1$ and the condition on $R$ of \cref{cor:a06d:C28}, neither involving $\eps_0$.

% a06_BALtorus_d4b: continuation of a06_BALtorus_d4 (input from sec/a06_BALtorus_c.tex): the k-uniformity of Ba2's unqualified O(1) lines.

\subsubsection*{The $k$-uniformity of the unqualified $O(1)$ lines of \cite{Ba2}}

Eight lines of \cite{Ba2} carry an $O(1)$ without a stated dependence: \cite[(2.61), (2.86), (2.91), (2.94), (2.95), (2.103), (2.117), (2.118)]{Ba2}. \cite[(2.61)]{Ba2} and \cite[(2.95)]{Ba2} inherit the $L=5$ constants of \cite[Lemmas~2.3, 2.4]{Ba2} (\cref{thm:a06:BALA}); \cite[(2.118)]{Ba2} is the geometric sum of \cref{prop:a06d:C211}; \cite[(2.117)]{Ba2}'s constants are $c_5$ (from \cite[Prop.~3.1$'$]{Ba4}) and $c_7=C_{2.57}+C_{2.78}$, with $C_{2.57}$ from \cref{cor:a06d:V13} and its field-dependent analogue and $C_{2.78}$ arithmetic; the remaining three are the following lemmas.

\begin{lemma}[{\cite[(2.86)]{Ba2} with an explicit $k$-free constant}]\label{lem:a06d:C61}
At scale $k$ let $\Lambda\defeq\Lambda^{(k)}_0$, $B$ the new block photon field, $A$ the scale-$k$ photon field, and assume \cite[(2.17)]{Ba2} at scale $k$ on $\Lambda'_{-1}$: $|A(x)-(Q^*B)(x)|\le2Ldp$, $|B(y)-B(y')|\le3Ldp$ for neighbouring $L$-blocks, $|B(y)|\le2p/(\lmu t)$. Then for $x\in\Lambda^{(k)}_1$, $|A'(x)|\defeq|A(x)-aL^{-2}(C^{(k)}_\Lambda Q^*B)(x)|\le C_{2.86}\,p$ with
\[
C_{2.86}\defeq2Ld+c_2S_{\theta_0}\Big[3dc_7\sum_ne^{-\delta_7|n|}(|n|_1+dL)+2\lmu+\frac{4c_7S_{\delta_7/2}E(1,\delta_7/4)}{\lmu}\Big]+\frac{2c_2c_7S_{\delta_7/2}S_{\theta''}e^{\theta''c_b}E(1,\theta'')}{\lmu},
\]
$\theta''\defeq\min(\theta_0,\delta_7/2)/2$, assuming $R\ge4(c_b+dL)$.
\end{lemma}

\begin{proof}
Photon operators are at zero background on the torus \cite[p.~608]{Ba1}, componentwise. $f\defeq\Lambda Q^*B$ lies in the range of $\Lambda P\Lambda$, so $aL^{-2}C_\Lambda f-f=-C_\Lambda\Lambda\Delta\Lambda f$ with $\Delta\defeq\Delta^{(k)}$, and since $f(x)=(Q^*B)(x)$ on $\Lambda$, $A'(x)=[A(x)-(Q^*B)(x)]+(C_\Lambda\Lambda\Delta\Lambda f)(x)$. $\Delta\mathbf 1=\delta_\mu\mathbf 1$ with $\delta_\mu=a_k\lmu^2t^2/(a_k+\lmu^2t^2)\le\lmu^2t^2$, so for $z\in\Lambda$, with $y(z)$ the $L$-block of $z$, $(\Delta\Lambda f)(z)=\sum_{z'\in\Lambda}\Delta(z,z')[B(y(z'))-B(y(z))]+B(y(z))[\delta_\mu-\sum_{z'\notin\Lambda}\Delta(z,z')]$. First sum: $|\Delta(z,z')|\le c_7e^{-\delta_7|z-z'|}$ (\cite[(2.27)]{Ba1}, \cref{prop:a06A:Ba1P22}); if $|z-z'|\le r/2$ the axis-ordered chain of $L$-blocks from $y(z)$ to $y(z')$ stays in $\Lambda_{-1}$ and \cite[(2.17)]{Ba2} gives $|B(y(z'))-B(y(z))|\le3dp(|z-z'|_1+dL)$; if $|z-z'|>r/2$ use $|B|\le2p/(\lmu t)$ and $e^{-\delta_7|n|}\le e^{-\delta_7r/4}e^{-\delta_7|n|/2}$, the tail being $\le4c_7S_{\delta_7/2}e^{-\delta_7r/4}p/(\lmu t)$. Second term: $|B|\delta_\mu\le2\lmu tp\le2\lmu p$, and $|B|\sum_{z'\notin\Lambda}|\Delta(z,z')|\le(2p/(\lmu t))c_7S_{\delta_7/2}e^{-(\delta_7/2)\dist(z,\Lambda^c)}$. Apply $|C_\Lambda(x,z)|\le c_2e^{-\theta_0|x-z|}$ (\cref{lem:a06A:blockinv}(c)): the bounded parts give $c_2S_{\theta_0}(\dots)p$, the boundary part gives $2c_2c_7S_{\delta_7/2}S_{\theta''}e^{-\theta''\dist(x,\Lambda^c)}p/(\lmu t)$ with $\dist(x,\Lambda^c)\ge r-c_b$ for $x\in\Lambda_1$, and $E(1,\theta'')$ absorbs $t^{-1}e^{-\theta''r}$.
\end{proof}

\begin{lemma}[{\cite[(2.91)]{Ba2}, value clause}]\label{lem:a06d:C62}
With $A'^{(k)}\defeq a_k\zeta^{(k)}G_kQ_k^*A'$ (the first term of \cite[(2.89)]{Ba2}), $|\zeta^{(k)}|\le1$, $\supp\zeta^{(k)}(\cdot,y)\subset\{|x-y|<r(t)-2M\}$ (\cref{thm:a06:BALA}(b)), and $|(G_kQ_k^*)(x,y)|\le C_1e^{-\delta_*|x-y|_1}$ uniformly in $k$ (\cref{thm:a06:A5}): for $x$ within $M$ of $B^k(\Lambda^{(k)}_2)$, $|\theta_{k+1}A'^{(k)}(x)|\le C_{2.91}\,p$ with $C_{2.91}\defeq aC_1S_*(C_{2.86}+C')$, $S_*\defeq\sum_ne^{-\delta_*|n|_1/2}$, $C'\defeq e^{\delta_*(c_b+M+1)/2}S_*(c_A+2aL^{-2}c_2S_{\theta_0})E(1,\delta_*/2)/\lmu$.
\end{lemma}

\begin{proof}
Split the $y$-sum: $y\in\Lambda^{(k)}_1$ uses \cref{lem:a06d:C61}; $y\in\Lambda^{(k)}_0\setminus\Lambda^{(k)}_1$, the only other $y$ reached since $2(r-c_b)>r-2M+M$, has $|A'(y)|\le M_v+2aL^{-2}c_2S_{\theta_0}p/(\lmu t)$ and $|x-y|\ge r-c_b-M-1$, which gives the $E(1,\delta_*/2)$ term.
\end{proof}

The derivative clause of \cite[(2.91)]{Ba2} is \cite[Lemma~2.3, (2.60)]{Ba2} at $L=5$ (\cref{thm:a06:BALA}).

\begin{lemma}[{\cite[(2.103), (2.94)]{Ba2}}]\label{lem:a06d:C63}
\textup{(i)} Let $n$ be odd and $I\defeq[1-\gamma_1/\gamma_0,1-\gamma_0/\gamma_1]$ with $\gamma_0\defeq\gamma_X$, $\gamma_1\defeq a+aL^{-2}$ (the two-sided bound $\gamma_X\le X_\Lambda\le a_k+aL^{-2}$ of \cref{lem:a06A:blockinv}(b), so $I$ does not depend on $j,k$). Then $-\frac12\ln(1-\lambda)\le\sum_{l\le n}\lambda^l/(2l)+c_n\lambda^{n+1}$ on $I$ with $c_n\defeq\max(0,\max_Ig)$, $g(\lambda)\defeq[-\frac12\ln(1-\lambda)-\sum_{l\le n}\lambda^l/(2l)]/\lambda^{n+1}$, $g(0)\defeq1/(2(n+1))$: $g$ is continuous on the compact $I$ ($1-\lambda\ge\gamma_0/\gamma_1>0$) and $\lambda^{n+1}\ge0$. \textup{(ii)} $O(1)(L^j\eta)^{d/2}p(L^k\eps)\le O(1)p(L^j\eps)$ for $j\le k$ with the same $O(1)$: $(L^j\eta)^{d/2}=L^{(j-k)d/2}\le1$ and $p(t_k)\le p(t_j)$.
\end{lemma}

% a06_BALtorus_d5: part of Step 3 of the proof of thm:a06:T1 (input from sec/a06_BALtorus_c.tex).
% Group (C), second half: the Z-factor expansions of Ba2 p. 578 (the terms containing B~^(k)) and the eps_0-freedom of the per-step large-field constants.

\subsubsection*{Step 3(C$_2$): the $Z^{(j)}$ expansions of \cite[p.~578]{Ba2}}

At step $k+1$ the fluctuation determinants $Z^{(j)}$, $j<k$, of \cite[(2.92)]{Ba2} are expanded in the new field. With $\tilde B$ the regular part of the configuration \cite[(2.93)--(2.94)]{Ba2}, $A_1\defeq\theta_{k+1}A'^{(k)}$ and $A_2\defeq\tilde B^{(k)}$, both rescaled to scale $j$, and $F_j(Y)\defeq\ln Z^{(j)}(Y)=-\frac12\ln\det X^{(j)}(Y)+\mathrm{const}$, $X^{(j)}(Y)\defeq\Lambda(\Delta^{(j)}(\Omega,Y)+aL^{-2}P(Y))\Lambda$ on $\ell^2(\Lambda;\R^N)$, $\Omega\defeq B^j(\Lambda^{(j)}_2)$, $\Lambda\defeq\Lambda^{(j)}_5$ \cite[(2.92), (2.99)]{Ba2}, exactly
\[
F_j(\tilde B+A_1+A_2)=F_j(\tilde B)+P_{\tilde B}[A_1]+R_j+D_j,\qquad R_j\defeq F_j(\tilde B+A_1)-F_j(\tilde B)-P_{\tilde B}[A_1],\quad D_j\defeq F_j(\tilde B+A_1+A_2)-F_j(\tilde B+A_1),
\]
$P_{\tilde B}[A_1]$ the order-$\bar n$ Taylor polynomial \cite[(2.104)]{Ba2}. $R_j$ (the order-$(\bar n+1)$ trace term) is bounded by the Cauchy estimate on the analytic family of \cref{lem:a06d:S1} below: $\sum_{j<k}|R_j|\le C_{\kappa}t_k^\kappa|\Lambda^{(k)}_2|/(L^{\bar n+1-d}-1)$ for $\kappa<(\bar n+1)/2$. The terms $D_j$ ``containing $\tilde B^{(k)}$'', which \cite[p.~578]{Ba2} bounds by $O((L^j\eps)^\kappa)|B^{k-j}(\Lambda^{(k)}_2)|$ without argument, are the content of this part: a per-site count at scale $j$ gives only $Ct_k^{\kappa+1/2}L^{(k-j)(d-1)}|\Lambda^{(k)}_2|$, which is not summable over $j<k$. The mechanism is that $A_2$ is smooth at scale $k$, so at scale $j$ it is almost pure gauge on every region of size $O(1)$ (its curvature is $L^{-2(k-j)}$ times smaller than at scale $k$, its size only $L^{-(k-j)}$ times smaller), and so is $B\defeq\tilde B+A_1$; $Z^{(j)}$ is exactly gauge invariant and even under charge conjugation, so after a local gauge fixing the first variation of each local piece of $\ln Z^{(j)}$ is bilinear in the two curvatures, $L^{-2(k-j)}\times L^{-2(k-j)}$ per scale-$j$ unit cube against $L^{3(k-j)}$ cubes per scale-$k$ cube, leaving $L^{-(k-j)}$, which sums. No identity on a single graph slot is used: the local pieces are exact gauge-invariant functions.

\paragraph{Setting.} Fix a scale $j$, $\Omega$ a union of unit blocks of the fine lattice $\xi\Z^3$, $\xi=L^{-j}$ (a box or a torus, Neumann bonds), $\Lambda$ a union of $L$-blocks in $\Omega^{(j)}$, $e_j\defeq et_j^{1/2}$, $U(z)\defeq\exp(qe_j\xi z)$ for $z\in\C$ ($q$ real antisymmetric, so $U$ is entire, $U(z)^T=U(-z)=U(z)^{-1}$, orthogonal for real $z$). For a real bond field $B$ on $\Omega$ put $X(B)\defeq X_\Lambda(\Omega,B)$ as above and $F(B)\defeq-\frac12\ln\det X(B)$. For a complex bond field $Z$ let $\hat w_\Delta(Z)\defeq\norm qe_j\sup\{|Z_b|:b_-\in\Delta\}$ for a unit cube $\Delta$ and $w(Z)\defeq\sup_\Delta\hat w_\Delta(Z)$; for a real $B$ let $f(B)\defeq\norm qe_j\sup|dB|$ be its curvature amplitude. Hypotheses: (Cur) $f(B)\le c_{\mathrm{cur}}$, with $c_{\mathrm{cur}}>0$ (depending on $d,\norm q,a,L$) such that the defects $2c_p\eps_g+d\kappa^2$ of the cube coercivity (\cref{lem:a06A:cube}) are at most half of $\min\{2\ell^{-2},c_p\}$ for $\ell\in\{1,L\}$; \cref{lem:a06A:cube} uses the field only through $|F|$, so under (Cur) it and \cref{lem:a06A:coer,lem:a06A:CT,lem:a06A:blockinv} and the first bound of \cref{prop:a06A:Ba1P22} hold for $B$ with no other condition, in particular $\gamma_X\le X(B)\le\gamma_1\defeq a+aL^{-2}$ and $|\Delta^{(j)}(\Omega,B;y,y')|\le c_7e^{-\delta_7|y-y'|}$ \cite[Prop.~2.2, (2.27)]{Ba1}, hence $|X(B)(y,y')|\le c_Xe^{-\delta_X|y-y'|}$ with $c_X\defeq c_7+aL^{-2}e^{\delta_7L}$, $\delta_X\defeq\delta_7$. All constants below depend only on $(d,N,\norm q,a,L,M)$ and the constants of \cref{thm:a06A:groupA}, never on $j,k,\eps,\eps_0,\Omega,\Lambda$, the torus or $B$.

\begin{lemma}[the analytic family of block operators]\label{lem:a06d:S1}
In the setting above let $B$ be real with $T_0\defeq T(B)\defeq D_B^TD_B+m_j^2+a_jQ_j(B)^TQ_j(B)\ge\gamma>0$ on $\Omega$, let $A$ be a real bond field, $s_A\defeq\sup_b|A_b|$, $w\defeq\norm qe_js_A$, $S_A$ the set of unit cubes containing an endpoint of a bond with $A_b\ne0$, and for $\tau\in\C$ let $T(\tau)$, $G(\tau)\defeq T(\tau)^{-1}$, $X(\tau)\defeq a_j-a_j^2Q_\tau G(\tau)Q_\tau^T+aL^{-2}P(\tau)$ be the operators in the background $B+\tau A$. Put $m_1\defeq5\sqrt d\gamma^{-1/2}+6a_+d\gamma^{-1}$, $c_\rho\defeq(72m_1dL)^{-1}$, $\theta_W\defeq\theta_0/4$ ($\theta_0$ the Combes--Thomas rate of group (A) at $(\gamma,1)$). For $|\tau|w\le c_\rho$: \textup{(i)} $T(\tau)$ is invertible, $G(\tau)$ analytic in $\tau$, and $\norm{T_0^{-1/2}(T(\tau)-T_0)T_0^{-1/2}}\le m_1|\tau|w\le1/72$; \textup{(ii)} $G(\tau)$ obeys the Combes--Thomas bounds of group (A) (plain, one derivative on either side, two derivatives) with rate $\theta_W$ and constant $c_{\mathrm{CT}}\defeq12\max(1,\gamma^{-1},\gamma^{-1/2})e^{\theta_W}$; \textup{(iii)} $\tilde W(\tau)\defeq X(\tau)-X(0)$ has $|\tilde W(\tau)(y,y')|\le C_W|\tau|w\sum_{\Delta\in S_A}e^{-\theta_W(\dist(y,\Delta)+\dist(\Delta,y'))}$ with $C_W=C_W(d,N,a,a_+,L,\gamma,\theta_0)$; \textup{(iv)} hence $\norm{\tilde W(\tau)}\le C_WS^2|\tau|w$ and $\sum_{y,y'}|\tilde W(\tau)(y,y')|\le C_WS^2|\tau|w|S_A|$, $S\defeq\sup_z\sum_ye^{-\theta_W\dist(y,z)}$. The same holds for every compression $\Lambda X(\tau)\Lambda$ and for $\Delta^{(j)}(\tau)$. The statement holds with $\tau A$ replaced by any complex bond field $Z$ with $w(Z)\le c_\rho$.
\end{lemma}

\begin{proof}
$U(B_b+\tau A_b)=U(B_b)U(\tau A_b)$, $\norm{U(B_b)}=1$ and $\norm{e^Z-1}\le\norm Ze^{\norm Z}$; since $|\tau|wdL\le1/72<\ln2$, $\norm{U(\tau A_b)-1}\le2\xi|\tau|w$ and $\norm{U(\tau A(\Gamma))-1}\le2\,\mathrm{len}(\Gamma)|\tau|w$. Hence $\norm{D_\tau-D_0}\le2\sqrt d|\tau|w$ (each site is the $b_+$ of at most $d$ bonds), $\norm{Q_\tau-Q_0}\le2d|\tau|w$, $\norm{Q_\tau}\le2$, $\norm{Q(\tau)-Q(0)}\le2dL|\tau|w$. (i): $T(\tau)-T_0=(D_\tau-D_0)^TD_\tau+D_0^T(D_\tau-D_0)+a_j[(Q_\tau-Q_0)^TQ_\tau+Q_0^T(Q_\tau-Q_0)]$; with $\norm{D_0T_0^{-1/2}}\le1$ and $\norm{T_0^{-1/2}}\le\gamma^{-1/2}$ the three groups contribute at most $m_1|\tau|w$, so $T(\tau)=T_0^{1/2}(1+K(\tau))T_0^{1/2}$ with $\norm K\le1/72$. (ii): conjugate by $e^{\theta\rho_f}$ as in the Combes--Thomas argument; the $\tau=0$ part contributes $\le1/8$ at $\theta\le\theta_0/4$, the $\tau$-part $V_\theta(\tau)\defeq e^{\theta\rho_f}(T(\tau)-T_0)e^{-\theta\rho_f}$ has every factor connecting points at distance $\le1$, so each norm grows by at most $e^\theta\le e$, and $\norm{e^{\theta\rho}D_\tau e^{-\theta\rho}T_0^{-1/2}}\le3$, giving $\le9m_1|\tau|w\le1/8$; so $\norm{K_\theta(\tau)}\le1/4$ and the three bounds follow with $\norm{T_\theta(\tau)^{-1}}\le\frac43\gamma^{-1}$. (iii): $\tilde W(\tau)=-a_j^2[(Q_\tau-Q_0)G_\tau Q_\tau^T+Q_0(G_\tau-G_0)Q_\tau^T+Q_0G_0(Q_\tau-Q_0)^T]+aL^{-2}[P(\tau)-P(0)]$ with $G_\tau-G_0=-G_\tau(T(\tau)-T_0)G_0$. $(Q_\tau-Q_0)^T\delta_y$ vanishes unless $B(y)\in S_A$, then is supported in $B(y)$ with norm $\le2d|\tau|w$, and $Q_\tau^T\delta_{y'}$ is supported in $B(y')$ with norm $\le2$; by (ii) the entry is $\le4dc_{\mathrm{CT}}e^{\theta_W}|\tau|we^{-\theta_W\dist(B(y),y')}$, and the third term is symmetric. $T(\tau)-T_0=\sum_{\Delta\in S_A}V_\Delta$, $V_\Delta$ collecting the bonds with $b_-\in\Delta$ and the pins of the block $\Delta$; for one $\Delta$, e.g.\ $\langle g,G_\tau(D_\tau-D_0)^T\mathbf 1_\Delta D_\tau G_0g'\rangle=\langle\mathbf 1_\Delta(D_\tau-D_0)G_\tau g,\mathbf 1_\Delta D_\tau G_0g'\rangle$, bounded by $2\sqrt d|\tau|w\norm{\mathbf 1_{\Delta^+}G_\tau g}\norm{\mathbf 1_\Delta D_\tau G_0g'}$ with $\Delta^+$ the union of the $\le3^d$ unit cubes touching $\Delta$, and the localised $L^2$ forms of (ii) bound the two factors by $3^dc_{\mathrm{CT}}e^{\theta_W}e^{-\theta_W\dist(\Delta,B(y))}\norm g$ and $c_{\mathrm{CT}}e^{-\theta_W\dist(\Delta,B(y'))}\norm{g'}$; with $g=Q_0^T\delta_y$, $g'=Q_\tau^T\delta_{y'}$ this is the $\Delta$-term. $P(\tau)-P(0)$ is block-diagonal over $L$-blocks, vanishes on $L$-blocks avoiding $S_A$, with entries $\le5dL|\tau|w\le5dLe^{2\theta_WL}|\tau|we^{-\theta_W(\dist(y,\Delta)+\dist(\Delta,y'))}$ for a $\Delta$ in the same $L$-block. Summing, $C_W\defeq a_+^2[8dc_{\mathrm{CT}}e^{2\theta_W}+4\cdot3^dc_{\mathrm{CT}}^2e^{\theta_W}(4\sqrt d+4a_+d)]+5adL^{-1}e^{2\theta_WL}$. (iv): Schur test and the sum of (iii). For a complex $Z$ the same entry bounds hold since $U$ is entire.
\end{proof}

\begin{lemma}[two kernel tools]\label{lem:a06d:G0}
\textup{(i)} Let $B$ be real with (Cur) and $Z,Z'$ complex with $w(Z),w(Z')\le c_\rho$. Then $|(X(B+Z)-X(B+Z'))(y,y')|\le C'_W\sum_\Delta\hat w_\Delta(Z-Z')e^{-\theta_W(d(y,\Delta)+d(\Delta,y'))}$, $C'_W\defeq3^de^{2\theta_W}C_W$. \textup{(ii)} Put $K_X\defeq c_X\sum_{n\in\Z^3}|n|e^{-\delta_X|n|/2}$, $S''\defeq e^2\sup_y\sum_\Delta e^{-(\theta_W/2)d(y,\Delta)}$, $\theta_2\defeq\min\{\delta_X/2,\gamma_X/(8K_X),\theta_W/2,1\}$ and $c_*\defeq\min\{c_\rho,\gamma_X/(8C'_WS''^2)\}$. For real $B$ with (Cur), complex $Z$ with $w(Z)\le c_*$, every $s\ge0$ and every $1$-Lipschitz $\rho$ on the unit lattice, $\norm{e^{\theta_2\rho}(s+X(B+Z))^{-1}e^{-\theta_2\rho}}\le2/(s+\gamma_X)$, hence $|(s+X(B+Z))^{-1}(y,y')|\le\frac2{s+\gamma_X}e^{-\theta_2|y-y'|}$.
\end{lemma}

\begin{proof}
(i) Assign each bond to the unit cube of its $b_-$, $Z-Z'=\sum_\Delta Y_\Delta$, order the cubes and set $Z_i\defeq Z'+\sum_{i'\le i}Y_{\Delta_{i'}}$; $Z_i$ agrees with $Z$ or $Z'$ bond by bond, so $w(Z_i)\le c_\rho$. Telescope $X(B+Z)-X(B+Z')=\sum_i[X(B+Z_i)-X(B+Z_{i-1})]$; the $i$-th term is $X(B+Z_{i-1}+\tau Y_{\Delta_i})-X(B+Z_{i-1})$ at $\tau=1$, bounded by \cref{lem:a06d:S1}(iii) at the complex base $B+Z_{i-1}+\tau Y_{\Delta_i}$, integrated in $\tau$; only the bonds of $Y_{\Delta_i}$ change, with amplitude $\le\hat w_{\Delta_i}(Z-Z')$, and the cubes containing their endpoints are the $\le3^d$ neighbours of $\Delta_i$, costing $e^{\theta_W}$ each. (ii) Schur test with $|e^{\theta|n|}-1|\le\theta|n|e^{\theta|n|}$: $\norm{e^{\theta\rho}X(B)e^{-\theta\rho}-X(B)}\le\theta K_X\le\gamma_X/8$; by (i) with $Z'=0$ and a Schur test, $\norm{e^{\theta\rho}(X(B+Z)-X(B))e^{-\theta\rho}}\le C'_WS''^2w(Z)\le\gamma_X/8$; $s+X(B)\ge s+\gamma_X$ is self-adjoint, so $e^{\theta\rho}(s+X(B+Z))e^{-\theta\rho}=(s+X(B))^{1/2}(1+K)(s+X(B))^{1/2}$ with $\norm K\le\frac12$. Entries: take $\rho\defeq d(\cdot,y')$.
\end{proof}

% a06_BALtorus_d5b: continuation of a06_BALtorus_d5 (input from sec/a06_BALtorus_c.tex): the local densities of ln Z^(j), the curvature form, Theorem G, and the eps_0-freedom of the per-step constants.

\begin{lemma}[gauge-invariant local densities of $\ln Z^{(j)}$]\label{lem:a06d:G1}
For $B$ real with (Cur) and $Z$ complex with $w(Z)\le c_*$ define, for $y\in\Lambda$,
\[
Q_y(B+Z)\defeq-\tfrac12\int_0^\infty\Tr_N\big[(1+s)^{-1}\delta_{yy}-((s+X(B+Z))^{-1})(y,y)\big]ds.
\]
Then: \textup{(a)} $F(B)=\sum_{y\in\Lambda}Q_y(B)$, and $\sum_yQ_y(B+Z)=-\frac12\ln\det X(B+Z)$ with the branch continuous from $Z=0$; \textup{(b)} (gauge invariance) $Q_y(B+d\lambda)=Q_y(B)$ for every real $\lambda$ on the fine sites of $\Omega$, $(d\lambda)_b\defeq\xi^{-1}(\lambda(b_+)-\lambda(b_-))$; \textup{(c)} (C-parity) $Q_y(-B)=Q_y(B)$; \textup{(d)} (single locality) for $w(Z),w(Z')\le c_*$, $Q_y(B+Z)$ is analytic and $|Q_y(B+Z)-Q_y(B+Z')|\le C_1\sum_\Delta\hat w_\Delta(Z-Z')e^{-\theta_2d(y,\Delta)}$, $C_1\defeq4NC'_WS^2_{\theta_2}/\gamma_X$, $S_\theta\defeq\sup_z\sum_ye^{-\theta|y-z|}$; \textup{(e)} (multi-point locality) for real $u_1,\dots,u_n$, $n\le\bar n+1$, $|D^nQ_y(B)[u_1,\dots,u_n]|\le C_{(n)}\sum_{\Delta_1,\dots,\Delta_n}\prod_i\hat w_{\Delta_i}(u_i)e^{-(\theta_2/n)\sum_id(y,\Delta_i)}$, $C_{(n)}\defeq2^{n-1}n^n3^de^{\theta_2}C_1c_*^{1-n}$; \textup{(f)} $DQ_y(0)=0$, and for real $Y$ with $sY$ under (Cur) for $s\in[0,1]$ and real $u$, $DQ_y(Y)[u]=\int_0^1D^2Q_y(sY)[Y,u]ds$.
\end{lemma}

\begin{proof}
(a) For a positive symmetric matrix, $\ln X=\int_0^\infty[(1+s)^{-1}-(s+X)^{-1}]ds$ (spectral calculus; the integrand is $O(s^{-2})$); take the trace over $\Lambda\otimes\R^N$, which is $\sum_y\Tr_N$ of the diagonal blocks; for complex $Z$ both sides are analytic on $w(Z)\le c_*$ by \cref{lem:a06d:G0}(ii) and agree on real $Z$. (b) Under $B\to B+d\lambda$, $U(B_b)\to g(b_+)U(B_b)g(b_-)^{-1}$ with $g(x)\defeq\exp(qe_j\lambda(x))$ orthogonal, all commuting; hence $D_B\to gD_Bg^{-1}$, $Q_j(B+d\lambda)=g_{\mathrm{unit}}Q_j(B)g^{-1}_{\mathrm{fine}}$ (the contours start at the unit point), $P(B+d\lambda)=g_{\mathrm{unit}}P(B)g^{-1}_{\mathrm{unit}}$, $G_j\to g_{\mathrm{fine}}G_jg^{-1}_{\mathrm{fine}}$, so $X(B+d\lambda)=g_{\mathrm{unit}}X(B)g^{-1}_{\mathrm{unit}}$ with $g_{\mathrm{unit}}$ block-diagonal over $y$, commuting with the compression to $\Lambda$; therefore $((s+X(B+d\lambda))^{-1})(y,y)=g(y)((s+X(B))^{-1})(y,y)g(y)^{-1}$ and $\Tr_N$ is invariant. (c) Write $q=O(\oplus_i\theta_iJ\oplus0)O^T$ and $c\defeq O(\oplus_i\sigma_z\oplus1)O^T$, orthogonal with $cqc^{-1}=-q$; then $U(-B_b)=cU(B_b)c^{-1}$ and $X(-B)=(1\otimes c)X(B)(1\otimes c)^{-1}$ (for $N=2$, $q=J$, $c=\sigma_z$: complex conjugation). (d) $(s+X')^{-1}-(s+X)^{-1}=(s+X')^{-1}(X-X')(s+X)^{-1}$ with $X\defeq X(B+Z)$, $X'\defeq X(B+Z')$; bound the $(y,y)$ block by \cref{lem:a06d:G0}(ii) for both resolvents and (i) for $X-X'$: $N(2/(s+\gamma_X))^2C'_W\sum_\Delta\hat w_\Delta\sum_{z,z'}e^{-\theta_2|y-z|}e^{-\theta_W(d(z,\Delta)+d(\Delta,z'))}e^{-\theta_2|z'-y|}$; since $\theta_2\le\theta_W/2$ and $|y-z|+d(z,\Delta)\ge d(y,\Delta)$ the double sum is $\le S^2_{\theta_2}e^{-\theta_2d(y,\Delta)}$, and $\int_0^\infty4(s+\gamma_X)^{-2}ds=4/\gamma_X$; analyticity by dominated convergence. (e) Decompose each $u_i=\sum_\Delta u_{i,\Delta}$; fix $(\Delta_1,\dots,\Delta_n)$, $\alpha_i\defeq\hat w_{\Delta_i}(u_i)>0$, and put $g(\tau)\defeq\sum_{S\subset\{1..n\}}(-1)^{n-|S|}Q_y(B+\sum_{i\in S}\tau_iu_{i,\Delta_i})$ on $|\tau_i|\le\rho_i\defeq c_*/(n\alpha_i)$, where $w(\sum\tau_iu_{i,\Delta_i})\le c_*$. For each $i$, pairing $S$ with $S\cup\{i\}$ gives differences with $Z-Z'=\tau_iu_{i,\Delta_i}$, each $\le C_13^de^{\theta_2}(c_*/n)e^{-\theta_2d(y,\Delta_i)}$ by (d); so $|g|\le2^{n-1}3^de^{\theta_2}C_1(c_*/n)\min_ie^{-\theta_2d(y,\Delta_i)}\le\dots e^{-(\theta_2/n)\sum_id(y,\Delta_i)}$. The mixed derivative $\partial_1\cdots\partial_ng(0)=D^nQ_y(B)[u_{1,\Delta_1},\dots,u_{n,\Delta_n}]$; Cauchy on the polydisc divides by $\prod\rho_i$; sum over the $\Delta_i$. (f) By (c), $DQ_y(0)=-DQ_y(0)$; then $\frac d{ds}DQ_y(sY)[u]=D^2Q_y(sY)[Y,u]$, each $sY$ satisfying (Cur).
\end{proof}

\begin{lemma}[curvature form of a background change at one density]\label{lem:a06d:G2}
Let $B,A$ be real bond fields on $\Omega$ such that $B+\sigma A$ satisfies (Cur) for $\sigma\in[0,1]$, $y\in\Lambda$, $R\ge2$ with $C_R(y)\defeq\{x:|x-y|_\infty\le R\}\subset\Omega$. Put $g_B\defeq\norm qe_j\sup\{|(dB)(p)|:p\subset C_R(y)\}$, $g_A$ likewise, $w_B\defeq\norm qe_j\sup_\Omega|B|$, $w_A$ likewise. Then
\[
|Q_y(B+A)-Q_y(B)|\le C_2\big[(g_B+g_A)g_A+e^{-\theta_2R/8}(w_B+w_A+Rg_B+Rg_A)(w_A+Rg_A)\big],\qquad C_2\defeq64C_{(2)}S^2_{\theta_2/4}.
\]
\end{lemma}

\begin{proof}
\emph{Local axial gauge.} For $x\in C_R(y)$ let $\gamma_{y,x}$ be the axis-ordered fine path from $y$ to $x$ (inside the box) and $\lambda_B(x)\defeq\sum_{b\in\gamma_{y,x}}\xi B_b$; extend to $\Omega$ by $\lambda_B(x)\defeq\lambda_B(\pi x)$, $\pi$ the coordinatewise clamp onto $C_R(y)$. $\lambda_B$ is linear in $B$; put $\beta\defeq B-d\lambda_B$, $\alpha\defeq A-d\lambda_A$, so $(B+\sigma A)-d\lambda_{B+\sigma A}=\beta+\sigma\alpha$. \emph{Inside bounds} ($b\subset C_R(y)$): $\xi\beta_b$ is the sum of $\xi^2(dB)(p)$ over the $\le|b-y|_1/\xi+1$ plaquettes of a strip inside $C_R(y)$, so $|\beta_b|\le(|b-y|_1+\xi)\sup_{C_R(y)}|dB|$, likewise for $\alpha$. \emph{Outside bounds}: $\pi$ maps the ends of a bond to one point or to the ends of a bond $b'\subset C_R(y)$, so $|(d\lambda_B)(b)|\le|B_{b'}|+|\beta_{b'}|\le\sup|B|+(3R+1)\sup_{C_R}|dB|$ and $|\beta_b|\le2\sup|B|+(3R+1)\sup_{C_R}|dB|$ for every $b$, likewise for $\alpha$. \emph{Gauge invariance and C-parity}: by \cref{lem:a06d:G1}(b), $Q_y(B+A)-Q_y(B)=Q_y(\beta+\alpha)-Q_y(\beta)=\int_0^1DQ_y(\beta+\sigma\alpha)[\alpha]d\sigma$, and by (f) this is $\int\!\!\int D^2Q_y(s(\beta+\sigma\alpha))[\beta+\sigma\alpha,\alpha]\,ds\,d\sigma$, every background real with (Cur). \emph{Estimate}: apply (e) with $n=2$; call $\Delta$ inside if $\Delta\subset C_{R-2}(y)$ ($d(y,\Delta)\le R$), otherwise $d(y,\Delta)\ge R-3$. Inside, $\hat w_\Delta(\beta+\sigma\alpha)\le(3d(y,\Delta)+5)(g_B+g_A)$ and $\hat w_\Delta(\alpha)\le(3d(y,\Delta)+5)g_A$; outside, $\hat w_\Delta(\beta+\sigma\alpha)\le2(w_B+w_A)+(3R+1)(g_B+g_A)$ and $\hat w_\Delta(\alpha)\le2w_A+(3R+1)g_A$. $\sum_\Delta(3d+5)e^{-(\theta_2/2)d}\le8S_{\theta_2/4}$ and $\sum_{\text{outside}}e^{-(\theta_2/2)d}\le S_{\theta_2/4}e^{-\theta_2(R-3)/4}$; expand the product of the two $\Delta$-sums.
\end{proof}

\begin{lemma}[{Ba{\l}aban's configuration near $\supp\theta_{k+1}$; \cite[pp.~567, 575--577]{Ba2}}]\label{lem:a06d:G3}
Let $0\le j<k$ and $N_k\defeq B^{k-1}(\Lambda^{(k-1)}_2)$ (rescaled to scale $j$), where $\theta_k=1$. \textup{(a)} On $N_k$ the configuration \cite[(2.93)]{Ba2} equals $A^{(k)}$, hence $B\defeq\tilde B+A_1=A^{(k)}-\tilde B^{(k)}$ on $N_k$. \textup{(b)} $\dist_j(\supp\theta_{k+1},N^c_k\cup\Omega^c)\ge D_{kj}\defeq L^{k-j}(r(t_k)-c_b-M)$ in scale-$j$ units.
\end{lemma}

\begin{proof}
(a) $\theta_{l+1}=1$ on $B^l(\Lambda^{(l)}_2)$ \cite[p.~567]{Ba2}, and $B^{k-1}(\Lambda^{(k-1)}_2)\subset B^l(\Lambda^{(l)}_2)$ for $l\le k-1$ by $\Lambda^{(i+1)}_2\subset\Lambda^{(i+1)}_0\subset\Lambda^{(i)'}_7\subset\Lambda^{(i)'}_2$ \cite[pp.~566, 570]{Ba2}; so every term $(1-\theta_{l+1})\theta_lA^{(l)}$, $l\le k-1$, of \cite[(2.93)]{Ba2} vanishes on $N_k$, where $\theta_k=1$; by \cite[(2.89)]{Ba2}, $\theta_{k+1}B^{(k+1)}+\theta_{k+1}A'^{(k)}+\tilde B^{(k)}=\theta_{k+1}A^{(k)}$, so the configuration is $(1-\theta_{k+1})A^{(k)}+\theta_{k+1}A^{(k)}=A^{(k)}$. (b) $\supp\theta_{k+1}$ lies in the $M$-neighbourhood of $B^k(\Lambda^{(k)}_2)$; $\Lambda^{(k)}_2\subset\Lambda^{(k-1)'}_7$, which lies five layers of width $\ge r(t_{k-1})-c_b$ (scale $k-1$) inside $\Lambda^{(k-1)'}_2$, i.e.\ $\ge r(t_k)-c_b$ in scale-$k$ units at $L=5$; multiply by $L^{k-j}$.
\end{proof}

The inputs read from \cite{Ba2}: (R1) $|\partial^\eta A^{(k)}|\le C_{Z2}p(t_k)$ on $B^k(\Lambda^{(k-1)'}_2)\supset N_k$ \cite[(2.60)]{Ba2}; (R2) $|\tilde B^{(k)}|\le C_\kappa t_k^\kappa$ and $|\partial^\eta\tilde B^{(k)}|\le C'_\kappa t_k^\kappa$ for every $\kappa$ \cite[(2.90)]{Ba2}; (R3) (Cur) at scale $j$ on $\Omega$ for $\tilde B+A_1+\sigma A_2$, $\sigma\in[0,1]$: by \cite[(2.94), (2.98)]{Ba2} the curvature amplitude is $\le4\norm qC_{98}et_j^{1/2}p(t_j)$, which is $\le c_{\mathrm{cur}}$ under the cap $\norm qC_{98}et^{1/2}p(t)\le c_{\mathrm{cur}}/4$ for $t\le\eps_0$; (R4) $w_B\le C_wt_j^{-1/2}p(t_j)$ on $\Omega$, from the restrictions \cite[(2.55)]{Ba2} at each scale $l\in(j,k+1]$ with \cite[(2.59)]{Ba2} and the rescaling $e_jA^{(l),(j)}=L^{-(l-j)}e_lA^{(l)}$.

\begin{theorem}[the terms containing $\tilde B^{(k)}$ are summable over $j<k$]\label{thm:a06d:G}
Let $1\le k\le K-1$ and $j<k$. Then for every $\kappa$,
\[
|D_j|\le C_G(\kappa)(1+(k-j))^{4+\bar p}L^{-(k-j)}t_k^{1+\kappa}p(t_k)|\Lambda^{(k)}_2|,\quad\text{hence}\quad\sum_{j<k}|D_j|\le C'_G(\kappa)t_k^\kappa|\Lambda^{(k)}_2|\ \text{ and }\ |D_j|\le C''_G(\kappa,\beta)t_k^\kappa|\Lambda^{(k)}_2|L^{(j-k)(1-\beta)}
\]
for every $\beta>0$, with $C_G,C'_G,C''_G$ depending only on $(d,N,\norm q,a,L,M,R,\varrho,\bar p,b_0,e,\lmu,c_1,C_{Z2},C_\kappa,C'_\kappa,C_{98},C_w)$ and the constants of \cref{lem:a06d:G1}, free of $j,k,\eps,\eps_0$, the torus and the history.
\end{theorem}

\begin{proof}
Fix $j<k$, $A\defeq A_2$, $B\defeq\tilde B+A_1$ at scale $j$, and $R\defeq R_{kj}\defeq4+(32/\theta_2)(k-j)\ln L$; require $R\ge R_*(\theta_2,L,c_b,M)$, fixed before $\eps_0$, so that $2R_{kj}+3\le D_{kj}$ for all $k-j\ge1$ (possible since $D_{kj}\ge L^{k-j}R/2$ grows faster than $R_{kj}$). \emph{Count}: $\supp A\subset\supp\theta_{k+1}$, and the set $S_A$ of scale-$j$ unit cubes meeting it has $|S_A|\le(2M+1)^3L^{3(k-j)}|\Lambda^{(k)}_2|$. \emph{Class I}, $y$ with $d(y,S_A)\le R$: by \cref{lem:a06d:G3}(b), $C_R(y)\subset\Omega\cap N_k$, and by (a) $B=A^{(k)}-\tilde B^{(k)}$ there, so by (R1), (R2) and the rescaling, $g_B\le2\norm qeL^{-2(k-j)}t_k^{1/2}(C_{Z2}p(t_k)+C'_\kappa t_k^\kappa)$, $g_A\le2\norm qeL^{-2(k-j)}t_k^{1/2}C'_\kappa t_k^\kappa$, $w_A\le\norm qeL^{-(k-j)}t_k^{1/2}C_\kappa t_k^\kappa$, and $w_B\le C_wt_j^{-1/2}p(t_j)=C_wL^{(k-j)/2}t_k^{-1/2}p(t_j)$ with $p(t_j)\le(1+(k-j)\ln L)^{\bar p}p(t_k)$ (R4). \cref{lem:a06d:G2} (hypotheses by (R3)): the main term is $\le CL^{-4(k-j)}t_k^{1+\kappa}p(t_k)$; the other is $\le Ce^{-\theta_2R/8}R^2L^{(k-j)/2}L^{-(k-j)}(1+(k-j)\ln L)^{\bar p}t_k^\kappa p(t_k)\le CR^2(1+(k-j)\ln L)^{\bar p}L^{-4.5(k-j)}t_k^\kappa p(t_k)$, since $e^{-\theta_2R/8}\le L^{-4(k-j)}$. The number of class-I $y$ is $\le(2R+3)^3|S_A|$. \emph{Class II}, $d(y,S_A)>R$: by \cref{lem:a06d:G1}(d) with $Z=A$, $Z'=0$ ($w(A)\le c_*$ for $t_k\le\eps_0$ by (R2), a cap), $|Q_y(B+A)-Q_y(B)|\le C_1w_A\sum_{\Delta\in S_A}e^{-\theta_2d(y,\Delta)}$, which summed over class II is $\le C_1w_A|S_A|S_{\theta_2/2}e^{-\theta_2R/2}\le CL^{3(k-j)}L^{-(k-j)}L^{-16(k-j)}t_k^{1/2+\kappa}(2M+1)^3|\Lambda^{(k)}_2|$. \emph{Sum}: by \cref{lem:a06d:G1}(a), $D_j=\sum_{y\in\Lambda}[Q_y(B+A)-Q_y(B)]$; class I gives $(2R+3)^3(2M+1)^3L^{3(k-j)}|\Lambda^{(k)}_2|\times[CL^{-4(k-j)}t_k^{1+\kappa}p+CR^2(1+(k-j)\ln L)^{\bar p}L^{-4.5(k-j)}t_k^\kappa p]$, whose first part is $\le C(1+(k-j))^3L^{-(k-j)}t_k^{1+\kappa}p|\Lambda^{(k)}_2|$ and whose second part is smaller by $L^{-1.5(k-j)}$ up to polynomials at the cost of one power of $t_k$ (take $\kappa$ one larger in (R2)); class II is smaller still. \emph{Sums over $j$}: $\sum_{m\ge1}(1+m)^{4+\bar p}L^{-m}<\infty$ and $t^{1+\kappa}p(t)\le P(1,1)t^\kappa$; for the per-$j$ form absorb $(1+m)^{4+\bar p}$ into $L^{\beta m}$.
\end{proof}

Per scale-$j$ cube the bound is $\sim L^{-4(k-j)}t_k^{1+\kappa}p=t_j^4t_k^{\kappa-3}p$, i.e.\ $O(t_j^{4-\beta})$ for $\kappa\ge3$: the printed claim ``$O((L^j\eps)^\kappa)|B^{k-j}(\Lambda^{(k)}_2)|$ with $\kappa>d$, $\kappa_0>0$'' holds with $\kappa=4-\beta$, $\kappa_0=1-\beta$, the double-curvature count.

\begin{corollary}[the $Z^{(j)}$ expansions, consumed form]\label{cor:a06d:G}
For $1\le k\le K-1$, $\sum_{j<k}|F_j(\tilde B+A_1+A_2)-F_j(\tilde B)-P_{\tilde B}[A_1]|\le(C_\kappa/(L^{\bar n+1-d}-1)+C'_G(\kappa))t_k^\kappa|\Lambda^{(k)}_2|$ for every $\kappa<(\bar n+1)/2$, with per-$j$ form $O(t_k^\kappa)|\Lambda^{(k)}_2|L^{(j-k)\kappa_0}$, $\kappa_0=\min(\bar n+1-d,1-\beta)>0$. This covers the three categories of the sentence of \cite[p.~578]{Ba2}: the order-$(\bar n+1)$ trace term, the R-vertices and the orders $>\bar n$ ($R_j$), and the terms containing $\tilde B^{(k)}$ ($D_j$); \cite[(2.104)]{Ba2} is unchanged.
\end{corollary}

Every gauge step above is the exact invariance of a full object, $Q_y$ or $F_j$, under a gauge transformation of the background on all of $\Omega$; the local axial gauge of \cref{lem:a06d:G2} is used only to bound fields, never to split an identity, and no Ward identity on a photon line enters.

\subsubsection*{The per-step large-field constants are free of $\eps_0$}

The order of constants in \cite[Sect.~3]{Ba2} requires that each per-step constant of the large-field chain be fixed before $\eps_0$, so that the final caps on $\eps_0$ (\cref{lem:s14:CW}) are consistent. A constant can depend on $\eps_0$ only through a cap ``for $e(t)$ sufficiently small'' or through a bound stated at the final scale; in the first case the constant is the sup over $t\in(0,1]$ of an explicit function, finite by $E(\kappa,\delta)$, $P(a,m)<\infty$, in the second case it is replaced by the scale-local bound. The constants of the chain are: $c_{3.29}$ of \cite[(3.29)]{Ba2}, from \cite[Prop.~3.1$'$, Sect.~4]{Ba4}, whose remainder $O(1)e^2p^2(e)(|\phi(x)|^2+|\phi(x')|^2)$ in \cite[(4.7)]{Ba4} is built from \cite[Lemma~2.1, (4.10), (4.13), (4.22)]{Ba4} with $\gamma''_0$ depending ``on $d$ and $a$ only'', so that $c_{3.29}(k)=C_{\mathrm{Ba4}}e^{2-\alpha}C_\phi^2\lambda^{-1/2}t_k^{(1-\alpha)/2}p_k^2$ is scale-local; $C_{2.57}$ of \cite[(2.57)]{Ba2}, which depends on $(d,L,M,\alpha_0,\bar n,a,c_1,b_0,\bar p,R,\varrho,e,\lambda,N,\beta)$ and on the $O(1)$ of Proposition~1, with $\eps_0$ entering only through the caps $p(t)^n\lambda(t)^{-m}e^{-\delta r(t)}\le t^\kappa$ (Ba{\l}aban's order of choices, \cite[p.~421]{Ba3}); the counterterm constants, $|\dm|\le c_m\eps^{-1}$ (per-graph power counting, no $\eps_0$) and the extensive bounds on the vacuum energy and on $\dm_{K,k}$ (Step~2 and \cref{thm:a06d:V12,thm:a06d:V14}), whose constants are built from renormalised power counting; and the $Z$-factor constant $C'_G$ of \cref{thm:a06d:G}, whose caps are preliminary. So every per-step constant is free of $\eps_0$, and the order statement of \cref{thm:a06:BALA} holds.

% a06_BALtorus_d6: part of Step 3 of the proof of thm:a06:T1 (input from sec/a06_BALtorus_c.tex).
% Group (D): Ba3 Prop. 1 at L = 5, first half: localisation conventions and the regrouping after separated lines are cut (p. 432).

\subsubsection*{Step 3(D$_3$): \cite[Prop.~1]{Ba3} at $L=5$, the regrouping after the cut}

The proof of \cite[Prop.~1, (1.33)]{Ba3} on \cite[pp.~420--445]{Ba3} consists of twenty sketched steps and seven further sketched places. Most of them are direct computations on the printed definitions (the slices \cite[(2.6), (2.10)--(2.12)]{Ba3}, the summation by parts \cite[(2.8)]{Ba3}, the induction \cite[(2.16)]{Ba3}, the generalised graphs of \cite[Prop.~2.2]{Ba3}, the estimate \cite[(2.17)]{Ba3}, the Ward--Takahashi identities of \cite[pp.~430--431]{Ba3} for full lines, the estimate \cite[(3.1)]{Ba3}, the cube steps of \cite[p.~433]{Ba3}, the norms \cite[(1.32)]{Ba3} at $B_0$ and at $\tilde B$), whose written form is a transcription with the $L=5$ constants of \cref{thm:a06:BALA}. Two steps need more than a transcription: the regrouping into renormalised classes after the separated lines are cut (\cite[p.~432]{Ba3}: ``if we remove a line in any graph of the class, then for the obtained graph we can find in this class all graphs necessary to form a new renormalized class''), and the renormalisation of the divergent classes (\cite[pp.~434--445]{Ba3}). They are proved here and in Step~3(D$_4$); the conventions are Ba{\l}aban's, made explicit.

\paragraph{Conventions.} (CV) Degrees are computed in Ba{\l}aban's count: each of \cite[(1.13), (1.14), (1.15)]{Ba3} is a vertex with its own degree (\cite[p.~423 (v), p.~414]{Ba3}); the two factors of one product \cite[(1.18)]{Ba3} share the block point $y$ and are joined in the bound function by a contact, a top-scale line of dimension $0$ whose kernel is the indicator that both summation points lie in $B^k(y)$. (Loc) Label boxes are closed grid boxes of side $\le2$ on the block-corner grid (\cite[p.~420]{Ba3}: ``for each leg of the vector field a smooth partition of unity \dots\ a support of each function is contained in a cube with sides of length 2''); a label is attached to each leg of a vector field (the support box of its partition function $H_y\subset y+[-1,1]^d$), to each \cite[(1.6)]{Ba3}/\cite[(1.7)]{Ba3} vertex (its unit cube) and to each block factor ($B^k(y)$), and every vertex lies in the intersection of the boxes of its labels. The number of label boxes of a graph is $V_{\mathrm{lab}}\le(2\bar n+1)V_{\mathrm{Bal}}$, $V_{\mathrm{Bal}}\le2\bar n(4N_A+5)$, constants of $\bar n$ only; put $l_N\defeq2V_{\mathrm{lab}}(\bar n)$. (Agr) The counterterm graph $T_\gamma G$ of a two-scalar-leg subgraph $\gamma$ is $G$ with $\gamma$'s two external legs attached at one point $x$ \cite[p.~417]{Ba3}, $x$ the leg point of the anchor vertex $v_a$ of $\gamma$; $\gamma$'s internal vertices keep their labels (\cite[p.~420]{Ba3}, the localisation is performed ``for the terms of $\delta m^2_k(\dots)$, $E_k(\dots)$ also''), so $G$ and $T_\gamma G$ have the same label multiset (\cref{conv:a06d:LOC}(c)). (R-W$'$) Two graphs are ``in agreement'' \cite[p.~432]{Ba3} if they carry the same label multiset; a \emph{per-label class} of a given type and order is the set of all graphs of that type and order with a given label multiset, equivalently (with the partition functions written as multipliers $t_yH_y$ on the vector legs and $t_Q\mathbf 1_Q$ on the \cite[(1.6)]{Ba3}/\cite[(1.7)]{Ba3} vertices) the coefficient of one monomial $\prod t$ in the expansion of \cite[(1.5)]{Ba3}; this is how \cite[p.~444]{Ba3} evaluates ``a product of the values of all the localization functions at the point $x'$'' and then takes the graph ``with the summations unrestricted''. (H$_A$) For an external field with a leg that can sit at several vertices of a class (the counterterm merge moves it to the anchor, leg exchanges move it too), the localisation $h$ of \cite[p.~420]{Ba3} (``$h=1$ on box$(v)$ and $h=0$ outside some neighborhood of box$(v)$'') is taken $=1$ on the label boxes of every vertex at which the leg sits in some term of the class and $0$ outside their $\frac14$-neighbourhood; for moves inside a cluster this is a neighbourhood of size $\le l_N$, and for moves across clusters the extra box is paid by \cref{lem:a06d:S25}(a). It costs nothing to the consumers, which sum over localisations \cite[p.~421]{Ba3} or use fields localised at one point.

\begin{definition}[clusters; the cluster cut rule]\label{def:a06d:CCR}
Let $\tau$ be a term of a per-label class and $\mathcal B$ its multiset of label boxes, common to the class. The \emph{clusters} are the connected components of the graph on $\mathcal B$ in which two boxes are joined when they intersect; $U_K$ is the union of the boxes of the cluster $K$. Each end of a line of $\tau$ lies at a vertex, all of whose labels contain it, so they intersect and lie in one cluster $K(\mathrm{end})$ (for a counterterm vertex, the cluster of $v_a$). The \emph{cut set} $S(\tau)$ is the set of lines whose two ends lie in different clusters. A cut line becomes the external field $h_KGh_{K'}$ (or its vector analogue), with $h_K$ smooth, $h_K=1$ on the $\eta$-neighbourhood of $U_K$, $\supp h_K$ in the $\frac14$-neighbourhood of $U_K$ and $|\partial^\eta h_K|,\mathrm{Lip}(\partial^\eta h_K)\le c(d)$ (e.g.\ $h_K=1-\prod_Q(1-\varphi_Q)$ over the unit grid cubes $Q$ of $U_K$, $\varphi_Q$ a fixed bump). This is Ba{\l}aban's threshold $1$ \cite[p.~432]{Ba3}, applied to clusters instead of single lines.
\end{definition}

\begin{lemma}[properties of the cut rule]\label{lem:a06d:S22}
\textup{(a)} If $K\ne K'$ then $\dist_\infty(U_K,U_{K'})\ge1$, and points of $\supp h_K$ and $\supp h_{K'}$ are at distance $\ge\frac12$. \textup{(b)} $\operatorname{diam}_\infty U_K\le2|K|\le l_N$. \textup{(c)} Every kept line joins two vertices of one cluster, so every kept-connected component lies in one $U_K$. \textup{(d)} Let $\gamma$ be a subgraph all of whose label boxes lie in one cluster $K$; any operation that changes only the vertices of $\gamma$ at which lines are attached (the counterterm merge of \cite[p.~417]{Ba3}, the exchange of two external legs of $\gamma$), or replaces $\gamma$ by another graph with the same labels, leaves $S$ unchanged. \textup{(e)} A connected subgraph whose vertices lie in at least two clusters contains a cut line and is disconnected by $S$, each piece lying in one cluster. \textup{(f)} A subgraph whose vertices lie in one cluster has no cut line.
\end{lemma}

\begin{proof}
(a) Two grid boxes either intersect or are at sup-distance $\ge1$; boxes of different clusters do not intersect; the collars add $\le\frac14$ on each side. (b) Two boxes of $K$ are joined by a chain of at most $|K|-1$ intersecting boxes of side $\le2$. (c) is the definition. (d) The moved ends stay at vertices of $\gamma$, whose labels are in $K$, and the other ends do not move; a replacement with the same label multiset does not change $\mathcal B$, hence not the clusters. (e) Kept lines are intra-cluster, so a kept path cannot leave a cluster. (f) is immediate.
\end{proof}

Let $G$ be a labelled graph and $\mathcal C=R(G)$ its class: the forest sum over $\mathcal D$, the partnered subgraphs of $G$ (the one-particle-irreducible two-scalar-leg subgraphs with mass counterterms, \cite[p.~430]{Ba3}, and the vacuum components with $E_1$, \cite[p.~432]{Ba3}), closed under leg exchange of $2\phi$ and $\phi\phi A$ subgraphs and under per-label completion of $2A$ and $\phi\phi A$ subgraphs. Split $\mathcal D=\mathcal D_0+\mathcal D_1$, $\mathcal D_0\defeq\{\gamma:\text{all label boxes of }\gamma\text{ in one cluster}\}$. For a forest $F=F_1+F_0$ ($F_1\subset\mathcal D_1$, $F_0\subset\mathcal D_0(F_1)$, the elements of $\mathcal D_0$ nested in or disjoint from every element of $F_1$), the term $\tau_F\defeq\prod_{\gamma\in F}(-T_\gamma)G$ is cut by its own cut rule.

\begin{theorem}[the regrouping after the cut]\label{thm:a06d:S23}
\textup{(i)} For fixed $F_1$ the cut set $S(\tau_{F_1+F_0})$ is the same for every $F_0$; call it $S(F_1)$. \textup{(ii)} $\sum_FE(\tau_F)=\sum_{F_1}R_{\mathcal D_0(F_1)}[E(\mathrm{Cut}_{S(F_1)}Y(F_1))]$, $Y(F_1)\defeq\prod_{\gamma\in F_1}(-T_\gamma)G$, each $F_0$-operation acting at its place (inside a counterterm coefficient when nested in an element of $F_1$). \textup{(iii)} Every element of $\mathcal D_1$ is disconnected by $S(\emptyset)$. \textup{(iv)} Every connected divergent subgraph of $\mathrm{Cut}\,Y(F_1)$ (kept lines only) is an element of $\mathcal D_0(F_1)$, renormalised by $R_{\mathcal D_0(F_1)}$; or a $2A$ or $\phi\phi A$ subgraph inside one cluster whose per-label partners are terms of the same block (\cref{thm:a06d:S28}); or a \cite[(2.4)]{Ba3} subgraph (no partner; treated by the summation by parts \cite[(2.8)]{Ba3}); or a subgraph containing a counterterm vertex of $F_1$, which by \cref{lem:a06d:S25} is convergent unless it is a vacuum graph or has one external vector leg (Ba{\l}aban's exception list, \cite[p.~432]{Ba3}). So each block is a renormalised class of its cut graph (\RENorb, \cref{hyp:k01:RENorb}, is the same regrouping with the slice decomposition in place of the cut).
\end{theorem}

\begin{proof}
(i) The $F_0$-operations move ends inside clusters (\cref{lem:a06d:S22}(d)). An element $\gamma_0\in\mathcal D_0(F_1)$ contains no line whose end is moved by an element $\gamma\in F_1$: such a line is external to $\gamma$ and attached at a vertex of $\gamma$, so $\gamma_0$ would share a vertex with $\gamma$, hence be nested with it by compatibility; $\gamma_0$ inside $\gamma$ contains no external line of $\gamma$, and $\gamma_0$ containing $\gamma$ would contain $\gamma$'s cut line, against $\gamma_0\in\mathcal D_0$. For vacuum elements of $F_1$ nothing moves: the vacuum graph and its counterterm graph (same graph, same labels, \cref{conv:a06d:LOC}(b)) are cut identically. (ii) Every forest of $\mathcal D$ splits uniquely as $F_1+F_0$ with $F_1$ a forest of $\mathcal D_1$ and $F_0$ a forest of $\mathcal D_0(F_1)$, and conversely; by (i) the cut is the same for all $F_0$ at fixed $F_1$, and $\mathrm{Cut}$ commutes with the $F_0$-operations because they act inside clusters and the elements of $\mathcal D_0$ have no cut line (\cref{lem:a06d:S22}(f)). (iii) \cref{lem:a06d:S22}(e). (iv) A divergent connected kept subgraph $H$ of $\mathrm{Cut}\,Y(F_1)$ without a counterterm vertex of $F_1$ lies in one cluster (\cref{lem:a06d:S22}(c)) and avoids the internal vertices of the elements of $F_1$ (these sit inside the counterterm coefficients). If $H$ is $2\phi$ it belongs to $\mathcal D$, lies in one cluster and is disjoint from $F_1$, so $H\in\mathcal D_0(F_1)$; if $H$ is $2A$ or $\phi\phi A$, \cref{thm:a06d:S28}; \cite[(2.4)]{Ba3} needs no partner; subgraphs through a counterterm vertex of $F_1$ are \cref{lem:a06d:S25}.
\end{proof}

\begin{corollary}[the cut kernels]\label{cor:a06d:S24}
In every term of every block, every cut line joins two different clusters, so its kernel is evaluated on $\supp h_K\times\supp h_{K'}$, at distance $\ge\frac12$. The estimate \cite[(3.1)]{Ba3} holds with $h_K,h_{K'}$ in place of $h,h'$: $\norm{h_KG_k(\Omega,\tilde B)h_{K'}}_{1,\alpha}\le Ce^{-(\delta_1/2)(\dist(U_K,U_{K'})-1/2)}$, its proof using only the slice bounds at the points of the supports, the separation, and $|\partial h|,\mathrm{Lip}(\partial h)\le c_h$; the points used lie within $\frac14+\eta$ of vertex label boxes, at the depth the hypotheses supply ($R\ge R_0+5$). The tree factor of \cite[(1.33)]{Ba3} satisfies $d(\{\square(v)\})\le\sum_{\text{spanning tree of cut kernels}}\dist(U_K,U_{K'})+\sum_K\operatorname{diam}U_K$, and the second sum is $\le l_N^2$, a constant.
\end{corollary}

With the printed per-line threshold $\dist(\square(v),\square(v'))\ge1$ and the counterterm leg move of \cite[p.~417]{Ba3}, a line from the moved leg of $\gamma$ to a vertex $w$ could be cut in $G$ and short in $T_\gamma G$, or conversely, and then either the cut kernel is not $O(1)$ or the cut does not commute with $T_\gamma$; the cluster rule removes this (the example of three boxes $[0,1],[1,2],[2,3]$ on one axis is one cluster, and the line is kept in $G$ and in $T_\gamma G$), with Ba{\l}aban's threshold $1$ and no larger cut threshold.

\begin{lemma}[disconnected counterterms are finite vertices]\label{lem:a06d:S25}
Let $\gamma\in\mathcal D_1$ (so disconnected by $S$), with anchor $v_a$, and let $c(x)$ be the coefficient of $T_\gamma$ in $Y(F_1)$ with $\gamma$'s labels fixed: $c(x)=\sum_{x'}$ of the amputated kernel of $\gamma$, renormalised by the inner forests, with its lines of $S(F_1)$ replaced by cut kernels. Assume \cite[(1.33)]{Ba3} for all classes with fewer vertices than $G$ (the induction on the number of vertices that the classes already use, \cite[p.~431]{Ba3}: ``our definition of the classes $G_{\mathrm{ren}}$ is in fact inductive''). Then \textup{(a)} $|c(x)|\le Ce_k^{d_v(\gamma)}\lambda_k^{d_s(\gamma)}\exp(-\delta'd(\text{labels of }\gamma))$ uniformly in $x,k,\eps,\eta,\Omega,\Omega_1,\Omega_2,\tilde B$ and the torus; \textup{(b)} $T_\gamma$ is therefore a \cite[(1.7)]{Ba3} vertex with a finite counterterm, of degree $D(v)=2$ (\cite[p.~423 (ii)]{Ba3}); every connected subgraph containing it has $D\ge1$ if $E\ge2$ and $D\ge0$ if $E=1$, with equality only for the tadpole graph, and vacuum graphs are treated by the vacuum rule of \cite[p.~432]{Ba3}; \textup{(c)} $c$ is an even polynomial in $q$, so the tadpole with this vertex vanishes after the gauge-out of $B_0$ \cite[p.~433]{Ba3}.
\end{lemma}

\begin{proof}
(a) The pieces of $\gamma|_{S(F_1)}$ lie in single clusters, each with fewer vertices than $G$ since $\gamma$ has at least two pieces; each piece with its inner renormalisation is a class of lower order whose external fields are the cut kernels, two-variable external fields in the sense of \cite[p.~419]{Ba3}; \cite[(1.33)]{Ba3} for each piece, the product over pieces, and \cref{cor:a06d:S24} for the cut kernels give (a), the decay coming from the cut kernels between the clusters of $\gamma$. (b) A \cite[(1.7)]{Ba3} vertex with a bounded coefficient has $D(v)=d-(d-2)=2$, and the degree count of \cite[(2.1)--(2.3)]{Ba3} applies with $\kappa=2$ at that vertex. (c) $\gamma$ has no external vector leg; every closed scalar index loop of $\gamma$ carries an even number of vector legs or vanishes ($\Tr q^{\mathrm{odd}}=0$), and the scalar path between $\gamma$'s two legs then carries an even number as well, since the total number of vector-leg ends of $\gamma$ is twice its number of photon lines.
\end{proof}

\begin{lemma}[per-label gauge invariance]\label{lem:a06d:S26}
Let $Y$ be a periodic box or $\xi\Z^d$, $X$ any bond field built from $A'$ and external fields (e.g.\ $X=\sum_yt_yH_yegA'+\tilde A$-type legs with their multipliers), $F$ any polynomial in the $|\phi(x)|^2$ with $x$-dependent coefficients (e.g.\ $\sum_Qt_Q\lambda\mathbf 1_Q|\phi|^4$, the \cite[(1.7)]{Ba3} vertices with any coefficients, the restriction to $\Omega_1$), and $Z(A;t)\defeq\int d\mu(A')\int d\phi\exp[-\frac12\langle\phi,(-\Delta_{A+X}+m^2)\phi\rangle-F(\phi)]$ with any Gaussian measure $d\mu(A')$. Then $Z(A+d\lambda;t)=Z(A;t)$ for every periodic (on $\xi\Z^d$: finitely supported) $\lambda$ and every $t$. Consequently every $t$-monomial coefficient of every $A$-Taylor coefficient of $\log Z$ satisfies the Ward--Takahashi identities of \cite[pp.~430--431]{Ba3}; in particular the two-vector-leg kernel $\Pi^{\mathrm{lab}}$ of a per-label class obeys $\sum_{y'}\Pi^{\mathrm{lab}}_{\mu\nu}(y,y')=0$.
\end{lemma}

\begin{proof}
$\phi=V_\lambda\psi$: the covariant Laplacian is covariant for the total field, $-\Delta_{A+d\lambda+X}=V_\lambda(-\Delta_{A+X})V_\lambda^*$ for every $X$, because only the total field enters $U$; $F$ is invariant, being a function of the $|\phi(x)|^2$; the $A'$-integral is untouched since the transformation acts on $A$ only. The identities follow coefficient by coefficient in $t$, and by (R-W$'$) the per-label class is one $t$-monomial.
\end{proof}

The identity is for complete per-label classes with full propagator lines, before any scale decomposition; it is never applied to a single ordering or scale assignment. In Step~3(D$_4$) it is applied only to the all-$C^\xi$ part of a regrouped local coefficient (full lines $C^\xi$ on $\xi\Z^3$), the remainders being bounded explicitly.

\begin{lemma}[{multipliers to the anchor; \cite[p.~444, (3.36)]{Ba3}}]\label{lem:a06d:S27}
For a node with anchor $\hat x$ and a smooth multiplier $H$ (partition function, $g_k$) at a vertex $x_i$ of the node, $H(x_i)=H(\hat x)+[H(x_i)-H(\hat x)]$ with $|H(x_i)-H(\hat x)|\le c|x_i-\hat x|$; the second term carries a factor $|x_i-\hat x|$ on a path inside the node, so its degree gains $1$. After all smooth multipliers of the node (including those of its external legs) are moved to $\hat x$, the remaining sums over the node's internal positions are unrestricted, except for sharp indicators (\cref{lem:a06d:S46}).
\end{lemma}

\begin{theorem}[the Ward--Takahashi and $\phi\phi A$ classes at fixed localisation]\label{thm:a06d:S28}
Let $H$ be a $2A$ or $\phi\phi A$ subgraph of a term of a per-label class. \textup{(a)} If the label boxes of $H$ lie in one cluster, every per-label partner $\tilde H$ (same type, order, label multiset) lies entirely in that cluster and has no cut line, and so does every leg-exchange partner (\cref{lem:a06d:S22}(d), (f)); the per-label completion is therefore a sum of terms of one block, and at the renormalisation step the local coefficient of the completed node is the per-label class sum, to which \cref{lem:a06d:S26} applies ($2A$ zeroth moment, after \cref{lem:a06d:S27} has moved the multipliers of $H$'s external legs to the anchor), as do \cref{lem:a06d:S27,lem:a06d:S46} with the reflection identities of \cite[(3.30)]{Ba3} ($\phi\phi A$ zeroth moment, $2A$ first moment). \textup{(b)} If they span at least two clusters, every partner is disconnected by $S$ (\cref{lem:a06d:S22}(e)) and no identity is needed.
\end{theorem}

\begin{proof}
\cref{lem:a06d:S22}.
\end{proof}

At fixed vertex localisations of all vertices of $H$ the zeroth-moment sum would be restricted to boxes and would not vanish (boundary layers of width $L^{j_0}\eta$); Ba{\l}aban localises legs, not vertices \cite[p.~420]{Ba3}, and moves the localisation functions to one point before the identity is used \cite[p.~444]{Ba3}, and with (R-W$'$) the identity holds label by label.

Consequences for the neighbouring steps: under the cluster rule the kept components lie in clusters (\cref{lem:a06d:S22}(c)), and the class-component cube of \cite[p.~432]{Ba3} has side $\le l_N$, common to the class; kept lines may join boxes at distance up to $l_N$, so \cite[(3.2)]{Ba3} for components carries the constant $e^{\delta l_N}$. The cut on the cube after the $\tilde B'$-expansion \cite[p.~433]{Ba3} is the cluster rule on the label boxes of the expanded class; every component with an old vertex lies in the cluster of an old vertex, of diameter $\le2V^\sharp_{\mathrm{lab}}$ ($V^\sharp_{\mathrm{lab}}$ the label count after the expansion), so the side of the localisation cube of a connected renormalised class after the expansion is $l^\sharp\defeq l_N+4V^\sharp_{\mathrm{lab}}(\bar n)$, which fixes $r^\flat$ in \cref{thm:a06:BALA}(a) through $r^\flat\ge3l^\sharp+R_0+12$. Chain components lie in one cluster, so their exponential factor is $\ge e^{-2\delta V^\sharp_{\mathrm{lab}}}$.

% a06_BALtorus_d7: part of Step 3 of the proof of thm:a06:T1 (input from sec/a06_BALtorus_c.tex).
% Group (D): Ba3 Prop. 1 at L = 5, second half: the renormalisation of the divergent classes (pp. 434-445) and the assembled statement; the one open row.

\subsubsection*{Step 3(D$_4$): the renormalisation of the divergent classes, \cite[pp.~434--445]{Ba3}}

After the cube steps of \cite[p.~433]{Ba3} a class component has scalar lines $G_k(0)$ on $\eta\Z^d$, sliced by \cite[(2.6)]{Ba3}, photon lines on the torus, vertices in the cube's inner part, all other objects external, and per-label full classes. Orderings are the lexicographic ones of the slice indices; for an ordering $\tilde l$ the components $G^{(\alpha)}_i$ along it form a laminar family. A \emph{dangerous node} is a component with $D\le0$; by the degree formula \cite[(2.1)--(2.3)]{Ba3} it is one of: a \cite[(2.4)]{Ba3} subgraph ($D=0$); $2\phi$ ($D\in\{-1,0\}$, $D=-1$ only at one loop, \cite[(3.6), (3.8)]{Ba3}); $2A$ ($D\in\{-1,0\}$, $D=-1$ only at one loop, \cite[(3.25)]{Ba3}); $\phi\phi A$ ($D=0$, one loop: $D\ge V/2+E/2-3$ with $E=3$ forces $V\le3$ and one loop); a vacuum component, cancelled by the vacuum counterterms \cite[(1.24)]{Ba3} with the same labels and lines, up to the torus-versus-$\eta\Z^d$ image terms. Subgraphs with no external scalar leg and an odd number of external vector legs vanish as factors at $B=0$ (their scalar index loops close inside them, $\Tr q^{\mathrm{odd}}=0$). A one-particle-reducible $2\phi$ node is not partnered (\cite[p.~430]{Ba3}: ``We consider one-particle-irreducible graphs only''); since $D(P_1+\text{line}+P_2)=D(P_1)+D(P_2)+2$, a dangerous reducible $2\phi$ chain consists of one-loop $D=-1$ parts, and after these are processed the chain is no longer dangerous. External vector fields inside vertices ($\tilde A$, and $\tilde B'$ after the cube step) count as external vector legs for the type of a node (an $\tilde A$-leg adds $\kappa=n'=1$ where an external leg adds $1$ to $E$), so the partnered $2\phi$ nodes are exactly the $2\phi$ nodes without such legs. A node $S$ is \emph{useful} along $\tilde l$ if its internal lines precede all its external lines, i.e.\ if it is a component. \emph{Processing} a useful dangerous node $S$ with an anchor $\hat x$ (a vertex of $S$): \cite[(2.4)]{Ba3} subgraphs are handled inside the smallest dangerous non-\cite[(2.4)]{Ba3} node containing them (\cref{lem:a06d:S44}), or alone by \cite[(2.8)]{Ba3} at either vertex; $2\phi$: the pairing of $G$ with its counterterm member, anchored at the vertex fixed by the counterterm convention (Agr), the exact identity \cite[(3.19)]{Ba3} for $D=0$ and \cite[(3.11)]{Ba3} for $D=-1$, which leaves a first-order local vertex with the regrouped coefficient and a remainder of degree $1+\alpha$; $2A$, $\phi\phi A$: Taylor expansion of the external fields at $\hat x$ (zeroth order for $D=0$, first order for $D=-1$), with \cref{lem:a06d:S27} for the multipliers and \cref{lem:a06d:S46} for sharp indicators, the local terms carrying the regrouped coefficients of \cref{thm:a06d:S48} and the remainders gaining $\alpha$ (resp.\ $1+\alpha$) on a path inside $S$. Non-\cite[(2.4)]{Ba3} inner nodes are processed before outer ones (\cite[p.~445]{Ba3}: ``assuming that it holds for all its subgraphs''), and the Taylor split of an outer node is applied to all terms produced by the inner processing (\cref{lem:a06d:S42}).

\begin{lemma}[exactness and degree preservation]\label{lem:a06d:S41}
\textup{(a)} Each processing step is an identity, at fixed external positions, of the sum over the orderings in which $S$ is useful: the Taylor formulas, \cite[(3.19), (3.26), (3.34), (2.8)]{Ba3} and the regrouping into classes are identities, and the counterterm pairing is an identity of class sums (\cref{thm:a06d:S23}). Hence the local coefficient of $S$, summed over $S$'s internal orderings, has the value of the regrouped \emph{unprocessed} coefficient; more precisely, the processed expression evaluated at external fields of $S$ that are polynomials of degree $\le1$ in the positions equals the unprocessed one, because every identity holds for every external field. \textup{(b)} Processing a useful node $S$ changes degrees only for subgraphs of $S$ containing a line of the Taylor/H\"older path ($+\alpha$, or $+1+\alpha$) and for subgraphs containing an external line of $S$ that carries the H\"older quotient ($-\alpha$, or $-1-\alpha$); every component along $\tilde l$ that contains such an external line contains $S$, so its net change is $0$. In the local term $S$ becomes a vertex $v_S$ at $\hat x$ with vertex power $D(S)$ at its vertex scale $s(v_S)=L^{j_0}\eta$ ($j_0$ the lowest external index), and every component containing $v_S$ has the degree of the original component with $S$, which is the convention \cite[(2.3)]{Ba3}.
\end{lemma}

\begin{proof}
(a) as listed; the inner processings are identities on complete groups of inner orderings, which are sub-orderings of $S$'s internal orderings. (b) A later component that contains a vertex of $S$ contains all of $S$, since $S$ is a component and components grow by merging along $\tilde l$; an earlier component contains no external line of $S$, these having higher index.
\end{proof}

\begin{lemma}[{inner processing acting on an outer node's external fields; \cite[p.~445]{Ba3}}]\label{lem:a06d:S42}
Let $S'$ be a dangerous non-\cite[(2.4)]{Ba3} node strictly inside a useful dangerous node $S$ along $\tilde l$, processed first with any anchor. In each term it produces, an external field $F$ of $S$ enters either unchanged or through a combination $\Delta[F]$ that vanishes for constant $F$ (a difference $F(x_i)-F(x_j)$, a Taylor remainder $F(x_i)-F(x_j)-(x_i-x_j)\cdot\partial F(x_j)$, or a derivative $\partial F(x_i)$); call the latter $\Delta$-terms. Then: \textup{(a)} $S$'s degree in a $\Delta$-term is $\ge D(S)+\alpha$: the gain of $S'$'s processing lies inside $S$ and the cost on $S$'s external line, which is not counted in $D(S)$; if $D(S)=0$, $S$ is not dangerous in a $\Delta$-term, and if $D(S)=-1$, $S$ is one-loop and contains no non-\cite[(2.4)]{Ba3} dangerous node. \textup{(b)} Apply $S$'s Taylor split (for $2\phi$, its counterterm pairing) to every term, $\Delta$-terms included: in a $\Delta$-term the local part is $\Delta$ evaluated at a constant field, i.e.\ $0$; for $2\phi$ the corresponding term of the counterterm member vanishes, since in $T_S(\cdot)$ the external fields of $S$ sit at $\hat x$ and $\Delta$ of them is $0$; the remainder of a $\Delta$-term is the $\Delta$-term itself. \textup{(c)} The local part of $S$, summed over all terms, is the processed expression evaluated at constant external fields of $S$, which by \cref{lem:a06d:S41}(a) equals the regrouped unprocessed coefficient; so \cref{thm:a06d:S48} bounds it, and the identities of \cref{lem:a06d:S26} and \cite[(3.30)]{Ba3} apply to it.
\end{lemma}

\begin{proof}
(a) is \cref{lem:a06d:S41}(b). (b) The external fields of $S$ are evaluated at $\hat x$ by the zeroth-order split, and every $\Delta$ annihilates constants; the counterterm member's statement is the definition of the counterterm graph \cite[p.~417]{Ba3}. (c) \cref{lem:a06d:S41}(a) at constant fields.
\end{proof}

So no anchor rule is needed: inner processings may act on outer fields; when they do, the outer node becomes convergent in those terms and its local coefficient is unchanged. Ba{\l}aban's choice (derivatives never on $G_0$'s legs) is one admissible choice among several. The anchors of $2A$ and $\phi\phi A$ nodes are free; the anchor of a $2\phi$ node is the vertex at which the counterterm graph of \cite[p.~417]{Ba3} merges the legs, the same in $G$ and in $T_SG$.

\begin{lemma}[{the \cite[(2.4)]{Ba3} subgraphs inside a node; \cite[p.~443]{Ba3}}]\label{lem:a06d:S44}
Let $S$ be a useful dangerous node of type $2\phi$, $2A$ or $\phi\phi A$ with anchor $\hat x$, and $\Sigma_1,\dots,\Sigma_r$ the fresh \cite[(2.4)]{Ba3} components inside $S$ along $\tilde l$ (vertex-disjoint, since a \cite[(1.8)]{Ba3} vertex has one derivative leg). There is an order of transport and a choice of H\"older paths such that in every term of the processed expression each $\Sigma_i$ either carries a Taylor/H\"older factor on its line, or has a vertex that carries no field external to $S$, at which \cite[(2.8)]{Ba3} is applied with no derivative on an external field of $S$.
\end{lemma}

\begin{proof}
Write $\Sigma_i=\{a_i,b_i\}$, with $a_i$ the anchor if the anchor is a vertex of $\Sigma_i$. Step~1: for each $i$ move every field external to $S$ that sits at $b_i$ onto $a_i$, to the Taylor order of $S$: $F(x_{b_i})=F(x_{a_i})[+(x_{b_i}-x_{a_i})\cdot\partial F(x_{a_i})]+$ remainder, the remainder carrying $|x_{b_i}-x_{a_i}|^\alpha$ (resp.\ $^{1+\alpha}$) on $\Sigma_i$'s line, which makes $\Sigma_i$'s degree $\ge\alpha>0$; the $\Sigma_i$ are vertex-disjoint, so these steps commute, and $\Delta$-fields are moved the same way. Step~2: transport all fields to $\hat x$ along any paths, as \cite[(3.19), (3.26), (3.34)]{Ba3} require (for $2\phi$ the two-step splitting $\phi_b(x')-\phi_b(x)=[\phi_b(x')-\phi_b(y)]+[\phi_b(y)-\phi_b(x)]$, $y$ the point of $a_i$). In every term produced by step~2, $b_i$ carries no field external to $S$: its fields were moved to $a_i$ in step~1, step~2 moves fields only from $a_i$ and other vertices, and fields at $\hat x$ are constants in $x_{b_i}$. \cite[(2.8)]{Ba3} at $b_i$ moves the derivative from $\Sigma_i$'s line onto $b_i$'s vector leg, onto $g$, or onto $b_i$'s other scalar leg, all internal to $S$ or carrying a constant field; the transformed component gets degree $1$ and later components through $b_i$ contain $\Sigma_i$'s line. If $\Sigma_i$ joins the two leg vertices of $S$ (the type \cite[(3.8)]{Ba3}, or the one-loop $2A$ bubble), step~1 is $S$'s own first-order Taylor expansion along $\Sigma_i$'s line. Inside a regrouped local coefficient there are no external fields at all (they are factored out at $\hat x$), so the integration by parts is free there.
\end{proof}

\begin{lemma}[{useless counterterms; \cite[p.~438]{Ba3}}]\label{lem:a06d:S45}
Let $\gamma$ be a $2\phi$ subgraph with external lines $l_a,l_b$, fix the indices of all lines outside $\gamma$, and let $j_*$ be the lexicographically smaller of the indices of $l_a,l_b$, $s_*\defeq L^{j_*}\eta$. Let $c_{\mathrm{use}}(x;j_*)$ be the counterterm coefficient of $T_\gamma$ restricted to the internal index assignments in which $\gamma$ is useless (maximal internal index $\ge j_*$), with $\gamma$'s inner divergences processed. Then $|c_{\mathrm{use}}(x;j_*)|\le Cs_*^{D(\gamma)}$ for $D(\gamma)=-1$ and $\le C(1+\log_L(1/s_*))\le C_{\alpha'}s_*^{-\alpha'}$ for $D(\gamma)=0$; consequently, in assignments where $\gamma$ is useless, $T_\gamma G$ is bounded by the bound function of $G/\gamma$ ($\gamma$ contracted to a vertex $v_\gamma$ at the anchor) with vertex power $D(\gamma)-\alpha'$ at $s(v_\gamma)=s_*$.
\end{lemma}

\begin{proof}
For a fixed top internal index $j_{\mathrm{top}}\ge j_*$ the bound function estimate of \cite[(2.16)]{Ba3}, with no sign condition on the degree and the anchor's cube fixed, bounds the internal sum by $Cs_{j_{\mathrm{top}}}^{D(\gamma)}$; it needs positive degrees for $\gamma$'s components along its orderings, which the inner processing supplies (induction on lines). Sum over $j_{\mathrm{top}}\in[j_*,k-1]$: for $D=-1$, $\sum_js_j^{-1}\le s_*^{-1}/(1-L^{-1})$; for $D=0$ the sum has $k-j_*=\log_L(1/s_*)$ terms. The lines of $G/\gamma$ at $v_\gamma$ are $l_a,l_b$, so $j(v_\gamma)=j_*$ and $s(v_\gamma)=s_*$, exactly the vertex scale of the bound function, and vertex powers that are powers of the vertex scale are allowed with any real exponent. In the useful assignments $T_\gamma G$ is paired with $G$ by \cite[(3.19)]{Ba3}/\cite[(3.11)]{Ba3} at the same indices (\cref{lem:a06d:S41}(a)); in the useless ones $G$ has no $\gamma$-divergence, since $\gamma$ is not a component. The loss $\alpha'\defeq\alpha/2$ is below the gains $\alpha$ of the H\"older remainders and below the half-integer gaps (degrees lie in $\frac12\Z$).
\end{proof}

\begin{lemma}[sharp indicators, boundary layers]\label{lem:a06d:S46}
Let $S$ be a dangerous node with $D(S)\le0$ whose local coefficient contains a sharp indicator $\mathbf 1_F(x_b)$ at an internal vertex $b$, $F$ the unit cube of a \cite[(1.6)]{Ba3}/\cite[(1.7)]{Ba3} vertex or $\Omega_1$. Write $c(\hat x)=\mathbf 1_F(\hat x)c^{\mathrm{unr}}(\hat x)+d(\hat x)$ with $c^{\mathrm{unr}}$ the coefficient with $\mathbf 1_F$ evaluated at $\hat x$. Then $|d(\hat x)|\le Cs_S^{D(S)}\exp(-\delta'\dist(\hat x,\partial F)/s_S)$, $s_S$ the top internal scale of $S$; summed over the cube $\Delta$ of side $s(v_S)$ of the local vertex, $\sum_{\hat x\in\Delta}\eta^d|d|\le Cs_S^{D(S)}(s_S/s(v_S))s(v_S)^d$; and after summing over $s_S\le s(v_S)$ the defect is bounded like a vertex of power $D(S)$ at $s(v_S)$, with no loss for $D(S)=0$ and better than \cite[(2.3)]{Ba3} for $D(S)=-1$.
\end{lemma}

\begin{proof}
$|\mathbf 1_F(x_b)-\mathbf 1_F(\hat x)|\le\mathbf 1\{\dist(\hat x,\partial F)\le|x_b-\hat x|\}$, the internal lines of $S$ at scales $\le s_S$ decay as $\exp(-\delta|x_b-\hat x|/s_S)$, and the bound function at fixed top scale gives the size; only the points of $\Delta$ within $O(s_S)$ of $\partial F$ contribute, a fraction $\le Cs_S/s(v_S)$ of $\Delta$; then a geometric sum.
\end{proof}

\begin{lemma}[position-space slices and remainder lines]\label{lem:a06d:S47}
\textup{(a)} Let $K$ be a kernel on $\xi\Z^3$ with $|\Delta_TD_SK(y,y')|\le C_K(\xi+\rho)^{-\omega-|S|-\alpha|T|}e^{-\delta\rho}$, $\rho\defeq|y-y'|$, for the derivative and H\"older operations used. With $\chi$ smooth, $=1$ on $[0,1]$, $0$ on $[L,\infty)$, $s_j=L^j\xi$, $\chi_0\defeq\chi(\rho/\xi)$, $\chi_j\defeq\chi(\rho/s_j)-\chi(\rho/s_{j-1})$ and the top piece $1-\sum_{j<J}\chi_j$ at $s_J=1$, the slices $K_j\defeq K\chi_j$ obey $|\Delta_TD_SK_j|\le C's_j^{-\omega-|S|-\alpha|T|}\exp(-\delta'\rho/s_j)$, so $K$ is a generalised line of the bound function with the derivatives acting on the sliced line. \textup{(b)} The remainder kernels $R=\breve G-C^\xi$ of the cube step ($|R|\le C$, $|\partial R|\le C(1+|\log(\xi+\rho)|)$, $|\partial R\partial^*|\le C(\xi+\rho)^{-1}$) have slices with bounds $s_j^{-(d-2)+1-|S|-\alpha|T|}$, one better than $C^\xi$-slices at every scale; for $|S|\ge1$ the logarithm is dominated by the terms where the derivative falls on $\chi_j$. \textup{(c)} For an inner node at scale $s_{j'}$ inside a coefficient on $\xi\Z^3$, the cumulative $C^\xi_{<j'}=C^\xi\chi_{<j'}$ splits as $C^{(j')}+R'_{j'}$ with $C^{(j')}\defeq(-\Delta^\xi+s_{j'}^{-2})^{-1}$, and after rescaling by $s_{j'}$, $R'_{j'}$ obeys the bounds of (b): $C^\xi-C^{(j')}=(s_{j'}^{-2}-1)C^\xi*C^{(j')}$ by the resolvent identity, and the cutoff part is supported at $\rho\sim s_{j'}$.
\end{lemma}

\begin{proof}
(a) On $\supp\chi_j$, $\rho$ is of order $s_j$ and $|\partial^k\chi_j|\le Cs_j^{-k}$; Leibniz. (b) Insert the bounds on $R$ into (a). (c) Resolvent identity and the convolution bound of two decaying kernels.
\end{proof}

% a06_BALtorus_d7b: continuation of a06_BALtorus_d7 (input from sec/a06_BALtorus_c.tex): the regrouped local coefficients, the induction of pp. 444-445, the assembled Proposition 1 at L = 5, and the one open row.

\begin{theorem}[bounds on the regrouped local coefficients]\label{thm:a06d:S48}
For every useful dangerous node $S$ of type $2A$ or $\phi\phi A$ the regrouped local coefficients satisfy: $2A$: $|c^{(0)}(\hat x)|\le Cs(v_S)^{D(S)}$ (zeroth moment), and for $D(S)=-1$ also $|c^{(1)}_\mu|\le C$ (first moment); $\phi\phi A$: $|c^{(0)}|\le C$; with $C$ independent of $\eta,k,\eps$, the labels, $\Omega$, $\tilde B$ and the torus. For $2\phi$ nodes the zeroth order cancels exactly against the counterterm member, and the $D=-1$ first moment has the bounded regrouped coefficient of \cite[(3.11)]{Ba3}.
\end{theorem}

\begin{proof}
Induction on the number of lines of the node. \emph{The value.} By \cref{lem:a06d:S41}(a), $c$ is the regrouped value; rescaled to $\xi=L^{-j_0}$ it is $s(v_S)^{D(S)}$ times the same expression on $\xi\Z^3$, with lines $\breve G_l=C^\xi+R_l$; smooth multipliers are moved to $\hat x$ (\cref{lem:a06d:S27}, gain $1$), sharp ones are treated by \cref{lem:a06d:S46}. \emph{The all-$C^\xi$ part} is the per-label full class with full lines $C^\xi$, internal sums unrestricted on $\xi\Z^3$, and every counterterm member re-expanded with its own lines (\cite[p.~445]{Ba3}: ``the same transformations in mass renormalization counterterms''), inner counterterm coefficients being fixed numbers, which \cref{lem:a06d:S26} allows: the $2A$ zeroth moment is $0$ by \cref{lem:a06d:S26}; the $2A$ first moment and the $\phi\phi A$ zeroth moment are $0$ by the reflection identities of \cite[(3.30)]{Ba3}, applicable because the coefficients are now constant; the torus-versus-$\xi\Z^3$ identification is the image-term argument of \cite[pp.~418--419]{Ba3}. \emph{Terms with at least one $R$-line, $D(S)=-1$}: these are one-loop nodes, bounded by the cancellation $\sum_{y'}(\partial C\partial^*)(y'-y)=0$ for the term $R\cdot\partial C\partial^*$, which uses only the bounds on $R$ and the forward-difference structure of the other line, so it applies also to $R'$-lines (\cref{lem:a06d:S47}(c)) and to slices on which the derivative acts. \emph{Terms with at least one $R$-line, $D(S)=0$}: re-slice every line (\cref{lem:a06d:S47}(a), (b)) and apply the bound function estimate \cite[(2.16)]{Ba3} to the coefficient (anchor fixed, external fields factored out). Along each ordering: a component containing an $R$-slice has degree $\ge D(H)+1$, which is $\ge\frac12$ unless $D(H)=-1$, and a $D(H)=-1$ component is a one-loop $2\phi$ or $2A$ node with an $R$-line, processed by the previous case ($2A$) or by its counterterm member with the same $R$-line ($2\phi$, exact; the first moment bounded as for \cite[(3.11)]{Ba3}); a component without an $R$-slice has its original degree, and the dangerous ones are inner nodes, processed by the induction hypothesis (local coefficients bounded, remainders gaining $\alpha$) and by \cref{lem:a06d:S44}; the whole coefficient has degree $\ge D(S)+1=1$; inner $2\phi$ nodes whose counterterm assignments are useless relative to the inner scale are bounded by \cref{lem:a06d:S45} (loss $\alpha'$). So every component has degree $\ge\min\{\frac12,\alpha\}-\alpha'>0$, and \cite[(2.16)]{Ba3} bounds the sum uniformly in $\xi$. \emph{$\phi\phi A$ remainder terms}: one-loop, so the inner components are trees, the only dangerous ones \cite[(2.4)]{Ba3}, handled by integration by parts inside the coefficient; the previous case applies with no inner nodes.
\end{proof}

\begin{theorem}[{\cite[(3.5)]{Ba3} for every class}]\label{thm:a06d:S49}
For every class component after the cube steps, $\sum_G\sum_{\tilde l}\sum_{j\in J(\tilde l)}E(G(j))$ is a finite sum of terms each of which is either a generalised bound function in which every component along every ordering of its support has degree $\ge\theta\defeq\alpha/2$, or is cancelled exactly (vacuum components, $2\phi$ zeroth orders). Hence \cite[(3.3)]{Ba3} holds by the bound function estimate \cite[(2.16)]{Ba3}.
\end{theorem}

\begin{proof}
Fix an ordering, i.e.\ an index assignment class $J(\tilde l)$, and process the dangerous nodes along $\tilde l$ innermost first. \emph{Order and anchors}: non-\cite[(2.4)]{Ba3} nodes innermost first with any anchors ($2\phi$: the counterterm anchor); each outer node's split applied to all terms produced by the inner processing, $\Delta$-terms included (\cref{lem:a06d:S42}); \cite[(2.4)]{Ba3} components treated inside the smallest dangerous non-\cite[(2.4)]{Ba3} node containing them (\cref{lem:a06d:S44}), or alone by \cite[(2.8)]{Ba3}. \emph{Local terms}: $2\phi$, the zeroth order cancels exactly against the counterterm member at the same indices in the useful assignments, and the $D=-1$ first moment has its bounded coefficient; $2A$ and $\phi\phi A$, the local vertex $v_S$ has coefficient bounded by \cref{thm:a06d:S48}, its value being the regrouped one (\cref{lem:a06d:S41}(a)); for $2A$ with $D=-1$ this vertex is the non-vanishing zero-momentum remainder vertex of \cite[(3.30)]{Ba3}, of power $-1$ at $s(v_S)$, i.e.\ degree $2+D(S)=1$ in the convention \cite[(2.3)]{Ba3}. \emph{Remainder terms} gain $\alpha$ or $1+\alpha$ inside $S$ and lose the same on $S$'s external lines (\cref{lem:a06d:S41}(b)); where an inner processing acts on an outer node's external field, the outer node gains at least $\alpha$ and has no local part in that term (\cref{lem:a06d:S42}); so after all nodes are processed every component that was dangerous has degree $\ge\alpha$, and every other component keeps its degree, which is $\ge\frac12$ since degrees lie in $\frac12\Z$. \emph{Useless counterterm assignments} are bounded by \cref{lem:a06d:S45} at a loss $\alpha'=\alpha/2$. \emph{Remainder vertices inside larger graphs}: a remainder vertex of a $2A$ node $S$ has power $D(S)$ at $s(v_S)$; components containing it have the degree of the same component with $S$ restored (\cref{lem:a06d:S41}(b)); if such a component is itself dangerous it is a node of the next level and is processed there, its local coefficient being the regrouped value with $S$ unprocessed (\cref{lem:a06d:S41}(a)), so the remainder vertex never needs a bound beyond its power. \emph{Induction}: this is Ba{\l}aban's ``for each ordering \dots\ the last divergent subgraph \dots\ apply (3.5) \dots\ then apply (3.19), (3.31) or (3.34) to $G_0$ itself'' \cite[pp.~444--445]{Ba3}, an induction on the number of lines, \cref{thm:a06d:S48} being applied inside it to coefficients of strictly smaller nodes. \emph{Vacuum and exception classes}: vacuum components cancel against \cite[(1.24)]{Ba3} up to image terms, and $1A$/$3A$ factors vanish at $B=0$.
\end{proof}

\paragraph{\cite[Prop.~1, (1.33)]{Ba3} at $L=5$.} With the steps above, every one of the twenty sketched steps and seven further places of \cite[pp.~420--445]{Ba3} is a written proof: the external-field readings of \cite[(1.16)]{Ba3} and of $\delta G_k(\Omega,\Omega_2,\tilde B)$ \cite[pp.~414, 424]{Ba3} (the latter at depth, as (I2a) of Step~3(D$_2$)), the slices and the summation by parts \cite[(2.6)--(2.12)]{Ba3}, the induction \cite[(2.16)]{Ba3}, \cite[Prop.~2.2, (2.17), Cor.~2.3]{Ba3}, the Ward--Takahashi identities for full lines and per label (\cref{lem:a06d:S26}), \cite[(3.1)]{Ba3} in the cluster form (\cref{cor:a06d:S24}), the regrouping (\cref{thm:a06d:S23}), the component cube of side $l_N$, the cube split with $l^\sharp$ (\cref{thm:a06:BALA}(a)), the constant field $B_0$ with the $k$-uniform constant of Step~3(C), the $\tilde B'$-expansion and its classes, the chains without old vertex (by an $L^2$ chain bound), gauge covariance and the gauge-out of $B_0$, $\delta G_k(\square,\eta\Z^d,0)$, the norms \cite[(1.32)]{Ba3} at $B_0$ and at $\tilde B$ (the constant carrying $(1+\operatorname{diam}\supp f)^{1-\alpha}$), and the renormalisation (\cref{thm:a06d:S49}). The inputs outside the proposition are the propagator estimates of group (A) (\cref{thm:a06A:groupA}) and \cref{thm:a06:A5}; the parameter rows, all fixed after $L$, $\bar n$, $R_0$ and before $\eps_0$, are $R\ge R_0+5$ (the depth of \cite[(3.1)]{Ba3}), $R\ge R_{\mathrm{flat}}+R_{\mathrm{Ba}}+2$ ($R_{\mathrm{flat}}$ the depth parameter of the cube step, $R_{\mathrm{Ba}}\ge R_0$ the depth of the H\"older clause of the background, $R_0$ the constant defined in \cite[(2.19)]{Ba4}; the H\"older depth of the cube step), $l^\sharp=l_N+4V^\sharp_{\mathrm{lab}}(\bar n)$ in \cref{thm:a06:BALA}(a), torus side $\ge\max\{6\sqrt3\,r(t_k)+2,8/\delta_Z\}$, the cap $e_k(\sigma_0+\sigma_1)(1+\operatorname{diam}\square)\le1$, and $\alpha_0<\frac12$ for the external-field bound of the kernel \cite[(1.16)]{Ba3} of order $m=n+n'\in\{4,5\}$, any $\alpha_0<1$ for $m\ge6$ (at the uses $m\ge\bar n+1\ge14$). The constants are independent of $\eps$, $k$, $\Omega$, $\Omega_1$, $\Omega_2$, $\tilde B$ and the torus, as \cite[(1.33)]{Ba3} asserts; they depend on $L$, which is fixed at $5$. The torus uniformity of the $O(1)$ (not in Ba{\l}aban's printed list, \cite[p.~421]{Ba3}) is the torus side of the propagator estimates of group (A): \cref{thm:a06:A5} contributes the periodic-box clause of \cite[Lemma~2.4]{Ba4}, \cref{lem:a06A:torus} gives \cite[Prop.~2.1]{Ba1} on a torus tiled by the large blocks with the constants of \cite{Ba4}, and the remaining estimates of \cref{thm:a06A:groupA} hold on the torus with the same constants, their proofs being local. With this, $[\mathrm B\text{-base}]_5$ of \cref{thm:a06:BALA}(iii) holds on every torus aligned with the blocks, with the single exception named next.

\paragraph{The walled reading.} The rotation part (\cref{sec:rot}) applies the bound \cite[(1.33)]{Ba3} once more, on tori carrying the smooth mass wall, to the two-point classes of the cut expansion of the finite order-$4$ mass counterterm and to the line-derivative families of the wall vertex: the scalar lines are $G^v_s(T_\eta,0)$ with a site potential $0\le v\le\bar v$ (in particular $v=(m^2+n\psi(t_s\cdot))t_s^2$), the photon lines are the unwalled torus propagator, the background is zero, the counterterms are the model's own, and in the two-point settings the root vertex is pinned. These classes are not per-label full classes: they are closed under the counterterm, two-vector-leg and bridgeless leg-exchange moves of \cref{def:a07:class}, but the exchange at a bridged $\phi\phi A$ subgraph carrying an external leg leaves the family. The statement needed, \cref{thm:a06:WN} below, is \cite[(1.33)]{Ba3} for such classes with walled scalar lines, with constants polynomial in $\bar v$ and free of $\eps$, $\eps_0$, the step and the torus. Its proof is the proof above run with walled lines, and it differs from the bulk at exactly two places: the zero-field kernel bounds of group (A) are replaced by their walled form (\cref{lem:a06:WN1}), and at the $\phi\phi A$ nodes the reflection identity over the full class is replaced by the clauses of \cref{thm:a06:WN}(iii), in which the exact identities of the bulk torus for the all-$C^\xi$ parts (the equality of the two leg-exchanged pinned coefficients of a bridgeless node, the vanishing of a bridged node's coefficient) are kept and the parts carrying the wall remainder $R^v$ become bounds, $|c_{S,1}(x)-c_{S,2}(x)|\le C^v_7$ for the bridgeless pair and $\le C^v_{\mathrm{BR}}$ for a bridged node paired with the counterterm member of its inner two-leg part.

% a06_WN_a: part 1 of 4 of the walled reading of Ba3 Prop. 1 (input from sec/a06_BALtorus_d7b.tex): the walled model and its zero-field kernels.
% Continues in sec/a06_WN_b.tex (classes, the two families, the statement), sec/a06_WN_c.tex (the phi phi A node clauses) and sec/a06_WN_d.tex (the bound and the constants).

\paragraph{The walled model.} Fix a step $s$, $1\le s\le K$, and work in scale-$s$ units: $\eta=L^{-s}$, $t_s\defeq L^s\eps$, $e_s\defeq e(L^s\eps)$ and $\lambda_s\defeq\lambda(L^s\eps)$ the running couplings of \cite[(1.33)]{Ba3}, and $T_\eta$ the torus, aligned with the blocks, with all sides $\ge1$. Let $v$ be a real function on the sites of $T_\eta$ with
\[
0\le v(x)\le\bar v\qquad\text{for every }x,
\]
and nothing else assumed about $v$: no smoothness, no constancy on any set, no symmetry. The walled scalar operator at zero background is
\[
H^v\defeq-\Delta^\eta_{T_\eta}+v+a_sP_s,\qquad G^v_s(T_\eta,0)\defeq(H^v)^{-1},
\]
with $a_sP_s$ the block term of the step-$s$ scalar propagator of \cite{Ba1}; at the constant $v=m^2t_s^2$ this is the bulk line $G_s(T_\eta,0)$ of Step~3(D). The wall of the rotation part is $v=(m^2+n\psi(t_s\,\cdot))t_s^2$ with a profile $0\le\psi\le1$, so that $\bar v=(m^2+n)t_s^2\le m^2+n$ once $t_s\le1$; the one-parameter family $v_\theta=(m^2+\theta n\psi(t_s\,\cdot))t_s^2$, $\theta\in[0,1]$, which interpolates between the bulk and the wall, satisfies the same bound. The \emph{walled model} is the perturbative scheme of \cite[Sect.~1]{Ba3} in which every scalar line $G_s(T_\eta,0)$, every slice \cite[(2.6)]{Ba3} and every zero-field object built from them is replaced by its $v$-analogue, while the vector lines are unchanged (the torus vector propagator carries no wall), the background is zero, and the counterterms are the model's own: a divergent two-scalar-leg subgraph $\gamma$ is renormalised by the vertex \cite[(1.7)]{Ba3} with the coefficient $\delta m^2_{s,\gamma}[v](x)$ given by Ba{\l}aban's rule \cite[p.~417]{Ba3} applied to the walled graph ($\gamma$ with both external legs at $x$ and its other leg point summed over the torus, with the same lines and the same labels, \cref{conv:a06d:LOC}, (Agr)), and the vacuum counterterm is $E_s[v]$, i.e.\ \cite[(1.24)]{Ba3} with walled lines. These counterterms depend on $x$, exactly as Ba{\l}aban's own do (\cite[p.~418]{Ba3}: the propagators ``are not translation invariant, so we have the dependence on $x$''), and the rule never uses translation invariance. Finite \cite[(1.7)]{Ba3} vertices with bounded $x$-dependent coefficients $c_v(x)$ are admitted as well; in the rotation part these are the non-renormalising parts of the inner mass counterterms and the \emph{wall vertex}, a two-scalar-leg vertex at $z$ with coefficient $-n\psi(t_sz)t_s^2$, which arises when a line of the bulk model is differentiated with respect to the wall height. The lines of a graph of the walled model are the full lines $G^v_s(T_\eta,0)$; the slices, orderings and regrouped lines of Step~3(D$_4$) are internal to the bounds.

\paragraph{Kernel notation.} For a kernel $K(x_1,x_2)$ on a lattice of spacing $\xi\in(0,1]$ and $S,T\subset\{1,2\}$, $D_S$ applies one forward difference $\partial^\xi_\mu$ in the variable $x_i$ for each $i\in S$, and $\Delta_T$ applies one H\"older quotient $|x_i-\tilde x_i|^{-\alpha}[K(\dots,x_i,\dots)-K(\dots,\tilde x_i,\dots)]$ for each $i\in T$ (in a fixed order; $\alpha\in(0,1)$ is the H\"older exponent of Step~3(D)); $\rho\defeq|x_1-x_2|$. Put $h_p(r)\defeq r^{-p}$ for $p>0$, $h_0(r)\defeq1+|\log r|$ and $h_p(r)\defeq1$ for $p<0$. In this notation the bounds of \cref{lem:a06d:S47}(b) on a remainder kernel $R$ read
\begin{equation}\label{eq:a06:WNR}
|\Delta_TD_SR(x_1,x_2)|\le C\,h_{|S|+\alpha|T|-1}(\xi+\rho)\,e^{-\delta\rho}\prod_{i\in T}|x_i-\tilde x_i|^\alpha,\qquad|S|,|T|\le2,
\end{equation}
and the zero-field, zero-mass form of the group-(A) estimates (\cite[Lemma~2.4, (2.35)--(2.36)]{Ba4}, \cref{thm:a06:A5}), extended to a difference in each variable by the same contour-shift argument, one factor $(e^{i\xi p_\mu}-1)/\xi$ per difference, for the pinned free operator $T\defeq-\Delta^\xi_X+bP_Y$ ($X$ a union of unit blocks with Neumann conditions, a periodic box or $\xi\Z^3$; $Y$ the unit blocks, or the $L$-blocks when $X$ is a union of $L$-blocks; $P_Y$ the block-average projection; $b\in[b_-,b_+]\subset(0,\infty)$) reads
\begin{equation}\label{eq:a06:WNZ1}
|\Delta_TD_ST^{-1}(x_1,x_2)|\le C_1(\xi+\rho)^{-p(S,T)}e^{-\delta\rho}\prod_{i\in T}|x_i-\tilde x_i|^\alpha,\qquad p(S,T)\defeq1+|S|+\alpha|T|,
\end{equation}
at every pair of points of $X$, with $T\ge\gamma'>0$ and $C_1,\delta,\gamma'$ depending only on $d=3$, $L$, $b_\pm$ and $\alpha$. (A quotient over two points at distance $>1$ is bounded by the two values, so \eqref{eq:a06:WNZ1} is needed only at $|x_i-\tilde x_i|\le1$; the same convention applies to every bound below.)

\begin{lemma}[lattice convolution]\label{lem:a06:WNK}
Let $0\le p_1,p_2<3$ and $\delta>0$. Then for all $x,y\in\xi\Z^3$ and $\xi\in(0,1]$,
\[
\sum_{z\in\xi\Z^3}\xi^3(\xi+|x-z|)^{-p_1}e^{-\delta|x-z|}(\xi+|z-y|)^{-p_2}e^{-\delta|z-y|}\ \le\ C_K\,h_{p_1+p_2-3}(\xi+|x-y|)\,e^{-\delta|x-y|/2},
\]
with $C_K=C_K(p_1,p_2,\delta)$. The same holds on a torus or a Neumann domain, the sum running over its sites.
\end{lemma}

\begin{proof}
Put $r\defeq|x-y|$. Since $|x-z|+|z-y|\ge r$, the exponentials are $\le e^{-\delta r/2}e^{-\delta(|x-z|+|z-y|)/2}$. On $\{|x-z|\le r/2\}$ one has $\xi+|z-y|\ge(\xi+r)/2$, and $\sum_{|x-z|\le r/2}\xi^3(\xi+|x-z|)^{-p_1}\le C(\xi+r)^{3-p_1}$ because $p_1<3$; the region $\{|z-y|\le r/2\}$ is symmetric; on the rest both factors are $\le C(\xi+|z-x|)^{-(p_1+p_2)}$ up to the exponentials, and $\sum_{|z|\ge r/2}\xi^3(\xi+|z|)^{-(p_1+p_2)}e^{-\delta|z|/2}$ is $\le C(\xi+r)^{3-p_1-p_2}$ if $p_1+p_2>3$, $\le C(1+|\log(\xi+r)|)$ if $p_1+p_2=3$ and $\le C$ otherwise. For $\xi+r\ge1$ the exponential factor dominates. On a torus or a Neumann domain the sites are a subset of $\xi\Z^3$ with the quotient (resp.\ restricted) distance, and the same sums bound the result.
\end{proof}

\begin{lemma}[walled zero-field kernels]\label{lem:a06:WN1}
Let $X$, $Y$, $b$ and $T=-\Delta^\xi_X+bP_Y$ be as in \eqref{eq:a06:WNZ1}, let $w\colon X\to[0,\bar v]$ and $T_w\defeq T+w$. Then:
\begin{enumerate}[label=\textup{(\alph*)},nosep,leftmargin=*]
\item\textup{(algebra)} $T_w\ge T\ge\gamma'$. The function $w$ commutes with every multiplication operator ($e^{\beta\rho}$ for a Lipschitz $\rho$, cutoffs, the gauge factors $V_\lambda$) and cancels in every difference of two operators carrying the same $w$; hence the gauge covariance $H^w(A^\lambda)=V_\lambda H^w(A)V_\lambda^*$, the commutator identities with cutoffs, the slice identity \cite[(2.6)]{Ba3} and its rescaling hold for the walled operator verbatim, the rescaled potential at scale $s_j$ being $s_j^2w(s_j\,\cdot)\in[0,\bar v s_j^2]\subset[0,\bar v]$, and the slices of the walled propagator are $G^{v,\eta}_{(j)}=a_j^2s_j^{-4}G^v_jQ_j^*C^{(j),v}Q_jG^v_j$ with $G^v_j\defeq(-\Delta+v+a_js_j^{-2}P_j)^{-1}$ and $C^{(j),v}$ the block operator built from $G^v_j$. Every walled zero-field scalar kernel is a scalar kernel times $\mathbf 1_N$, because $w$ is a scalar.
\item\textup{(the walled kernel bound)} $T_w^{-1}$ satisfies \eqref{eq:a06:WNZ1} at every pair of points, with $C_1$ replaced by $C_1^w\defeq C(d,L,\alpha,b_+,\gamma',\delta,C_1)(1+\bar v)^2$ and $\delta$ replaced by $\delta^w\defeq\min\{\delta/4,\beta'/\sqrt3\}/4$, where $\beta'\le1$ is the Combes--Thomas rate of \cref{lem:a07:CT} for $T_w$ (it depends only on $\gamma'$, $b_+$ and $L$, not on $w$).
\item\textup{(block sums)} For $|S|,|T|\le1$ and every block $B(y)$: $|\Delta_TD_S(T_w^{-1}\mathbf 1_{B(y)})(x)|\le c^we^{-\delta^w|x-y|}\prod_{i\in T}|x-\tilde x|^\alpha$.
\item\textup{(the remainder)} On $\xi\Z^3$ or a periodic box let $C^\xi\defeq(-\Delta^\xi+1)^{-1}$ be the free unit-mass kernel of Step~3(D$_4$) and $R^w\defeq T_w^{-1}-C^\xi$. Then $R^w$ satisfies \eqref{eq:a06:WNR} at every pair of points with $C=C_R(1+\bar v)^3$ and $\delta=\delta_R$, $C_R$ and $\delta_R$ depending only on $d,L,\alpha,b_\pm,\gamma'$ and the constants of \eqref{eq:a06:WNZ1}. In particular $R^w$ is a remainder line in the sense of \cref{lem:a06d:S47}(b), with constant $C_R(1+\bar v)^3$, and so are its slices and the kernels of \cref{lem:a06d:S47}(c) built from it.
\item\textup{(slices, cube kernels)} The slice bounds \cite[(2.10)--(2.12)]{Ba3} and the kernel bounds of the cube step of Step~3(D) hold for the walled slices $G^{v,\eta}_{(j)}$, including the two-ended entries and the mixed second difference, with constants depending polynomially on $\bar v$ and otherwise as in the bulk; the Combes--Thomas bound of \cref{lem:a07:CT} holds for $H^v$ with the $v$-free constants.
\item\textup{(periodisation)} $G^v_s(T_\eta,0)(x,x')=\sum_{n}G^{v_{\mathrm{per}}}_s(\eta\Z^3,0)(\tilde x,\tilde x'+n)$, the sum over the period lattice of $T_\eta$, $\tilde x,\tilde x'$ lifts of $x,x'$ and $v_{\mathrm{per}}$ the periodic extension of $v$; the sum converges absolutely, and the same identity holds slice by slice.
\end{enumerate}
\end{lemma}

\begin{proof}
(a) $w$ is a real scalar multiplication operator: it is self-adjoint and $\ge0$, it commutes with every other multiplication operator, and it cancels in $H^w(A)-H^w(A')$ and in $[H^w,\chi]$. The slice identity is the resolvent algebra of $-\Delta+w$ with the pin terms, which uses no property of the zeroth-order operator beyond coercivity; rescaling by $s_j$ multiplies a potential by $s_j^2$ and composes it with the dilation. Since $v$ acts as $v\mathbf 1_N$, the inverse of $H^v$ at zero field is a scalar kernel times $\mathbf 1_N$.

(b) \emph{Step 1 (Combes--Thomas).} $T_w=D^*D+w+bP_Y$ is of the form of \cref{lem:a07:CT} with $V=w\ge0$, $\Pi_B$ the block-average projections and $T_w\ge\gamma'$; so $\norm{e^{\beta\rho}T_w^{-1}e^{-\beta\rho}}\le2/\gamma'$ for every $1$-Lipschitz $\rho$ and $0\le\beta\le\beta'$, $\rho$ being $\sqrt3$-Lipschitz for the sup-metric, so that $\beta\le\beta'/\sqrt3$ is imposed, with $\beta'$ depending only on $\gamma'$, $b_+$ and the block size.
\emph{Step 2 (identity).} $T_w^{-1}=T^{-1}-T^{-1}wT^{-1}+T^{-1}wT_w^{-1}wT^{-1}$, which combines $T_w^{-1}=T^{-1}-T^{-1}wT_w^{-1}$ with $T_w^{-1}=T^{-1}-T_w^{-1}wT^{-1}$.
\emph{Step 3 (no operations).} Put $r\defeq|x-y|$. The first term is $\le C_1(\xi+r)^{-1}e^{-\delta r}$ by \eqref{eq:a06:WNZ1}. The second is $\le\bar vC_1^2C_Kh_{-1}(\xi+r)e^{-\delta r/2}=\bar vC_1^2C_Ke^{-\delta r/2}$ by \cref{lem:a06:WNK} with $p_1=p_2=1$. For the third write $a(z)\defeq T^{-1}(x,z)w(z)$, $b(z')\defeq w(z')T^{-1}(z',y)$ and $\rho\defeq|\cdot-x|$; then $|\ip a{T_w^{-1}b}|\le\norm{e^{\beta\rho}a}_2\,(2/\gamma')\,\norm{e^{-\beta\rho}b}_2$, and since $|z'-x|\ge r-|z'-y|$, $\norm{e^{-\beta\rho}b}_2\le e^{-\beta r}\norm{e^{\beta|\cdot-y|}b}_2$; both weighted norms are $\le\bar vC_1J$ with $J^2\defeq\sup_\xi\sum_z\xi^3(\xi+|z|)^{-2}e^{-(\delta-2\beta)|z|}<\infty$ (here $d=3$ and $\beta\le\delta/4$). Hence $|T_w^{-1}(x,y)|\le C^{(1)}(1+\bar v)^2(\xi+r)^{-1}e^{-\delta'r}$ with $\delta'\defeq\min\{\delta/2,\beta\}$ (for $\xi+r\le1$ the constant terms are $\le(\xi+r)^{-1}$, for $\xi+r\ge1$ the exponential absorbs the power).
\emph{Step 4 (general $S,T$).} Apply $\Delta_TD_S$ to the identity of step 2. The operations at $x_1$ act only on the first factor $T^{-1}(x_1,\cdot)$ and those at $x_2$ only on the last factor $T^{-1}(\cdot,x_2)$; $w$ and the middle factor $T_w^{-1}$ are never differentiated. Put $p_i\defeq1+[i\in S]+\alpha[i\in T]\in[1,2+\alpha]$, so that $p(S,T)=p_1+p_2-1$. The first term is \eqref{eq:a06:WNZ1}. The second is, by \eqref{eq:a06:WNZ1} for the two factors and \cref{lem:a06:WNK}, $\le2^{|T|}\bar vC_1^2C_Kh_{p_1+p_2-3}(\xi+\rho)e^{-\delta\rho/2}\prod|x_i-\tilde x_i|^\alpha$. For the third insert the bound of step 3 for the middle factor and apply \cref{lem:a06:WNK} twice: the convolution of $(\xi+|x_1-z|)^{-p_1}$ with $(\xi+|z-z'|)^{-1}$ gives $h_{p_1-2}(\xi+|x_1-z'|)$, which is $\le C_\alpha(\xi+|x_1-z'|)^{-q}$ with $q\defeq\max\{p_1-2,0\}+\alpha/2<3$ (this covers the logarithm at $p_1=2$), and the second convolution with $(\xi+|z'-x_2|)^{-p_2}$ gives $h_{q+p_2-3}(\xi+\rho)$. Since $p_1+p_2-3<p(S,T)$ and $q+p_2-3<p(S,T)$ with $p(S,T)\ge1$, both $h$-factors are $\le C(\xi+\rho)^{-p(S,T)}$ for $\xi+\rho\le1$, and for $\xi+\rho\ge1$ the exponentials dominate. Each application of \cref{lem:a06:WNK} halves the rate: the second term decays at the rate $\delta/2$ and the third at $\delta'/4$. Taking $\beta\defeq\min\{\delta/4,\beta'/\sqrt3\}$, so that $\delta'=\beta$, and collecting the constants gives $C_1^w=C(1+\bar v)^2$ and the rate $\delta^w=\min\{\delta/2,\beta/4\}=\beta/4$.

(c) Sum (b) over $x'\in B(y)$: since $p(S,T)\le2+\alpha<3$, $\sum_{x'\in B(y)}\xi^3(\xi+|x-x'|)^{-p}e^{-\delta^w|x-x'|}\le C(p)e^{-\delta^w\dist(x,B(y))}$.

(d) Since $(C^\xi)^{-1}=-\Delta^\xi+1$ and $T_w^{-1}$ exists, $R^w=T_w^{-1}\big((C^\xi)^{-1}-T_w\big)C^\xi=T_w^{-1}(1-w)C^\xi-b\,T_w^{-1}P_YC^\xi$. This is the identity \cite[p.~437]{Ba3} ($G^\xi(0)=C^\xi+G^\xi(0)(1-m^2-aP)C^\xi$) with the constant $1-m^2$ replaced by the function $1-w(z)$, which is never differentiated: the operations at $x_1$ fall on $T_w^{-1}(x_1,\cdot)$, those at $x_2$ on $C^\xi(\cdot,x_2)$. In the first term, (b) for $T_w^{-1}$, the free bound \eqref{eq:a06:WNZ1} for $C^\xi$ (it is the case $b=0$, $m^2=1$ of the group-(A) estimates), $|1-w|\le1+\bar v$ and \cref{lem:a06:WNK} give $C(1+\bar v)^3h_{p_1+p_2-3}(\xi+\rho)e^{-\delta_R\rho}\prod|x_i-\tilde x_i|^\alpha$ with $p_1+p_2-3=|S|+\alpha|T|-1$. The second term is $b\sum_Y|Y|^{-1}(T_w^{-1}\mathbf 1_Y)(x_1)(\mathbf 1_Y^{\mathsf T}C^\xi)(x_2)$, and (c) for $T_w^{-1}\mathbf 1_Y$ together with the block-sum bound for $C^\xi\mathbf 1_Y$ give $Ce^{-\delta_R\rho}\prod|x_i-\tilde x_i|^\alpha$ after the sum over $Y$, which is $\le h_{|S|+\alpha|T|-1}(\xi+\rho)$ for $\xi+\rho\le1$. The slices of $R^w$ and the kernels of \cref{lem:a06d:S47}(c) are obtained from these bounds exactly as in \cref{lem:a06d:S47}, which uses only the bounds \eqref{eq:a06:WNR}.

(e) The slice $G^{v,\eta}_{(j)}$ is the composition of (a) with $G^v_j$ bounded by (b) at scale $j$ (unit-block units after rescaling, with the potential $s_j^2v(s_j\,\cdot)\in[0,\bar v]$) and with the block-operator kernel $C^{(j),v}$, whose exponential bound uses only the coercivity $H^v\ge H^0\ge\gamma_0$ and the Combes--Thomas argument of \cref{lem:a07:CT}, so its constants are $v$-free; the bulk proof of \cite[(2.10)--(2.12)]{Ba3} in Step~3(D) is a composition of these three bounds and of lattice sums, and the mixed second difference and the two-ended entries are the cases $|S|=2$ of (b). The kernel bounds of the cube step are sums of slice bounds over the scales. The Combes--Thomas bound is \cref{lem:a07:CT} with $V=v$.

(f) The right side is well defined because the image sum converges absolutely by (b) on $\eta\Z^3$ with the periodic potential $v_{\mathrm{per}}$ (which lies in $[0,\bar v]$); it is periodic, and it solves the torus equation because the $\eta\Z^3$ operator with $v_{\mathrm{per}}$ commutes with the period translations; the torus inverse is unique. The slice identity (a) is compatible with the periodisation term by term.
\end{proof}

% a06_WN_b: part 2 of 4 of the walled reading of Ba3 Prop. 1 (input from sec/a06_BALtorus_d7b.tex): classes closed under the moves, the two families of the rotation part, and the statement.

\paragraph{Classes closed under the moves.} Recall the moves (P), (P$'$), (W), (X) and the classes of \cref{def:a07:class}, and the cut terms and orbits of \cref{def:a07:cut}. A \emph{decorated graph} of the walled model is a graph of the model (lines $G^v_s(T_\eta,0)$ and the torus vector propagator, vertices \cite[(1.6)--(1.15)]{Ba3} with their labels (Loc), counterterm vertices with the model's own coefficients, finite \cite[(1.7)]{Ba3} vertices with bounded coefficients) whose external legs carry external fields which may be kernels of several variables \cite[p.~419]{Ba3}, in particular the two ends of a cut line; these legs are the \emph{slots} of the graph. A finite set $\cC$ of decorated graphs with a common slot structure (the same slots, with the same kinds, differentiations, labels and fields), a common order, a common label multiset and a common multiset of finite-vertex coefficients is \emph{closed under the moves} if every move of the set $M_0$, consisting of (P) and (P$'$) at a divergent two-scalar-leg subgraph of a member, (W) at a divergent two-vector-leg subgraph, and (X) at a \emph{bridgeless} divergent $\phi\phi A$ subgraph (at the member itself included, when the member is of one of these types), gives a member. The set $M_0$ is a subset of the moves of \cref{def:a07:class} and has the properties (A1)--(A5) of that definition, which are properties of the individual moves; so \cref{lem:a07:orb} applies with $M=M_0$. Per-label full classes are closed under the moves (\cref{def:a07:class}), and so is every class $\Cl_{M}(K)$, $M\supset M_0$, of a connected component $K$ of a cut graph in \cref{lem:a07:orb}(c), with the cut ends as slots. The restriction to bridgeless subgraphs in (X) is forced by the two-point settings: the exchange at a bridged divergent $\phi\phi A$ subgraph carrying an external leg produces a one-particle-reducible graph, which is not a term of a family of \cite[(1.23)]{Ba3} graphs (\cref{thm:a06:WN}(iii)(c) shows that no partner is needed there). In a \emph{two-point setting} the members have two external scalar legs, undifferentiated, one of them, the \emph{root}, at a vertex pinned at the point $x$, the other at a vertex that is summed; the coefficient is a function of $x$, and the bounds below hold uniformly in $x$. A $\phi\phi A$ subgraph $S$ is \emph{bridgeless} if no internal line of $S$ disconnects it (a self-loop is not a bridge), and \emph{bridged} otherwise.

\paragraph{The two families of the rotation part.} (a) \emph{The cut expansion of the finite order-$4$ mass counterterm.} The order-$4$ two-point coefficient of the walled model is computed once with the $\eps$-level lines $C^v_0$, $C^\gamma$ and once with the step-$K$ lines $G^v_K$, $G^\gamma_K$; the difference, with the scalar lines decomposed as $C^v_0=G^v_K+R^v_K$ and the vector lines as $C^\gamma=G^\gamma_K+R^\gamma_K$, is the sum \eqref{eq:a07:tel} of \cref{def:a07:cut} over the cut terms with at least one cut line, the cut ends carrying the kernels $R^v_K$, $R^\gamma_K$. By \cref{lem:a07:orb}(c),(d) each orbit is a product of classes $\Cl_M(K)$ of the connected components $K$ of its cut graph, each closed under the moves, with the cut ends as slots; the component carrying the root is a two-point class with the root pinned at $x$, the others are vacuum components with slots. (b) \emph{The line-derivative family of the wall vertex.} Let $\cF^{(4)}$ be the set of two-point graphs of \cite[(1.23)]{Ba3} of coupling order $4$: one-particle-irreducible, two undifferentiated external scalar legs, the root at $x$, the inner \cite[(1.7)]{Ba3} vertices carrying the model's own step-$s$ counterterm of their order and their non-renormalising parts as finite vertices. For $\Gamma\in\cF^{(4)}$ and a scalar line $l$ of $\Gamma$ let $\Gamma_{l\to v}$ be $\Gamma$ with $l$ subdivided by a two-scalar-leg vertex $v$ at a point $z$ with a bounded coefficient $c(z)$ (the wall vertex), and
\[
(\Gamma_{l\to v})_{\mathrm{ren}}\ \defeq\ \sum_{F}\prod_{\gamma\in F}(-T_\gamma)\,\Gamma_{l\to v},
\]
the sum over the forests $F$ of divergent two-scalar-leg subgraphs of $\Gamma$ avoiding $l$ (i.e.\ not containing $l$), $T_\gamma$ the counterterm operation of (Agr). The family is $\sum_{\Gamma\in\cF^{(4)}}\sum_l(\Gamma_{l\to v})_{\mathrm{ren}}$, with no cut line.

\begin{lemma}[the line-derivative family]\label{lem:a06:WNNC}
\textup{(a)} Every connected subgraph of $\Gamma_{l\to v}$ that contains $v$ has degree $D\ge1$ or vanishes. Hence the divergent subgraphs of $\Gamma_{l\to v}$ are exactly the divergent subgraphs of $\Gamma$ avoiding $l$, and $(\Gamma_{l\to v})_{\mathrm{ren}}$ is the renormalised forest sum of $\Gamma_{l\to v}$. \textup{(b)} A forest of divergent two-scalar-leg subgraphs of $\Gamma\in\cF^{(4)}$ avoiding $l$ has at most one element. \textup{(c)} Group the terms $T_F\Gamma_{l\to v}\defeq\prod_{\gamma\in F}(-T_\gamma)\Gamma_{l\to v}$ ($\Gamma\in\cF^{(4)}$, $l$ a scalar line of $\Gamma$, $F$ as above) by label multiset and by the multiset of finite-vertex coefficients, the wall vertex included. Each group is a two-point class closed under the moves, with the root pinned at $x$.
\end{lemma}

\begin{proof}
(a) The vertex $v$ is a \cite[(1.7)]{Ba3} vertex with a bounded coefficient, so it has $D(v)=2$ and every connected subgraph $H$ containing it has $D(H)\ge1$ if $H$ has at least two slots (\cref{lem:a06d:S25}(b), \cite[Cor.~2.3]{Ba3}). A connected subgraph with exactly one slot vanishes: with a scalar slot by the symmetry $\phi\to-\phi$ of every vertex (its amputated kernel $K_a$ with one open scalar index satisfies $K_a=-K_a$), with a vector slot at zero background because its scalar index loops close inside it and $\Tr q^{\mathrm{odd}}=0$ (Step~3(D$_4$)); and a connected subgraph of $\Gamma_{l\to v}$ with no slot is $\Gamma_{l\to v}$ itself, which has two. The divergent subgraphs of $\Gamma_{l\to v}$ therefore avoid $v$, i.e.\ are subgraphs of $\Gamma$ not containing $l$, with the same vertices, lines, slots and degree (\cref{lem:a07:orb}(a)); conversely a divergent subgraph of $\Gamma$ avoiding $l$ is a subgraph of $\Gamma_{l\to v}$ with the same data.

(b) A proper divergent two-scalar-leg subgraph of $\Gamma$ has coupling order $2$ (\cite[(1.23)]{Ba3}; the order-$4$ one is $\Gamma$ itself, which contains $l$). Two such subgraphs cannot be nested, and two disjoint ones $\gamma,\gamma'$ would exhaust the vertices of $\Gamma$: each has exactly two slots, the two external legs of $\Gamma$ are among the four slots, so at most one line joins $\gamma$ to $\gamma'$; if exactly one does it is a bridge of $\Gamma$, and if none does $\Gamma$ is disconnected; both contradict $\Gamma\in\cF^{(4)}$.

(c) Fix a group. \emph{(P), (P$'$).} A divergent two-scalar-leg subgraph $\gamma$ of a member $T_F\Gamma_{l\to v}$ avoids $v$ and $l$ by (a); the move replaces the term with forest $F$ by the term with forest $F\cup\{\gamma\}$ (or $F\setminus\{\gamma\}$), a forest avoiding $l$ by (b), of the same $\Gamma$ and $l$; the counterterm vertex carries $\gamma$'s labels (Agr) and the model's own coefficient, so the labels and the finite-vertex coefficients of the group are kept. \emph{(W).} A divergent two-vector-leg subgraph $H$ of a member avoids $v$ and $l$ by (a). Replacing $H$ by a graph $H'$ of its per-label two-vector-leg class (same two vector slots, same order, same label multiset, its counterterm members included) changes nothing outside $H$: the line $l$ with the vertex $v$ on it and the two external scalar legs are untouched, the order and the labels are kept, and the result is one-particle-irreducible: since $T_F\Gamma_{l\to v}$ is one-particle-irreducible, the two attachment points of $H$ are connected outside $H$ (else $H$'s vector line to the rest would be a bridge), so an internal bridge of $H'$ separating its two slots is not a bridge of the new graph, and an internal bridge of $H'$ with a slot-free side bounds a one-slot subgraph, which vanishes by (a). Smoothing $v$ back into a line gives a graph of $\cF^{(4)}$ with the counterterm vertices of $F$ inside or outside $H$ as the case may be (a forest element inside $H$ is a counterterm member of $H$'s class), so the result is a member of the group. \emph{(X).} The exchange at a bridgeless $\phi\phi A$ subgraph $P$ moves only the attachments at $P$'s two scalar slots and no vertex: the line $l$ keeps its ends up to this re-attachment and $v$ stays on it, the external legs stay undifferentiated (both slots of $P$ carry no difference, since $P$ is divergent), and the result is one-particle-irreducible: for a line $\ell\notin P$, $P$ is connected in both graphs minus $\ell$ with the same multiset of outer attachment points, and for $\ell\in P$, $P\setminus\ell$ is connected because $P$ is bridgeless. The forest element, if any, is nested with $P$ or disjoint from it (a divergent two-scalar-leg subgraph crossing a bridgeless order-$3$ $\phi\phi A$ subgraph would meet it in a single one-vector-leg vertex whose two internal lines of $P$ would be its two scalar slots, leaving its vector line as a bridge of the forest element); a forest element sharing a vertex with $P$ but no line (the tadpole at the \cite[(1.6)]{Ba3} vertex of the double-line shape) is untouched, since \textup{(X)} moves attachments only; so $F$ is carried along. \emph{Bridged $\phi\phi A$ subgraphs.} No move of $M_0$ acts there. (The exchange at a bridged divergent $\phi\phi A$ subgraph $\{u\}\cup S_2$ of a member, with the external leg of the member at $u$, would move that leg to $S_2$, which then hangs on the bridge alone: the result is one-particle-reducible and is not a term of the family. Such a member exists: $u$, $u'$ one-vector-leg vertices \cite[(1.8)]{Ba3} joined by a vector line, $w$ a \cite[(1.6)]{Ba3} vertex with a self-loop, scalar lines $u$--$w$ and $w$--$u'$, the root at $u$ and the wall vertex on $w$--$u'$; the node $\{u,w\}$ is a dangerous bridged $\phi\phi A$ node of degree $0$.) Finally every member has its root pinned at $x$, since no move moves a vertex.
\end{proof}

\begin{theorem}[the walled form of {\cite[Prop.~1]{Ba3}} for classes closed under the moves]\label{thm:a06:WN}
In the walled model at a step $s\le K$ (potential $0\le v\le\bar v$ on the scalar lines only, vector lines unwalled, zero background, the model's own counterterms, finite \cite[(1.7)]{Ba3} vertices with bounded coefficients, torus aligned with the blocks with sides $\ge1$) let $\cC$ be a class closed under the moves, of order $\le\bar n$, with slot fields $\Phi_{\mathrm{ext}}$, $A_{\mathrm{ext}}$ and finite-vertex coefficients $c_v$. Then:
\begin{enumerate}[label=\textup{(\roman*)},nosep,leftmargin=*]
\item\textup{(regrouping)} \Cref{lem:a07:orb} (with $M=M_0$ in the two-point settings, under which the families of \cite[(1.23)]{Ba3} graphs are closed) and \cref{thm:a06d:S23} hold verbatim with walled lines; in particular the cut expansion of the finite order-$4$ mass counterterm is a finite sum over orbits, each orbit the contraction through the kernels $R^v_K$, $R^\gamma_K$ of per-component classes closed under the moves, with the invariants of \cref{lem:a07:orb}(b),(f); and the line-derivative family of the wall vertex is a finite disjoint union of two-point classes closed under the moves (\cref{lem:a06:WNNC}).
\item\textup{(the bound)} With $e_s=e(L^s\eps)$ and $\lambda_s=\lambda(L^s\eps)$,
\[
\big|E^v(\cC;\{\square(v)\},\Phi_{\mathrm{ext}},A_{\mathrm{ext}})\big|\ \le\ O(1)^v\,e_s^{\,d_v}\lambda_s^{\,d_s}\prod_{v\ \mathrm{finite}}\sup_z|c_v(z)|\;e^{-\delta^v_0d(\{\square(v)\})}\prod\norm{h\Phi_{\mathrm{ext}}}_{1,\alpha_0}\prod\norm{h'A_{\mathrm{ext}}}_{1,\alpha_0},
\]
the form \cite[(1.33)]{Ba3} ($E^v$ the class sum of \cref{lem:a07:orb}(d) with walled lines): $d_v$, $d_s$ count the couplings of the action vertices, a counterterm vertex counting as the exact counterterm of the order it replaces; the localisations $h$, $h'$ are those of (H$_A$); the slot fields are kernels of several variables, each measured by its norm \cite[(1.32)]{Ba3} in the variable at the slot; in a two-point setting the left side is the coefficient at the pinned root $x$, and the bound is uniform in $x$. The constants $O(1)^v$ and $\delta^v_0$ depend only on the order, on the parameters of Step~3(D) ($d=3$, $L=5$, $a$, $N$, $\norm q$, $\bar n$, $\alpha_0$, $m$, $\lmu$, $M$, $R$) and polynomially on $\bar v$; they do not depend on $\eps$, on $\eps_0\in(0,\eps_0^*]$ below the caps of Step~3, on $s$, $K$, the labels, the torus side, $x$, or on $v$ beyond $\bar v$.
\item\textup{(the $\phi\phi A$ node clauses)} Let $S$ be a dangerous $\phi\phi A$ node of a member $G$ of $\cC$ along an ordering, and let $c$ be its regrouped zeroth-order local coefficient in the form of \cref{thm:a06d:S48}: rescaled to $\xi\Z^3$, lines $\breve G_l=C^\xi+R_l$ with $R_l=R^v_l$ of \cref{lem:a06:WN1}(d) on the scalar lines and the bulk remainder on the vector lines, multipliers at the anchor, and the decomposition $c=c^C+c^R$ into the all-$C^\xi$ part and the terms with at least one $R$-line.
\begin{enumerate}[label=\textup{(\alph*)},nosep,leftmargin=*]
\item $S$ bridgeless, $G^{\mathrm{sw}}$ the member obtained by \textup{(X)} at $S$, a common anchor in $S$: $c(S\text{ in }G)+c(S\text{ in }G^{\mathrm{sw}})=0$ exactly, for any line kernels.
\item $S$ bridgeless with its two scalar slots at distinct vertices $u_1\ne u_2$, the pinned root at $u_1$ in $G$ and at $u_2$ in $G^{\mathrm{sw}}$, anchors at the root; $c_{S,1}(x)$, $c_{S,2}(x)$ the two coefficients: $c^C_{S,1}(x)=c^C_{S,2}(x)$, and $|c_{S,1}(x)-c_{S,2}(x)|\le C^v_7$.
\item $S$ bridged. Then $S=\{u\}\cup S_2$, where $u$ is a one-vector-leg vertex \cite[(1.8)]{Ba3} carrying one scalar slot and the vector slot of $S$, its difference acting on the bridge $l_0=(u,w)$, $w\in S_2$, and $S_2$ is a two-scalar-leg subgraph of order $2$ (the exchange bubble, the \cite[(1.6)]{Ba3} tadpole, a seagull with a vector loop, or a \cite[(1.7)]{Ba3} vertex); these are the shapes \cite[(2.21b)--(g)]{Ba3}. If $u$ is summed, then with the anchor in $S_2$, $c=0$ for every member and any line kernels. If $u$ carries the pinned root (anchor $u$ at $x$), then $c^C=0$ for every member, and the $R$-part of the pair $(S\text{ with }S_2)+(S\text{ with }T_{S_2})$, both members of $\cC$ by \textup{(P)}, is $\le C^v_{\mathrm{BR}}$ in absolute value; the $R$-part of a single member is in general not bounded uniformly in $\xi$.
\end{enumerate}
$C^v_7$ and $C^v_{\mathrm{BR}}$ are the $\phi\phi A$ constant of \cref{thm:a06d:S48} times a polynomial in $\bar v$, free of $\xi$, $j_0$, $s$, $\eps$, the labels, the torus and $x$.
\end{enumerate}
\end{theorem}

The rotation part consumes (ii) and (iii) at the step $s=K$ for the per-component classes of the cut expansion in (i), and at every step $s=a\le K$ for the line-derivative family, both with $v$ in the family $v_\theta$ and at coupling order $\le4$. The proof of (ii) is by induction on the number of vertices of the members, as the proof of \cite[(1.33)]{Ba3} is (\cref{lem:a06d:S25}); the induction runs over classes closed under the moves, the completed pieces of a cut partnered subgraph being such classes of lower order.

\begin{proof}[Proof of \textup{(i)}]
\Cref{lem:a07:orb} uses the lines only through the properties (A1)--(A5) of the moves, the multilinearity of the expansion \eqref{eq:a07:tel} and the symmetry of the decomposed kernels; $R^v_K=C^v_0-G^v_K$ is a difference of two symmetric resolvents, hence symmetric. In the two-point setting the family of \cite[(1.23)]{Ba3} graphs is closed under $M_0$: (P) contracts or expands a one-particle-irreducible two-leg subgraph (a bridge of the new graph would be one of the old), and (W), (X) at bridgeless subgraphs keep one-particle-irreducibility and the undifferentiated external legs by the arguments given in the proof of \cref{lem:a06:WNNC}(c). \Cref{thm:a06d:S23} and \cref{lem:a06d:S22} are statements about labels and the cut rule. Degrees count lines (CV). So (i) is the bulk statement, read with walled lines, and the second half is \cref{lem:a06:WNNC}(c).
\end{proof}

% a06_WN_c: part 3 of 4 of the walled reading of Ba3 Prop. 1 (input from sec/a06_BALtorus_d7b.tex): the proof of thm:a06:WN(iii), the phi phi A node clauses with walled lines.

\begin{proof}[Proof of \cref{thm:a06:WN}\textup{(iii)}]
\emph{The shapes.} A dangerous $\phi\phi A$ node has order $3$ and one loop (Step~3(D$_4$)), and a node containing a finite \cite[(1.7)]{Ba3} vertex is not dangerous (\cref{lem:a06d:S25}(b)); by \cite[p.~429]{Ba3} its vertices are of the types \cite[(1.6)--(1.8), (1.10)]{Ba3} or divergent-order counterterm vertices. With an odd number of vector legs the vertex multisets are three one-vector-leg vertices \cite[(1.8)]{Ba3}; one one-vector-leg vertex with one two-vector-leg vertex of type \cite[(1.8)]{Ba3} or \cite[(1.10)]{Ba3}; one three-vector-leg vertex \cite[(1.8)]{Ba3}; or one one-vector-leg vertex with a \cite[(1.6)]{Ba3} vertex or with an order-$2$ counterterm vertex. The line count gives one loop in each case, so the internal lines form one cycle through all vertices (a triangle or a double line) in the bridgeless case, and a bridge otherwise. The single three-vector-leg vertex is excluded (a vector self-loop leaves a differentiated scalar slot, a scalar self-loop leaves three vector slots). The bridgeless shapes are: \emph{T}, three one-vector-leg vertices in a triangle, the apex carrying the vector slot and the two base vertices the scalar slots, which are their undifferentiated legs, the base vertices' differences acting on their lines to the apex, and the apex's difference on one of its two lines; \emph{B8}, resp.\ \emph{B10}, a one-vector-leg vertex $u$ and a two-vector-leg vertex $w$ of type \cite[(1.8)]{Ba3}, resp.\ \cite[(1.10)]{Ba3}, joined by a vector line and a scalar line, the slots being $u$'s undifferentiated scalar leg, $w$'s undifferentiated scalar leg and $w$'s second vector leg; \emph{BL}, a one-vector-leg vertex $u$ and a \cite[(1.6)]{Ba3} vertex $w$ joined by two scalar lines ending at legs of different index pairs of $w$, the slots being the two remaining legs of $w$ and $u$'s vector leg. In each shape every scalar slot is an undifferentiated leg and the vector slot carries no difference, as the degree $D=0$ requires. The bridged shapes are \cite[(2.21a)--(g)]{Ba3}; \cite[(2.21a)]{Ba3}, a vector bridge whose other side is a one-vector-leg vertex with a scalar loop, vanishes identically, and in the remaining shapes: a side of the bridge with no slot would be a one-slot subgraph, which vanishes (\cref{lem:a06:WNNC}(a)); a vector bridge is excluded: the side carrying the vector slot has two vector legs (the slot and the bridge), hence order $\ge2$, so the other side has order $1$ and is a single one-vector-leg vertex with no loop, whose two scalar legs are then both external, i.e.\ both scalar slots lie on it, and its difference acts on a slot, against $D=0$; so the bridge $l_0$ is scalar, each side carries one scalar slot, the side with the vector slot has order $\ge1$ and the other, a two-scalar-leg subgraph, order $\ge2$, whence orders $1+2$: the first side is a single one-vector-leg vertex $u$ whose difference, not allowed on its scalar slot or its vector leg, acts on $l_0$, and $S_2$ is one of the order-$2$ two-scalar-leg graphs, none of which has a bridge. This is the structure asserted in (c), and it is what Ba{\l}aban's arrows in \cite[(2.21b)--(g)]{Ba3} show.

\emph{(a): the common-anchor cancellation.} Let $\sigma_i\in u_i$ be the two scalar slots of $S$, on outer lines or external fields $F_i$. $G^{\mathrm{sw}}$ differs from $G$ only in the attachments at $\sigma_1,\sigma_2$, so $S$ has the same vertices, internal lines, differences, vertex localisations and anchor $\hat x$ in both, and the external lines of $S$ keep their indices, so that usefulness, the regrouped lines of the internal and external slices (which depend on the external indices only) and the transport of the multipliers to $\hat x$ (\cref{lem:a06d:S27}) agree. Hence the regrouped position kernel $K_S(y_1,y_2;\dots)$ of $S$, including its $R$-lines and the walled kernels, is the same function in both members, whatever the line kernels are. Its index structure is a matrix $M_S$ on the open index path from $\sigma_1$ to $\sigma_2$: every closed index cycle inside $S$ carries an even power of $q$ (else $\Tr q^{\mathrm{odd}}=0$ kills the term), and the total number of vector-leg ends at the vertices of $S$ is $2(\text{internal vector lines})+1$, odd, so the open path carries $q^{r}$ with $r$ odd times even polynomials in $q$, and $M_S^{\mathsf T}=-M_S$. At zeroth order the external fields are constants at $\hat x$ and the two local terms are $c(\hat x)F_1(\hat x)^{\mathsf T}M_SF_2(\hat x)$ and $c(\hat x)F_2(\hat x)^{\mathsf T}M_SF_1(\hat x)=-c(\hat x)F_1(\hat x)^{\mathsf T}M_SF_2(\hat x)$, with the same $c(\hat x)$; they cancel. Walled lines change nothing: they are scalar kernels times $\mathbf 1_N$ (\cref{lem:a06:WN1}(a)), so $M_S$ is unchanged; for slot fields which are kernels of several variables (the ends of $R^v_K$, $R^\gamma_K$), the zeroth-order term at $\hat x$ is their value at $\hat x$ in the slot variable, controlled by the norm \cite[(1.32)]{Ba3}. The remainders of the Taylor expansion are terms of the processing of Step~3(D$_4$), bounded term by term. This is Ba{\l}aban's reason for \cite[(2.18)]{Ba3}, ``the same expressions but with the opposite signs'' \cite[p.~444]{Ba3}, with the two members anchored at one vertex.

\emph{(b): the pinned root.} Here $u_1\ne u_2$ (the shapes T, B8, B10; in BL both scalar slots sit at $w$, so $u_1=u_2$ and (a) applies with the anchor at $w$). The local terms are $c_{S,1}(x)\Phi$ and $-c_{S,2}(x)\Phi$, with $c_{S,i}$ the regrouped coefficient with $u_i$ at $x$ and the anchor at $u_i$ (the sign from (a)). \emph{Split.} By the steps of \cref{thm:a06d:S48}: rescale to $\xi\Z^3$ ($\xi=L^{-j_0}$, $j_0$ the lowest external index), where by \cref{lem:a06:WN1}(f) and the rescaling of \cref{lem:a06:WN1}(a) every regrouped scalar line is $\breve G^v_l=C^\xi+R^v_l$ with a potential in $[0,\bar v]$ and $R^v_l$ as in \cref{lem:a06:WN1}(d), and every vector line is $C^\xi+R_l$ with the bulk remainder; move the multipliers to the anchor (\cref{lem:a06d:S27}); expand multilinearly, $c_{S,i}=c^C_{S,i}+c^R_{S,i}$. The shapes T, B8, B10 contain no \cite[(1.6)]{Ba3} or \cite[(1.7)]{Ba3} vertex, so no sharp indicator occurs and the internal sums are unrestricted on $\xi\Z^3$; and they contain no divergent proper subgraph (a two-scalar-leg or two-vector-leg subgraph inside a one-loop triangle or double line would need a bridge), so no inner counterterm occurs. \emph{The all-$C^\xi$ parts.} $c^C_{S,i}(x)$ is the same expression as in the bulk model: the same graph, vertex tensors and differences, the multipliers at the anchor, lines $C^\xi$, unrestricted sums. It is the zero-momentum value of one translation-invariant kernel on $\xi\Z^3$ with $u_i$ pinned: with $K^C_S(y_{u_1},y_{u_2},\dots)$ the all-$C^\xi$ kernel, $c^C_{S,1}(x)=\sum K^C_S(x,y_{u_2},\dots)=\sum K^C_S(0,y,\dots)$ and $c^C_{S,2}(x)=\sum K^C_S(y_{u_1},x,\dots)=\sum K^C_S(0,-y,\dots)$, both the total sum of $K^C_S(0,\cdot,\dots)$ over $\xi\Z^3$ (absolutely convergent at fixed $\xi$), so $c^C_{S,1}(x)=c^C_{S,2}(x)$, independent of $x$. The wall does not enter this term. \emph{The $R$-parts.} Each term of $c^R_{S,i}$ has at least one $R$-line. \Cref{lem:a06:WN1}(d) gives the bounds \eqref{eq:a06:WNR} for $R^v_l$ at every pair of points, with constant $C_R(1+\bar v)^3$, so $R^v_l$ is a remainder line of \cref{lem:a06d:S47}(b), and the step of \cref{thm:a06d:S48} for ``terms with at least one $R$-line, $D(S)=0$'' applies: the node is one-loop, its inner components are trees, the only dangerous ones are the \cite[(2.4)]{Ba3} subgraphs, treated by the summation by parts inside the coefficient (\cref{lem:a06d:S44}), every component after re-slicing has positive degree, and the bound function estimate \cite[(2.16)]{Ba3} bounds each $c^R_{S,i}$ by $C(1+\bar v)^p$ uniformly in $\xi$, $j_0$, $s$, $\eps$, the torus and $x$ (the kernel bounds are pointwise, and the anchor is fixed at $x$). The exponent $p$ counts the $R$-lines and the factors of $\bar v$ in their bounds, at most three lines here. \emph{Conclusion.} $|c_{S,1}(x)-c_{S,2}(x)|=|c^R_{S,1}(x)-c^R_{S,2}(x)|\le2C(1+\bar v)^p=: C^v_7$. This is the form of the $\phi\phi A$ clause of \cref{thm:a06d:S48}, $|c^{(0)}|\le C$, for the pair.

\emph{(c): bridged nodes.} Write $\Sigma_2(w,w')$ for the amputated kernel of $S_2$ between its leg point $w$ on the bridge and its other leg point $w'$, with $S_2$'s lines (walled or not, at the regrouped scale), and $F(y)\defeq(\partial^\xi_\mu\breve G_{l_0})(x_u,y)$ for the bridge factor, $u$'s difference acting on the bridge.

\emph{$u$ summed.} Choose the anchor $\hat x$ in $S_2$ (anchors of $\phi\phi A$ nodes are free, Step~3(D$_4$)). At zeroth order every external field of $S$ is a constant at $\hat x$, and so are $u$'s multipliers, which are smooth (\cref{lem:a06d:S27}; a one-vector-leg vertex carries no sharp indicator). The dependence of the regrouped kernel on $x_u$ is then the single factor $F$, and $\sum_{x_u}\xi^3(\partial^\xi_\mu\breve G_{l_0})(x_u,y)=0$ for every summable kernel: a forward difference in the summed variable telescopes, on $\xi\Z^3$ and on the torus. So $c=0$ for every member and any line kernels, before any $C^\xi/R$ split, and the same holds for the member with $S_2$ replaced by its counterterm vertex. No translation invariance of the lines is used.

\emph{$u$ pinned.} The anchor is $u$, at $x$. \emph{All-$C^\xi$ part.} On $\xi\Z^3$ with lines $C^\xi$ and unrestricted sums, $\sum_{w'}\Sigma^C_2(w,w')$ is a constant $g$ by translation invariance, and $\sum_w\xi^3(\partial^\xi_\mu C^\xi)(x,w)\,g=g\,\partial^\xi_\mu(C^\xi\mathbf 1)(x)=0$ since $C^\xi\mathbf 1\equiv1$. The counterterm member is the same computation with $\Sigma_2$ replaced by its own row sum. A sharp indicator inside $S_2$ occurs only at the \cite[(1.6)]{Ba3} tadpole, which cancels exactly against its counterterm member label by label ((Agr); the counterterm carries the same unit-cube label), so the pair has no indicator defect; the exchange bubble and the seagull loop carry only the smooth labels of their vector legs. So $c^C=0$ per member (and per pair). This is Ba{\l}aban's ``by translation invariance it can be written as a derivative of a constant'' \cite[p.~444]{Ba3} for \cite[(2.21c)]{Ba3}, and the exact cancellations he asserts for \cite[(2.21b), (2.21d), (2.21f), (2.21g)]{Ba3}, applied to the all-$C^\xi$ parts. \emph{Why the $R$-part is not bounded per member.} The term in which the bridge carries $R^v_{l_0}$ and $S_2$ carries lines $C^\xi$ equals $g\sum_w\xi^3(\partial^\xi_\mu R^v_{l_0})(x,w)=g\,\partial^\xi_\mu(R^v\mathbf 1)(x)$, where $g\sim\xi^{-1}$ (the node $S_2$ has $D(S_2)=-1$, and the $R$-line is not inside it, so it gains nothing for $S_2$). In the bulk $R\mathbf 1$ is constant ($-\Delta+m^2+aP$ maps constants to constants), and the term vanishes; under a wall $\breve G^v\mathbf 1=\mathbf 1-\breve G^v v$ is not constant, so the term is of order $\xi^{-1}$ and is not bounded uniformly in $\xi$. \emph{The pair.} Both members are in $\cC$ by (P). Let $f$ be the external field at $S_2$'s far slot $w'$ (constant at the anchor at zeroth order, Taylor-expanded at first order), and take the anchor convention (Agr) for $S_2$ with the anchor at the bridge end $w$, so that the counterterm member carries $\delta m^2_{S_2}(w)=\sum_{w'}\xi^3\Sigma_2(w,w')$, the row sum of the same kernel (the model's own counterterm, \cite[p.~417]{Ba3}). Then
\[
\sum_{w,w'}\xi^6F(w)\Sigma_2(w,w')f(w')-\sum_w\xi^3F(w)\,\delta m^2_{S_2}(w)f(w)\ =\ \sum_{w,w'}\xi^6F(w)\Sigma_2(w,w')\big[f(w')-f(w)\big],
\]
an identity for any kernels, with no symmetry used. Its zeroth-order term ($f$ constant) is $0$ exactly, including the terms in which an $R$-line sits inside $S_2$ (these have $D(S_2)+1=0$ and are paired with the counterterm member carrying the same $R$-line, as in the step of \cref{thm:a06d:S48} for $D(H)=-1$ components). The first-order term is $\sum_w\xi^3F(w)c_\mu(w)\,\partial_\mu f$ with $c_\mu(w)\defeq\sum_{w'}\xi^3\Sigma_2(w,w')(w'-w)_\mu$ the first moment of $S_2$; its all-$C^\xi$ part is bounded by \cite[(3.16)]{Ba3} (the cancellation $\sum_{y'}(\partial C^\xi\partial^*)(y'-y)=0$ of \cref{thm:a06d:S48}), and its $R^v$-part by $|\partial R^v\partial^*|\le C_R(1+\bar v)^3(\xi+\rho)^{-1}e^{-\delta_R\rho}$ (\cref{lem:a06:WN1}(d)) against the other line of the bubble, a sum convergent in $d=3$; so $|c_\mu(w)|\le C(1+\bar v)^3$, and $\sum_w\xi^3|F(w)|\le C(1+\bar v)^3$ by the bounds on $\partial C^\xi$ and $\partial R^v$. The H\"older remainder $f(w')-f(w)-(w'-w)\cdot\partial f(w)$ gives the node $S$ a differentiated external leg, after which every component has positive degree and \cite[(2.16)]{Ba3} bounds the sum. If the bridge attaches at $S_2$'s other leg and the anchor convention puts the counterterm at the far leg, the same identity holds with $f$ and $F$ exchanged, by the symmetry $\Sigma_2(w,w')=\Sigma_2(w',w)$ of the real symmetric scalar kernels at zero field (\cref{lem:a06:WN1}(a)). Altogether the $R$-part of the pair is $\le C(1+\bar v)^p=: C^v_{\mathrm{BR}}$, uniformly in $\xi$, $j_0$, $s$, $\eps$, the labels, the torus and $x$. The bridged shapes without a counterterm member (\cite[(2.21a), (2.21c), (2.21e)]{Ba3}: the one whose expression vanishes identically \cite[p.~442]{Ba3}, and those in which $S_2$ carries a difference on a leg or an extra power of $\eta$) need none: for them the all-$C^\xi$ part vanishes as above, and the $R$-part is convergent per member, since the difference on $S_2$'s leg, or the extra power of $\eta$, makes $D(S_2)\ge0$ before the gain of the $R$-line.
\end{proof}

\begin{remark}[reducible two-scalar-leg chains]\label{rem:a06:WNCH}
A one-particle-reducible two-scalar-leg chain of one-loop $D=-1$ parts is renormalised by the counterterms of its parts (Step~3(D$_4$)). The identity behind it, $T_{P_2}(b',c)=\delta_{b'c}\sum_{c'}K_{P_2}(b',c')$ with the same lines, indices and labels, holds for any kernels because the counterterm is computed with the model's own lines; so it holds verbatim in the walled model. The alternative route through translation invariance is not used.
\end{remark}

% a06_WN_d: part 4 of 4 of the walled reading of Ba3 Prop. 1 (input from sec/a06_BALtorus_d7b.tex): the proof of thm:a06:WN(ii), the constants, and two remarks on scope.

\begin{proof}[Proof of \cref{thm:a06:WN}\textup{(ii)}]
The proof of \cite[(1.33)]{Ba3} at $L=5$ for the bulk model, Steps~3(D$_3$)--3(D$_4$) with the twenty sketched steps listed after \cref{thm:a06d:S49}, is run for the walled model and for a class $\cC$ closed under the moves. It is organised as follows: first, where the scalar lines enter and what replaces each entry; second, where the class enters and what supplies each partner; third, the assembly and the induction.

\emph{Step 1: the scalar lines.} At zero background the scalar lines of every graph, slice, counterterm and cut kernel are zero-field kernels, and the proof uses them in exactly four ways. \emph{(1) Kernel bounds}: the group-(A) estimates (the slice bounds \cite[(2.6), (2.10)--(2.12)]{Ba3}, the kernel bounds of the cube step, the two-domain differences $\delta G_s(\square,\eta\Z^3,0)$, the cut-kernel estimate \cite[(3.1)]{Ba3} in the cluster form of \cref{cor:a06d:S24}, the Combes--Thomas bounds of the chains), each a composition of the pointwise bounds on the zero-field kernels with lattice sums; for walled lines these are \cref{lem:a06:WN1}(b),(c),(e), with constants polynomial in $\bar v$. The two-domain difference is bounded by the same resolvent identity as in the bulk, with \cref{lem:a06:WN1}(b) on each domain, since at zero field the bounds hold at every point of a Neumann domain and of the torus. \emph{(2) Algebra}: the slice identity, its rescaling, the gauge covariance used in the gauge-out of $B_0$, the cancellation of the operator in the differences of the $\tilde B'$-expansion and of the domain changes; for walled lines these are \cref{lem:a06:WN1}(a), and at zero background the $\tilde B'$-expansion, the constant field $B_0$ and its gauge-out are void. \emph{(3) The regrouped lines of the renormalisation}: in Step~3(D$_4$) the cumulative regrouped line of a node is, by the slice identity, the inverse of the operator pinned at the lowest external scale, which after rescaling to $\xi\Z^3$ is $\breve G_l=C^\xi+R_l$ (\cref{thm:a06d:S48}, \cref{lem:a06d:S47}); for walled lines the rescaled operator carries the potential $s^2v_{\mathrm{per}}(s\,\cdot)\in[0,\bar v]$ (\cref{lem:a06:WN1}(a),(f)), so $\breve G^v_l=C^\xi+R^v_l$ with the \emph{same} $C^\xi$ and with $R^v_l$ a remainder line of \cref{lem:a06d:S47}(b) with constant $C_R(1+\bar v)^3$ (\cref{lem:a06:WN1}(d)); hence every step of \cref{thm:a06d:S48} that concerns terms with at least one $R$-line (the one-loop $D=-1$ cancellation, which uses only the bounds on the $R$-line and the forward-difference structure of the $C^\xi$-line; the re-slicing of \cref{lem:a06d:S47}; the bound function estimate \cite[(2.16)]{Ba3}) runs verbatim, and every all-$C^\xi$ expression, with the inner counterterm members re-expanded with their own lines, is literally the same expression as in the bulk model \cite[pp.~437, 441]{Ba3}. \emph{(4) The counterterm values}: the model's own, by definition (the walled model). The torus-versus-$\xi\Z^3$ identification of \cite[pp.~418--419]{Ba3} is \cref{lem:a06:WN1}(f), with the image terms bounded, as in the bulk, by \cite[(1.33)]{Ba3} for the classes of lower order into which they decompose (the induction below). The sharp-indicator lemma \cref{lem:a06d:S46} uses only the decay of the internal lines, supplied by \cref{lem:a06:WN1}(b); the useless-counterterm lemma \cref{lem:a06d:S45} and the \cite[(2.4)]{Ba3} lemma \cref{lem:a06d:S44} are kernel-free given the bounds; \cref{lem:a06d:S41,lem:a06d:S42} are identities. Subgraphs with no scalar slot and an odd number of vector slots vanish at $B=0$ also with walled lines, which are scalar kernels times $\mathbf 1_N$ (\cref{lem:a06:WN1}(a)). A finite \cite[(1.7)]{Ba3} vertex with an $x$-dependent coefficient never sits inside a dangerous node (\cref{lem:a06d:S25}(b), \cite[Cor.~2.3]{Ba3}), so such coefficients enter only through their suprema, as the statement counts them.

\emph{Step 2: the class.} The class enters the bulk proof only at the dangerous nodes, as the sum over a node's partners with the rest of the graph fixed (\cref{lem:a06d:S41}(a): the counterterm pairing is an identity of class sums; \cref{thm:a06d:S48}: the local coefficient is the class sum at the node; \cref{thm:a06d:S28}), at the cut counterterms of partnered subgraphs (\cref{lem:a06d:S25}, which needs \cite[(1.33)]{Ba3} for the completed pieces of lower order), at the class-component cube (one label multiset), and in the regrouping after the cut (\cref{thm:a06d:S23}, kernel-free, (i)). Node by node, with walled lines: \emph{two-scalar-leg nodes} are paired with their counterterm member (P), the exact identities \cite[(3.19)]{Ba3} ($D=0$) and \cite[(3.11)]{Ba3} ($D=-1$) against the model's own counterterm with the same lines, the same localisation and (Agr); the $D=-1$ first moment has its all-$C^\xi$ part bounded by \cite[(3.16)]{Ba3} and its $R^v$-part by $|\partial R^v\partial^*|\le C(1+\bar v)^3(\xi+\rho)^{-1}e^{-\delta\rho}$, so $|c_\mu|\le C(1+\bar v)^3$; reducible chains by \cref{rem:a06:WNCH}. \emph{Two-vector-leg nodes}: the per-label two-vector-leg class at the node, which (W) supplies; the zeroth moment of its all-$C^\xi$ part vanishes by \cref{lem:a06d:S26}, applied as in \cref{thm:a06d:S48} to the full lines $C^\xi$ of the class sum (the lemma allows any $x$-dependent $|\phi|^2$ coefficients, so the inner counterterm and finite-vertex coefficients, and the potential itself, are admissible coefficients; \cref{lem:a06d:S26} is thus verbatim in the walled model), the first moment of its all-$C^\xi$ part by the reflection identity \cite[(3.30)]{Ba3} on $C^\xi$, which is model-free; the $R^v$-parts by the $R$-line steps of \cref{thm:a06d:S48}, with the zero-momentum remainder vertex kept. \emph{$\phi\phi A$ nodes}: in the bulk, \cref{thm:a06d:S48} bounds the zeroth moment by the reflection identity \cite[(3.37)]{Ba3} over the per-label full $\phi\phi A$ class; in a class closed under the moves this is replaced by (iii): at a bridgeless node the (X)-partner is in $\cC$ and (iii)(a) cancels the pair exactly at a common anchor, or, when the node carries the pinned root at one of its two slot vertices, (iii)(b) bounds the pair by $C^v_7$; at a bridged node no partner is needed, (iii)(c), the counterterm member (P) of $S_2$ being in $\cC$; and the pair bounds are of the form $|c^{(0)}|\le C$ that \cref{thm:a06d:S49} consumes for a $\phi\phi A$ local vertex. \emph{\cite[(2.4)]{Ba3} subgraphs}: summation by parts, \cref{lem:a06d:S44}, kernel-free. \emph{Vacuum components}: cancelled by the model's own $E_s[v]$ with the same labels and lines, up to image terms (Step~1). \emph{One- and three-vector-leg subgraphs}: vanish (Step~1). The cut counterterms of partnered subgraphs (\cref{lem:a06d:S25}) are classes of lower order of the completed pieces, which are closed under the moves (a move at a divergent subgraph of a piece is a move of the uncut member, and the completion of a partnered subgraph is by the moves); and the class-component cube depends on the label multiset only, common to $\cC$. Per-label full classes are closed under the moves, so the bulk statement for them is the special case $\cC=$ full class; and every partner that the walled node mechanisms use is in $\cC$ by the closure, while every other step is a bound per term.

\emph{Step 3: assembly and induction.} With Steps~1 and 2, \cref{thm:a06d:S48,thm:a06d:S49} hold for $\cC$ with walled lines, with the $\phi\phi A$ clause in the pair form of (iii); the induction on the number of lines inside \cref{thm:a06d:S48} and on the number of vertices in \cref{lem:a06d:S25} runs over classes closed under the moves of lower order (the completed pieces of a cut partnered subgraph, and the pieces into which an image term decomposes), to which the statement applies by the induction hypothesis; the base is the one-vertex graphs. The bound function estimate \cite[(2.16)]{Ba3} then gives \cite[(3.3)]{Ba3} for every class component, and the assembly of \cite[pp.~420--445]{Ba3} listed after \cref{thm:a06d:S49}, with the cluster cut and the cube steps, gives \cite[(1.33)]{Ba3} in the stated form, with $d_v$, $d_s$ counting the action couplings, the counterterm vertices as exact counterterms, and the finite vertices through their suprema (they have $D(v)=2$ whatever their coefficient, \cref{lem:a06d:S25}(b)); the localisation of an external field whose leg moves under the moves is (H$_A$); the slot fields of several variables are measured by \cite[(1.32)]{Ba3} in the slot variable, as the two-variable external fields of \cite[p.~419]{Ba3}; and in a two-point setting the root vertex is fixed at $x$ throughout, every kernel bound used being pointwise, so the bound is uniform in $x$.

\emph{Constants.} The walled inputs carry $(1+\bar v)^2$ (\cref{lem:a06:WN1}(b)), $(1+\bar v)^3$ (\cref{lem:a06:WN1}(d)) and the constants of \cref{lem:a06:WN1}(e), polynomial in $\bar v$; a graph of order $\le\bar n$ has a number of lines bounded by $\bar n$, so $O(1)^v$ is a polynomial in $\bar v$ of degree bounded in terms of $\bar n$, and $\delta^v_0$ is a fixed fraction of $\delta^w$, $\delta_R$ and the bulk rates. The parameter rows of Step~3(D) are kept, with $R_0$ replaced by $R^v_0\defeq R_0+10$ in every depth row for a scalar propagator bound and the torus side $\ge\max\{6\sqrt3\,r(t_s)+2,8/\delta^v_0\}$, all lower bounds on $R$, $r^\flat$ and the torus side or caps on $\eps_0$ of Ba{\l}aban's kind, fixed after $R_0$ and before $\eps_0$; nothing in the order of the choices changes. The constants are therefore independent of $\eps$, of $\eps_0\in(0,\eps^*_0]$ below the caps, of $s$, $K$, the labels, the torus side, $x$, and of $v$ beyond $\bar v$, and they depend on $L=5$, on the parameters of Step~3(D) and polynomially on $\bar v$.
\end{proof}

\begin{remark}[non-circularity of the uses in the rotation part]\label{rem:a06:WNcirc}
\Cref{thm:a06:WN}(ii) is a bound in which the suprema of the finite-vertex coefficients are parameters. Its uses in the rotation part at coupling order $4$ carry, as finite vertices, only the non-renormalising parts of the order-$2$ inner counterterms, whose sizes come from the order-$2$ statement, and the wall vertex, whose coefficient is $\le nt_s^2$; the sizes of the order-$4$ finite vertices are then \emph{consequences} of (ii) at order $4$, not inputs to it. The induction of the proof runs on the number of vertices inside the theorem.
\end{remark}

\begin{remark}[a bounded wall is not a boundary]\label{rem:a06:WNbdry}
\Cref{thm:a06:WN} is stated for a bounded site potential on the scalar lines with the vector lines unwalled, on a torus. It does not cover scalar or vector lines with Dirichlet or Neumann boundary conditions on a box, and the proof of (iii)(b) does not extend to them: a boundary changes the kernel by image terms $C^\xi(x-y^*)$ which, for $x$ and $y$ in the boundary layer, are as singular as $C^\xi$ itself and so are not remainder lines in the sense of \eqref{eq:a06:WNR}; with a Neumann vector line the odd part of the pinned pair is left uncancelled at every scale, and nothing bounds $|c_{S,1}-c_{S,2}|$ uniformly in $\xi$ in the layer. A bounded potential changes the kernel by $R^v$, which satisfies \eqref{eq:a06:WNR} at \emph{every} pair of points (\cref{lem:a06:WN1}(d); the multiplier $1-w$ is never differentiated, and no smoothness of $v$ is used), so there is no boundary layer, and the vector line is the translation-invariant torus propagator. This is why the rotation part compares King's cubes with walled tori and never applies the renormalisation theorem on a cube.
\end{remark}


% a06_BALtorus_ct1: Step 3(E), first part: the per-graph bound on the mass counterterm, |dm^2| <= c_m eps^{-1}
% (Ba2 p. 559, Ba3 (1.21)-(1.23)), by momentum-space power counting on the torus. Input from sec/a06_BALtorus_c.tex.

\paragraph{Step 3(E): the sizes of the mass counterterm.} Three statements about the mass counterterm $\dm$ are asserted in \cite{Ba2,Ba3} without proof and are consumed by every step of the chain: $\dm=O(\eps^{-1})$ \cite[p.~559]{Ba2}, and the two finite-part statements of \cite[p.~418]{Ba3}, $\dm_{K,k}(x)=O((L^k\eps)^{-1})$ and ``$\dm_{\mathrm{fin}}$ convergent as $\eps\to0$'' (``Future considerations will imply''), which enter the proof of Proposition~1 only through the size of the non-renormalising part of the \cite[(1.7)]{Ba3} vertex \cite[p.~421]{Ba3}. This step proves the first in the form stated, the second in the form stated, and the third in the form in which it is consumed: in scale-$k$ units the non-renormalising coefficient $c_k\defeq t_k^2(\dm_{\mathrm{fin}}+\dm_{K,k})$ is $O(t_k(1+\ln t_k^{-1})^5)$ with a constant free of $\eps_0$ (\cref{thm:a06e:T0}), which every consumer absorbs (\cref{lem:a06e:T1}); $\dm_{\mathrm{fin}}$ itself is bounded by a constant depending on $\eps_0$ (\cref{cor:a06e:T0p}).

\subsubsection*{The per-graph bound $\dm=O(\eps^{-1})$}

\emph{Setting} \cite[pp.~605--606]{Ba1}, \cite[pp.~413, 416--417]{Ba3}. $d=3$, $T_\eps\defeq\prod_\mu(\eps\Z/\ell_\mu\Z)$ with all sides $\ell_\mu\ge1$, $\ell_\mu/\eps\in\N$ (write $\ell_T^{-3}$ for $\prod_\mu\ell_\mu^{-1}$); functions carry the $\eps^3$-weighted product and $\hat f(p)\defeq\sum_x\eps^3f(x)e^{-ipx}$ for $p\in T^*\defeq\prod_\mu(2\pi/\ell_\mu)\Z\cap B_\eps$, $B_\eps\defeq[-\pi/\eps,\pi/\eps)^3$. For $p\in\R^3$, $[p]\in B_\eps$ is its representative modulo $(2\pi/\eps)\Z^3$. $(\partial_\mu f)(x)\defeq\eps^{-1}(f(x+\eps e_\mu)-f(x))$ has symbol $\hat\partial_\mu(p)=\eps^{-1}(e^{i\eps p_\mu}-1)$, and $\hat\Delta(p)\defeq\sum_\mu4\eps^{-2}\sin^2(\eps p_\mu/2)=\sum_\mu|\hat\partial_\mu(p)|^2$. The free covariances are $C_0\defeq(-\Delta+m^2)^{-1}$ (scalar, $m>0$) and $C\defeq(-\Delta+\lmu^2)^{-1}$ (vector, Feynman gauge, componentwise \cite[(1.11)]{Ba1}). By \cite[(1.20)--(1.21)]{Ba3}, $\dm=\sum_{2\le\alpha+2\beta\le4}e^\alpha\lambda^\beta\dm_{(\alpha,\beta)}$ is fixed graph by graph: $\dm_G\defeq\sum_{x'\in T_\eps}\eps^3\Sigma^\eps_G(x-x')$, $\Sigma^\eps_G$ the amputated one-particle-irreducible two-point graph $G$ (no difference on its external legs) with the vertices of \cite[(1.6)--(1.8), (1.10)]{Ba3} at $\eta=\eps$, $L^k\eps=1$, $\tilde B=0$, $g_k=1$ \cite[(1.22)--(1.23)]{Ba3}; $\dm$ is an $N\times N$ matrix polynomial in $q^2$, treated as a number. The vertices, the Taylor terms of $\frac12|D_A\phi|^2-\frac12|\partial\phi|^2=-\eps^{-2}\phi(b_-)\cdot(U(A_b)-1)\phi(b_+)$, are: $(V_\lambda)$ $-\lambda\sum_x\eps^3|\phi|^4$, order $2$; $(V_\delta)$ $-\frac12\sum_x\eps^3\dm_1|\phi|^2$ with $\dm_1$ the order-$2$ part of $\dm$, order $2$; $(V^8_n)$ $e^n(-1)^n\eps^{n-1}(n!)^{-1}\sum_b\eps^3[(\partial\phi)(b)\cdot q^n\phi(b_-)](A_b)^n$, $n\ge1$, order $n$; $(V^{10}_n)$ $e^n\eps^{n-2}(n!)^{-1}\sum_b\eps^3[\phi(b_-)\cdot q^n\phi(b_-)](A_b)^n$, $n\ge2$ even (for odd $n$, $q^n$ is antisymmetric and the term vanishes), order $n$. The order of $G$ is $\omega_G\defeq\sum_v\mathrm{ord}(v)=\alpha+2\beta$. The $\eps$-exponent of a vertex is $s_v\defeq0$ $(V_\lambda)$, $n-1$ $(V^8_n)$, $n-2$ $(V^{10}_n)$, and $s_v\defeq-1$ for $V_\delta$, the last used only after $|\dm_1|\le C\eps^{-1}$ is known; $s(G)\defeq\sum_vs_v$.

\begin{lemma}[symbols]\label{lem:a06e:M0}
$\frac4{\pi^2}|[p]|^2\le\hat\Delta(p)\le|[p]|^2$, $|\hat\partial_\mu(p)|\le|[p]|$, and $\hat C_0(p)=(\hat\Delta(p)+m^2)^{-1}\le\frac{\pi^2}4(|[p]|^2+m^2)^{-1}$, the same for $\hat C$ with $\lmu$.
\end{lemma}

\begin{proof}
All three are periodic, so take $p=[p]$; use $|\sin t|\ge2|t|/\pi$ for $|t|\le\pi/2$ and $|\sin t|\le|t|$; for the last claim, $\frac4{\pi^2}|p|^2+m^2\ge\frac4{\pi^2}(|p|^2+m^2)$.
\end{proof}

\begin{lemma}[Riemann domination]\label{lem:a06e:M1}
Let $F\colon B_\eps^{\,n}\to[0,\infty)$ satisfy $F(k)\le C_1F(k')$ whenever $|k_{i,\mu}-k'_{i,\mu}|\le\pi/\ell_\mu$ for all $i,\mu$. Then $C_1^{-1}(2\pi)^{-3n}\int_{B_\eps^n}F\le\ell_T^{-3n}\sum_{k\in(T^*)^n}F(k)\le C_1(2\pi)^{-3n}\int_{B_\eps^n}F$. For $F=\prod_cg_c(q_c(k))$ with $g_c(q)=(|[q]|^2+m_c^2)^{-\beta_c}$, $\beta_c\ge0$, each $q_c$ one of $\pm k_i$, $\pm(k_1+k_2)$, the hypothesis holds with $C_1=\prod_c(2+242/m_c^2)^{\beta_c}$.
\end{lemma}

\begin{proof}
The boxes of sides $2\pi/\ell_\mu$ centred at the points of $T^*$ tile $B_\eps$, and $\ell_T^{-3}=(2\pi)^{-3}|\text{box}|$; compare $F$ at the centre with $F$ on the cube. For the product form, $|[q_c(k)]-[q_c(k')]|\le2\sqrt3\pi/\min_\mu\ell_\mu\le11$, and for $|\delta|\le11$, $(|q'|+|\delta|)^2+m^2\le2|q'|^2+242+m^2\le(2+242/m^2)(|q'|^2+m^2)$.
\end{proof}

\begin{lemma}[radial integrals]\label{lem:a06e:M2}
For $\beta\in\R$, $m>0$, $0<\eps\le1$, $r>0$ let $I(\beta)\defeq\int_{B_\eps}(|p|^2+m^2)^{-\beta}dp$ and $H_\beta(r)\defeq\int_{B_\eps\cap\{|p|\le r\}}(|p|^2+m^2)^{-\beta}dp$. With $C=C(\beta,m)$: if $\beta<\frac32$, $I(\beta)\le C\eps^{2\beta-3}$ and $H_\beta(r)\le C(1+r)^{3-2\beta}$; if $\beta=\frac32$, $I\le C(1+\ln\eps^{-1})$ and $H\le C(1+\ln(1+r))$; if $\beta>\frac32$, $I,H\le C$.
\end{lemma}

\begin{proof}
$B_\eps\subset\{|p|\le\sqrt3\pi/\eps\}$; in polar coordinates $(s^2+m^2)^{-\beta}\le\max(m^{-2\beta},(1+m^2)^{-\beta})$ on $s\le1$ and $\le Cs^{-2\beta}$ on $s\ge1$; $\int_1^{\max(\rho,1)}s^{2-2\beta}ds$ gives the three cases.
\end{proof}

\begin{lemma}[structure and degree]\label{lem:a06e:M3}
Let $G$ be a graph of $\dm$ of order $\omega_G\le4$. Then: \textup{(a)} the loop number is $\mathcal L\defeq I-V+1=(\omega_G-\omega_{\mathrm{ins}})/2$, $\omega_{\mathrm{ins}}$ the total order of the $V_\delta$ vertices; so $\mathcal L\le2$, and a $V_\delta$ vertex occurs only in one-loop graphs of order $4$, exactly one, of order $2$. \textup{(b)} The internal-line graph $G'$ is connected and bridgeless; for $\mathcal L=1$ it is a cycle, for $\mathcal L=2$ a theta (two vertices of degree $3$ joined by three internally disjoint chains) or a figure-eight (one vertex of degree $4$ on two cycles); all other vertices have degree $2$ in $G'$. \textup{(c)} At zero external momentum all lines of a chain (or cycle) carry one momentum $q_c$ ($k_1,k_2,k_1+k_2$ for a theta; $k_1,k_2$ for a figure-eight; $k_1$ for a cycle), and $|\dm_G|\le C_G\eps^{s(G)}\ell_T^{-3\mathcal L}\sum_{k\in(T^*)^{\mathcal L}}\prod_cg_c(q_c)$ with $g_c(q)=(|[q]|^2+m_c^2)^{-\beta_c}$, $\beta_c\defeq\sum_{l\in c}(1-d_l/2)\ge0$, $d_l\in\{0,1,2\}$ the number of differences acting on the line $l$, $m_c\in\{m,\lmu\}$ on a chain of one line type and $m_c\defeq\min(m,\lmu)$ on a mixed chain. \textup{(d)} With $\omega(G)\defeq3\mathcal L-2\sum_l(1-d_l/2)-s(G)$ and, for a cycle $H$ of $G'$, $\omega(H)\defeq3-2\sum_{l\in H}(1-d_l/2)-s(H)$: $\omega(G)=2-\omega_G/2$, and every cycle $H\ne G'$ of a two-loop $G'$ has $\omega(H)\le1$, with equality only if $H$ is a two-point subgraph of order $2$.
\end{lemma}

\begin{proof}
(a) $\sum_vh_v=2I+2$ ($h_v$ the number of legs), so $\mathcal L=\sum_v(h_v/2-1)$, and $h_v/2-1=\mathrm{ord}(v)/2$ for $V_\lambda$ ($h=4$), $V^8_n$, $V^{10}_n$ ($h=n+2$) and $0$ for $V_\delta$ ($h=2$). (b) $G$ connected with external legs at vertices gives $G'$ connected; 1PI gives no bridge, so no vertex of degree $\le1$; $\sum_v(\deg_{G'}v-2)=2(\mathcal L-1)$. For $\mathcal L=2$ with two vertices $u,w$ of degree $3$, a path leaving $u$ and returning to $u$ would make the connection to $w$ a bridge, so all three paths end at $w$. (c) Momentum conservation at the vertices with zero external momenta; per line $|\hat C(q)||\hat\partial(q)|^{d_l}\le\frac{\pi^2}4(|[q]|^2+m_c^2)^{-1+d_l/2}$ by \cref{lem:a06e:M0}; vertex phases have modulus $1$, $\norm q\le1$. (d) With $\mathcal L=I-V+1$ and $2I=\sum h_v-2$, $\omega(G)=\sum_v(h_v/2-3-s_v+d_v)+2$, $d_v=1$ for $V^8$ vertices (the difference acts on an internal line) and $0$ otherwise; the bracket is $-\mathrm{ord}(v)/2$ for every type ($-1$, $-1$, $(n+2)/2-3-(n-1)+1=-n/2$, $(n+2)/2-3-(n-2)=-n/2$). For a cycle $H$ (as many lines as vertices), $\omega(H)=3-\mathrm{ord}(H)/2-E_H/2-\#\{V^8\text{ vertices of }H\text{ whose difference acts off }H\}$, $E_H\defeq\sum_{v\in H}(h_v-2)$: the bracket per vertex is $-1-e_v/2$ ($V_\lambda$, $e_v=2$), $-n/2-e_v/2-(1-d_{v,H})$ ($V^8_n$), $-n/2-e_v/2$ ($V^{10}_n$), $-1-e_v/2$ ($V_\delta$). For a proper cycle of a theta or a figure-eight, $E_H\ge2$ (the third chain, resp.\ the other cycle, leaves $H$). $\mathrm{ord}(H)\ge2$: a cycle through $\ge2$ vertices has order $\ge2$; a self-loop can sit only at the junction of a figure-eight, which has $\ge4$ internal legs, so it is $V_\lambda$, $V^8_n$ or $V^{10}_n$ with $n\ge2$. Hence $\omega(H)\le1$, with equality iff $\mathrm{ord}(H)=2$, $E_H=2$ and no difference leaves $H$.
\end{proof}

\begin{lemma}[one loop]\label{lem:a06e:M4}
If $\mathcal L(G)=1$, then $|\dm_G|\le C\eps^{s(G)}I(\beta)$ with $\beta$ the $\beta_c$ of the cycle, and $|\dm_G|\le C\eps^{-1}$.
\end{lemma}

\begin{proof}
The first bound is \cref{lem:a06e:M1,lem:a06e:M3}(c). Then $\omega(G)=3-2\beta-s$ and \cref{lem:a06e:M2}. Order $2$ (no $V_\delta$, $s\ge0$, $\omega=1$): $s=0$ gives $\beta=1$, $I\le C\eps^{-1}$; $s\ge1$ gives $\beta=1-s/2<\frac32$ and $\eps^sI\le C\eps^{s+2\beta-3}=C\eps^{-1}$. Order $4$ with one $V_\delta$ ($s=s'-1$, $s'\ge0$ the sum over the other vertices, $\omega=0$, $3-2\beta=s'-1$): $s'=0$ gives $\beta=2$ and $C\eps^{-1}$; $s'=1$ gives $\beta=\frac32$ and $C(1+\ln\eps^{-1})$; $s'\ge2$ gives $\eps^{s'-1}\eps^{2\beta-3}=1$.
\end{proof}

\begin{lemma}[theta]\label{lem:a06e:M5}
If $G'$ is a theta with chains $c=1,2,3$, then $|\dm_G|\le C\eps^{-1}$.
\end{lemma}

\begin{proof}
By \cref{lem:a06e:M1,lem:a06e:M3}(c), $|\dm_G|\le C\eps^s\,T$ with $T\defeq\iint_{B_\eps^2}g_1(k)g_2(l)g_3(k+l)\,dk\,dl$; at order $4$ with $\mathcal L=2$ there is no $V_\delta$, so $s\ge0$. (1) For $a\in\{1,2,3\}$ let $S_a$ be the set where $|[q_a]|$ is the smallest of $|[k]|,|[l]|,|[k+l]|$, $b,c$ the other indices and $q_*$ the one of $q_b,q_c$ of smaller norm; each $g$ decreases in the norm, so on $S_a$ the integrand is $\le g_a(q_a)g_b(|q_*|)g_c(|q_*|)$ with $|[q_a]|\le|[q_*]|$. (2) Any two of $k,l,k+l$ are coordinates with unit Jacobian, so $\iint_{S_a}\le2\int_{B_\eps}g_b(q)g_c(q)H_{\beta_a}(|q|)dq$. (3) Put $\omega^0\defeq6-2(\beta_1+\beta_2+\beta_3)$, $\omega^0_{bc}\defeq3-2(\beta_b+\beta_c)$; then $\omega(G)=\omega^0-s=0$ and $\omega(H_{bc})=\omega^0_{bc}-s(H_{bc})\le1$ by \cref{lem:a06e:M3}(d), with $s(H_{bc})\le s$. (4) $\beta_a<\frac32$: $H\le C(1+|q|)^{3-2\beta_a}\le C(|q|^2+m^2)^{(3-2\beta_a)/2}$, so the sector is $\le C\,I(\beta')$ with $3-2\beta'=\omega^0=s$, and $\eps^s$ times it is $\le C(1+\ln\eps^{-1})$. (5) $\beta_a=\frac32$: $H\le C(1+\ln\eps^{-1})$ and $\omega^0_{bc}=\omega^0=s$, so the sector is $\le C(1+\ln\eps^{-1})^2$. (6) $\beta_a>\frac32$: $H\le C$ and $\eps^sI(\beta_b+\beta_c)\le C\max(\eps^{s-\omega^0_{bc}},1+\ln\eps^{-1})$ with $\eps^{s-\omega^0_{bc}}\le\eps^{s(H_{bc})-\omega^0_{bc}}=\eps^{-\omega(H_{bc})}\le\eps^{-1}$. Sum the six pieces.
\end{proof}

\begin{lemma}[figure-eight]\label{lem:a06e:M6}
If $G'$ is a figure-eight with cycles $1,2$, then $|\dm_G|\le C\eps^sI(\beta_1)I(\beta_2)\le C\eps^{-1}$.
\end{lemma}

\begin{proof}
The cycles carry independent momenta; $\omega(G)=6-2\beta_1-2\beta_2-s=0$ and $\omega(H_i)\le1$. Both $\beta_i<\frac32$: $\eps^{s+2\beta_1-3+2\beta_2-3}=\eps^{-\omega(G)}=1$. $\beta_1<\frac32=\beta_2$: $\eps^{-\omega(G)}(1+\ln\eps^{-1})$. $\beta_1<\frac32<\beta_2$: $\eps^{s-(3-2\beta_1)}\le\eps^{-\omega(H_1)}\le\eps^{-1}$, using $s\ge s(H_1)$. Both $\ge\frac32$: $\le C(1+\ln\eps^{-1})^2$.
\end{proof}

\begin{theorem}[the mass counterterm per graph]\label{thm:a06e:CT1}
For every graph $G$ of $\dm$ (\cite[(1.22)--(1.23)]{Ba3}: all graphs of order $\alpha+2\beta\le4$, $d=3$), $|\dm_G|\le C_G\eps^{-1}$ with $C_G=C_G(m,\lmu,N)$, uniformly in $0<\eps\le1$ and in the torus (all sides $\ge1$). Hence $|\dm|\le c_m\eps^{-1}$ with $c_m\defeq\sum_Ge^\alpha\lambda^\beta C_G$, a finite sum; $c_m$ does not depend on $\eps_0$.
\end{theorem}

\begin{proof}
Order $2$: by \cref{lem:a06e:M3}(a) these are one-loop graphs without $V_\delta$, so \cref{lem:a06e:M4} gives $|\dm_1|\le C\eps^{-1}$, which justifies $s_v=-1$ for the $V_\delta$ vertices at order $4$. Order $4$: by \cref{lem:a06e:M3} each graph is a cycle with one $V_\delta$ (\cref{lem:a06e:M4}), a theta (\cref{lem:a06e:M5}) or a figure-eight (\cref{lem:a06e:M6}). At order $\le4$, $n\le4$ and $V\le4$, so there are finitely many graphs.
\end{proof}

The proof uses the momentum-space form: in $x$-space the photon--scalar bubble line $\partial C_0\partial^*$ has a logarithmically divergent $\ell^1$ norm, while its symbol $|\hat\partial|^2\hat C_0$ is bounded (\cref{lem:a06e:M0}). Not every order-$4$ graph is $O(\ln\eps^{-1})$: graphs with an order-$2$ two-point subgraph (the last graph of \cite[(1.22)]{Ba3}, the tadpole on a tadpole of \cite[(1.23)]{Ba3}) are of size $\eps^{-1}$. The third graph of \cite[(1.22)]{Ba3} is \cref{lem:a06e:M6} at the junction $V^{10}_4$: $\eps^2C_0(0)^2=O(1)$, Ba{\l}aban's ``except the third''.

\begin{corollary}[{\cite[(2.11)--(2.13)]{Ba2}}]\label{cor:a06e:C1}
The bounds \cite[(2.11)--(2.13)]{Ba2} on the counterterm vertices hold for Ba{\l}aban's torus $\dm$ with constants free of $\eps$, $\eps_0$ and the torus.
\end{corollary}

\begin{proof}
They are stated in Step~3 as consequences of $|\dm|\le c\eps^{-1}$ by exact computation; \cref{thm:a06e:CT1} supplies the hypothesis.
\end{proof}

% a06_BALtorus_ct2: Step 3(E), second part: the block kernels at zero background (slice ladder, remainder, order-2
% infrared remainders), the multi-cut expansion, the cut-line mechanism, and dm^2_{K,k} = O((L^k eps)^{-1}) (Ba3 (1.28)).

\subsubsection*{The block kernels at zero background}

\emph{Setting} \cite[(1.25)--(1.28), p.~418]{Ba3}. On a torus $T_\eps$ aligned with the blocks (all sides $\ge1$), at $A=0$ and in $\eps$-units, put $l_j\defeq L^j\eps=t_j$, $G_j\defeq(C_0^{-1}+a_jl_j^{-2}P_j)^{-1}$ for $1\le j\le K$, with $P_j$ the orthogonal projection onto the functions constant on the blocks of side $l_j$ (so $G_K$ is the operator of \cref{lem:a06:P}) and $G_0\defeq0$; $S_j\defeq G_{j+1}-G_j$ (the slice of index $j$; $S_0=G_1$) and $R_j\defeq C_0-G_j$. Then $G_K=G_k+G_{K,k}$ with $G_{K,k}=\sum_{j=k}^{K-1}S_j$, which is \cite[(1.27)]{Ba3}, and $R_K$ is the ``operator with a regular kernel'' of \cite[(1.25)]{Ba3}: $R_K=G_K(a_Kl_K^{-2}P_K)C_0=C_0(a_Kl_K^{-2}P_K)G_K$. The vector kernels $G^v_j$, $S^v_j$, $R^v_j$ are the same with $\lmu$ in place of $m$, componentwise (Feynman gauge; \cite[p.~608]{Ba1}: at $A=0$ the vector transformations are the scalar ones with $N=d$). $\dm_K$, $E_K$, $\dm_k$, $E_k$ are ``the same expressions'' \cite[(1.26), (1.28)]{Ba3} with $G_K$, resp.\ $G_k$, in place of $C_0$ and $C$: the \cite[(1.7)]{Ba3} vertices inside them keep their numerical coefficients $\dm_1$, $\dm_{(4)}$, and only the lines are replaced (\cref{conv:a06d:LOC}; in \cite[(1.23)]{Ba3} $\dm_1$ is a symbol for a number). Write $\dm[G,G^v]$, $E[G,G^v]$ for the class sums \cite[(1.23)--(1.24)]{Ba3} with the lines $(G,G^v)$, so that $\dm=\dm[C_0,C]$, $\dm_K=\dm[G_K,G^v_K]$, $\dm_{\mathrm{fin}}=\dm-\dm_K$, $\dm_{K,k}=\dm_K-\dm_k$, and likewise for $E$. Superscripts $(2)$, $(4)$ denote the parts of order $\alpha+2\beta=2$, $4$ (the orders $(3,0)$, $(1,1)$ vanish, $\alpha$ odd).

\begin{lemma}[operator monotonicity]\label{lem:a06e:B0}
For $1\le k\le K$ and every $x$: \textup{(a)} $C_{M_k}\le G_k\le C_0$ as operators, $C_M\defeq(-\Delta+m^2+M^2)^{-1}$, $M_k^2\defeq al_k^{-2}$; \textup{(b)} $0\le R_k(x,x)\le c_Rl_k^{-1}$ and $|G_{K,k}(x,x)|\le c_Rl_k^{-1}$, $c_R\defeq C(m)a^{1/2}$; \textup{(c)} the same for the vector kernels.
\end{lemma}

\begin{proof}
(a) $0\le a_kl_k^{-2}P_k\le M_k^2$ and inversion is operator-monotone decreasing. Diagonal entries are monotone ($K(x,x)=\ip{e_x}{Ke_x}$, $e_x\defeq\eps^{-3}\mathbf 1_x$), so $C_{M_k}(0)\le G_k(x,x)\le C_0(0)$ and $0\le R_k(x,x)\le C_0(0)-C_{M_k}(0)=\ell_T^{-3}\sum_pM_k^2\hat C_0(p)(\hat\Delta(p)+m^2+M_k^2)^{-1}\le CM_k^2\int_{\R^3}((p^2+m^2)(p^2+M_k^2))^{-1}dp\le C'M_k$ (\cref{lem:a06e:M1}). (b) $G_{K,k}(x,x)=G_K(x,x)-G_k(x,x)$ lies between $C_{M_K}(0)-C_0(0)\ge-c_Rl_K^{-1}$ and $C_0(0)-C_{M_k}(0)$, and $l_K\ge l_k$.
\end{proof}

\begin{lemma}[the slice ladder]\label{lem:a06e:SL}
Let $\sigma,\sigma'\in\{0,1\}$ count forward differences in the first and second variable, $y\defeq x-x'$ (torus distance). \textup{(a)} $G_k=\sum_{j<k}S_j$ and $R_k=R_K+\sum_{j=k}^{K-1}S_j$. \textup{(b)} For $0\le j<K$, $|\partial^\sigma\partial'^{\sigma'}S_j(x,x')|\le C_Kl_j^{-1-\sigma-\sigma'}e^{-\delta_K|y|/l_j}$, with $C_K$, $\delta_K$ free of $j$, $\eps$, $K$ and the torus. \textup{(c)} $|\partial^\sigma\partial'^{\sigma'}G_k(x,x')|\le C(\eps+|y|)^{-1-\sigma-\sigma'}e^{-\delta_K|y|/(2l_k)}$, and $|\partial^\sigma\partial'^{\sigma'}C_0(x,x')|\le C(\eps+|y|)^{-1-\sigma-\sigma'}e^{-2c_*|y|}$ with $c_*=c_*(m)>0$ (several differences in one variable allowed, $\sigma+\sigma'\le3$), on every torus with sides $\ge1$. \textup{(d)} The same for the vector kernels.
\end{lemma}

\begin{proof}
(a) is telescoping. (b) For $j\ge1$ the slice identity \cite[(1.27)]{Ba3}: $S_j=a_j^2l_j^{-4}G_jQ_j^*C^{(j)}Q_jG_j$, with $Q_j$ the block average and $C^{(j)}=((\Delta^{(j)}+aL^{-2}P)|_{T^{(j)}})^{-1}$ the step-$j$ fluctuation covariance. In scale-$j$ units the legs $G_jQ_j^*\delta_{\hat y}$, their differences and H\"older quotients (over separations $\le2$) obey $c_0e^{-\delta_0\dist(\hat x,B_j(\hat y))}$, with $c_1(\alpha)e^{2\delta_0}$ in place of $c_0$ for the H\"older quotients of exponent $\alpha$ (the printed bound decays in $\dist(\{\hat x,\hat x'\},\hat y)\ge\dist(\hat x,B_j(\hat y))-2$), by \cite[Lemma~2.4, (2.35)--(2.36)]{Ba4} on the periodic box $T^{(j)}$ (a union of unit blocks, the torus being aligned with the blocks), in the form of \cref{thm:a06:A5}: the constants are free of $j$ and of the torus, and no depth condition arises (this is \cite[Prop.~2.1, (2.24)--(2.25)]{Ba1} at $A=0$ with the source $Q_j^*\delta_{\hat y}$, $\norm{Q_j^*\delta_{\hat y}}_\infty=1$); and $|C^{(j)}(\hat y,\hat y')|\le c_2e^{-\theta_0|\hat y-\hat y'|}$ by \cite[(2.37)]{Ba4} with the $j$-free constants of \cref{thm:a06:A5}, i.e.\ \cref{lem:a06A:blockinv}(c) at $A=0$ with $\Lambda=T^{(j)}$, whose proof uses, of the geometry in (H1), only that $T^{(j)}$ is a union of $L$-blocks (the blocks of scale $j+1\le K$); keep half of each exponent for the two block sums (\cref{lem:a06d:V11}) and use $|\hat x-\hat x'|\le\dist(\hat x,B(\hat y))+|\hat y-\hat y'|+\dist(\hat x',B(\hat y'))+2$; rescaling to $\eps$-units gives $l_j^{-1}$ for the kernel and $l_j^{-1}$ per difference, and $\delta_K\defeq\min(\delta_0,\theta_0)/2$. For $j=0$, $S_0=G_1$ in scale-$1$ units has the fixed spacing $1/L$: the $\ell^2$ decay of $G_1$ at rate $\theta_0$ (the Combes--Thomas bound of group (A), \cref{lem:a06A:CT} for $T=G_1^{-1}\ge\gamma_\emptyset$ by \cref{lem:a06A:coer}, at the fixed spacing $1/L$ and with no depth restriction) gives the pointwise bounds with a constant depending on $L$, and the rescaling is as before. (c) Sum (b) over $j<k$: the slices with $l_j\ge|y|$ form a geometric sum $\le C\max(\eps,|y|)^{-1-\sigma-\sigma'}$, and those with $l_j<|y|$ give $l_j^{-a}e^{-\delta_K|y|/l_j}\le C_a|y|^{-a}(l_j/|y|)$, again geometric; split $e^{-\delta_K|y|/l_j}\le e^{-\delta_K|y|/(2l_j)}e^{-\delta_K|y|/(2l_k)}$ for $j<k$. For $C_0$: on $\eps\Z^3$ shift the contour of the Fourier integral of $(\hat\Delta(p)+m^2)^{-1}$ by $i\tau$, $\tau\defeq m/4$, in the direction of the largest coordinate of $y$ (for $\eps\tau\le1$, $\Re\hat\Delta(p+i\tau e)\ge\hat\Delta(p)-1.1\tau^2$), together with the short-distance bound $\int_{B_\eps}|p|^{\sigma+\sigma'}(|p|^2+m^2)^{-1}dp\le C\eps^{-1-\sigma-\sigma'}$; the torus kernel is the periodisation, and the images at distance $\ge\ell_T/2\ge\frac12$ are summable by the mass. (d) is the same with $\lmu$.
\end{proof}

\begin{lemma}[the remainder]\label{lem:a06e:RK}
On every torus aligned with the blocks with all sides $\ge1$, for $1\le a\le K$ and $\sigma,\sigma'\in\{0,1\}$: $|\partial^\sigma\partial'^{\sigma'}R_a(x,x')|\le C_{R}\max(l_a,|y|)^{-1-\sigma-\sigma'}e^{-c_*|y|}$, $C_R=C(m,L,a,C_K,\delta_K)$, free of $\eps$, $a$, $K$ and the torus; the same for $R^v_a$.
\end{lemma}

\begin{proof}
Case $|y|\ge l_a$: $R_a=C_0-G_a$, and \cref{lem:a06e:SL}(c) bounds both kernels by $C(\eps+|y|)^{-1-\sigma-\sigma'}$ times $e^{-2c_*|y|}$, resp.\ $e^{-\delta_K|y|/(2l_a)}\le e^{-\delta_K|y|/2}$ ($l_a\le1$); take $c_*\le\delta_K/4$. Case $|y|<l_a$: $R_a=G_a(a_al_a^{-2}P_a)C_0$ \cite[(1.25)]{Ba3}, whose kernel is $a_al_a^{-5}\sum_B(\int_B\partial^\sigma G_a(x,\cdot))(\int_BC_0\partial'^{\sigma'*}(\cdot,x'))$, the sum over the blocks $B$ of side $l_a$; by \cite[(2.35)]{Ba4} at scale $a$ in the form of \cref{thm:a06:A5} ($\mathbf 1_B=Q_a^*\delta_{\hat y}$ in scale-$a$ units; the torus is a periodic box that is a union of the blocks of side $l_a$), $|\int_B\partial^\sigma G_a(x,\cdot)|=|(\partial^\sigma G_a\mathbf 1_B)(x)|\le Cl_a^{2-\sigma}e^{-\delta_0\dist(x,B)/(4l_a)}$, by \cref{lem:a06e:SL}(c) $\int_B|C_0\partial'^{\sigma'*}(\cdot,x')|\le Cl_a^{2-\sigma'}$, and $\sum_Be^{-\delta_0\dist(x,B)/(4l_a)}\le C$; so the kernel is $\le Cl_a^{-1-\sigma-\sigma'}$, and $e^{-c_*|y|}\ge e^{-c_*}$.
\end{proof}

\begin{proposition}[the order-$2$ infrared remainders]\label{prop:a06e:DM1}
For $a,b\in\{1,\dots,K\}$ with $|a-b|\le1$ let $\Delta_{a,b}(x)\defeq\dm_1-\dm_1[G_a,G^v_b](x)$. Then $|\Delta_{a,b}(x)|\le C_Dl_{\min(a,b)}^{-1}$, $C_D=C(N,m,\lmu,e,\lambda)$, uniformly in $x$, $\eps$ and the torus. In particular ${\dm}^{(2)}_{\mathrm{fin}}=\Delta_{K,K}=O(l_K^{-1})$ and ${\dm}^{(2)}_{K,k}=\Delta_{k,k}-\Delta_{K,K}=O(l_k^{-1})$.
\end{proposition}

\begin{proof}
By \cite[(1.23)]{Ba3}, $\dm_1[G,G^v](x)=-4(N+2)\lambda G(x,x)+e^2dq^2G^v(x,x)+b(G,G^v)(x)$ with $b(G,G^v)(x)\defeq-e^2q^2\sum_{x'}\eps^3\sum_\mu(\partial_\mu G\partial'^*_\mu)(x,x')G^v(x,x')$. (1) Tadpoles: $-4(N+2)\lambda R_a(x,x)$ and $e^2dq^2R^v_b(x,x)$ are bounded by \cref{lem:a06e:B0}. (2) $b(C_0,C)-b(G_a,G^v_b)=b(R_a,C)+b(G_a,R^v_b)$ by bilinearity. (3) $b(R_a,C)$: by \cref{lem:a06e:RK} $|\partial\partial'R_a|\le C\max(l_a,|y|)^{-3}$ and $|C(x,x')|\le C(\eps+|y|)^{-1}$ (\cref{lem:a06e:SL}(c), no additive constant), so $\sum_{x'}\eps^3|\cdot|\le C[l_a^{-3}\int_{|y|\le l_a}(\eps+|y|)^{-1}dy+\int_{|y|\ge l_a}|y|^{-4}dy]\le Cl_a^{-1}$. (4) $b(G_a,R^v_b)$, zero sum: $\sum_{x'}\partial'^*_{x'}F(x')=0$ on the torus, so $R^v_b(x,x')$ may be replaced by $R^v_b(x,x')-R^v_b(x,x)$; integrating $|\partial'R^v_b|\le C\max(l_b,s)^{-2}$ along a lattice path gives $|R^v_b(x,x')-R^v_b(x,x)|\le C\min(|y|l_b^{-2},l_b^{-1})$. For $|y|\le l_b$, $|\partial\partial'G_a|\le C(\eps+|y|)^{-3}$ gives $Cl_b^{-2}\int_{|y|\le l_b}|y|^{-2}dy\le Cl_b^{-1}$; for $|y|>l_b$, $|\partial\partial'G_a|\le C(\eps+|y|)^{-3}e^{-\delta_K|y|/(2l_a)}$ gives $Cl_b^{-1}\int_{|y|>l_b}|y|^{-3}e^{-\delta_K|y|/(2l_a)}dy\le Cl_b^{-1}(1+\ln_+(l_a/l_b))\le C(1+\ln L)l_{\min(a,b)}^{-1}$. (Without the zero-sum step the bound would carry $\ln(l_b/\eps)$.)
\end{proof}

\subsubsection*{The multi-cut expansion and the cut-line mechanism}

Let $F[G]$ denote a class sum of \cite[(1.23)]{Ba3} (two-point, rooted at $x$) or \cite[(1.24)]{Ba3} (vacuum) with all lines in a kernel pair $G=(G,G^v)$, written as Proposition~1 writes it: the sum over the forest families $\{G_{\mathcal F}\}$ of its graphs, $\mathcal F$ a forest of divergent subgraphs, each maximal member contracted to its counterterm vertex, where a subgraph is a set of lines with the vertices it touches and its degree \cite[(2.1)]{Ba3} counts only its own lines, and the coefficients of the counterterm vertices are fixed numbers (\cref{conv:a06d:LOC}).

\begin{lemma}[the multi-cut expansion]\label{lem:a06e:MC}
Let $G=H+D$ be kernel pairs. \textup{(a)} $F[H+D]-F[H]=\sum_\Gamma\sum_{\emptyset\ne\Lambda\subset\mathrm{lines}(\Gamma)}(\Gamma\setminus\Lambda)_{\mathrm{ren}}[H;D\text{ on }\Lambda]$, where $(\Gamma\setminus\Lambda)_{\mathrm{ren}}\defeq\{G_{\mathcal F}:\mathcal F$ a forest of divergent subgraphs of $\Gamma$ none of which contains a line of $\Lambda\}$, every line not in $\Lambda$ carries $H$, and every line of $\Lambda$ is cut and carries $D$ as a two-variable external field. \textup{(b)} $(\Gamma\setminus\Lambda)_{\mathrm{ren}}$ is exactly the forest family of the cut graph $\Gamma\setminus\Lambda$, and it factorises over the connected components (clusters) of $\Gamma\setminus\Lambda$. \textup{(c)} With $H=G_{j+1}$, $D=-S_j$: $m_j\defeq F[G_{j+1}]-F[G_j]=\sum_\Gamma\sum_{\Lambda\ne\emptyset}(-1)^{|\Lambda|+1}(\Gamma\setminus\Lambda)_{\mathrm{ren}}[G_{j+1};S_j\text{ on }\Lambda]$; every uncut line carries $G_{j+1}$.
\end{lemma}

\begin{proof}
In $G_{\mathcal F}$ the lines inside members of $\mathcal F$ are contracted, so the remaining lines are those of $\Gamma\setminus\bigcup\mathcal F$; expanding each as $H+D$ gives $\sum_{\Lambda\subset\mathrm{lines}(G_{\mathcal F})}G_{\mathcal F}[H\text{ off }\Lambda,D\text{ on }\Lambda]$, the term $\Lambda=\emptyset$ being $G_{\mathcal F}[H]$; counterterm vertices are vertices, not lines, and are not expanded. For fixed $\Lambda$, the forests that occur are those disjoint from $\Lambda$. (b) For a subgraph containing no line of $\Lambda$ the degree counts only its own lines, so it is the same in $\Gamma$ and in $\Gamma\setminus\Lambda$; hence the divergent subgraphs of $\Gamma\setminus\Lambda$ are exactly those of $\Gamma$ avoiding $\Lambda$. A divergent subgraph is connected, so it lies in one cluster; forests and values factorise over clusters, which are joined only through the $D$-fields on $\Lambda$. (c) is (a) with $G_j=G_{j+1}-S_j$.
\end{proof}

Proposition~1 is multilinear in the coefficients of the \cite[(1.7)]{Ba3} vertices: a \cite[(1.7)]{Ba3} vertex of order $\mathrm{ord}(v)=\alpha+2\beta$ whose coefficient is a fixed number $c_v$ (in the units of the step, couplings included) enters the bound \cite[(1.33)]{Ba3} through $\sup|c_v|\,t^{-\mathrm{ord}(v)/2}$ in place of the factor $t^{\mathrm{ord}(v)/2}$ that the coupling count $e_k^{d_v}\lambda_k^{d_s}$ assigns to it (\cite[p.~421]{Ba3}: ``the vertices (1.7) give also a positive power of $L^k\eps$ depending on a term in (1.29)''). So a localised family of order $o$ whose \cite[(1.7)]{Ba3} vertices have coefficients $c_v$ is bounded by $O(1)t^{o/2}\prod_v(\sup|c_v|t^{-\mathrm{ord}(v)/2})e^{-\delta_0d(\Lambda)}\prod\norm{h\Phi}_{1,\alpha_0}$. This is the form in which Proposition~1 is used below and in Steps~3(D$_1$)--3(D$_2$).

\begin{lemma}[the cut-line mechanism]\label{lem:a06e:CUT}
Fix $0\le j<K$ and a class $F$ of graphs of order $o$ of \cite[(1.23)]{Ba3} or \cite[(1.24)]{Ba3}, and let $m_j=F[G_{j+1}]-F[G_j]$. \textup{(a)} Each term of \cref{lem:a06e:MC}(c) is a product of clusters, renormalised families at step $j+1$ with all lines $G_{j+1}$, joined through the cut lines, which are two-variable external fields: in scale-$(j+1)$ units $|\partial^\sigma\partial'^{\sigma'}S_j(x,x')|\le C_KL^{1+\sigma+\sigma'}e^{-\delta_KL|x-x'|}$, and the two-variable norm $\norm{\cdot}_{1,\alpha_0}$ of $hS_jh'$ for localisations at boxes at distance $r$ is $\le C_KL^{2+\alpha_0}e^{-\delta_KLr}$ (\cref{lem:a06e:SL}(b); for $j=0$ its $L$-dependent constant). \textup{(b)} The only \cite[(1.7)]{Ba3} vertices in a two-point class of order $\le4$ or a vacuum class are $\dm_1$ and $\dm_{(4)}$ vertices; in a cluster with lines $G_{j+1}$ the non-renormalising part of a $\dm_1$ vertex is $\Delta_{j+1,j+1}=O(l_{j+1}^{-1})$ (\cref{prop:a06e:DM1}), so in scale-$(j+1)$ units its coefficient is $O(t_{j+1})$: it enters Proposition~1 with its printed power $t^{\mathrm{ord}/2}=t$. \textup{(c)} Proposition~1 per cluster, the clusters joined through the cut fields as in \cref{lem:a06d:V10}, bounds each localised term by $O(1)C_K^{|\Lambda|}L^{c|\Lambda|}t_{j+1}^{o/2-\sum_v(\mathrm{ord}(v)/2-1)^+}e^{-\delta''d(\{\square\})}$, the sum over the counterterm vertices $v$ of the class (a $\dm_{(4)}$ vertex carries $t\,q(t)$, $q(t)=(1+\ln t^{-1})^5$, i.e.\ $t^{1-\beta}$ for every $\beta>0$, not $t^2$: \cref{thm:a06e:T0}). \textup{(d)} Summing over localisations (\cref{lem:a06d:V11}), with the root cube fixed: for a two-point class without $\dm_{(4)}$ vertices, $|m_j(x)|\le Cl_j^{o/2-2}$; for a vacuum class, $|m_j|\le Ct_j^{o/2}|T^{(j)}|$; with the exponent of \textup{(c)}, times $q(t_{j+1})$ per $\dm_{(4)}$ vertex, when $\dm_{(4)}$ vertices occur. The constants are free of $\eps$, $\eps_0$, $j$, $K$ and the torus.
\end{lemma}

\begin{proof}
(a) is \cref{lem:a06e:MC}(c) and the rescaling of \cref{lem:a06e:SL}(b) to scale $j+1$, where $S_j$ has side $l_j/l_{j+1}=1/L$. (b) In a two-point class of order $\le4$ the \cite[(1.7)]{Ba3} vertices have order $\le2$ (a $\dm_{(4)}$ vertex alone is the whole order); in a cluster with lines $G_{j+1}$, every $\dm_1$ vertex is paired with the order-$2$ two-point subgraphs of the same class with lines $G_{j+1}$, and the pair sums to $\dm_1-\dm_1[G_{j+1},G^v_{j+1}]=\Delta_{j+1,j+1}$ at that site, plus renormalised two-point insertions, which Proposition~1 bounds. (c) Proposition~1 with the external fields of (a), in the multilinear form above; the cut fields tie the clusters, and the tree of the whole term has length at most the sum of the cluster trees and the cut-line distances. (d) In scale-$(j+1)$ units the two-point coefficient is $l_{j+1}^2m_j(x)$, bounded by $Ct_{j+1}^{o/2}$ after the sum over localisations with the root fixed; for a vacuum class the root ranges over the $|T^{(j+1)}|$ boxes. The factor $t_{j+1}^{(\cdot)}=L^{(\cdot)}t_j^{(\cdot)}$ changes constants only; the number of terms is $\le2^{I}$, $I$ the number of lines.
\end{proof}

\begin{theorem}[{$\dm_{K,k}=O((L^k\eps)^{-1})$, \cite[(1.28)]{Ba3}}]\label{thm:a06e:DM}
For $0\le k<K$ and every $x$: \textup{(a)} $\dm_{K,k}=\sum_{j=k}^{K-1}m_j$ with $m_j\defeq\dm_{j+1}-\dm_j$, the graphs of \cite[(1.23)]{Ba3} with all lines in $\{G_j,S_j\}$ and at least one $S_j$ line. \textup{(b)} Order $2$: for $j\ge1$, $m^{(2)}_j=\Delta_{j,j}-\Delta_{j+1,j+1}$, so $|m^{(2)}_j|\le2C_Dl_j^{-1}$; for $j=0$, $m^{(2)}_0=\dm_1-\Delta_{1,1}$ and $|\dm_1|\le c_ml_0^{-1}$ (\cref{thm:a06e:CT1}). \textup{(c)} Order $4$: $|m^{(4)}_j(x)|\le C_4$ (\cref{lem:a06e:CUT} with $o=4$, $l_j^{o/2-2}=1$; the only \cite[(1.7)]{Ba3} vertices are $\dm_1$ vertices). \textup{(d)} Hence $|\dm_{K,k}(x)|\le\sum_{j\ge k}(2C_Dl_j^{-1}+C_4)\le Cl_k^{-1}+C_4(K-k)\le C'l_k^{-1}$, since $K-k\le\log_L(1/l_k)\le l_k^{-1}$ ($l_K\le\eps_0\le1$); so $\dm_{K,k}(x)(L^k\eps)^2=O(L^k\eps)$, exactly as \cite[p.~418]{Ba3} asserts, with a constant free of $\eps$, $\eps_0$, $k$, $K$ and the torus. \textup{(e)} ${\dm}^{(2)}_{\mathrm{fin}}=\Delta_{K,K}=O(l_K^{-1})$. The order-$(0,2)$ part of $\dm_{K,k}$ obeys $|\cdot|\le C\lambda^2(1+\ln l_k^{-1})$ without Proposition~1.
\end{theorem}

\begin{proof}
(a) is telescoping under \cref{conv:a06d:LOC}; (b) and (e) are \cref{prop:a06e:DM1}; (c) is \cref{lem:a06e:CUT}(d); (d) is the sum. For the last clause: the order-$\lambda^2$ graphs are the sunset, the tadpole on a tadpole and the tadpole with a $\dm_{(0,1)}$ vertex, the last two forming the class ``$\lambda$-tadpole whose loop line carries the blob $4(N+2)\lambda(C_0(0)-G(x',x'))$''; the sunset part is $16(2N+4)\lambda^2\sum_{x'}\eps^3[G_K^3-G_k^3](x,x')$, bounded by \cref{lem:a06e:SL} through $|G_{K,k}(x,x')|\le\sum_{j=k}^{K-1}Cl_j^{-1}e^{-\delta_K|y|/l_j}$ and $|G_K|,|G_k|\le C(\eps+|y|)^{-1}$ by $\sum_{j=k}^{K-1}Cl_j^{-1}\int(\eps+|y|)^{-2}e^{-\delta_K|y|/l_j}dy\le C(K-k)$; the tadpole class is $C_N\lambda^2\sum_{x'}\eps^3[(G_K^2-G_k^2)R_K-G_k^2G_{K,k}(x',x')]$, with $|R_K|\le c_Rl_K^{-1}$, $|G_{K,k}(x',x')|\le c_Rl_k^{-1}$ (\cref{lem:a06e:B0}), bounded by $Cl_K^{-1}\sum_{j<K}l_j^{-1}l_j^2+Cl_k^{-1}l_k\le C$.
\end{proof}

% a06_BALtorus_ct3: Step 3(E), third part: the finite part dm^2_fin (Ba3 (1.26)) in the form in which it is consumed,
% and the size of the non-renormalising (1.7) vertex at every step (the vertex tolerance).

\subsubsection*{The finite part $\dm_{\mathrm{fin}}$}

By \cref{lem:a06e:MC}(a) with $H=(G_K,G^v_K)$ and $D=(R_K,R^v_K)$, $\dm_{\mathrm{fin}}=\dm[C_0,C]-\dm[G_K,G^v_K]$ is a finite sum of terms in which every cluster is a renormalised family at step $K$ with lines $G_K$ and every cut line is an $R_K$ line. The order-$2$ part is \cref{thm:a06e:DM}(e). The order-$4$ part is bounded in \cref{thm:a06e:DMf} below, up to a polylogarithm; the proof is Proposition~1 at step $K$ per cluster, with the clusters glued along the $R_K$ lines by an infrared power count.

\begin{lemma}[the remainder in scale-$K$ units]\label{lem:a06e:RKp}
Let $\alpha_0\in(0,1)$ and $t\defeq t_K=l_K$. \textup{(a)} For $\sigma,\sigma'\in\{0,1\}$, $|\partial^\sigma\partial'^{\sigma'}R_K(x,x')|\le C_R\max(l_K,|y|)^{-1-\sigma-\sigma'}e^{-c_*|y|}$ (\cref{lem:a06e:RK}), and the H\"older quotients of exponent $\alpha_0$ (in units of $l_K$, over separations $\le2l_K$) of $\partial^\sigma\partial'^{\sigma'}R_K$, $\sigma+\sigma'\le2$, in either variable obey the same bound with $\max(l_K,|y|)^{-1-\sigma-\sigma'}$ (so the mixed quotient, $\sigma=\sigma'=1$, carries $\max(l_K,|y|)^{-3}e^{-c_*|y|}$). \textup{(b)} In scale-$K$ units, $\hat R\defeq l_KR_K(l_K\cdot,l_K\cdot)$: for localisations $h,h'$ at boxes of side $\le2$ at distance $r$, $\norm{(h\otimes h')\partial^\sigma\partial'^{\sigma'}\hat R}_{0,\alpha_0}\le C_R\max(1,r)^{-1-\sigma-\sigma'}e^{-c_*tr}$, and the two-variable norm $\norm{(h\otimes h')\hat R}_{1,\alpha_0}$ is $\le C_R\max(1,r)^{-1}e^{-c_*tr}$. \textup{(c)} The same for $R^v_K$. The constants are free of $\eps$, $K$ and the torus.
\end{lemma}

\begin{proof}
(a) The sup and first-difference bounds are \cref{lem:a06e:RK}. For the H\"older quotients use $R_K=G_K(a_Kl_K^{-2}P_K)C_0$ \cite[(1.25)]{Ba3}: with the block average, $R_K(x,x')=a_Kl_K^{-5}\sum_B(G_K\mathbf 1_B)(x)u_B(x')$, $u_B(x')\defeq\sum_{z\in B}\eps^3C_0(z,x')$, the sum over the blocks of side $l_K$. Since $G_K\mathbf 1=(m^2+a_Kl_K^{-2})^{-1}\mathbf 1$ on the torus, $\sum_B\partial_x(G_K\mathbf 1_B)(x)=0$ and $\partial_xR_K(x,x')=a_Kl_K^{-5}\sum_B\partial_x(G_K\mathbf 1_B)(x)[u_B(x')-u_{B_x}(x')]$, $B_x$ the block of $x$. By \cite[Lemma~2.4, (2.35)--(2.36)]{Ba4} at scale $K$ in the form of \cref{thm:a06:A5} ($\mathbf 1_B=Q_K^*\delta_{\hat y}$ in scale-$K$ units; the torus is a periodic box that is a union of the blocks of side $l_K$, so no depth condition arises and the constants are free of the torus) rescaled to $\eps$-units, $|G_K\mathbf 1_B|(x)\le c_0l_K^2e^{-\delta_0\dist(x,B)/l_K}$, $|\partial_xG_K\mathbf 1_B|\le c_0l_Ke^{-\delta_0\dist(x,B)/l_K}$, and the H\"older quotient of $\partial_xG_K\mathbf 1_B$ (in units of $l_K$) obeys the same bound with $c_1(\alpha_0)e^{2\delta_0}$ in place of $c_0$, $c_1(\alpha_0)$ the constant of \cite[(2.36)]{Ba4} (which decays in $\dist(\{\hat x,\hat x'\},\hat y)\ge\dist(\hat x,B)-2$ for separations $\le2l_K$). By \cref{lem:a06e:SL}(c) along a path from $B_x$ to $B$, for $|y|\ge4l_K+2\dist(B,B_x)$: $|u_B(x')-u_{B_x}(x')|\le Cl_K^3(\dist(B,B_x)+l_K)|y|^{-2}e^{-c_*|y|}$, each difference in $x'$ adding $|y|^{-1}$; for the remaining $B$ ($\dist(B,B_x)\ge|y|/4-l_K$) bound the two terms separately by $|u_B|\le Cl_K^3\max(l_K,\dist(B,x'))^{-1}$. The sum over $B$ converges by $e^{-\delta_0\dist/l_K}$, which absorbs the polynomial factors and makes the remaining-$B$ part $\le Ce^{-\delta_0|y|/(8l_K)}$. This gives (a) for the H\"older quotient of $\partial_xR_K$ at $|y|\ge4l_K$; for $|y|<4l_K$ the same sum with $|u_B|\le Cl_K^2$ gives $Cl_K^{-2-\sigma'}$. The $x'$-variable is symmetric ($R_K=C_0a_Kl_K^{-2}P_KG_K$). The mixed quotient (the H\"older quotient of $\partial_x\partial_{x'}R_K$ in either variable) is the same sum with one more difference: \cref{lem:a06e:SL}(c) at $\sigma+\sigma'\le3$ for the differences of $u_B$, and \cite[(2.36)]{Ba4} (in the form of \cref{thm:a06:A5}) for $\partial_xG_K\mathbf 1_B$. (b) is (a) rescaled: a two-point kernel scales by $l_K$, each difference by $l_K$, and $e^{-c_*|y|}=e^{-c_*tr}$; the norm $\norm{\cdot}_{1,\alpha_0}$ is the maximum of the sup, the first differences and their H\"older quotients.
\end{proof}

\begin{lemma}[derivative credit]\label{lem:a06e:DC}
Let a localised family at step $K$ (zero background, torus) have an external leg $l$ at a vertex $v$ of type \cite[(1.8)]{Ba3}, with the difference of $v$ acting on $l$, and let $f_l$ be the external field on $l$. Then \cite[(1.33)]{Ba3} holds with $\norm{hf_l}_{1,\alpha_0}$ replaced by $C_{\mathrm{DC}}\norm{\hat h\,\partial f_l}_{0,\alpha_0}$, $\hat h$ a localisation equal to $1$ on the $1$-neighbourhood of $\supp h$, $C_{\mathrm{DC}}$ depending on $(d,\alpha_0)$ and the localisation family only; the same in each variable of a two-variable external field.
\end{lemma}

\begin{proof}
Vertices are localised by weights on their summation variables (unit-cube indicators for \cite[(1.6)--(1.7)]{Ba3}, smooth partitions of unity on the vector legs for the others, \cite[p.~420]{Ba3}), so $f_l$ enters only through $(\partial f_l)(b)$ at the bonds of the localisation of $v$, weighted by the partition, and the value is unchanged under $f_l\to f_l-f_l(z_0)$, $z_0\in\supp h$ (at zero background constants have zero difference). Apply \cite[(1.33)]{Ba3} to $f_l-f_l(z_0)$: on $\supp h$, $|f_l(z)-f_l(z_0)|\le3d\sup|\partial f_l|$ along an axis path, and $\norm{h(f_l-f_l(z_0))}_{1,\alpha_0}\le C(\sup_{\supp\hat h}|f_l-f_l(z_0)|+\norm{\hat h\,\partial f_l}_{0,\alpha_0})$.
\end{proof}

\begin{lemma}[infrared power counting for the gluing]\label{lem:a06e:IRC}
Let $Q$ be a finite connected multigraph (self-edges allowed) with a root node $c_0$; each non-root node $c$ has a position $z_c\in\Z^3$, $z_{c_0}\defeq0$ (or positions on a torus, with the torus distance in the weights, torus balls in the count and the torus triangle inequality; every step below holds verbatim). A non-self edge $e$ between $c,c'$ has weight $w_e(z_c-z_{c'})$ with $0\le w_e(z)\le\max(1,|z|)^{-a_e}e^{-\mu|z|}$, $a_e\ge0$, $0<\mu\le\frac12$; self-edges have weight $\le1$. For a partition $P$ of the nodes into groups each inducing a connected subgraph put $\mathrm{IR}(P)\defeq3(|P|-1)-\sum_{e\text{ between groups}}a_e$ and $D\defeq\max_P\mathrm{IR}(P)\ge0$. Then $\sum_{(z_c)}\prod_ew_e\le C_Q(1+\log_2\mu^{-1})^{|E|}\mu^{-D}$, $C_Q$ depending on $Q$ and $(a_e)$ only.
\end{lemma}

\begin{proof}
For a non-self edge put $j_e\defeq0$ if $|z_e|\le1$ and $j_e\defeq j$ if $2^{j-1}<|z_e|\le2^j$; $J\defeq\lceil\log_2\mu^{-1}\rceil$, $\hat j_e\defeq\min(j_e,J)$. Kruskal's rule with the edges in non-decreasing $j_e$ gives a spanning tree $\tau$; the configuration is fixed by the tree-edge vectors, so the number of configurations with a given scale vector is $\le\prod_{e\in\tau}C2^{3j_e}$. Weights: $w_e\le2^{a_e}2^{-a_ej_e}e^{-\mu2^{j_e-1}[j_e\ge1]}$. For $e\in\tau$ with $j_e>J$: $2^{(3-a_e)j_e}e^{-\mu2^{j_e-1}}\le2^{(3-a_e)\hat j_e}g(j_e-J)$ with $g(n)\defeq2^{3n}\exp(-2^{n-2})$ summable; for $e\notin\tau$, $2^{-a_ej_e}\le2^{-a_e\hat j_e}$. Telescoping: $\sum_{e\in\tau}3\hat j_e-\sum_ea_e\hat j_e=\sum_{i=1}^J[3\#\{e\in\tau:j_e\ge i\}-\sum_{e:j_e\ge i}a_e]$; for level $i$ let $P_i$ be the partition into the components of the edges with $j_e<i$; by Kruskal's rule $\#\{e\in\tau:j_e\ge i\}=|P_i|-1$, every edge between groups of $P_i$ has $j_e\ge i$, and the other edges with $j_e\ge i$ only lower the bracket; so the bracket is $\le\mathrm{IR}(P_i)\le D$ and the total is $\le JD$, i.e.\ $2^{JD}\le(2/\mu)^D$. A rejected edge $e$ closes a cycle of tree edges $f$ with $j_f\le j_e$ (Kruskal's rule), and the triangle inequality along that path gives $2^{j_e-1}<|z_e|\le|Q|2^{\max_fj_f}$, so $\max_fj_f\le j_e\le\max_fj_f+\log_2|Q|+1$; hence each tree scale vector extends to at most $(\log_2|Q|+2)^{|E|-|\tau|}$ full scale vectors, there are $\le(J+1)^{|\tau|}$ truncated tree scale vectors, and the excess sums converge; the stated exponent $|E|\ge|\tau|$ is an upper bound.
\end{proof}

\begin{lemma}[degrees]\label{lem:a06e:ID}
For a connected subgraph $g$ of a graph of Step~3(E) (vertex set and all lines among them) let $E(g)$ be the number of its legs not on its lines, $x(g)$ the number of its $V^8$ vertices whose difference acts on such a leg, and $\omega(g)\defeq3\mathcal L(g)-2\sum_{l\in g}(1-d_l/2)-\sum_{v\in g}s_v$ with $s(V_\delta)\defeq-1$. \textup{(a)} $\omega(g)=3-E(g)/2-\mathrm{ord}(g)/2-x(g)$, an integer. \textup{(b)} (Gluing identity.) Let $\Gamma$ be a graph of $\dm_{(4)}$ at zero external momentum, rooted at $x$, $\Lambda$ a set of its lines, the clusters the components of $\Gamma\setminus\Lambda$; give an $R_K$ line $e$ the exponent $a_e\defeq1+d_e$, $d_e\in\{0,1,2\}$ the number of \cite[(1.8)]{Ba3} differences acting on its ends. For a partition $P$ of the clusters into groups connected through $\Lambda$, with $g_1,\dots,g_n$ the subgraphs induced by the groups, $\mathrm{IR}(P)=\sum_i\omega(g_i)-\omega(\Gamma)$. \textup{(c)} (Subgraph degrees.) If $\Gamma$ is a graph of $\dm_{(4)}$ (two-point, order $4$, one-particle irreducible, no difference on its external legs), then $\omega(\Gamma)=0$, every connected subgraph $g\ne\Gamma$ has $\omega(g)\le1$, with $\omega(g)=1$ only if $g$ is blob-shaped ($\mathrm{ord}(g)=2$, $E(g)=2$, $x(g)=0$) and $\omega(g)\le0$ otherwise, and no two vertex-disjoint connected subgraphs of $\Gamma$ are both blob-shaped.
\end{lemma}

\begin{proof}
(a) With $\mathcal L=I-V+1$ and $2I=\sum_vh_v-E$: $\omega=\sum_v(h_v/2-3-s_v+d_{v,g})-E/2+3$, $d_{v,g}\in\{0,1\}$ indicating a $V^8$ difference acting on a line of $g$; the bracket is $-\mathrm{ord}(v)/2$ for $V_\lambda$, $V^{10}_n$, $V_\delta$, and $-n/2-(1-d_{v,g})$ for $V^8_n$ (\cref{lem:a06e:M3}(d)). Parity: $h_v\equiv\mathrm{ord}(v)\pmod2$ for every type, so $E(g)\equiv\mathrm{ord}(g)$ and $\omega(g)\in\Z$. (b) Lines between groups are $R_K$ lines; $\mathcal L(\Gamma)=\sum_i\mathcal L(g_i)+(E_{\mathrm{inter}}-n+1)$, so $\omega(\Gamma)=\sum_i\omega(g_i)+3(E_{\mathrm{inter}}-n+1)-\sum_{\mathrm{inter}}(2-d_e)=\sum_i\omega(g_i)-[3(n-1)-\sum_{\mathrm{inter}}(1+d_e)]$. (c) $\omega(\Gamma)=3-1-2-0=0$. For a proper connected $g$: $E(g)\ge2$ ($E=0$ forces $g=\Gamma$; $E=1$ makes the single leg a bridge or disconnects $\Gamma$). The only order-$1$ vertex is $V^8_1$, whose single-vertex subgraph has $E=3$ and $x=1$. So $E+\mathrm{ord}+2x\ge4$ with equality only for $(E,\mathrm{ord},x)=(2,2,0)$: $\omega=3-(E+\mathrm{ord}+2x)/2\le1$, with equality exactly for blob-shaped $g$ and $\omega\le0$ otherwise (integrality). Two vertex-disjoint blob-shaped $g_1,g_2$ carry all of the order, hence all vertices, and each has exactly two external legs; the two external legs of $\Gamma$ are among these four: if both sit on $g_1$, the two legs of $g_2$ must join $g_1$, which has none left; if one on each, the remaining two form a bridge, against one-particle irreducibility; if neither, $\Gamma$ has no external leg.
\end{proof}

\begin{theorem}[the order-$4$ finite part]\label{thm:a06e:DMf}
For every $x$ and every torus aligned with the blocks with sides $\ge1$, $|{\dm}^{(4)}_{\mathrm{fin}}(x)|\le C_{\mathrm{DMf}}\,l_K^{-1}(1+\ln l_K^{-1})^5$, with $C_{\mathrm{DMf}}$ free of $\eps$, $\eps_0$, $K$ and the torus.
\end{theorem}

\begin{proof}
(1) By \cref{lem:a06e:MC}(a), ${\dm}^{(4)}_{\mathrm{fin}}(x)=\sum_\Gamma\sum_{\Lambda\ne\emptyset}(\Gamma\setminus\Lambda)_{\mathrm{ren}}[G_K;R_K\text{ on }\Lambda](x)$, a finite sum (at most $4$ vertices and $5$ lines per $\Gamma$). (2) Each cluster is a renormalised family at step $K$ with lines $G_K$; its \cite[(1.7)]{Ba3} vertices are $\dm_1$ vertices (no order-$4$ counterterm occurs in $\dm_{(4)}$), whose non-renormalising part is $\Delta_{K,K}=O(l_K^{-1})$ (\cref{prop:a06e:DM1}), of the printed size; in scale-$K$ units Proposition~1 gives per localised cluster family $O(1)t^{o_c/2}e^{-\delta'd(\{\square\})}$ times the norms of its external fields, $\sum_co_c=4$. (3) The external fields are the ends of the $R_K$ lines; an end at which a \cite[(1.8)]{Ba3} difference acts is measured by the differenced kernel (\cref{lem:a06e:DC}); by \cref{lem:a06e:RKp}(b) a line $e$ between boxes at distance $r$ contributes $\le C_R\max(1,r)^{-a_e}e^{-c_*tr}$, $a_e=1+d_e$, and a line with both ends in one cluster $\le C_R$; clusters joined by $R_K$ lines are joined as in \cref{lem:a06d:V10}. (4) For each cluster sum its boxes around a root box with half of its tree decay (\cref{lem:a06d:V11}); the distance of two attached boxes differs from that of the roots by at most $d_c+d_{c'}$, and $(1+d_c+d_{c'})^{a_e}e^{c_*t(d_c+d_{c'})}$ is absorbed by the other half ($c_*t\le\delta'/8$). (5) \Cref{lem:a06e:IRC} with $\mu\defeq c_*t$ bounds the sum over the cluster roots by $C(1+\log t^{-1})^{|\Lambda|}t^{-D}$, and by \cref{lem:a06e:ID}(b),(c), $\mathrm{IR}(P)=\sum_i\omega(g_i)-\omega(\Gamma)\le1$ for every partition (at most one blob-shaped group, every other group $\le0$), so $D\le1$. (6) In scale-$K$ units the coefficient is $\dm(x)l_K^2$, so $|{\dm}^{(4)}_{\mathrm{fin}}(x)|l_K^2\le Ct^2t^{-1}(1+\log t^{-1})^5$ with $t=l_K$.
\end{proof}

The logarithm is the count of scale vectors in \cref{lem:a06e:IRC}; the statement without it is not claimed, and nothing below needs it.

\subsubsection*{The non-renormalising vertex at every step}

For $0\le j\le K$ put $c_j(x)\defeq l_j^2(\dm_{\mathrm{fin}}(x)+\dm_{K,j}(x))$, $\dm_{K,K}\defeq0$: the coefficient, in scale-$j$ units, of the part of the step-$j$ vertex \cite[(1.7), (1.29)]{Ba3} that is not absorbed into renormalised pairs ($c_0=\eps^2\dm$, since $\dm_0=0$); the renormalising piece $\dm_j(e_j,\lambda_j,\Omega_1,x)$ of \cite[(1.29)]{Ba3} is of size $l_j^2/\eps$-type by \cref{thm:a06e:CT1} rescaled and is cancelled inside the renormalised classes. Let $q(t)\defeq(1+\ln t^{-1})^5$ and $Q_\beta\defeq\sup_{0<t\le1}t^\beta q(t)=(5/\beta)^5e^{\beta-5}$ for $\beta\in(0,5]$.

\begin{theorem}[the vertex tolerance]\label{thm:a06e:T0}
Let $C_T\defeq2\max(C_D,c_m)+2C_4/\ln L+C_{\mathrm{DMf}}$. For every $0\le j\le K$ and $x$: \textup{(a)} $|c_j(x)|\le C_Tl_jq(l_j)$; \textup{(b)} for $\beta\in(0,5]$, $|c_j(x)|\le C_TQ_\beta l_j^{1-\beta}$; \textup{(c)} $|c_j(x)|\le C_TQ_1$; \textup{(d)} $|c_j(x)|\le C_T(1+\ln(L/\eps_0))^5l_j$; \textup{(e)} only the order-$4$ finite piece carries $q$: $|c^{(2)}_j|\le\max(C_D,c_m)l_j$, $|l_j^2{\dm}^{(4)}_{K,j}|\le(C_4/\ln L)l_j^2(1+\ln l_j^{-1})\le Cl_j$, $|l_j^2{\dm}^{(4)}_{\mathrm{fin}}|\le C_{\mathrm{DMf}}l_jq(l_j)$. $C_T$ is free of $\eps$, $\eps_0$, $j$, $K$ and the torus.
\end{theorem}

\begin{proof}
Split $c_j=c^{(2)}_j+c^{(4)}_j$ by order. Order $2$: for $j\ge1$, ${\dm}^{(2)}_{\mathrm{fin}}+{\dm}^{(2)}_{K,j}=\Delta_{K,K}+(\Delta_{j,j}-\Delta_{K,K})=\Delta_{j,j}$ (\cref{thm:a06e:DM}(b),(e)) and $|\Delta_{j,j}|\le C_Dl_j^{-1}$; for $j=0$ it is $\dm_1$ with $|\dm_1|\le c_ml_0^{-1}$; so $|c^{(2)}_j|\le\max(C_D,c_m)l_j$. Order $4$, $(K,j)$ piece: $|{\dm}^{(4)}_{K,j}|\le C_4(K-j)$ (\cref{thm:a06e:DM}(c)) and $K-j\le\log_L(1/l_j)\le(1+\ln l_j^{-1})/\ln L$, so $l_j^2|{\dm}^{(4)}_{K,j}|\le(C_4/\ln L)l_j^2(1+\ln l_j^{-1})\le(C_4/\ln L)l_j$, as $t(1+\ln t^{-1})\le1$ on $(0,1]$. Order $4$, finite piece: $f(t)\defeq t^{-1}q(t)$ is decreasing on $(0,1]$ ($f'/f=-t^{-1}(1+5/(1+\ln t^{-1}))<0$), and $l_j\le l_K$, so \cref{thm:a06e:DMf} gives $l_j^2|{\dm}^{(4)}_{\mathrm{fin}}|\le C_{\mathrm{DMf}}l_j^2f(l_K)\le C_{\mathrm{DMf}}l_j^2f(l_j)=C_{\mathrm{DMf}}l_jq(l_j)$. Sum, with $q\ge1$: (a). (b), (c) are $q(t)\le Q_\beta t^{-\beta}$, $tq(t)\le Q_1$. (d): $l_j^2f(l_K)=l_j(l_j/l_K)(1+\ln l_K^{-1})^5\le l_j(1+\ln(L/\eps_0))^5$, since $l_j\le l_K$ and $l_K>\eps_0/L$; the other parts are $\le Cl_j$ already.
\end{proof}

\begin{corollary}[$\dm_{\mathrm{fin}}$ at fixed $\eps_0$]\label{cor:a06e:T0p}
$|\dm_{\mathrm{fin}}(x)|\le C_TL\eps_0^{-1}(1+\ln(L/\eps_0))^5$, uniformly in $\eps$, $x$ and the torus: Ba{\l}aban's ``convergent as $\eps\to0$'' \cite[(1.26)]{Ba3}, with a constant depending on $\eps_0$.
\end{corollary}

\begin{proof}
\Cref{thm:a06e:T0}(d) at $j=K$ ($c_K=l_K^2\dm_{\mathrm{fin}}$) and $l_K>\eps_0/L$.
\end{proof}

\begin{lemma}[family tolerance]\label{lem:a06e:T1}
Let $F$ be a localised renormalised family at a step $j\le K$, in any application of Proposition~1 of Steps~3(B)--3(D), and let $n_7(F)\le\bar n/2$ be the number of its non-renormalising \cite[(1.7)]{Ba3} vertices. If a bound $B_F$ on $F$ was obtained with $|c_j|\le Cl_j$ at each such vertex, then \cref{thm:a06e:T0} gives the bound $(C_T/C)^{n_7}q(l_j)^{n_7}B_F$, with constants free of $\eps_0$. Consequently: \textup{(i)} a bound $O(1)t^\kappa p(t)^Ne^{-\delta d}$ becomes $O(1)t^\kappa p(t)^Nq(t)^{\bar n/2}e^{-\delta d}$, the same polylogarithmic class ($p(t)^Nq(t)^n\le C_{N,n,\beta}t^{-\beta}$); \textup{(ii)} a conclusion ``$C_\beta t^{\kappa-\beta}$ for every $\beta>0$'' holds with the same $\kappa$ and a new $\eps_0$-free $C_\beta$; \textup{(iii)} a bound that used only $|c_j|\le C$ holds with $C_TQ_1$; \textup{(iv)} a geometric step sum $\sum_{i=k}^{K-1}t_i^aq(t_i)^n$ is $\le(1-L^a)^{-1}t_k^aq(t_k)^n$ for $a<0$ and $\le(1-L^{-a/2})^{-1}Q^n_{a/(2n)}$ for $a>0$, $n\ge1$. Every consumer of the counterterm sizes in Steps~3(B)--3(D) uses the \cite[(1.7)]{Ba3} vertex only as a positive power of $t$ (\cite[p.~421]{Ba3}; the exponents $\kappa_0>0$, $\frac12-\beta$, $\frac14-\beta$ of \cite[(3.58)]{Ba1}, \cite[(2.57), (2.116)--(2.117)]{Ba2}), so each tolerates the form \textup{(a)} of \cref{thm:a06e:T0}; where such a bound was written with an exact factor $t$ it holds with $tq(t)$.
\end{lemma}

\begin{proof}
Proposition~1 is multilinear in the coefficients $c_j$ (Step~3(E), before \cref{lem:a06e:CUT}); replace $\sup|c_j|\le Cl_j$ by $C_Tl_jq(l_j)$ at each of the $n_7$ vertices; $n_7\le\bar n/2$ since each \cite[(1.7)]{Ba3} vertex has order $\ge2$ and the family has order $\le\bar n$. (i), (ii): $p(t)=b_0(1+\ln t^{-1})^{\bar p}$, so $p^Nq^n$ is a power of $1+\ln t^{-1}$, bounded by $C_\beta t^{-\beta}$. (iii) is \cref{thm:a06e:T0}(c). (iv): for $a<0$, $q(t_i)\le q(t_k)$ for $i\ge k$ and the series is geometric; for $a>0$, $t^{a/2}q(t)^n\le Q^n_{a/(2n)}$ and $\sum_it_i^{a/2}\le(1-L^{-a/2})^{-1}$.
\end{proof}

% a06_BALtorus_ct4: Step 3(E), fourth part: the vacuum counterterm pieces E_{K,k} at every step (Ba3 (1.28), (1.30)),
% Ba2 (2.14) for the constant the chain carries, and Ba1 Props. 3.1-3.2 and Ba2 p. 579 ("Prop. I.3.1 can also be applied").

\subsubsection*{The vacuum counterterm pieces}

The vacuum counterterm $E_1$ is split by \cite[(1.26), (1.28), (1.30)]{Ba3} as $E_1=\Efin+E_{K,k}+E_k$, $E_K=E[G_K,G^v_K]$, $E_k=E[G_k,G^v_k]$ (class sums \cite[(1.24)]{Ba3} with the counterterm coefficients kept as numbers, \cref{conv:a06d:LOC}), $E_{K,k}\defeq E_K-E_k$ and $\Efin\defeq E_1-E_K$. $\Efin$ is a global constant, factored out of the chain by \cref{lem:a06:CF} and bounded by \cref{thm:a06:F}; the pieces $E_{K,k}$ are what the step-$k$ action carries, and the following statements bound them with constants free of $\eps_0$.

\begin{lemma}[the $c$-number of the step-$k$ polynomial]\label{lem:a06e:VK0}
At zero background on the torus, the field-independent part of the step-$k$ polynomial \cite[(1.4)--(1.5)]{Ba3} with vertex domain $\Omega_1$ is $c_k(\Omega_1)=-(\Efin+E_{K,k})+a_k(\Omega_1)$, where $a_k(\Omega_1)$ is the sum of the connected vacuum families containing at least one block-spin pair vertex \cite[(1.12)--(1.14)]{Ba3}. On the chain run with $E_0+E_K$ (\cref{lem:a06:CF}) it is $-E_{K,k}+a_k(\Omega_1)$.
\end{lemma}

\begin{proof}
By \cite[(1.4)]{Ba3} the vacuum graphs of $-\log$ of the fluctuation integral at zero background without pair vertices have the vertices of the action localised in $\Omega_1$ (\cite[p.~419]{Ba3}: the vertices \cite[(1.6)--(1.7)]{Ba3} by $\sum_{x\in\Omega_1}$, the others by $g_k$, and ``we perform this localization for the terms of $\dm_k$, $E_k$ also'' \cite[p.~420]{Ba3}), the full $\dm$ of \cite[(1.29)]{Ba3} in the vertex \cite[(1.7)]{Ba3}, and lines $G_k(T,0)$; order by order these are the graphs of $E_k(e'e_kg_k,\lambda'\lambda_k,\Omega_1)$ (\cite[p.~418]{Ba3}: ``the same form, only the propagators are replaced by $G^\eps_k(0)$, $G^\eps_k$''), while $\Efin$ and $E_{K,k}$ carry no $\Omega_1$ \cite[(1.30), p.~419]{Ba3}. So the $\Omega_1$-dependent vacuum part cancels against $-E_k(\Omega_1)$ in \cite[(1.5)]{Ba3}, and the pair-vertex families remain.
\end{proof}

\begin{theorem}[the vacuum pieces at step $k$]\label{thm:a06e:VK}
For $0\le k<K$, on every torus aligned with the blocks with sides $\ge1$: \textup{(a)} the pair-vertex vacuum families $a_k(\Omega_1)$ are $O(t_k)$ per localisation cube; \textup{(b)} $|E_{K,k}|\le C_{\mathrm{VK}}t_k|T^{(k)}|=C_{\mathrm{VK}}\ell_T^3l_k^{-2}$ for the classes without $\dm_{(4)}$ vertex, and $\le C_\beta t_k^{1-\beta}|T^{(k)}|$ for every $\beta>0$ for the classes with one; \textup{(c)} at order $\lambda$, without Proposition~1: $E^{(0,1)}_{K,k}=\lambda N(N+2)\sum_x\eps^3G_{K,k}(x,x)[R_K(x,x)+R_k(x,x)]$, so $|E^{(0,1)}_{K,k}|\le2\lambda N(N+2)c_R^2\ell_T^3l_k^{-2}$, and $0\le\Efin^{(0,1)}=\lambda N(N+2)\sum_x\eps^3R_K(x,x)^2\le\lambda N(N+2)c_R^2\ell_T^3l_K^{-2}$. The constants are free of $\eps$, $\eps_0$, $k$, $K$ and the torus.
\end{theorem}

\begin{proof}
(a) Proposition~1 with no external fields gives $O(1)e_k^{d_v}\lambda_k^{d_s}e^{-\delta_0d}$ per localised family; a pair vertex contains a \cite[(1.14)]{Ba3} member, of order $\ge1$ in $e$, and all $A'$-legs are contracted in pairs, so $d_v\ge2$ and the family carries $t_k$ (as in the proof of \cref{lem:a06d:V9}); sum over localisations with the root fixed (\cref{lem:a06d:V11}). (b) $E_{K,k}=\sum_{j=k}^{K-1}\epsilon_j$ with $\epsilon_j\defeq E[G_{j+1}]-E[G_j]$, and \cref{lem:a06e:CUT}(d) for vacuum classes of order $o$ gives $|\epsilon^{(o)}_j|\le Ct_j^{o/2}|T^{(j)}|=C\ell_T^3l_j^{o/2-3}$; summed over $j\ge k$: $o=2$ gives $C\ell_T^3l_k^{-2}$, $o=4$ gives $C\ell_T^3l_k^{-1}$, $o=6$ gives $C\ell_T^3(K-k)$, $o\ge8$ gives $C\ell_T^3$, all $\le C\ell_T^3l_k^{-2}$ since $l_k^2(K-k)\le1$. The \cite[(1.7)]{Ba3} vertices of the vacuum families have orders $2$ or $4$; the non-renormalising part of a $\dm_1$ vertex at step $j+1$ is $\Delta_{j+1,j+1}=O(l_{j+1}^{-1})$ (\cref{prop:a06e:DM1}), of the printed size, and that of a $\dm_{(4)}$ vertex is $\sum_{i>j}m^{(4)}_i+{\dm}^{(4)}_{\mathrm{fin}}$, whose coefficient in scale-$(j+1)$ units is $\le C_Tl_{j+1}q(l_{j+1})$ (\cref{thm:a06e:T0}(e)), so such a vertex carries $t_j^{1-\beta}$ in place of $t_j^2$ and the classes containing one obey $C_\beta t_j^{o/2-1-\beta}|T^{(j)}|$; summed over $j\ge k$ with $o\ge4$ this is $\le C_\beta\ell_T^3l_k^{-2-\beta}=C_\beta t_k^{-\beta}t_k|T^{(k)}|$. (c) $E^{(0,1)}[G]=-\lambda N(N+2)\sum_x\eps^3G(x,x)[G(x,x)-2C_0(0)]$ (the figure-eight and the cycle with the $\dm_{(0,1)}=-4(N+2)C_0(0)$ vertex); subtracting at $G_K$ and $G_k$, resp.\ at $C_0$ and $G_K$, gives the two displayed identities, and \cref{lem:a06e:B0} the bounds.
\end{proof}

\begin{corollary}[{\cite[(2.14)]{Ba2} for the constant the chain carries}]\label{cor:a06e:C1b}
On the chain run with $E_0+E_K$ the vacuum constant at step $0$ is $E_K=E_{K,0}$ ($G_0=0$), and $|E_K|\le C\eps|T_\eps|$ for the classes without $\dm_{(4)}$ vertex, $\le C_\beta\eps^{1-\beta}|T_\eps|$ for the others, with constants free of $\eps$, $\eps_0$ and the torus. Read as a density, $E_K(X)\defeq(E_K/|T_\eps|)|X|$: $E_K=E_K(\Lambda_7)+E'_K(\Lambda_7^c)$ with $|E'_K(\Lambda_7^c)|\le C_\beta\eps^{1-\beta}|\Lambda_7^c|$, which is \cite[(2.14)]{Ba2} with $\kappa_0=1-\beta$ for every $\beta>0$ ($\kappa_0=1$ for the classes without $\dm_{(4)}$ vertex), for the constant to which the chain applies it.
\end{corollary}

\begin{proof}
\Cref{thm:a06e:VK}(b) at $k=0$, $t_0|T^{(0)}|=\eps|T_\eps|$; \cref{lem:a06:CF} for the booking.
\end{proof}

This bound holds for the class sums only. The vacuum graph made of a scalar ring through three $\lambda$-vertices, each with one scalar self-loop, is a term of $E_1$ of order $\lambda^3$ whose value per $\eps$-site is $\ge c_N\lambda^3m^{-3}$ for small $\eps$ (its three insertions are $4(N+2)\lambda C_0(0)\sim\eps^{-1}$ each and $\ell_T^{-3}\sum_p\hat C_0(p)^3\ge c\,m^{-3}$); it cancels within the order-$\lambda^3$ coefficient against the rings with $\dm_{(0,1)}$ insertions.

\begin{proposition}[{\cite[Props.~3.1--3.2]{Ba1} and \cite[p.~579]{Ba2}}]\label{prop:a06e:B12}
Given Proposition~1 at $L=5$ (Steps~3(D$_3$)--3(D$_4$)) with the vertex sizes of \cref{thm:a06e:T0}: \textup{(i)} (\cite[Prop.~3.1]{Ba1}) at step $k$ of the lower bound, the sum of the $R$-terms (the remainders of the Taylor expansion in the background, \cite[(1.9), (1.11), (1.15)]{Ba3}) and the sum of the field-dependent terms at whose vertices the powers of $L^k\eps$ add up to more than $d$ are $O(1)(L^k\eps)^\kappa|T^{(k)}_1|$ with $\kappa>d$; the $c$-number part of the step-$k$ expansion is $-E_{K,k}+a_k(\Omega_1)$, bounded by \cref{thm:a06e:VK}, and the total $c$-number the chain collects is $E_K=O(\eps|T_\eps|)$, which the final bound absorbs. \textup{(ii)} (\cite[Prop.~3.2, (3.57)--(3.58)]{Ba1}) the coefficients $v^{(k)}_{j;\mu}(B^{(k+1)},\psi;x,y)$ of $V^{(k)}$ with $n+m\ge1$ obey $|v^{(k)}|\le O(1)(L^k\eps)^{\kappa_0}\exp(-\delta_0d(x_1,\dots,y_m))$ with $\kappa_0>0$ independent of $k$; the coefficient with $n=m=0$ is a function of the new fields $(B^{(k+1)},\psi)$, whose field-dependent graphs obey the same bound through their external legs, and whose pure $c$-number part is that of \textup{(i)}: $O(t_k^{1-\beta})$ per localisation cube. \textup{(iii)} (\cite[p.~579]{Ba2}) at step $k$ of the upper bound, on $\Omega=B^k(\Lambda^{(k)}_0)$ with the background $\tilde B$ regular by \cite[(2.98)]{Ba2}, the terms with $R$-vertices and the terms of order $>\bar n$ are $O((L^k\eps)^\kappa)|\Lambda^{(k)}_2|$, the terms containing $\tilde B^{(k)}$ are $O(t_k^\kappa)|\Lambda^{(k)}_2|$ for every $\kappa$ by \cref{cor:a06d:G}, and the field-dependent terms of $Q^{(k)}$ localised in $B^k(\Lambda^{(k)c}_7)$ are bounded by \cref{prop:a06f:H}, the corresponding terms of $P^{(k)}$ being the subject of \cref{rem:a06f:HP}; the $c$-number part is as in \textup{(i)}.
\end{proposition}

\begin{proof}
(i) Every $R$-term carries $e(L^k\eps)^{\bar n+1}$ with $\bar n+1$ background legs, and Proposition~1 with the field restrictions of \cite[Sect.~3]{Ba1} bounds it; a field-dependent term with more than $d$ powers of $L^k\eps$ at its vertices is bounded by Proposition~1 with that power; the exponent count $\kappa>d$ is arithmetic on the printed orders and needs $\bar n\ge6$ (\cref{thm:a06:BALA}). The $c$-number part is \cref{lem:a06e:VK0} with \cref{lem:a06:CF}, bounded by \cref{thm:a06e:VK}(a),(b); the total is \cref{cor:a06e:C1b}. For a $c$-number vacuum class of order $>6$ no graph-by-graph bound of the printed kind holds (a ring of $n$ step-$k$ tadpole insertions $\lambda_kG_k(x,x)\sim\lambda t_k/\eta$ has size $\sim(\lambda t_k/\eta)^n$ per unit block), which is why the class sums of \cref{thm:a06e:VK} are used. (ii) For $n+m\ge1$ this is Proposition~1's localised bound at fixed external fields, with $\kappa_0$ its uniform exponent (\cref{lem:a06e:T1} for the \cite[(1.7)]{Ba3} vertices); the coefficient with $n=m=0$ is free of the fluctuation fields but depends on $(B^{(k+1)},\psi)$, its graphs with external legs fall under Proposition~1, and its pure $c$-number part is read per localisation cube as the density of $-E_{K,k}+a_k(\Omega_1)$. (iii) is (i)'s argument on $\Omega=B^k(\Lambda^{(k)}_0)$: Proposition~1 is stated for a general $\Omega$ with a regular background \cite[p.~412]{Ba3}, and the depth and H\"older clauses of the background are those of \cite[(2.94)--(2.98)]{Ba2}; the $\tilde B^{(k)}$ terms are \cref{cor:a06d:G}; the $Q^{(k)}$ terms are \cref{prop:a06f:H} and the $P^{(k)}$ terms \cref{rem:a06f:HP}.
\end{proof}

% a06_BALtorus_ct5: Step 3(F), first part: Ba2 p. 568, (2.46)-(2.47) (the restrictions and the Neumann cut at the
% conditional integral), with the Neumann-cut lemma. Input from sec/a06_BALtorus_c.tex.

\paragraph{Step 3(F): the remaining asserted estimates of \cite{Ba2}.} The estimates of \cite[Sect.~2]{Ba2} asserted at the conditional integral \cite[p.~568]{Ba2}, at the quadratic-form step \cite[(2.108)--(2.109)]{Ba2}, at the collar \cite[Prop.~2.1, (2.57)]{Ba2} (its field-dependent part), at the passage from \cite[(3.5)]{Ba2} to \cite[(3.21)]{Ba2}, and the region-entropy estimate \cite[(3.42)]{Ba2} are proved in this step, at block factor $5$.

\subsubsection*{The conditional integral: \cite[p.~568, (2.46)--(2.47)]{Ba2}}

\emph{Setting.} The step from scale $k-1$ to scale $k$, in scale-$(k-1)$ units, $t\defeq L^{k-1}\eps$; $\Lambda_i\defeq\Lambda^{(k-1)}_i$, $\Omega\defeq B^{k-1}(\Lambda_2)$, $\Lambda'$ the set of coarse points (the $L$-blocks of the unit lattice) whose block lies in $\Lambda$; $\varphi,A$ the new fields on coarse points, $\varphi',A'$ the old fields on unit points; $A$ as an operator argument is the background $A^{(k),\eps}$. The exponential of \cite[(2.46)]{Ba2} has ten terms: $T_1$, $T_2$ (old-field vector forms on $\Lambda^{(k-2)'}_6\cap\Lambda_5^c$ with $\Delta^{(k-1)}$ and $\Delta^{(k-1)}C^{(k-1)}_{\Lambda_5}\Delta^{(k-1)}$), $T_3=-a(L^k\eps)^{-2}\ip{(\Lambda^{(k-2)'}_6\cap\Lambda_5^c)A'}{\Delta^{(k-1)}C^{(k-1)}_{\Lambda_5}Q^*A}$, $T_4=-\frac12\ip A{\Delta^{(k),L^k\eps}_{\Lambda_5}A}$, $T_5,T_6$ (old-field forms on $\Lambda^{(k-2)'}_6\cap\Lambda_3^c$, not touched on p.~568), $T_7=-\frac12\ip{(\Lambda_3\cap\Lambda_5^c)\varphi'}{\Delta^{(k-1)}(\Omega,A)(\Lambda_3\cap\Lambda_5^c)\varphi'}$, $T_8$ the same with $\Delta C_{\Lambda_5}\Delta$, $T_9=-a(L^k\eps)^{-2}\ip{(\Lambda_3\cap\Lambda_5^c)\varphi'}{\Delta^{(k-1)}C^{(k-1)}_{\Lambda_5}Q^*(A)\varphi}$, $T_{10}=-\frac12\ip\varphi{\Delta^{(k),L^k\eps}_{\Lambda_5}(\Omega,A)\varphi}$. The coarse forms are \cite[(2.27)]{Ba2}: $\ip\psi{\Delta^{(k)}_\Lambda\psi}=aL^{d-2}\sum_{y\in\Lambda'}|\psi(y)|^2-a^2L^{-4}\ip\psi{QC^{(k-1)}_\Lambda Q^*\psi}$. \cite[p.~568]{Ba2} asserts: (2.47), the localisation of $T_{10}$ to the three forms on $\Lambda'^c_6$, $\Lambda'_6\cap\Lambda'^c_7$ and $\Lambda'_6$ with the operators on $B^{k-1}(\Lambda_2\cap\Lambda_7^c)$ (``Neumann boundary conditions'') and on $\Omega$, up to $O((L^{k-1}\eps)^\kappa)|\Lambda_6|$; the same for $T_4$ ``except that Neumann boundary conditions are not introduced'' (vector forms on the torus); the restriction of the linear terms $T_3$, $T_9$ to $\Lambda'_5\cap\Lambda'^c_6$; and the Neumann cut in $T_7$--$T_9$. Readings: $S\defeq\partial B^{k-1}(\Lambda_7)$, $\Omega_1\defeq B^{k-1}(\Lambda_2\cap\Lambda_7^c)$, $\Omega_{\mathrm{cut}}$ is $\Omega$ with the bonds crossing $S$ removed, so $\Omega_{\mathrm{cut}}=\Omega_1\sqcup B^{k-1}(\Lambda_7)$; ``Neumann conditions on $\partial B^{k-1}(\Lambda_2\cap\Lambda_7^c)$'' is $\Omega\to\Omega_{\mathrm{cut}}$, and the printed $\Delta_{\Lambda_5}(B^{k-1}(\Lambda_2\cap\Lambda_7^c),A)$ is the $\Omega_1$-block of $\Delta_{\Lambda_5}(\Omega_{\mathrm{cut}},A)$ (\cref{lem:a06f:F2}(d); read literally the printed operator is undefined, since \cite[(2.27)]{Ba2} needs $\Lambda\subset\Omega^{(k-1)}$); the third line of (2.47) reads $B^{k-1}(\Lambda_2)$; the vector forms are the instance $A=0$, $\lmu$ for $m$, componentwise. Inputs: (G) the geometry of \cite[(2.7)--(2.8)]{Ba2}: with $d_S(z)\defeq\dist(B^{k-1}(z),S)$, $z\notin\Lambda_6$ gives $d_S(z)\ge r-c_d-1$, $z\in\Lambda_3\cap\Lambda_5^c$ gives $d_S(z)\ge2(r-c_d)-1$, and $\dist(\Lambda_6,\Lambda_5^c)$, $\dist(\Lambda_6^c,\Lambda_7)\ge r-c_d$, $r=r(t)$; (FR) the field restrictions \cite[(2.16)--(2.17)]{Ba2} at scale $k-1$ \cite[p.~560]{Ba2}: $|\varphi|\le M_\varphi\defeq cp(t)\lambda(t)^{-1/4}$, $|A|\le M_A\defeq cp(t)t^{-1}$ on $\Lambda'_5$, the old fields the same on $\Lambda_3\cap\Lambda_5^c$, and the old $A'$ on $\Lambda^{(k-2)'}_6\cap\Lambda_5^c$ by the printed factor $\chi_{k-1,\Lambda^{(k-1)c}_5}$ of \cite[(2.45)--(2.46)]{Ba2}; $M\defeq\max(M_\varphi,M_A,1)\le Ct^{-2}$; (RG) $A^{(k),\eps}$ regular on $\Omega$ (\cite[p.~562, (2.98)]{Ba2}) and $e\le e_0$; (A) the propagator estimates of group (A), i.e.\ Step~3(A) at the step index $k-1$, under its hypotheses ($\Omega$ a union of big blocks, every pin set $\Lambda\subset\Omega^{(k-1)}$ a union of $L$-blocks (the proofs of Step~3(A) use of (H1) only this), $A$ regular on $\Omega$ by (RG), $e_{k-1}\le e_1$): for every pin set $\Lambda$ the operator $T_\Lambda(\Omega)\defeq-\Delta^{\eta,N}_{A,\Omega}+m^2_{k-1}+a_{k-1}P_{k-1}(\Omega\setminus B\Lambda)+a_kL^{-2}P_k(B\Lambda)$ is $\ge\gamma_L$ (\cref{lem:a06A:coer}, which is \cite[(1.8)]{Ba4} on a general region); the Combes--Thomas bound $\norm{\mathbf 1_QT^{-1}f},\norm{\mathbf 1_QD_AT^{-1}f}\le Ce^{-\theta_0\dist(Q,\supp f)}\norm f$ for unit cubes $Q$ and $T=D_A^*D_A+V+\sum_Bc_B\Pi_B\ge\gamma$ ($V\ge0$, $\Pi_B$ projections on disjoint sets of diameter $\le\ell_*$, $0\le c_B\le a$), $\theta_0=\theta_0(\gamma,\ell_*)$ (\cref{lem:a06A:CT}(i)--(ii), paired with the functions, resp.\ bond functions, supported in $Q$; $C\defeq\max\{2\gamma^{-1},(2\gamma^{-1/2}+6\sqrt d\,\theta_0\gamma^{-1})e^{\theta_0}\}$; this contains the $G_k$ clause of \cite[Cor.~2.3]{Ba4}); the decay $|\Delta^{(k-1)}(\Omega,A;z,z')|\le c_7e^{-\delta_7|z-z'|}$ (the first bound of \cref{prop:a06A:Ba1P22}, which is \cite[(2.27)]{Ba1}; its proof gives the constant $a+2a^2\gamma_\emptyset^{-1}e^{\theta_0(\gamma_\emptyset,1)}$ and the rate $\theta_0(\gamma_\emptyset,1)$, both depending on $(d,N,\norm q,a,L)$ only) and $|C_\Lambda(y,y')|\le c_2e^{-\theta_0|y-y'|}$ (\cref{lem:a06A:blockinv}(c), which is \cite[(1.16)]{Ba4}); the block inversion $X_\Lambda^{-1}=W^*G_{k-1}(\Omega,\Lambda,A)W+a_{k-1}^{-1}-(a_kL^{-2}/a_{k-1}^2)P_\Lambda$, $X_\Lambda\defeq(\Delta^{(k-1)}+aL^{-2}P)|_\Lambda$, $W\defeq Q_{k-1}^*\Lambda(1-(a_k/a_{k-1})L^{-2}P_\Lambda)$ (\cref{lem:a06A:blockinv}(a), which is \cite[(3.8)]{Ba4}); and the restriction difference $|C_\Lambda(y,y')-C^{(k-1)}(\Omega,A;y,y')|\le c_4\exp[-\frac{\theta_0}6(|y-y'|+\dist(y,\Omega^{(k-1)}\setminus\Lambda)+\dist(y',\Omega^{(k-1)}\setminus\Lambda))]$ (\cref{lem:a06A:CLdiff}, which is \cite[(1.18)]{Ba4}). On a torus these hold with the same constants (\cref{lem:a06A:torus}).

\begin{lemma}[Neumann cut]\label{lem:a06f:F1}
Let $\Omega$ be a union of big blocks with $A$ regular on $\Omega$, $Y\subset\Omega^{(k-1)}$ a union of large blocks, $S\defeq\partial B^{k-1}(Y)$, $\Omega_{\mathrm{cut}}$ as above, $\Lambda\subset\Omega^{(k-1)}$ a pin set, $T\defeq T_\Lambda(\Omega)$ and $T_{\mathrm{cut}}\defeq T_\Lambda(\Omega_{\mathrm{cut}})$ (the Neumann Laplacian of $\Omega_{\mathrm{cut}}$, same pins). Then $T_{\mathrm{cut}}\ge\gamma_L$, and for $f,f'\in\ell^2(\Omega)$ supported in sets of diameter $\le s$, both at distance $\ge R\ge4$ from $S$: $|\ip f{(T_{\mathrm{cut}}^{-1}-T^{-1})f'}|\le c_F(s)\norm f\norm{f'}\exp[-\frac{\theta_0}4(R+\dist(\supp f,\supp f'))]$, $c_F$ depending on $(d,L,a,\gamma_L,s)$ only. (For $R<4$ the bound follows from the Combes--Thomas bound applied to each resolvent.)
\end{lemma}

\begin{proof}
(1) The proof of \cref{lem:a06A:coer} drops all bonds between cells, and the cut bonds join cells since $Y$ is a union of large blocks, so $T_{\mathrm{cut}}\ge\gamma_L$; \cref{lem:a06A:CT} holds for $T_{\mathrm{cut}}$ with the bond set of $\Omega_{\mathrm{cut}}$ (its proof uses only that each site is $b_+$ of at most $d$ bonds). Equivalently: $\Omega_{\mathrm{cut}}=(\Omega\setminus B^{k-1}(Y))\sqcup B^{k-1}(Y)$ is a disjoint union of two unions of big blocks, $T_{\mathrm{cut}}=T_{\Lambda\setminus Y}(\Omega\setminus B^{k-1}(Y))\oplus T_{\Lambda\cap Y}(B^{k-1}(Y))$, and the two lemmas apply to each summand. (2) Let $\zeta(x)\defeq\min\{1,\max\{0,\frac4R\dist(x,S)-1\}\}$: $\zeta=0$ where $\dist(\cdot,S)\le R/4$, $\zeta=1$ where $\dist(\cdot,S)\ge R/2$, $|\zeta(x)-\zeta(x')|\le\frac4R|x-x'|$; so $\zeta=1$ on $\supp f\cup\supp f'$ and $\zeta$ vanishes at both ends of every cut bond (they lie within $\eta\le1\le R/4$ of $S$). (3) Put $u\defeq T^{-1}f'$; $T(\zeta u)=f'+[T,\zeta]u$. With $T-T_{\mathrm{cut}}=D_J^*D_J$, $D_J$ the covariant difference on the cut bonds, $D_J(\zeta u)=0$, so $T_{\mathrm{cut}}(\zeta u)=T(\zeta u)$ and $\zeta u=T_{\mathrm{cut}}^{-1}(f'+[T,\zeta]u)$; pairing with $f$: $\ip f{T_{\mathrm{cut}}^{-1}f'}-\ip f{T^{-1}f'}=-\ip f{T_{\mathrm{cut}}^{-1}[T,\zeta]T^{-1}f'}$. (4) With $\zeta$ acting on bond functions by $\zeta(b_-)$, $D_A(\zeta\phi)=\zeta D_A\phi+K\phi$, $(K\phi)(b)\defeq\eta^{-1}(\zeta(b_+)-\zeta(b_-))U(A_b)\phi(b_+)$, $\norm K\le4\sqrt d/R$; $[D_A^*D_A,\zeta]=D_A^*K-K^*D_A$; the potential commutes with $\zeta$; $\norm{[\Pi_B,\zeta]}\le2\operatorname{osc}_B\zeta\le8\sqrt dL/R$. Every term of $[T,\zeta]$ is supported in both variables in the layer $\mathcal L\defeq\{R/4-c\le\dist(\cdot,S)\le R/2+c\}$, $c=\sqrt dL$, and the bonds where $K\ne0$ are not cut bonds. (5) Cover $\mathcal L$ by unit cubes $Q$; the Combes--Thomas bounds for $T_{\mathrm{cut}}$ and $T$ give $\norm{\mathbf 1_QD_{A,\mathrm{cut}}T_{\mathrm{cut}}^{-1}f},\norm{\mathbf 1_QT_{\mathrm{cut}}^{-1}f}\le Ce^{-\theta_0\dist(Q,\supp f)}\norm f$ and the same for $T^{-1}f'$; writing the pairing as $\ip{D_{A,\mathrm{cut}}T_{\mathrm{cut}}^{-1}f}{KT^{-1}f'}-\ip{KT_{\mathrm{cut}}^{-1}f}{D_AT^{-1}f'}+\sum_Bc_B\ip{T_{\mathrm{cut}}^{-1}f}{[\Pi_B,\zeta]T^{-1}f'}$, each term couples neighbouring cubes only, so its modulus is $\le CR^{-1}\sum_{Q\subset\mathcal L}e^{-\theta_0(\dist(Q,\supp f)+\dist(Q,\supp f'))}\norm f\norm{f'}$. (6) For $Q\subset\mathcal L$ both distances are $\ge R/2-2c$ and their sum is $\ge\dist(\supp f,\supp f')-\sqrt d$; so $\theta_0(a+a')\ge\frac{\theta_0}4(R-4c)+\frac{\theta_0}4(\dist(\supp f,\supp f')-\sqrt d)+\frac{\theta_0}2a$, and $\sum_Qe^{-(\theta_0/2)\dist(Q,\supp f)}\le C(1+s)^d$ uniformly in the torus.
\end{proof}

\begin{lemma}[{the operators of \cite[(2.27), (2.46)]{Ba2} under the cut}]\label{lem:a06f:F2}
Take $Y=\Lambda_7$, $\Omega=B^{k-1}(\Lambda_2)$. Write $\Delta\defeq\Delta^{(k-1)}(\Omega,A)$, $\Delta^{\mathrm{cut}}\defeq\Delta^{(k-1)}(\Omega_{\mathrm{cut}},A)$; $C_\Lambda$, $C^{\mathrm{cut}}_\Lambda$ for $\Lambda\subset\Lambda_2$; $\Delta^{(k)}_\Lambda$, $\Delta^{(k),\mathrm{cut}}_\Lambda$ the coarse forms built from them, and $\Delta^{(k)}\defeq\Delta^{(k)}_{\Lambda_2}$. There are $c,\theta_*>0$, depending on $(d,N,\norm q,a,L,\gamma_L)$ only, such that for unit points $z,z'$ and coarse points $y,y'$: \textup{(a)} $|\Delta|,|\Delta^{\mathrm{cut}}|(z,z')\le ce^{-\theta_*|z-z'|}$, the same for $C_\Lambda$, $C^{\mathrm{cut}}_\Lambda$ and for the coarse forms; \textup{(b)} $|\Delta-\Delta^{\mathrm{cut}}|(z,z')\le ce^{-\theta_*(d_S(z)+d_S(z')+|z-z'|)}$, the same for $C_\Lambda-C^{\mathrm{cut}}_\Lambda$ and $\Delta^{(k)}_\Lambda-\Delta^{(k),\mathrm{cut}}_\Lambda$ (with $d_S(y)\defeq\dist(B^{k-1}B(y),S)$); \textup{(c)} for $y,y'\in\Lambda'$, $|\Delta^{(k)}_\Lambda-\Delta^{(k)}|(y,y')\le ce^{-\theta_*(|y-y'|+\dist(y,\Lambda_2\setminus\Lambda)+\dist(y',\Lambda_2\setminus\Lambda))}$; \textup{(d)} block structure: $G_{k-1}(\Omega_{\mathrm{cut}},\Lambda,A)=G_{k-1}(\Omega_1,\Lambda\cap\Lambda_7^c,A)\oplus G_{k-1}(B(\Lambda_7),\Lambda\cap\Lambda_7,A)$, and likewise $\Delta^{\mathrm{cut}}$, $C^{\mathrm{cut}}_\Lambda$, $\Delta^{(k),\mathrm{cut}}_\Lambda$ split into their $\Omega_1$- and $B(\Lambda_7)$-blocks; \textup{(e)} for products $K_1K_2K_3$ of up to three kernels of \textup{(a)} (and $Q^*$, block-local), $|K_1K_2K_3-K^{\mathrm{cut}}_1K^{\mathrm{cut}}_2K^{\mathrm{cut}}_3|(z,z')\le c'e^{-(\theta_*/3)(d_S(z)+d_S(z')+|z-z'|)}$.
\end{lemma}

\begin{proof}
(a) The first bound of \cref{prop:a06A:Ba1P22} (\cite[(2.27)]{Ba1}) for $\Delta$ and \cref{lem:a06A:blockinv}(c) (\cite[(1.16)]{Ba4}) for $C_\Lambda$: on $\Omega$ directly, and on $\Omega_{\mathrm{cut}}$ because their proofs use only \cref{lem:a06A:coer,lem:a06A:CT} and the algebra of \cref{lem:a06A:blockinv}(a), all valid on $\Omega_{\mathrm{cut}}$ by \cref{lem:a06f:F1}(1) (equivalently, by (d), whose proof is algebraic: each operator on $\Omega_{\mathrm{cut}}$ is the direct sum of its blocks on the two unions of big blocks $\Omega_1$ and $B^{k-1}(\Lambda_7)$); for the coarse kernels, $|(QKQ^*)(y,y')|\le L^d\max_{z\in B(y),z'\in B(y')}|K(z,z')|$ and the diagonal term of \cite[(2.27)]{Ba2} is local. (b) For $\Delta$: $(\Delta-\Delta^{\mathrm{cut}})(z,z')=a_{k-1}^2\ip{q(z)}{(G_{k-1}(\Omega_{\mathrm{cut}},A)-G_{k-1}(\Omega,A))q(z')}$, $q(z)\defeq Q_{k-1}^*\delta_z$ of diameter $\le\sqrt d$ and norm $1$; \cref{lem:a06f:F1} with $\Lambda=\emptyset$ and $R=\min(d_S(z),d_S(z'))$, and $\min(d_S(z),d_S(z'))+|z-z'|\ge\frac12(d_S(z)+d_S(z')+|z-z'|)-c''$. For $C_\Lambda$: by the block inversion on $\Omega$ and on $\Omega_{\mathrm{cut}}$ (algebra, needing only $T_\Lambda>0$), $W$ depends only on the pins, so $(C_\Lambda-C^{\mathrm{cut}}_\Lambda)(y,y')=\ip{W\delta_y}{(G_{k-1}(\Omega,\Lambda)-G_{k-1}(\Omega_{\mathrm{cut}},\Lambda))W\delta_{y'}}$ with $W\delta_y$ supported in the $L$-block of $y$, norm $\le1$; \cref{lem:a06f:F1} with $s=\sqrt dL$. The coarse forms follow. (c) \cref{lem:a06A:CLdiff} (\cite[(1.18)]{Ba4}; here $\Omega^{(k-1)}=\Lambda_2$, so $C^{(k-1)}(\Omega,A)=C_{\Lambda_2}$ and the distances are to $\Lambda_2\setminus\Lambda$), then the coarse step. (d) The Neumann Laplacian of $\Omega_{\mathrm{cut}}$ is the direct sum over $\Omega_1$ and $B(\Lambda_7)$; the pins are block-local and every block lies in one part, since $\Lambda_7$ is a union of large blocks; so $T_\Lambda(\Omega_{\mathrm{cut}})$ is block diagonal, $\Delta^{\mathrm{cut}}(z,z')$ vanishes unless $B(z),B(z')$ lie in the same part, $X^{\mathrm{cut}}_\Lambda$ is block diagonal and so is its inverse. (e) Telescope $K_1K_2K_3-K^{\mathrm{cut}}_1K^{\mathrm{cut}}_2K^{\mathrm{cut}}_3=(K_1-K^{\mathrm{cut}}_1)K_2K_3+K^{\mathrm{cut}}_1(K_2-K^{\mathrm{cut}}_2)K_3+K^{\mathrm{cut}}_1K^{\mathrm{cut}}_2(K_3-K^{\mathrm{cut}}_3)$; in the first term, with $X\defeq d_S(z)+d_S(w)+|z-w|$, $Y\defeq|w-w'|$, $Z\defeq|w'-z'|$ and the target $T\defeq d_S(z)+d_S(z')+|z-z'|$, the $1$-Lipschitz property of $d_S$ gives $T\le X+2Y+2Z$, so $\theta_*(X+Y+Z)\ge\frac{\theta_*}3T+\frac{\theta_*}3(Y+Z)$, and $e^{-(\theta_*/3)(Y+Z)}$ pays the sums over $w,w'$ uniformly in the torus; the other terms are the same with the roles exchanged. (A target rate $\theta_*/2$ is false: nothing would be left for the $w'$-sum.)
\end{proof}

\begin{proposition}[{\cite[p.~568]{Ba2}}]\label{prop:a06f:C10}
Assume (G), (FR), (RG), (A). For every $\kappa>0$ there is $t_0(\kappa)>0$ such that for $t=L^{k-1}\eps\le t_0$ each of the following replacements changes the exponential of \cite[(2.46)]{Ba2} by at most $Ct^\kappa|\Lambda^{(k-1)}_6|$, $C$ independent of $k$, $\eps$, the torus and the fields: \textup{(i)} $T_{10}\to$ the right side of \cite[(2.47)]{Ba2} without its $O$-term; \textup{(ii)} $T_4\to$ its analogue with $A$ in place of $\varphi$, every vector form the torus form; \textup{(iii)} $T_3\to T_3$ with $A$ replaced by $(\Lambda'_5\cap\Lambda'^c_6)A$; \textup{(iv)} $T_9\to T_9$ with $\varphi$ replaced by $(\Lambda'_5\cap\Lambda'^c_6)\varphi$ and $\Delta^{(k-1)}$, $C^{(k-1)}_{\Lambda_5}$ taken on $\Omega_{\mathrm{cut}}$, i.e.\ (by \cref{lem:a06f:F2}(d)) $\Delta^{(k-1)}(\Omega_1)$ and $C^{(k-1)}_{\Lambda_5\cap\Lambda_7^c}(\Omega_1)$; \textup{(v)} $T_7,T_8\to$ the same with every operator on $\Omega_{\mathrm{cut}}$. In particular $\kappa>d$ is allowed, as \cite[(2.48)]{Ba2} requires.
\end{proposition}

\begin{proof}
$\theta_*$ from \cref{lem:a06f:F2}; all sums are over unit or coarse points with torus-free constants. (1) Counting: for $z\notin\Lambda_6$, $d_S(z)\ge r-c$ by (G), and every unit cube meeting $S$ is adjacent to a cube of $\Lambda_7$, so for any set $X$ with $d_S\ge r-c$ on $X$, $\sum_{z\in X}e^{-\theta d_S(z)}\le e^{-\theta(r-c)/2}\sum_{Q\cap S\ne\emptyset}\sum_ze^{-\theta\dist(z,Q)/2}\le C_\theta e^{-\theta r/2}|\Lambda_7|$. (2) (i): write $\varphi=\varphi_o+\varphi_m+\varphi_i$ supported in $\Lambda'_5\setminus\Lambda'_6$, $\Lambda'_6\setminus\Lambda'_7$, $\Lambda'_7$; with $\Delta_5\defeq\Delta^{(k)}_{\Lambda_5}(\Omega)$ (symmetric), $\ip\varphi{\Delta_5\varphi}=\ip{\varphi_o}{\Delta_5\varphi_o}+2\ip{\varphi_o}{\Delta_5\varphi_m}+2\ip{\varphi_o}{\Delta_5\varphi_i}+\ip{\varphi_m+\varphi_i}{\Delta_5(\varphi_m+\varphi_i)}$; by \cref{lem:a06f:F2}(d) the first two terms of (2.47) are $\frac12\ip{\varphi_o}{\Delta^{(k),\mathrm{cut}}_{\Lambda_5}\varphi_o}+\ip{\varphi_o}{\Delta^{(k),\mathrm{cut}}_{\Lambda_5}\varphi_m}$ and the third is $\frac12\ip{\varphi_m+\varphi_i}{\Delta^{(k)}(\varphi_m+\varphi_i)}$. The difference is $E_1+E_2+E_3$: $E_1\defeq\frac12\ip{\varphi_o}{(\Delta_5-\Delta^{(k),\mathrm{cut}}_{\Lambda_5})\varphi_o}+\ip{\varphi_o}{(\Delta_5-\Delta^{(k),\mathrm{cut}}_{\Lambda_5})\varphi_m}$, by \cref{lem:a06f:F2}(b) and (1) $|E_1|\le cM^2\sum_{y\in\supp\varphi_o}e^{-\theta_*d_S(y)}\sum_{y'}e^{-\theta_*|y-y'|}\le CM^2e^{-\theta_*r/2}|\Lambda_7|$; $E_2\defeq\ip{\varphi_o}{\Delta_5\varphi_i}$, with $|y-y'|\ge r-c$ between $\supp\varphi_o$ and $\Lambda'_7$ (G), $|E_2|\le CM^2e^{-\theta_*r/2}|\Lambda_7|$; $E_3\defeq\frac12\ip{\varphi_m+\varphi_i}{(\Delta_5-\Delta^{(k)})(\varphi_m+\varphi_i)}$, by \cref{lem:a06f:F2}(c) with $\dist(y,\Lambda_2\setminus\Lambda_5)\ge r-c$ on $\Lambda_6$, $|E_3|\le CM^2e^{-\theta_*r}|\Lambda_6|$. (3) (ii): as (i) with $E_1$ absent (no cut) and the torus operators; $E_3$ by \cref{lem:a06A:CLdiff} (\cite[(1.18)]{Ba4}) on the torus, at $A=0$ with $\Omega=T_\eps$ and $\Lambda=\Lambda_5$ (\cref{lem:a06A:torus}). (4) (iii): the error is $a(L^k\eps)^{-2}\ip{A'_{\mathrm{out}}}{\Delta C_{\Lambda_5}Q^*(\Lambda'_6A)}$ with $A'_{\mathrm{out}}$ supported in $\Lambda_5^c$; by \cref{lem:a06f:F2}(a) $|(\Delta C_{\Lambda_5}Q^*)(x,y')|\le ce^{-\theta_*|x-y'|/2}$ and $|x-y'|\ge r-c$ (G), so the error is $\le Ct^{-2}M^2e^{-\theta_*r/4}|\Lambda_6|$. (5) (iv): restricting $\varphi$ is as (4); replacing $\Omega$ by $\Omega_{\mathrm{cut}}$ in $\Delta$, $C_{\Lambda_5}$ changes the kernel at $(x,y')$ by $\le ce^{-(\theta_*/3)(d_S(x)+d_S(y')+|x-y'|)}$ (\cref{lem:a06f:F2}(e)) with $d_S(x)\ge2(r-c)-1$ on $\Lambda_3\cap\Lambda_5^c$, hence by (1) $\le Ct^{-2}M^2e^{-\theta_*r/6}|\Lambda_7|$; then \cref{lem:a06f:F2}(d): $\varphi'$ is supported off $B(\Lambda_7)$ and $Q^*\varphi$ lies in $\Lambda_5\setminus\Lambda_6\subset\Lambda_7^c$, so only the $\Omega_1$-blocks act. (6) (v): as (5) with both outer points in $\Lambda_3\cap\Lambda_5^c$. (7) Every error is $\le Ct^{-4}M^2e^{-\theta_*r(t)/6}(|\Lambda_6|+|\Lambda_7|)\le C't^{-8}e^{-\theta_*r(t)/6}|\Lambda_6|$ ($M\le Ct^{-2}$, $\Lambda_7\subset\Lambda_6$), and $e^{-\theta_*r(t)/6}\le t^{\kappa+9}$ for $t\le t_0(\kappa)$ since $\varrho>1$ (the bound of Step~3(D$_2$), Geometry).
\end{proof}

The proof uses neither the Theorem of \cite{Ba4} nor the two-sided depth bounds; ``Proposition I.2.1'' of the print is not needed for these terms. The terms omitted at \cite[p.~568]{Ba2} (``described more precisely later, in (2.109)'') are the quadratic-form step treated next.

% a06_BALtorus_ct6: Step 3(F), second part: the quadratic-form step Ba2 (2.108)-(2.109) (the half of the stage of
% Ba2 Prop. 2.6 that concerns no interaction term), and the passage from Ba2 (3.5) at k = 0 to (3.21) and (3.42).

\subsubsection*{The quadratic-form step \cite[(2.108)--(2.109)]{Ba2}}

\emph{Setting.} Step $k$ in scale-$k$ units, $t\defeq t_k$, $\eta\defeq L^{-k}$; $\Omega_A\defeq B^k(\Lambda^{(k-1)'}_2)$ (the propagator domain of \cite[(2.53), (2.56)]{Ba2}), $\Omega_B\defeq B^k(\Lambda^{(k-1)'}_2\cap\Lambda^{(k)c}_5)$ and $\Omega_C\defeq B^k(\Lambda^{(k)}_2)$ (the domains of \cite[(2.108)]{Ba2}), all unions of big blocks; $\tilde B\defeq A^{(k)}$ on $\Omega_A$. On $\Omega_C$, $\tilde B=B^{(k+1),\eta}$: $\theta_{k+1}=1$ on $B^k(\Lambda^{(k)}_2)$ \cite[pp.~567, 576]{Ba2}, $\Omega_C\subset B^k(\Lambda^{(k)}_0)$, and by \cite[p.~579]{Ba2} the configuration there is $(1-\theta_{k+1})A^{(k)}+\theta_{k+1}B^{(k+1),\eta}$; the operators $G_k(\Omega_C,\cdot)$, $\Delta^{(k)}(\Omega_C,\cdot)$, $C^{(k)}_\Lambda(\Omega_C,\cdot)$ depend on the field only through the bonds of $\Omega_C$ \cite[(1.3)--(1.4)]{Ba4}. So the replacement of \cite[Prop.~2.6]{Ba2} is the pure domain change $\Omega_A\to\Omega_C$: its interaction-term half is \cref{thm:a07:dom}, and its quadratic-form half is the following. Layer distances \cite[(2.8)]{Ba2}: $\dist(\Lambda^{(k)}_{i+j},\Lambda^{(k)c}_i)\ge j(r-1)$ and $\dist(\Lambda^{(k-1)'}_7,\Lambda^{(k-1)'c}_2)\ge r-1$ ($L=5$). Standing lower bound: $R\ge R_0+3$, so that $r-2\ge R_0+1$ ($R$ may be enlarged, \cite[(2.7)]{Ba2}).

\begin{lemma}[discrete trace inequality]\label{lem:a06f:T}
Let $u$ be a function on a box $Y$ of $\eta\Z^d$ with all its bonds, $z$ a site, $e$ a lattice direction and $S_z\defeq\{z+m\eta e:0\le m<\eta^{-1}\}\subset Y$. Then $|u(z)|^2\le2\sum_{v\in S_z}\eta|u(v)|^2+2\sum_{b'\subset S_z}\eta|(D_Au)(b')|^2$.
\end{lemma}

\begin{proof}
$|U|=1$ gives $||u(b'_+)|-|u(b'_-)||\le\eta|D_Au(b')|$, so $|u(z)|\le|u(v)|+\sum_{b'\subset[z,v]}\eta|D_Au(b')|$ for every $v\in S_z$; average over $v$ with weight $\eta$ (total weight $1$), apply Cauchy--Schwarz and square.
\end{proof}

\begin{lemma}[domain difference of $\Delta^{(k)}$, decaying from the new boundary]\label{lem:a06f:Q2}
Let $\Omega'\subset\Omega$ be unions of big blocks, $A$ regular on $\Omega$, $\Sigma$ the set of bonds of $\Omega$ with exactly one end in $\Omega'$, and $\dist(\Sigma,\Omega^c)\ge R_0+1$. Then for $y,y'\in\Omega'^{(k)}$: $|\Delta^{(k)}(\Omega',A;y,y')-\Delta^{(k)}(\Omega,A;y,y')|\le c_Q\exp[-\delta_Q(|y-y'|+\dist(B(y),\Sigma)+\dist(B(y'),\Sigma))]$, $\delta_Q\defeq\min(\delta_0,\theta_0)/4$, $c_Q$ depending on $(c_0,\delta_0,\gamma_\emptyset,\theta_0,a,d)$. (This is stronger than \cite[(2.29)]{Ba1}, which decays from all of $\Omega'^c$.)
\end{lemma}

\begin{proof}
(1) With $H_\Omega\defeq-\Delta^{\eta,N}_{A,\Omega}+m^2t^2+a_kP_k(A)$: $H_\Omega=(H_{\Omega'}\oplus H_{\Omega\setminus\Omega'})+D_\Sigma^*D_\Sigma$, $D_\Sigma$ the covariant difference on the bonds of $\Sigma$ ($P_k$ and the mass are block diagonal and both pieces are unions of blocks). For $f'$ on $\Omega'$, $u'\defeq G_k(\Omega')f'$ and $u^-\defeq u'\oplus0$: $H_\Omega u^-=f'+D_\Sigma^*D_\Sigma u^-$, hence $(G_k(\Omega)f')|_{\Omega'}-u'=-(G_k(\Omega)D_\Sigma^*D_\Sigma u^-)|_{\Omega'}$. (2) With $q(y)\defeq Q_k^*(A)\delta_y$: $\delta\Delta(y,y')=a_k^2\ip{q(y)}{(G_k(\Omega)-G_k(\Omega'))q(y')}=-a_k^2\ip{D_\Sigma G_k(\Omega)q(y)}{D_\Sigma u^-}$, $u'=G_k(\Omega')q(y')$. (3) For $b=(z,w)\in\Sigma$ with $z\in\Omega'$, $|(D_\Sigma u^-)(b)|=\eta^{-1}|u'(z)|$. (4) By \cite[(2.25)]{Ba1} at $b$, the restricted form of the theorem of \cite{Ba4}, on the torus by \cref{lem:a06A:torus} (its hypotheses hold: $\Omega$ is a union of big blocks, $A$ is regular on $\Omega$, $\norm{q(y)}_\infty=1$, and $\dist(b,\Omega^c)\ge R_0$), $|D_bG_k(\Omega)q(y)|\le c_0e^{-\delta_0\dist(b,B(y))}$. (5) For a unit block $Y$ of $\Omega'$ let $\Sigma_Y$ be the bonds of $\Sigma$ with $z_b\in Y$ ($\#\Sigma_Y\le2d\eta^{-(d-1)}$); each $z_b$ lies on a face of $Y$, the segments away from $w$ stay in $Y$ and those from one face partition $Y$, so \cref{lem:a06f:T} over the $\le2d$ faces gives $\sum_{b\in\Sigma_Y}\eta^{d-1}|u'(z_b)|^2\le4d(\norm{\mathbf 1_Yu'}^2+\norm{\mathbf 1_YDu'}^2)$. (6) \Cref{lem:a06A:CT} on $\Omega'$ with $T=H_{\Omega'}\ge\gamma_\emptyset$ (\cref{lem:a06A:coer} with $\Lambda=\emptyset$; $\Omega'$ is a union of big blocks and $A$ is regular on it) gives $\norm{\mathbf 1_Yu'}\le2\gamma_\emptyset^{-1}e^{-\theta_0\dist(Y,B(y'))}$ and $\norm{\mathbf 1_YDu'}\le Ce^{-\theta_0\dist(Y,B(y'))}$, no depth restriction. (7) Cauchy--Schwarz on $\Sigma_Y$ with weights $\eta^{d-1}$: $|\delta\Delta|\le a^2\sum_Y\sum_{\Sigma_Y}\eta^d|D_bGq(y)|\eta^{-1}|u'(z_b)|\le C\sum_Ye^{-\delta_0\dist(Y,B(y))-\theta_0\dist(Y,B(y'))}$ over the blocks $Y$ adjacent to $\Sigma$. (8) For such $Y$, $\dist(Y,B(y))+\dist(Y,B(y'))\ge\frac12(\dist(B(y),B(y'))+\dist(B(y),\Sigma)+\dist(B(y'),\Sigma))-2$; keep half of the exponent for the sum over $Y$ and use $|y-y'|\le\dist(B(y),B(y'))+1$.
\end{proof}

\begin{lemma}[{\cite[(2.108)]{Ba2} is an identity}]\label{lem:a06f:C53}
For $\varphi$ on $\Lambda^{(k-1)'}_6$ put $\varphi_o\defeq(\Lambda^{(k-1)'}_6\cap\Lambda^{(k)c}_3)\varphi$, $\varphi_i\defeq\Lambda^{(k)}_3\varphi$, $\varphi_e\defeq(\Lambda^{(k)}_3\cap\Lambda^{(k)c}_4)\varphi$, $\varphi_4\defeq\Lambda^{(k)}_4\varphi$ ($\varphi_i=\varphi_e+\varphi_4$), $\Delta_A\defeq\Delta^{(k)}(\Omega_A,\tilde B)$, $\Delta_B\defeq\Delta^{(k)}(\Omega_B,\tilde B)$, $\Delta_C\defeq\Delta^{(k)}(\Omega_C,B^{(k+1),\eta})=\Delta^{(k)}(\Omega_C,\tilde B)$. Then $\frac12\ip{\Lambda'_6\varphi}{\Delta_A\Lambda'_6\varphi}=\frac12\ip{\varphi_o}{\Delta_B\varphi_o}+\ip{\varphi_o}{\Delta_B\varphi_e}+\frac12\ip{\varphi_i}{\Delta_C\varphi_i}+\frac12\ip{\Lambda'_6\varphi}{H_k\Lambda'_6\varphi}$ with $\frac12\ip\varphi{H_k\varphi}\defeq\frac12\ip{\varphi_o}{(\Delta_A-\Delta_B)\varphi_o}+\ip{\varphi_o}{(\Delta_A-\Delta_B)\varphi_e}+\ip{\varphi_o}{\Delta_A\varphi_4}+\frac12\ip{\varphi_i}{(\Delta_A-\Delta_C)\varphi_i}$. With the block-spin term $\frac12aL^{d-2}\sum|\psi-Q(\theta_k\tilde B)\varphi|^2$, unchanged on both sides, this is \cite[(2.108)]{Ba2} with the factor $\frac12$ of its fourth term restored.
\end{lemma}

\begin{proof}
$\Lambda^{(k)}_3\subset\Lambda^{(k)}_0\subset\Lambda^{(k-1)'}_7\subset\Lambda^{(k-1)'}_6$, so $\Lambda'_6\varphi=\varphi_o+\varphi_i$; the $\Delta$'s are real symmetric ($U\in SO(N)$, $q$ real antisymmetric); expand and subtract. All forms are defined: $B^k(\Lambda'_6\cap\Lambda^c_3)$ and $B^k(\Lambda_3\cap\Lambda_4^c)$ lie in $\Omega_B$, $B^k(\Lambda_3)$ in $\Omega_C$.
\end{proof}

\begin{lemma}[{\cite[(2.109)]{Ba2}}]\label{lem:a06f:C54}
For $y,y'\in\Lambda^{(k-1)'}_6$ the kernel of $H_k$ obeys $|h_k(y,y')|\le C_He^{-\delta_H(r-2)}e^{-\delta_H|y-y'|}$, $\delta_H\defeq\min(\delta_Q,\delta_7/2)/2$ ($\delta_7$ the rate of \cref{prop:a06A:Ba1P22}, the form of \cite[(2.27)]{Ba1} proved in Step~3(A)); and under \cite[(2.55)]{Ba2}, for $t\le t_0(\kappa)$, $\frac12|\ip{\Lambda'_6\varphi}{H_k\Lambda'_6\varphi}|\le Cc_1^2p(t/L)^2\lambda(t/L)^{-1/2}e^{-(\delta_H/2)(r-2)}|\Lambda^{(k)}_3|\le C_\kappa t^\kappa|\Lambda^{(k)}_3|$, the printed error of \cite[(2.108)]{Ba2}.
\end{lemma}

\begin{proof}
Four pieces. (oo), (oe): \cref{lem:a06f:Q2} with $(\Omega,\Omega')=(\Omega_A,\Omega_B)$; $\Omega_A\setminus\Omega_B=B^k(\Lambda^{(k)}_5)$, so $\Sigma$ consists of the bonds crossing $\partial B^k(\Lambda^{(k)}_5)$, with $\dist(\Sigma,\Omega_A^c)\ge\dist(\Lambda^{(k)}_5,\Lambda^{(k-1)'c}_2)-1\ge r-2\ge R_0+1$; for $y\in\Lambda_3^c$, $\dist(B(y),\Sigma)\ge2(r-1)-2$, for $y'\in\Lambda_3\cap\Lambda_4^c$, $\dist(B(y'),\Sigma)\ge r-2$; so $|(\Delta_A-\Delta_B)(y,y')|\le c_Qe^{-(\delta_Q/2)(r-2)}e^{-(\delta_Q/2)(|y-y'|+\dist(B(y),\Sigma)+\dist(B(y'),\Sigma))}$. (o4): for $y\in\Lambda_3^c$, $y'\in\Lambda_4$, $|y-y'|\ge r-1$, and \cref{prop:a06A:Ba1P22} (\cite[(2.27)]{Ba1}, under the hypotheses of \cref{lem:a06f:Q2} on $\Omega_A$) gives $|\Delta_A(y,y')|\le c_7e^{-(\delta_7/2)(r-1)}e^{-(\delta_7/2)|y-y'|}$. (ii): \cref{lem:a06f:Q2} with $(\Omega_A,\Omega_C)$, $\Sigma$ the bonds crossing $\partial\Omega_C$, $\dist(\Sigma,\Omega_A^c)\ge r-2$, and $\dist(B(y),\Sigma)\ge r-2$ for $y\in\Lambda_3$. Assembling the kernels gives the kernel bound. Error: $\sup|\varphi|\le c_1p(t/L)\lambda(t/L)^{-1/4}$ on $\Lambda'_6\subset\Lambda'_{-1}$ by \cite[(2.55)]{Ba2}; the (oo) sum over $y,y'\in\Lambda_3^c$ is $\le Ce^{-(\delta_Q/2)(r-2)}\sum_ye^{-(\delta_Q/4)\dist(B(y),\Sigma)}\le Ce^{-(\delta_Q/2)(r-2)}|\Lambda^{(k)}_5|$ (at most $|\Lambda_5|$ unit blocks are adjacent to $\Sigma$ from inside); in (oe), (o4), (ii) the second index runs over $\Lambda_3$ or $\Lambda_4$, giving $\le C|\Lambda^{(k)}_3|$ times the exponential; then $|\Lambda_5|\le|\Lambda_3|$, $p(t/L)\le(1+\ln L)^{\bar p}p(t)$, $\lambda(t/L)=\lambda_k/L$, and $p^2\lambda^{-1/2}e^{-\delta r}\le t^\kappa$ for $t\le t_0(\kappa)$.
\end{proof}

With \cite[(2.29)]{Ba1} in place of \cref{lem:a06f:Q2} the error would be $O(t^\kappa)|\Lambda^{(k-1)'}_6|$, still summable at $\kappa>d$ but not the printed $|\Lambda^{(k)}_3|$.

\subsubsection*{From \cite[(3.5)]{Ba2} at $k=0$ to \cite[(3.21)]{Ba2}, and the payment of \cite[Sect.~3C]{Ba2}}

\begin{lemma}[monotone nesting]\label{lem:a06f:C141}
Every renormalisation transformation $T^t_{a,L}$, $T^t_{a,L,\tilde A}$ of \cite[Sect.~3]{Ba2} is $F\mapsto\int t(\cdot,F\text{-arguments})F$ with a kernel $t\ge0$ normalised by \cite[(3.32)]{Ba2}; hence $F\le F'$ pointwise implies $T[F]\le T[F']$, and nested brackets preserve inequalities.
\end{lemma}

\begin{lemma}[{\cite[(3.21)]{Ba2}}]\label{lem:a06f:C142}
\cite[(3.5)]{Ba2} holds at $k=0$ (\cref{prop:a06d:C211}, the finite downward induction from the base case). At $k=0$ the innermost bracket of \cite[(3.5)]{Ba2} is the bare Gaussian density $\exp(-\frac12\ip{A_0}{(-\Delta^\eps+\lmu^2)A_0})\exp(-\frac12\ip{\phi_0}{(-\Delta^\eps_{\tilde A^\eps}+m^2)\phi_0})$ with the bare characteristic functions \cite[(2.34)--(2.42), (2.47)--(2.49)]{Ba2}; each characteristic function factorises into a vector condition and a scalar condition with the vector configuration as parameter, $\zeta$ placed inside the scalar bracket; integrating first over the scalar fields at fixed vector fields (Tonelli) and \cref{lem:a06f:C141} give \cite[(3.21)]{Ba2}, whose last factor is $\exp(\sum_{k=1}^KC_\kappa t_k^{\kappa_0}|T_\eps|)\le\exp(C_\kappa\eps_0^{\kappa_0}(1-L^{-\kappa_0})^{-1}|T_\eps|)$, i.e.\ \cite[(2.118)]{Ba2}.
\end{lemma}

\begin{lemma}[{the residue of \cite[Sect.~3B]{Ba2}}]\label{lem:a06f:C143}
In \cite[(3.30)--(3.40)]{Ba2} the only factors left besides $\prod\zeta''$ and $\exp(E_0)$ are: from \cite[(3.31)]{Ba2}, $\exp(c_{3.29}t_k^{\kappa_0}|\Lambda_k|)$ with $c_{3.29}$ the remainder constant of \cite[Prop.~3.1, (3.29)]{Ba2} (\cref{thm:a06:BALA}, the $\eps_0$-freedom paragraph of Step~3(C$_2$)); from \cite[(3.33)]{Ba2}, $2^{N/2}\exp(-\frac14ap(t_k)^2)\le\exp(-\frac18ap(t_k)^2)$ per exceptional block, valid once $b_0^2\ge4N\ln2/a$ (a condition on $b_0$, fixed before $\eps_0$); from \cite[(3.36), (3.40)]{Ba2}, the Jacobians $\exp(\frac N2\ln2\sum_k|\Lambda_k|)$ and $\exp(\frac d2\ln2\sum_k|\Lambda_k|)$; and from \cite[(3.38)]{Ba2} the diamagnetic inequality and the normalisations \cite[(3.32)]{Ba2} (equalities). Here $\Lambda_k\defeq\Lambda^{(k-1)'}_5\cap\Lambda^{(k)c}_5$, counted in points of the scale-$k$ unit lattice. So the $O(1)$ of \cite[(3.39)--(3.40)]{Ba2} is $c_5\defeq\frac{N+d}2\ln2+c_{3.29}\eps_0^{\kappa_0}\le\frac{N+d}2\ln2+c_{3.29}$, and the $O(1)$ in front of $(L^k\eps)^{\kappa_0}|\Lambda^{(k-1)'}_7\cap\Lambda^{(k)c}_7|$ in \cite[(2.117), (3.42)]{Ba2} is $c_7\defeq C_{2.57}+C_{2.78}$ (Step~3(C$_2$), the $O(1)$ lines). \cite[(3.40)]{Ba2} prints $\sum_{k=1}^K|\Lambda_k|$ where \cite[(3.36), (3.42)]{Ba2} print $\sum_{k=0}^K$; the constants use $\sum_{k=0}^K$.
\end{lemma}

\begin{lemma}[the payment]\label{lem:a06f:C144}
The residue $\exp(c_5\sum_k|\Lambda_k|)$ of \cref{lem:a06f:C143} is exactly the second factor of the left side of \cite[(3.42)]{Ba2} with $O(1)=c_5$, and the first factor carries $O(1)=c_7$. So \cite[(2.117)]{Ba2}, hence the upper bound, follows from \cite[(3.42)]{Ba2} with the constants $(c_5,c_7)$, which are free of $\eps$, $K$, the torus, the history and $\eps_0$; \cite[(3.42)]{Ba2} itself is \cref{thm:a06f:R7} below, which imposes its own cap on $\eps_0$ after $c_5$, $c_7$ are fixed.
\end{lemma}

\begin{proof}[Proof of \cref{lem:a06f:C141,lem:a06f:C142,lem:a06f:C143,lem:a06f:C144}]
\Cref{lem:a06f:C141}: positivity of the Gaussian kernels. \Cref{lem:a06f:C142}: as stated, the $k=0$ bracket being the bare density by the definitions of the first step (\cite[p.~588]{Ba2}: ``using inequality (3.5) for $k=0$''), and $t_k=L^{k-K}t_K\le L^{k-K}\eps_0$. \Cref{lem:a06f:C143}: arithmetic on the printed steps, with $p(t)\ge b_0$ for the $b_0$-condition. \Cref{lem:a06f:C144}: $|\Lambda_k|=|\Lambda^{(k-1)'}_5\cap\Lambda^{(k)c}_5|$ by definition (\cite[p.~592]{Ba2}: ``$|\cdot|$ means the number of elements''), an identity of terms.
\end{proof}

% a06_BALtorus_ct7: Step 3(F), third part: the region-entropy estimate Ba2 (3.42) at block factor 5 (Sect. 3C of Ba2
% completed), with the counterterm moves of Ba3 p. 422 paid in it, and the g_k-localisation of Q^(k) (Ba2 p. 579).

\subsubsection*{The region sum \cite[(3.42), Sect.~3C]{Ba2} at block factor $5$}

\emph{Setting.} $d=3$, $L=5$, $M\ge1$, $s_b\defeq LM$ the side of a large block in lattice units of $T^{(k)}$ \cite[p.~607]{Ba1}; $R>R_0$, $\theta\defeq s_b/R\le\frac17$; $\varrho\ge1$ and $\bar p$ with $2\bar p>\varrho d+2$ (Ba{\l}aban's $\varrho=2$, $\bar p=6$: exponent $2\bar p-2-\varrho d=4$); $\eps_0\le e^{-1}$; $\hat l_k\defeq1+\ln t_k^{-1}$, $r_k\defeq r(t_k)=R\hat l_k^\varrho$, $p_k\defeq p(t_k)=b_0\hat l_k^{\bar p}$. The metric is the torus sup-metric \cite[(1.3)]{Ba1}; $N_w(X)$ is the closed $w$-neighbourhood; boxes are images of products of integer intervals. The torus $T^{(k)}=\Z^d/\Gamma_k$, $\Gamma_k=L\Gamma_{k+1}$, is aligned with the blocks: $\Gamma_k=\bigoplus_\mu P_{k,\mu}\Z$ with $s_b\mid P_{k,\mu}$ for $k\le K-1$ (a condition on the torus: its sides are multiples of the side of a large block at scale $K-1$, hence at every scale up to $K-1$; cf.\ \cref{def:s10:torus}, whose sides lie in $M\N$, and \cref{def:a07:btorus}); so the large blocks and the $L$-blocks tile $T^{(k)}$ and the prime $X\mapsto X'$ of \cite[(1.16)]{Ba1} is defined. Families and regions are as printed: at scale $k$ a family $F_k=(P_v,Q_v,R_v,P_s,Q_s,R_s)$ inside $\Lambda^{(k-1)'}_7$ ($\Lambda^{(-1)'}_7\defeq T^{(0)}$); each element $e$ has a point set $\mathrm{pts}(e)$ (a $P$-element its $L$-block; a bond or plaquette its ends or corners; an $R$-element its point) and a charge point $x_e\in\mathrm{pts}(e)$ (the lowest corner), with $\mathrm{pts}(e)\subset N_{L-1}(x_e)$; $S_k\defeq\bigcup_e\mathrm{pts}(e)$; $\Lambda^{(k)}_0$ is the union of the large blocks of $\Lambda^{(k-1)'}_7$ at distance $>r_k$ from $S_k$ \cite[p.~570]{Ba2}, $\Lambda^{(k)}_i$ by \cite[(2.8)]{Ba2}, ``admissible'' as on \cite[p.~566]{Ba2}. Write $X_k\defeq\Lambda^{(k)c}_0$, $\mathrm{Col}_k\defeq\Lambda^{(k-1)'}_7\cap\Lambda^{(k)c}_7$, $\mathrm{Sh}_k\defeq\Lambda^{(k-1)'}_5\cap\Lambda^{(k)c}_5$, $N_k\defeq|T^{(k)}|$, $I_k\defeq\{z\in T^{(k)}:\dist(z,T^{(k)}\setminus\Lambda^{(k-1)'}_7)\ge s_b\}$ ($I_0=T^{(0)}$). Constants: $c_{\mathrm{box}}\defeq\min\{a3^{-d}/8,\gamma_03^{-d}/4,3^{-d}\}$ (Ba{\l}aban's $c_0$ of \cite[p.~593]{Ba2}); $c_5,c_7\ge0$, $\kappa_0>0$ the $O(1)$'s and the exponent of \cite[(3.42)]{Ba2} (\cref{lem:a06f:C143}); $n_d\defeq4+2\max\{d,d(d-1)/2\}$ ($n_3=10$); $\nu_L\defeq14+(16+2L)\theta$, $\nu'_L\defeq\nu_L+2+3\theta$, $S_L\defeq\sum_{m\ge0}L^{-m}(1+m\ln L)^\varrho$, $C^0_D\defeq\nu'_L(S_L-1)$, $C_D\defeq C^0_D+14(1+\theta)$, $C_1\defeq n_d\ln2+d+1+c_5+c_7$, $\Theta_L\defeq C_1C_D^dR^d/(c_{\mathrm{box}}b_0^2\ln L)$, $C_{\mathrm{fin}}\defeq L^d(c_7\eps_0^{\kappa_0}+c_5)+e^{-(d+1)}\eps_0^{d+1}/(L-1)$. At $L=5$, $\varrho=2$, $\theta=\frac17$: $S_5=3.470$, $C^0_D=49.76$, $C_D=65.76$.

The printed \cite[Sect.~3C]{Ba2} is not followed: its partition \cite[(3.43)]{Ba2} (cubes with side in $(r,2r)$) need not exist at $L=5$, and the count \cite[(3.45)]{Ba2} with the first inequality of \cite[(3.46)]{Ba2} is false as printed: at $k=1$, with one $R_s$-element at a point $x_*$ of $T^{(0)}$ and $F_1=\emptyset$, $\Lambda^{(0)c}_7$ is a cube of large blocks whose prime is a union of scale-$1$ large blocks only if its corners are divisible by $L$, which fails for $L-1$ of every $L$ positions of $x_*$; then a scale-$1$ large block straddling the face lies in $\Lambda^{(1)c}_0$ and meets $\Lambda^{(0)'}_7$, so the printed $C_1$ is non-empty while $3^d(|P^{(1)}_v|+\dots)=0$, and the first inequality of \cite[(3.46)]{Ba2} reads $1\le\exp(-c_0p^2|C_1|)<1$. The count below is on the interior $I_k$ (\cref{lem:a06f:RC}), and the straddling layer is covered by the inherited boxes (\cref{lem:a06f:RD}).

\begin{lemma}[consequences of the printed rules]\label{lem:a06f:RE}
For every admissible sequence and $k\le K-1$: \textup{(a)} every point of every element of $F_k$ lies in $X_k$; \textup{(b)} $\Lambda^{(k)c}_i\subset\Lambda^{(k)c}_{i+1}$ and $\Lambda^{(k)c}_i\subset N_{i(r_k+s_b-1)}(X_k)$, $i=0,\dots,8$; \textup{(c)} if $z\in X_k\cap I_k$, the large block $\beta\ni z$ lies in $\Lambda^{(k-1)'}_7$ and $\dist(\beta,S_k)\le r_k$; \textup{(d)} $T^{(k)}\setminus\Lambda^{(k-1)'}_7=(\Lambda^{(k-1)c}_7)'\subset X_k$; \textup{(e)} $X_k,\mathrm{Col}_k,\mathrm{Sh}_k\subset\Lambda^{(k)c}_7$; \textup{(f)} if $X_i=\emptyset$ for all $i<j$ and $X_j\ne\emptyset$, then $\Lambda^{(j-1)'}_7=T^{(j)}$ and $X_j\cap I_j=X_j$.
\end{lemma}

\begin{proof}
(a) The large block $\beta$ of a point of $\mathrm{pts}(e)$ has $\dist(\beta,S_k)=0\le r_k$, so it is not one of the blocks of \cite[p.~570]{Ba2}; the large blocks tile $T^{(k)}$ and $\Lambda^{(k)}_0$ is a union of them. (b) \cite[(2.8)]{Ba2} adds the large blocks at distance $\le r_k$ (blocks of $\Lambda_i^c$ included); a point of such a block is within $(s_b-1)+r_k$ of $\Lambda^c_i$; induct. (c) $\beta$ has sup-diameter $s_b-1<s_b$, so $\beta\subset\Lambda^{(k-1)'}_7$, and being a candidate not chosen, $\dist(\beta,S_k)\le r_k$. (d) $\Lambda^{(k-1)}_7$ is a union of $L$-blocks, so $(X^c)'=T^{(k)}\setminus X'$; $\Lambda^{(k)}_0\subset\Lambda^{(k-1)'}_7$. (e) by (b). (f) $X_i=\emptyset$ gives $\Lambda^{(i)c}_7=\emptyset$ by (b), so $\Lambda^{(i)'}_7=T^{(i+1)}$.
\end{proof}

\begin{lemma}[number of families]\label{lem:a06f:RN}
For an admissible $\Lambda^{(k)}_0$ the number of families $F_k$ admissible for it is $\le2^{n_d|X_k|}$.
\end{lemma}

\begin{proof}
By \cref{lem:a06f:RE}(a) every element has all its points in $X_k$; charge each to its charge point; a point is the charge point of at most one $P_v$-, one $P_s$-, one $R_v$-, one $R_s$-element and of at most $\max\{d,d(d-1)/2\}$ bonds or plaquettes of each $Q$-type; a family is a subset of these $\le n_d|X_k|$ candidates.
\end{proof}

\begin{lemma}[the box partition]\label{lem:a06f:RP}
Let $q_k\defeq\lceil(r_k+s_b)/s_b\rceil$. For each $k\le K-1$ there is a partition $\mathcal Q_k$ of $T^{(k)}$ into boxes and a symmetric relation $Q\sim Q'$ (equal or adjacent) such that: (Q1) every box is a union of large blocks; (Q2) every box is the image of a product of $d$ integer intervals with at most $2r_k+3s_b\le(2+3\theta)r_k$ points each; (Q3a) every box has at most $3^d$ boxes $Q'\sim Q$; (Q3b) if $\dist(z,x)\le r_k+s_b$ then $Q(z)\sim Q(x)$.
\end{lemma}

\begin{proof}
In each direction $\mu$ let $m_\mu\defeq P_{k,\mu}/s_b$. If $m_\mu\ge q_k$, cut $\Z/P_{k,\mu}\Z$ into consecutive arcs, each a union of between $q_k$ and $2q_k-1$ consecutive large-block intervals $[js_b,(j+1)s_b)$ (write $m_\mu=nq_k+\rho'$, $0\le\rho'<q_k$, and give the $\rho'$ extra blocks to the $n$ arcs, at most $q_k-1$ each); if $m_\mu<q_k$ use one arc. The product of the lifted arc partitions is $\Gamma_k$-invariant, so its images form a partition of $T^{(k)}$; $Q\sim Q'$ if some lifts have arc indices differing by at most one in every coordinate. (Q1): arcs are unions of large-block intervals. (Q2): an arc has at most $(2q_k-1)s_b<2(r_k+s_b)+s_b$ points, or $P_{k,\mu}<q_ks_b\le r_k+2s_b$ in the one-arc case; $s_b\le\theta r_k$. (Q3a): the boxes $\sim Q$ are the images of the $3^d$ neighbours of one lift. (Q3b): take lifts with $|\tilde z-\tilde x|_\infty=\dist(z,x)\le r_k+s_b$; in a direction with $m_\mu\ge q_k$ every arc has $\ge q_ks_b\ge r_k+s_b$ points, so the two arcs are equal or consecutive; in a one-arc direction the indices agree.
\end{proof}

\begin{lemma}[the cube count on the interior]\label{lem:a06f:RC}
Put $\mathcal C_k\defeq\{Q\in\mathcal Q_k:Q\cap X_k\cap I_k\ne\emptyset\}$. For every family $F_k$ admissible for $\Lambda^{(k)}_0$, every $Q\in\mathcal C_k$ satisfies $Q\sim Q(x_e)$ for some element $e$ of $F_k$; hence $|\mathcal C_k|\le3^d|F_k|$, $|F_k|\defeq|P^{(k)}_v|+\dots+|R^{(k)}_s|$. At $k=0$ this is \cite[(3.44)]{Ba2}.
\end{lemma}

\begin{proof}
Let $z\in Q\cap X_k\cap I_k$; its large block $\beta$ lies in $Q$ (Q1); by \cref{lem:a06f:RE}(c) there are $z'\in\beta$, $y\in S_k$ with $\dist(z',y)\le r_k$, $y\in\mathrm{pts}(e)$ with $\dist(y,x_e)\le L-1<s_b$; so $\dist(z',x_e)\le r_k+s_b$ and $Q=Q(z')\sim Q(x_e)$ by (Q3b); count with (Q3a).
\end{proof}

\begin{lemma}[the bound on $\zeta''$]\label{lem:a06f:RZ}
$\zeta''_{\Lambda^{(k)}_0}\le2^{n_d|X_k|}\exp(-c_{\mathrm{box}}p_k^2|\mathcal C_k|)$.
\end{lemma}

\begin{proof}
The six coefficients $a/8,\gamma_0/4,\gamma_0/4,a/8,\gamma_0/4,1$ of \cite[(3.41)]{Ba2} are each $\ge3^dc_{\mathrm{box}}$, so each term of \cite[(3.41)]{Ba2} is $\le\exp(-3^dc_{\mathrm{box}}p_k^2|F_k|)\le\exp(-c_{\mathrm{box}}p_k^2|\mathcal C_k|)$ by \cref{lem:a06f:RC} ($\mathcal C_k$ depends on the region sequence only); the number of terms is \cref{lem:a06f:RN}.
\end{proof}

\begin{lemma}[weighted region entropy]\label{lem:a06f:RW}
For $\theta'\ge0$, $\sum_Xe^{-\theta'|X|}\le\exp(N_ke^{-\theta'})$, the sum over all unions $X$ of large blocks of $T^{(k)}$. With $\theta_k\defeq(d+1)\hat l_k$: $\sum_{k<K}N_ke^{-\theta_k}\le e^{-(d+1)}\eps_0^{d+1}|T^{(K)}|/(L-1)$.
\end{lemma}

\begin{proof}
The large blocks partition $T^{(k)}$ into $N_k/s_b^d$ blocks, so the sum is $(1+e^{-\theta's_b^d})^{N_k/s_b^d}\le\exp(N_ks_b^{-d}e^{-\theta's_b^d})\le\exp(N_ke^{-\theta'})$. Next $e^{-\theta_k}=e^{-(d+1)}t_k^{d+1}$ and $N_kt_k^d=|T^{(K)}|t_K^d$, so $\sum_{k<K}N_ke^{-\theta_k}=e^{-(d+1)}|T^{(K)}|t_K^d\sum_{k<K}t_k$ and $\sum_{k<K}t_k=t_K/(L-1)\le\eps_0/(L-1)$. (A $k$-independent price would give $\sum_kN_ke^{-\theta'}\sim\eps^{-d}|T|$: the logarithm in \cite[(3.47)]{Ba2} is exactly the price $\hat l_k$.)
\end{proof}

\begin{proposition}[from the sum to the supremum]\label{prop:a06f:RS}
Let $\mathcal F(\Lambda)$ be the summand of \cite[(3.42)]{Ba2} for the admissible sequence $\Lambda=(\Lambda^{(0)}_0,\dots,\Lambda^{(K-1)}_0)$ and $\Xi(\Lambda)\defeq\sum_{k<K}[(n_d\ln2+\theta_k)|X_k|+c_7t_k^{\kappa_0}|\mathrm{Col}_k|+c_5|\mathrm{Sh}_k|]$. For every set $\mathcal A$ of admissible sequences, $\sum_{\Lambda\in\mathcal A}\mathcal F(\Lambda)\le\exp(C_{\mathrm{fin}}|T^{(K)}|)\sup_{\Lambda\in\mathcal A}\exp(-\sum_{k<K}c_{\mathrm{box}}p_k^2|\mathcal C_k|+\Xi(\Lambda))$.
\end{proposition}

\begin{proof}
The $k=K$ terms of the two exponents of \cite[(3.42)]{Ba2} are $\le(c_7t_K^{\kappa_0}+c_5)L^d|T^{(K)}|$. By \cref{lem:a06f:RZ}, $\mathcal F(\Lambda)\le e^{L^d(\dots)|T^{(K)}|}\Psi(\Lambda)e^{-\sum_{k<K}\theta_k|X_k|}$ with $\Psi(\Lambda)\defeq\exp(-\sum_{k<K}c_{\mathrm{box}}p_k^2|\mathcal C_k|+\Xi(\Lambda))$; bound $\Psi$ by its supremum; $\Lambda\mapsto(X_0,\dots,X_{K-1})$ is injective and every $X_k$ is a union of large blocks, so $\sum_{\Lambda\in\mathcal A}e^{-\sum\theta_k|X_k|}\le\prod_{k<K}\sum_Xe^{-\theta_k|X|}\le\exp(\sum_{k<K}N_ke^{-\theta_k})$ by \cref{lem:a06f:RW}; collect into $C_{\mathrm{fin}}$.
\end{proof}

\begin{lemma}[{covering families; \cite[(3.48)--(3.52)]{Ba2} repaired}]\label{lem:a06f:RD}
Set $\mathcal D_0\defeq\mathcal C_0$ and, for $1\le k\le K-1$, $\mathcal D_k\defeq\{P_k(Q):Q\in\mathcal D_{k-1}\}\cup\mathcal C_k$, $P_k(Q)$ the prime of the union of the $L$-blocks of $T^{(k-1)}$ meeting $N_{w_{k-1}}(Q)$, $w_{k-1}\defeq7(r_{k-1}+s_b)+Ls_b$ (constructed on lifts; a box). Then for $k\le K-1$: \textup{(i)} $X_k\subset\bigcup\mathcal D_k$; \textup{(ii)} $|\mathcal D_k|\le\sum_{j\le k}|\mathcal C_j|$; \textup{(iii)} every box of $\mathcal D_k$ has side $\le C^0_Dr_k$; \textup{(iv)} $|\Lambda^{(k)c}_i|\le(C_Dr_k)^d\sum_{j\le k}|\mathcal C_j|$ for $i\le7$, and $|\Lambda^{(k)c}_8|\le(C^{(8)}_Dr_k)^d\sum_{j\le k}|\mathcal C_j|$ with $C^{(8)}_D\defeq C^0_D+16(1+\theta)$.
\end{lemma}

\begin{proof}
(ii) is immediate. (i) Induction on $k$; at $k=0$, $I_0=T^{(0)}$. Let $z\in X_k$. If $z\in I_k$, $z\in Q(z)\in\mathcal C_k$. Otherwise $\dist(z,T^{(k)}\setminus\Lambda^{(k-1)'}_7)\le s_b-1$, so by \cref{lem:a06f:RE}(d) there is $z_1$ with $B(z_1)\subset\Lambda^{(k-1)c}_7$ and $\dist(z,z_1)\le s_b-1$; by \cref{lem:a06f:RE}(b) at scale $k-1$ and the induction hypothesis every point $u$ of $B(z_1)$ lies in $N_{7(r_{k-1}+s_b)}(Q)$ for some $Q\in\mathcal D_{k-1}$; lifting, a point of $B(z)$ is within $L(s_b-1)$ of $u$, so $B(z)$ meets $N_{w_{k-1}}(Q)$ and $z\in P_k(Q)$. (iii) A box born at scale $j$ has side $s_j\le(2+3\theta)r_j$ (Q2); the union of the $L$-blocks meeting an interval of $n$ points has at most $n+2(L-1)$ points, so $s_{i+1}\le(s_i+2w_i+2L-2)/L$, and $2w_i+2L-2\le14r_i+(16+2L)s_b\le\nu_Lr_i$; iterating from $j<k$, $s_k\le L^{-(k-j)}s_j+\nu_L\sum_{i=j}^{k-1}L^{-(k-i)}r_i\le\nu'_L\sum_{i=j}^{k-1}L^{-(k-i)}r_i$, and $r_i\le(1+(k-i)\ln L)^\varrho r_k$ gives $s_k\le\nu'_L(S_L-1)r_k$; at $k=j$, $2+3\theta\le C^0_D$. (iv) By \cref{lem:a06f:RE}(b) and (i), $\Lambda^{(k)c}_i\subset\bigcup_{Q\in\mathcal D_k}N_{i(r_k+s_b)}(Q)$, each of at most $(C^0_Dr_k+2i(r_k+s_b))^d$ points; use (ii), with $i\le7$, resp.\ $i=8$.
\end{proof}

\begin{theorem}[{\cite[(3.42)]{Ba2} at $L=5$}]\label{thm:a06f:R7}
Let $c_5$, $c_7$, $\kappa_0$ be as in \cref{lem:a06f:C143}, let $c_8\ge0$ be free of $k$, $\eps$, $K$, the torus and the history, and add to the left side of \cite[(3.42)]{Ba2} the factor $\exp(\sum_{k<K}c_8|\Lambda^{(k)c}_8|)$ (the counterterm moves, \cref{prop:a06f:D2C}). Put $C_1^{(8)}\defeq C_1+c_8$, $\Theta^{(8)}_L\defeq C^{(8)}_1(C^{(8)}_D)^dR^d/(c_{\mathrm{box}}b_0^2\ln L)$ and assume
\[
(\mathrm{E.ent})\qquad(1+\ln(L/\eps_0))^{2\bar p-2-\varrho d}\ \ge\ \Theta^{(8)}_L.
\]
\textup{(a)} The (augmented) left side of \cite[(3.42)]{Ba2} is $\le\exp(C_{\mathrm{fin}}|T^{(K)}|)$. \textup{(b)} If (E.ent) holds with $2\Theta^{(8)}_L$, the part of the sum over the admissible sequences with some $X_k\ne\emptyset$ is $\le\exp(-c_{\mathrm{box}}p(t_{K-1})^2/2)\exp(C_{\mathrm{fin}}|T^{(K)}|)$. \textup{(c)} Both bounds are uniform in $\eps$, $K$, the torus and the history; (E.ent) is a cap on $\eps_0$ fixed after $d,a,L,M,R,b_0$ and after $c_5,c_7,c_8$, in Ba{\l}aban's order of choices. With $c_8=0$ this is \cite[(3.42)]{Ba2} as printed.
\end{theorem}

\begin{proof}
(a) By \cref{prop:a06f:RS} (with $\Xi$ augmented by $\sum_{k<K}c_8|\Lambda^{(k)c}_8|$, which is carried through its proof unchanged) it suffices that $\Xi(\Lambda)\le\sum_{j<K}c_{\mathrm{box}}p_j^2|\mathcal C_j|$. By \cref{lem:a06f:RE}(e), $t_k^{\kappa_0}\le1\le\hat l_k$ and $\Lambda^{(k)c}_7\subset\Lambda^{(k)c}_8$, every coefficient in $\Xi$ is $\le C^{(8)}_1\hat l_k$, so by \cref{lem:a06f:RD}(iv) and an exchange of sums, $\Xi\le C^{(8)}_1\sum_{k<K}\hat l_k|\Lambda^{(k)c}_8|\le C^{(8)}_1(C^{(8)}_D)^dR^d\sum_{j<K}|\mathcal C_j|\sum_{k=j}^{K-1}\hat l_k^{1+\varrho d}$. $\hat l_k$ decreases in $k$ and $K-j\le\log_L(\eps_0/t_j)\le\hat l_j/\ln L$, so the inner sum is $\le\hat l_j^{2+\varrho d}/\ln L$ and $\Xi\le c_{\mathrm{box}}b_0^2\Theta^{(8)}_L\sum_{j<K}\hat l_j^{2+\varrho d}|\mathcal C_j|$. For $j\le K-1$, $t_j\le\eps_0/L$, so $\hat l_j\ge1+\ln(L/\eps_0)$ and (E.ent) gives $\Theta^{(8)}_L\le\hat l_j^{2\bar p-2-\varrho d}$; hence $\Xi\le\sum_jc_{\mathrm{box}}b_0^2\hat l_j^{2\bar p}|\mathcal C_j|=\sum_jc_{\mathrm{box}}p_j^2|\mathcal C_j|$. (b) With $2\Theta^{(8)}_L$ the same chain gives $\Xi\le\frac12\sum_jc_{\mathrm{box}}p_j^2|\mathcal C_j|$; for a non-trivial sequence let $j$ be the first index with $X_j\ne\emptyset$; by \cref{lem:a06f:RE}(f), $X_j\cap I_j=X_j$, so $|\mathcal C_j|\ge1$, and $p_j\ge p_{K-1}$. (c) No step used a bound on $L$, $\eps$, $K$, the torus or the history beyond the setting.
\end{proof}

\begin{proposition}[the counterterm moves in the ledger]\label{prop:a06f:D2C}
Add at each step $k\le K-1$ the bounds of \cref{lem:a06d:D22,lem:a06d:D23} to the exponent that \cite[(2.43), (2.115)]{Ba2} carries into \cite[(2.117), (3.42)]{Ba2}. Then the left side of \cite[(3.42)]{Ba2} is bounded by \cref{thm:a06f:R7}(a), and its non-trivial part by \textup{(b)}, with $c_8\defeq Q_1(C_{\mathrm{D22}}+C_{\mathrm{D23}})$, $C_{\mathrm{D22}}$, $C_{\mathrm{D23}}$ the constants of those lemmas and $Q_1$ absorbing the factor $q(t_k)$ of a $\dm_{(4)}$ vertex (\cref{lem:a06e:T1}), and the $t^\kappa$ parts in the bulk sum \cite[(2.118)]{Ba2}: $C_\kappa\sum_kt_k^\kappa|T^{(k)}|\le C_\kappa\eps_0^{\kappa-d}(1-L^{d-\kappa})^{-1}|T|$ for $\kappa>d$.
\end{proposition}

\begin{proof}
$\Lambda^{(k)}_7\setminus\Lambda^{(k)}_8\subset\Lambda^{(k)c}_8$ and $T^{(k)}_1\setminus\Lambda^{(k)}_8=\Lambda^{(k)c}_8$, and $t_k^{1/2-\beta}\le1$, $t_kq(t_k)\le Q_1$, so the moves are $\le c_8|\Lambda^{(k)c}_8|+C_\kappa t_k^\kappa|T^{(k)}|$: the first part is a term of the augmented $\Xi$, the second a bulk term; they enter \cite[(2.117)]{Ba2} with the same sign and in the same place as the collar terms, and \cref{lem:a06f:C141} carries them through \cite[Sect.~3A]{Ba2}.
\end{proof}

\subsubsection*{The $g_k$-localisation of $Q^{(k)}$ \cite[p.~579]{Ba2}}

\begin{proposition}[the terms of $Q^{(k)}$ removed by $g_k$]\label{prop:a06f:H}
Let $1\le k\le K-1$, $\Omega_c\defeq B^k(\Lambda^{(k)}_0)$, $\Phi\defeq\theta_{k+1}A'^{(k)}$, and $Q^{(k)}[Y]\defeq\sum_{j<k}\sum_{n=1}^{\bar n}\frac1{n!}\frac{d^n}{d\tau^n}F_j(\Omega_c;\tilde B+\tau Y)|_{\tau=0}$ \cite[(2.105)--(2.106)]{Ba2}, in the setting of \cref{thm:a06d:G}. Let $c_g\defeq\sup|\partial^\eta g_k|$ ($\supp g_k\subset\Lambda^{(k)}_7$, $g_k=1$ on $\Lambda^{(k)}_8$). Then $|Q^{(k)}[\Phi]-Q^{(k)}[g_k\Phi]|\le F_H(t_k)(2M+1)^3|\Lambda^{(k)c}_8|$ with $F_H(t)\defeq C_Htp(t)^2(1+c_g)^2$, bounded on $(0,1]$ and free of $\eps$, $k$, $\eps_0$, the torus and the history. The other two replacements of \cite[(2.107)]{Ba2} (the background $\tilde B\to B^{(k+1),\eta}$ and the domain $\Omega_c\to B^k(\Lambda^{(k)}_2)$) are $O(t^\kappa)|\Lambda^{(k)}_7|$ for every $\kappa$.
\end{proposition}

\begin{proof}
Fix $j$ and $y$; put $\Phi_1\defeq g_k\Phi$, $\Phi_2\defeq(1-g_k)\Phi$, $\supp\Phi_2\subset\supp\theta_{k+1}\setminus B^k(\Lambda^{(k)}_8)$. For each $n$, $D^nQ_y(\tilde B)[\Phi^n]-D^nQ_y(\tilde B)[\Phi_1^n]$ is a sum of $n$ terms with at least one slot $\Phi_2$ (multilinearity). Gauge-fix at $y$ as in \cref{lem:a06d:G2} ($Q_y$ is gauge invariant, \cref{lem:a06d:G1}(b), and the axial map is linear). For $n=1$ use \cref{lem:a06d:G1}(f), which gives $D^2Q_y(s\beta)[\beta,\alpha_{\Phi_2}]$, bilinear in the curvature amplitudes of $\tilde B$ and $\Phi_2$; for $n\ge2$ use \cref{lem:a06d:G1}(e): each term is bounded by products of $n$ local amplitudes, one of them that of $\Phi_2$, plus terms suppressed by $e^{-\theta_2R/(4n)}$ with $R\defeq R_{kj}$ as in \cref{thm:a06d:G}. Curvatures in scale-$k$ units: $|\partial\Phi|\le Cp$ (the derivative clause of \cite[(2.91)]{Ba2}, Step~3(C$_1$): $\partial\Phi=(\partial\theta_{k+1})A'^{(k)}+\theta_{k+1}\partial A'^{(k)}$ with $|A'^{(k)}|\le C_{2.91}p$ by \cref{lem:a06d:C62}, $|\partial\theta_{k+1}|\le C$, and $\partial A'^{(k)}=\partial A^{(k)}-\partial B^{(k+1)}-\partial\tilde B^{(k)}$ by \cite[(2.89)]{Ba2}, bounded by (R1), \cite[(2.98)]{Ba2} and (R2) of Step~3(C$_2$)), $|\partial\Phi_2|\le|\partial\Phi|+2c_g|\Phi|\le C(1+c_g)p$, and on $N_k$, $\tilde B=A^{(k)}-\Phi-\tilde B^{(k)}$ (\cref{lem:a06d:G3}(a)) so $|\partial\tilde B|\le Cp$ by (R1), (R2); at scale $j$ all three amplitudes are $\le C(1+c_g)L^{-2(k-j)}t_k^{1/2}p$. Per $y$ the main terms are $\le C(1+c_g)^2L^{-4(k-j)}t_kp^2$, and the remaining terms are as in the proof of \cref{thm:a06d:G}; the $y$ within $R$ of $\supp\Phi_2$ number $\le(2R+3)^3(2M+1)^3L^{3(k-j)}|\Lambda^{(k)c}_8|$; sum over $y$ and over $j<k$ as there. The other two replacements act at distance $\ge5L^{k-j}r_k$ (scale-$j$ units) from $\supp g_k\Phi$, and \cref{lem:a06d:G1}(e) with one slot carrying the background or domain change gives $\le Ce^{-\theta_2L^{k-j}r_k/(2\bar n)}$ per $y$ (for the domain change, compare $F_j$ on the two domains through the common sub-box; the difference is localised at the domain boundary by \cite[(2.29)]{Ba1} and \cref{lem:a06d:G1}(d)).
\end{proof}

% a06_BALtorus_ct8: Step 3(F), fourth part: the field-dependent part of Ba2 Prop. 2.1, (2.57) (the collar quartic
% and the remainder O((L^k eps)^{kappa_0})|Col|); the g_k-gated terms of P^(k) of Ba2 p. 579 and the final step (2.116)
% are in the part files sec/a06_pgate_a.tex, _b.tex, _c.tex, input at the end of this file.

\subsubsection*{The collar: the field-dependent part of \cite[Prop.~2.1, (2.57)]{Ba2}}

\emph{Setting.} Step $k$ in scale-$k$ units ($t\defeq t_k$, $\eta\defeq L^{-k}$, $e_k$, $\lambda_k$, $p=p(t)$, $r=r(t)$); $\Omega_A\defeq B^k(\Lambda^{(k-1)'}_2)$, $\Omega_1\defeq B^k(\Lambda^{(k-1)'}_7)$, $\Omega'_1\defeq B^k(\Lambda^{(k)}_7)$, $\mathrm{Col}\defeq\Omega_1\setminus\Omega'_1$, $|\mathrm{Col}|=|\Lambda^{(k-1)'}_7\cap\Lambda^{(k)c}_7|$ unit cubes; $\tilde B\defeq A^{(k)}$, which equals $\theta_kA^{(k)}$ on $\Omega_A$; $\Phi\defeq\phi^{(k)}=a_kG_k(\Omega_A,\tilde B)Q_k^*(\tilde B)\Lambda^{(k-1)'}_6\varphi$ \cite[(2.56)]{Ba2}; $c_k(\Omega')$ the field-independent part of $P^{(k)}(\Omega',\tilde B,\varphi)$. Under \cite[(2.55)]{Ba2}, \cite[Prop.~2.1, (2.57)]{Ba2} asserts $P^{(k)}(\Lambda^{(k-1)'}_7,\theta_kA^{(k)},\varphi)=-\lambda_k\sum_{x\in B^k(\mathrm{Col})}\eta^d|\Phi(x)|^4+P^{(k)}(\Lambda^{(k)}_7,\theta_kA^{(k)},\varphi)+O((L^k\eps)^{\kappa_0})|\mathrm{Col}|$, ``a corollary of the analysis of the perturbation expansions'' (the sketch is \cite[p.~421]{Ba3}). Inputs: Proposition~1 with the vertex sizes of \cref{thm:a06e:T0}; \cite[(2.24)--(2.25)]{Ba1} at points of $\Omega_1$ (depth $\ge r-2\ge R_0$); \cite[(2.55), (2.60)]{Ba2} and the H\"older clause of \cite[p.~412]{Ba3} for $\tilde B$; \cite[Lemma~2.4, (2.65)--(2.66)]{Ba2} at every $y\in\Lambda^{(k-1)'}_7$ (\cref{thm:a06:BALA}(c)); the slice kernels of \cref{lem:a06e:SL} at $A=0$; the layer distances $\dist(\Lambda^{(k)}_{i+j},\Lambda^{(k)c}_i)\ge j(r-1)$, $\dist(\Omega_1,\Omega_A^c)\ge r-2$.

\begin{theorem}[the collar quartic and its remainder]\label{thm:a06f:C2}
$P^{(k)}(\Omega_1)-P^{(k)}(\Omega'_1)=-\lambda_k\sum_{x\in\mathrm{Col}}\eta^3|\Phi(x)|^4+R_{\mathrm{field}}+[c_k(\Omega_1)-c_k(\Omega'_1)]$ with $|R_{\mathrm{field}}|\le C_\beta t^{1/2-\beta}|\mathrm{Col}|$ for every $\beta>0$ and $|c_k(\Omega_1)-c_k(\Omega'_1)|\le Ct(1+\ln t^{-1})^5|\mathrm{Col}|\le C_\beta t^{1-\beta}|\mathrm{Col}|$, under the caps of Step~3(D$_2$) (Geometry). So \cite[(2.57)]{Ba2} holds with $\kappa_0=\frac12-\beta$, which implies the printed $\frac{4-d}4-\beta=\frac14-\beta$. The constants depend on $(d,L,M,\alpha_0,\bar n,a,c_1,b_0,\bar p,R,\varrho,e,\lambda,N,\beta)$, on the $O(1)$ of Proposition~1 and on $c_0,\delta_0,R_0$ of \cite[(2.24)]{Ba1}; they are free of $\eps$, $k$, $\eps_0$, the sets and $\varphi$.
\end{theorem}

The proof is \cref{lem:a06f:C21,lem:a06f:C22,lem:a06f:C23,lem:a06f:C24,lem:a06f:C25,lem:a06f:C26,lem:a06f:C27}.

\begin{lemma}[how $\Omega_1$ enters]\label{lem:a06f:C21}
In \cite[(1.4)--(1.5)]{Ba3}, $\Omega_1$ enters only through: \textup{(i)} $\mathbf 1_{\Omega_1}$ in the sums of the vertices \cite[(1.6)--(1.7)]{Ba3}; \textup{(ii)} $g_k$ ($\supp g_k\subset\Omega_1$, $g_k=1$ at distance $\ge M$ from $\partial\Omega_1$, derivatives up to order $2$ bounded by $1$) on every $A'$-leg, through $e'g_kA'+A^{(k)}$, the $A'$-legs of the pair vertices \cite[(1.14)]{Ba3} included; \textup{(iii)} $\dm_k(e'e_kg_k,\lambda'\lambda_k,\Omega_1,x)$ and $E_k(\dots,\Omega_1)$ of \cite[(1.29)--(1.30)]{Ba3}, sums of graphs whose internal vertices are localised by \textup{(i)}--\textup{(ii)} (\cite[pp.~419--420]{Ba3}: ``We localize the vertices (1.6) and (1.7) restricting summations to the set $\Omega_1$. The other vertices are localized by functions $g_k$''; ``We perform this localization for the terms of $\dm_k$, $E_k$ also''). The Gaussian measures, the propagators $G_k(\Omega_A,\tilde B)$ and the external fields $\Phi$, $\varphi$ are $\Omega_1$-free.
\end{lemma}

\begin{lemma}[marking]\label{lem:a06f:C22}
Put $D_{\mathrm{good}}\defeq\{x\in\Omega'_1:\dist(x,\partial\Omega'_1)\ge M+2\}$ and $\widetilde{\mathrm{Col}}\defeq\Omega_1\setminus D_{\mathrm{good}}$, and localise every graph in unit cubes \cite[p.~420]{Ba3}. \textup{(a)} A localised family all of whose cubes (counterterm-internal vertices included) lie in $D_{\mathrm{good}}$ has the same value for $\Omega_1$ and $\Omega'_1$. \textup{(b)} $P(\Omega_1)-P(\Omega'_1)$ is the sum over the localised families with a cube meeting $\widetilde{\mathrm{Col}}$ of $E(\cdot;\Omega_1)-E(\cdot;\Omega'_1)$. \textup{(c)} $|\widetilde{\mathrm{Col}}|\le(1+(2M+7)^d)|\mathrm{Col}|$.
\end{lemma}

\begin{proof}
$\dist(\Omega'_1,\Omega_1^c)\ge7(r-1)-1>M+3$, so $\partial\Omega'_1$ faces only $\mathrm{Col}$. On $D_{\mathrm{good}}$ (with the side-$2$ supports of the leg partitions) $\mathbf 1_{\Omega_1}=\mathbf 1_{\Omega'_1}=1$ and $g^{\Omega_1}=g^{\Omega'_1}=1$; \cref{lem:a06f:C21} for the counterterms. (c): every unit block of $\Omega'_1\setminus D_{\mathrm{good}}$ is within $M+3$ of a block of $\mathrm{Col}$, and each block of $\mathrm{Col}$ is so near to at most $(2M+7)^d$ of them.
\end{proof}

\begin{lemma}[external fields]\label{lem:a06f:C23}
For localisations $h$ at cubes in $\Omega_1$: $\norm{h\Phi}_{1,\alpha_0}\le C_\Phi p\lambda_k^{-1/4}$, and the block fields \cite[(1.12)]{Ba3} obey $|\varphi(y)|\le C_\Phi p\lambda_k^{-1/4}$ for $y\in\Lambda^{(k-1)'}_7$, with $C_\Phi\defeq3C_hac_0c_1L^{1/4}(1+\ln L)^{\bar p}$ ($C_h$ the localisation constant, $3$ for the three parts of \cite[(1.32)]{Ba3}). There are no external vector fields ($n'=0$ \cite[p.~414]{Ba3}; every $A'$-leg is contracted).
\end{lemma}

\begin{proof}
$\Phi=a_kG_k(\Omega_A,\tilde B)f$ with $f\defeq Q_k^*(\tilde B)\Lambda'_6\varphi$, $\norm f_\infty\le\sup_{\Lambda'_6}|\varphi|\le c_1p(t/L)\lambda(t/L)^{-1/4}$ by \cite[(2.55)]{Ba2} ($\Lambda'_6\subset\Lambda'_{-1}$); $p(t/L)\le(1+\ln L)^{\bar p}p(t)$, $\lambda(t/L)=\lambda_k/L$; apply \cite[(2.24)--(2.25)]{Ba1} at points of $\Omega_1$, at depth $\ge r-2\ge R_0$ in $\Omega_A$, $\tilde B$ regular by \cite[(2.60)]{Ba2}.
\end{proof}

\begin{lemma}[census and parity]\label{lem:a06f:C24}
The expressions of $P^{(k)}$ are the vertices \cite[(1.6), (1.7), (1.8), (1.10)]{Ba3} and the pairs from \cite[(1.12)--(1.14)]{Ba3}, with $n'=0$ (\cite[p.~414]{Ba3}); a pair is the scalar product of two of \cite[(1.12)--(1.14)]{Ba3}, at least one of the form \cite[(1.14)]{Ba3}, with factor $-a_k$ or $-\frac12a_k$, and \cite[(1.12)]{Ba3} is an external field; the $R$-vertices do not occur (\cite[(1.5)]{Ba3} has order $\le\bar n$). Let $s(v)$ be the number of scalar legs of $v$ (\cite[(1.12)]{Ba3} counted) and $X(v)\defeq2d_v(v)+4d_s(v)-s(v)$: $X=0$ for \cite[(1.6)]{Ba3}; $X\ge2$ for \cite[(1.7)]{Ba3} (coefficient of order $\ge2$); $X=2n-2$ for \cite[(1.8)]{Ba3} and \cite[(1.10)]{Ba3} at order $n$; $X=2m-2$ for a pair, $m\ge1$ the total order of its \cite[(1.14)]{Ba3} members. For a connected graph $G$ with $l_s$ internal scalar lines and $n_{\mathrm{ext}}$ external scalar legs, $W(G)\defeq2d_v+4d_s-n_{\mathrm{ext}}=\sum_vX(v)+2l_s$ is even and $\ge0$, and $W=0$ exactly for the lone \cite[(1.6)]{Ba3} vertex with four external legs and for two order-$1$ objects (\cite[(1.8)]{Ba3} with $n=1$, or a pair \cite[(1.12)]{Ba3}$\cdot$\cite[(1.14)]{Ba3}$_1$, \cite[(1.13)]{Ba3}$\cdot$\cite[(1.14)]{Ba3}$_1$ with both scalar legs external) joined by one vector line.
\end{lemma}

\begin{proof}
$\sum_vs(v)=n_{\mathrm{ext}}+2l_s$. $W=0$ forces $l_s=0$ and $X(v)=0$ for every $v$; the $X=0$ objects carry $0$ or $1$ $A'$-leg, all $A'$-legs are contracted and $P^{(k)}$ is connected, and a connected graph with vector lines only and at most one $A'$-leg per vertex has at most two vertices.
\end{proof}

\begin{lemma}[the families with $W\ge2$]\label{lem:a06f:C25}
For a localised family with $W(G)\ge2$, at $\Omega'\in\{\Omega_1,\Omega'_1\}$: $|E(G_{\mathrm{ren}},\{\square\};\Omega')|\le O(1)e^{d_v}\lambda^{d_s-n_{\mathrm{ext}}/4}C_\Phi^{n_{\mathrm{ext}}}p^{n_{\mathrm{ext}}}t^{W/4}e^{-\delta_0d(\{\square\})}$; summed over the marked families of \cref{lem:a06f:C22}, this part of $R_{\mathrm{field}}$ is $\le C_\beta t^{1/2-\beta}|\mathrm{Col}|$.
\end{lemma}

\begin{proof}
Proposition~1 with the external scalar fields of \cref{lem:a06f:C23} and no external vector field; the background satisfies the hypotheses of \cite[p.~412]{Ba3} ($|\partial\tilde B|\le O(p)$ by \cite[(2.60)]{Ba2}, its H\"older clause), there is no $\tilde A$, and its $O(1)$ is independent of $\eps$, $k$, the domains and $\tilde B$ \cite[p.~421]{Ba3}. In $d=3$, $e_k^{d_v}\lambda_k^{d_s}(p\lambda_k^{-1/4})^{n_{\mathrm{ext}}}=e^{d_v}\lambda^{d_s-n_{\mathrm{ext}}/4}p^{n_{\mathrm{ext}}}t^{W/4}$, and $W\ge2$ gives $t^{1/2}$; $n_{\mathrm{ext}}\le C(\bar n)$, so $p^{n_{\mathrm{ext}}}\le C_\beta t^{-\beta}$; the \cite[(1.7)]{Ba3} vertices with the finite parts $\dm_{\mathrm{fin}}$, $\dm_{K,k}$ enter through \cref{thm:a06e:T0} and \cref{lem:a06e:T1}. Root the tree at a marked cube ($\le|\widetilde{\mathrm{Col}}|$ choices, including the localised internal vertices of $\dm_k$, $E_k$, \cref{lem:a06f:C21}(iii)) and sum the rest by \cref{lem:a06d:V11}; there are finitely many families of order $\le\bar n$.
\end{proof}

\begin{lemma}[the terms with $W=0$]\label{lem:a06f:C26}
\textup{(a)} The lone \cite[(1.6)]{Ba3} vertex contributes to $P(\Omega_1)-P(\Omega'_1)$ exactly $-\lambda_k\sum_{x\in\mathrm{Col}}\eta^3|\Phi(x)|^4$ (its sum is sharp over $\Omega_1$, \cite[p.~413]{Ba3}). \textup{(b)} The single-line terms contribute $\le C_{\mathrm{SL}}e^2\lambda^{-1/2}p^4t^{1/2}$ per marked cube, hence $\le Ct^{1/2-\beta}|\mathrm{Col}|$ in all.
\end{lemma}

\begin{proof}[Proof of \textup{(b)}]
The order-$1$ objects (both scalar legs external) are $V_8=-e_k\sum_b\eta^3[(D_{\tilde B}\Phi)(b)\cdot q\Phi(b_-)]g_kA'_b$ from \cite[(1.8)]{Ba3}; $P_{12}(y)=a_ke_k\varphi(y)\cdot\sum_{x\in B^k(y)}\eta^3A'(\Gamma^{(k)}_{y,x})qU(\tilde B(\Gamma^{(k)}_{y,x}))\Phi(x)$ from \cite[(1.12)]{Ba3}$\cdot$\cite[(1.14)]{Ba3}$_1$; and $P_{13}(y)$, which is $P_{12}$ with $\varphi(y)$ replaced by $-(Q_k(\tilde B)\Phi)(y)$. The weight $g_k$ forces $y\in\Lambda^{(k-1)'}_7$. By \cite[Lemma~2.4, (2.65)]{Ba2} at every $y\in\Lambda^{(k-1)'}_7$, $\Phi(x)=U(\tilde B(\Gamma^{(k)}_{x,y}))\varphi(y)+\mathrm e(x)$ with $|\mathrm e(x)|\le c_{65}p$; $U(\tilde B(\Gamma_{y,x}))U(\tilde B(\Gamma_{x,y}))=1$ and $\varphi(y)\cdot q\varphi(y)=0$ ($q$ real antisymmetric), so the leading part of $P_{12}$ vanishes and its coefficient is $\le ae_kC_\Phi p\lambda_k^{-1/4}c_{65}p$; the same for $P_{13}$, since $(Q_k(\tilde B)\Phi)(y)=\varphi(y)+O(p)$; for $V_8$, \cite[(2.66)]{Ba2} gives $|D_{\tilde B}\Phi|\le c_{66}p$ on $B^k(\Lambda^{(k-1)'}_7)$, so its coefficient is $\le e_kc_{66}pC_\Phi p\lambda_k^{-1/4}$. The vector line: $A'=\sum_jA'^{(j),\eta}$ with independent Gaussian $A'_j$ of covariance $C^{(j)}$ \cite[p.~414]{Ba3}, and $A'(\Gamma^{(k)}_{y,x})=\sum_jA'^{(j),\eta}(\Gamma^{(j+1)}_{x_{j+1},x})$ \cite[(1.3), p.~412]{Ba3}; its covariance is slice-diagonal, $G^{v}_1+\sum_{j\ge1}S^v_j$ in scale-$k$ units, with $|S^v_j(b,b')|\le Cl_j^{2-d}e^{-\delta_K|b-b'|/l_j}$, $l_j\defeq L^{j-k}$ (\cref{lem:a06e:SL}(b),(d) rescaled; the $j=0$ term with its $L$-dependent constant), so $\sum_{b'}\eta^3|S^v_j(b,b')|\le Cl_j^2$, and the slice-$j$ contour $\Gamma^{(j+1)}$ has $\eta$-length $\le1.25dLl_j$. The three pairings per marked cube: $(V_8,V_8)$: $\le C(e_kp^2\lambda_k^{-1/4})^2\sum_jl_j^2\le Ce^2\lambda^{-1/2}p^4t^{1/2}$; $(V_8,P)$: $\sum_j(1.25dLl_j)Cl_j^2\le C$, the same order; $(P,P')$: per pair of blocks $\sum_j(1.25dLl_j)^2Cl_j^{2-d}e^{-\delta_K\dist/l_j}=\sum_jC'l_je^{-\dots}$ ($d=3$), and the sum over the second block with the decay is $\le C$ (the slice-$j$ field is integrated only over a contour of length $O(l_j)$: no logarithm; the $P$-side bonds are summed through the $\eta^3$-weighted $x\in B^k(y')$, $x\mapsto b'$ being $O(1)$-to-one). Sum the marked cubes over $\widetilde{\mathrm{Col}}$.
\end{proof}

\begin{lemma}[the $c$-number clause]\label{lem:a06f:C27}
\textup{(i)} At zero background on the torus, $c_k(\Omega_1)=-(\Efin+E_{K,k})+a_k(\Omega_1)$ (\cref{lem:a06e:VK0}), so the collar difference within the vertex set of $E_1$ is exactly $0$ at every order $\le\bar n$. \textup{(ii)} At $(\Omega_A,\tilde B)$: $|c_k(\Omega_1)-c_k(\Omega'_1)|\le Ct(1+\ln t^{-1})^5|\mathrm{Col}|\le C_\beta t^{1-\beta}|\mathrm{Col}|$.
\end{lemma}

\begin{proof}
(ii) By \cref{lem:a06f:C22} the difference is a sum over localised vacuum families with a cube meeting $\widetilde{\mathrm{Col}}$ (the vacuum graphs with their $E$-partners, \cite[pp.~420, 422]{Ba3}, and the pair-vertex vacuum graphs). Proposition~1 with no external fields gives $O(1)e_k^{d_v}\lambda_k^{d_s}e^{-\delta_0d}$; a vacuum graph has $d_v\ge2$ or $d_s\ge1$ or contains a \cite[(1.7)]{Ba3} vertex of order $\ge2$ (every expression of $P^{(k)}$ has order $\ge1$ in $e$ unless it is \cite[(1.6)]{Ba3} or \cite[(1.7)]{Ba3}; each $A'$-leg carries one power of $e$; $A'$-legs are contracted in pairs), so $W=2d_v+4d_s\ge4$ and $t^{W/4}\le t$, times $q(t)=(1+\ln t^{-1})^5$ per $\dm_{(4)}$ vertex for the finite counterterm parts (\cref{thm:a06e:T0}, \cref{lem:a06e:T1}); no external field, so no power of $p$. Sum by \cref{lem:a06d:V11} over roots in $\widetilde{\mathrm{Col}}$. At zero background this is \cref{cor:a06d:V13}.
\end{proof}

\begin{proof}[Proof of \cref{thm:a06f:C2}]
\Cref{lem:a06f:C22}(b) splits the difference into the marked families; those with $W\ge2$ are \cref{lem:a06f:C25}, those with $W=0$ are \cref{lem:a06f:C26}, and the field-independent ones are \cref{lem:a06f:C27}. $\eps_0$ enters only through the caps of Step~3(D$_2$) and the ``$e(t)$ sufficiently small'' of \cite[(2.24)]{Ba1} and of Proposition~1; the torus uniformity of Proposition~1's $O(1)$ is the torus side of the propagator estimates of group (A).
\end{proof}

\begin{remark}\label{rem:a06f:HP}
The second sentence of \cite[p.~579]{Ba2} bounds the terms localised in $B^k(\Lambda^{(k)c}_7)$ by $O((L^k\eps)^{\kappa_0})|\Lambda^{(k)}_2\cap\Lambda^{(k)c}_7|$: for $Q^{(k)}$ this is \cref{prop:a06f:H}; for the terms of $P^{(k)}(\Lambda^{(k)}_7,\theta_kA^{(k)},\varphi)$ whose vector leg is gated by $g_k$ there it is \cref{prop:a06g:HP} below, with $\kappa_0=\frac12-\beta$ and the part of the gate inside $\Lambda^{(k)}_7\setminus\Lambda^{(k)}_8$ booked as in \cref{cor:a06g:ledger}; the field-dependent part of the final-step bound \cite[(2.116)]{Ba2} is \cref{prop:a06g:FS}. Both run the mechanism of \cref{thm:a06f:C2}: the gated leg, resp.\ the whole final domain, takes the place of the collar vertex.
\end{remark}

% a06_pgate_a: Step 3(F), fifth part (input from sec/a06_BALtorus_ct8.tex): the new fluctuation field a of Ba2 (2.88)-(2.89)
% at step k (value, derivative and Hoelder bounds) and the geometry of the g_k-gate. Continued in a06_pgate_b (the gated terms
% of P^(k)) and a06_pgate_c (the final step, Ba2 (2.116)).

\subsubsection*{The $g_k$-gated terms of $P^{(k)}$ and the final step \cite[p.~579, (2.116)]{Ba2}}

\emph{Setting.} Step $k$, $1\le k\le K-1$, in scale-$k$ units, with the Setting of \cref{thm:a06f:C2} ($\Omega_A$, $\Omega_1$, $\Omega'_1=B^k(\Lambda^{(k)}_7)$, $\mathrm{Col}=\mathrm{Col}_k$, $\Phi=\phi^{(k)}$, $p=p(t_k)$, $r=r(t_k)$). At step $k+1$ \cite{Ba2} writes $\theta_{k+1}A^{(k)}=\theta_{k+1}A'^{(k)}+\tilde B^{(k)}+\theta_{k+1}B^{(k+1),\eta}$ \cite[(2.88)--(2.89)]{Ba2}, with $A'^{(k)}=a_k\zeta^{(k)}G_kQ_k^*A'$ (\cref{lem:a06d:C62}), $\theta_{k+1}$ the smooth cutoff of $B^k(\Lambda^{(k)}_2)$ changing in a slice of thickness $<M$ \cite[p.~567]{Ba2}, and every term of $\tilde B^{(k)}$ carrying the factor $\theta_{k+1}$. Put
\[
\tilde B\defeq(1-\theta_{k+1})\theta_kA^{(k)}+\theta_{k+1}B^{(k+1),\eta},\qquad a\defeq\theta_{k+1}A'^{(k)}+\tilde B^{(k)},\qquad\text{so that}\qquad\theta_kA^{(k)}=\tilde B+a,
\]
since $\theta_k=1$ on $\supp\theta_{k+1}\subset B^k(\Lambda^{(k-1)'}_7)\subset B^{k-1}(\Lambda^{(k-1)}_2)$ (\cref{lem:a06d:G3}); $\tilde B$ is the background of \cite[p.~582]{Ba2}, regular by \cite[(2.94)--(2.98)]{Ba2}, and $\supp a\subset\supp\theta_{k+1}$, which lies within $M$ of $B^k(\Lambda^{(k)}_2)$, hence in $B^k(\Lambda^{(k)}_1)$ ($\dist(\Lambda^{(k)}_2,\Lambda^{(k)c}_1)\ge r-1>M+1$ under the caps of Step~3(D$_2$)). The gate is $g_k$ ($\supp g_k\subset\Lambda^{(k)}_7$, $g_k=1$ on $\Lambda^{(k)}_8$, $|\partial^\eta g_k|\le c_g$, as in \cref{prop:a06f:H}). The cutoffs are taken, as the printed ``smooth'' allows, from fixed classes: $\theta_{k+1}$ a fixed smooth function of the mollified distance to $B^k(\Lambda^{(k)}_2)$, so that $|\partial^\eta\theta_{k+1}|\le c_\theta$ and $|\partial^\eta\partial^\eta\theta_{k+1}|\le c_{\theta2}$ uniformly in $k$ and $\eta$ (a smooth function of the lattice sup-distance would have second differences of order $\eta^{-1}$ along the edges of the slice); and $\zeta^{(k)}$ with $|\zeta^{(k)}|\le1$, $|\partial^\eta_x\zeta^{(k)}|\le1$, $|\partial^\eta_x\partial^\eta_x\zeta^{(k)}|\le2$ (the cut-off of \cref{thm:a06:BALA}(b) is so chosen: a fixed smooth profile of the mollified distance to $y$, as for $\theta_{k+1}$, over its transition region of width $\frac15r-2M$, which is $\ge2$ for $r\ge10M+10$, a cap on $\eps_0$ free of $k$). Put $\delta'\defeq\min(\delta_*,\delta_0)$, $S'_*\defeq\sum_{n\in\Z^3}e^{-\delta'|n|/2}$, and $C'(\delta')$ for the constant $C'$ of \cref{lem:a06d:C62} with $\delta'$ in place of $\delta_*$.

\begin{lemma}[the field $a$]\label{lem:a06g:A}
On $\supp\theta_{k+1}$, for $t_k\le t_0$, where $t_0\in(0,1]$ is fixed by $C_\kappa t_0^\kappa\le C_{2.91}b_0$, $C'_\kappa t_0^\kappa\le b_0$ and $C^H_\kappa t_0^\kappa\le b_0$ ($C_\kappa$, $C'_\kappa$ the constants of (R2) of Step~3(C$_2$), $C^H_\kappa$ that of (c) below; a cap on $\eps_0$, free of $k$; $p(t)\ge b_0$): \textup{(a)} $|a|\le C_{2.91}p+C_\kappa t_k^\kappa\le2C_{2.91}p$; \textup{(b)} for every fine bond $b=(x,\mu)$ with $x$ or $x+\eta e_\mu$ in $\supp\theta_{k+1}$, $|\partial^\eta_\mu a(x)|\le C'_{2.91}p+C'_\kappa t_k^\kappa\le(C'_{2.91}+1)p$, with
\[
C'_{2.91}\defeq c_\theta C_{2.91}+a(C_1e^{\delta_*}+c_0e^{\delta_0})S'_*(C_{2.86}+C'(\delta'));
\]
\textup{(c)} for parallel bonds $b,b'$ within distance $M$ of $\supp\theta_{k+1}$ with $|b-b'|\le1$, $|\partial^\eta a(b)-\partial^\eta a(b')|\le C^H_{2.91}p\,|b-b'|^{\alpha_0}$. The constants are free of $k$, $\eps$, $\eps_0$, the torus and the history. So $a$ satisfies the hypotheses on an external vector field of \cite[p.~412]{Ba3} with $\norm{\cdot}_{1,\alpha_0}$-norm $\le C_ap$, $C_a\defeq2C_{2.91}+C'_{2.91}+1+C^H_{2.91}$.
\end{lemma}

\begin{proof}
(a) is \cref{lem:a06d:C62} for $\theta_{k+1}A'^{(k)}$ and (R2) of Step~3(C$_2$) for $\tilde B^{(k)}$ \cite[(2.90)]{Ba2}; $C_\kappa t^\kappa\le C_{2.91}p$ for $t\le t_0$.

(b) \emph{Lattice distances.} Fine points are at mutual distances in $\eta\Z$, and the slice has thickness $<M$, so $\supp\theta_{k+1}$ lies within $M-\eta$ of $B^k(\Lambda^{(k)}_2)$ and both $x$ and $x+\eta e_\mu$ lie within $M$ of it; \cref{lem:a06d:C62} and its proof apply at both points (the proof bounds $A'^{(k)}$ itself there, without the factor $\theta_{k+1}$). \emph{Leibniz rule.} $\partial_\mu(\theta_{k+1}A'^{(k)})(x)=(\partial_\mu\theta_{k+1})(x)A'^{(k)}(x+\eta e_\mu)+\theta_{k+1}(x)(\partial_\mu A'^{(k)})(x)$. The first term is $\le c_\theta C_{2.91}p$. In the second,
\[
\partial_\mu A'^{(k)}(x)=a_k\sum_y\big[(\partial_{\mu,x}\zeta^{(k)})(x,y)(G_kQ_k^*)(x+\eta e_\mu,y)+\zeta^{(k)}(x,y)(\partial_\mu G_kQ_k^*)(x,y)\big]A'(y);
\]
both kernels are $\le(C_1e^{\delta_*}+c_0e^{\delta_0})e^{-\delta'|x-y|}$, by $|\partial\zeta|\le1$, \cref{thm:a06:A5} ($|\cdot|_1\ge|\cdot|$ and one bond shift) and the derivative clause of \cite[Prop.~2.2, (2.58)]{Ba2} (\cite[Thm., (1.9)--(1.10)]{Ba4}, at zero background on the torus, where the depth restriction is void, \cref{lem:a06A:torus}); both vanish unless $|x-y|<r-2M+\eta$, and such $y$ lie in $\Lambda^{(k)}_0$ by the layer count of \cref{lem:a06d:C62} ($2(r-c_b)>r-2M+M+1$). Split the $y$-sum as there: on $\Lambda^{(k)}_1$, \cref{lem:a06d:C61} gives $\le(\dots)S'_*C_{2.86}p$; on $\Lambda^{(k)}_0\setminus\Lambda^{(k)}_1$, $|x-y|\ge r-c_b-M-1$ and the factor $e^{-\delta'|x-y|}$ with $E(1,\delta'/2)$ turn $p/(\lmu t)$ into $C'(\delta')p$. With $|\theta_{k+1}|\le1$ and $a_k\le a$ the second term is $\le a(C_1e^{\delta_*}+c_0e^{\delta_0})S'_*(C_{2.86}+C'(\delta'))p$. For $\tilde B^{(k)}$ the derivative bound is (R2).

(c) Write $\partial(\theta_{k+1}A'^{(k)})=(\partial\theta_{k+1})\,\mathsf SA'^{(k)}+\theta_{k+1}\partial A'^{(k)}$ ($\mathsf S$ the bond shift) and telescope the difference at $b,b'$ term by term. The factors $\partial\theta_{k+1}$, $\theta_{k+1}$ and the kernel factor $\partial\zeta^{(k)}$ are Lipschitz on the $\eta$-lattice with constants $c_{\theta2}$, $c_\theta$ and $2$; at separation $|b-b'|\ge\eta$ a Lipschitz bound is $\le$ the constant times $|b-b'|^{\alpha_0}$, since $|b-b'|\le1$; their cofactors ($A'^{(k)}$, $\partial A'^{(k)}$ and the kernel sums) are bounded by \cref{lem:a06d:C62} and by (b), the kernel $(G_kQ_k^*)(\cdot+\eta e_\mu,y)$ of the $\partial\zeta^{(k)}$ term being itself Lipschitz in $x$ with constant $c_0e^{\delta_0}e^{-\delta_0\dist(\cdot,y)}$ by the derivative clause of \cite[(2.58)]{Ba2} (a difference along a lattice path of length $|b-b'|\le1$), summed over $y$ as in (b). The $\theta$-factors are telescoped first, so each cofactor is either evaluated at a bond of $\supp\theta_{k+1}$ or is a difference of two kernel sums of the type of (b) at points within $M+1$ of $B^k(\Lambda^{(k)}_2)$; the bounds of (b) hold there, since the kernel bounds are global on the torus and the layer count with $M+1$ in place of $M$ keeps the $y$-range in $\Lambda^{(k)}_0$. The term $\zeta^{(k)}(\partial G_kQ_k^*)A'$: its variation in $b$ is the H\"older clause of \cite[Prop.~2.1, (2.24)]{Ba1} (\cite[Thm., (1.10)]{Ba4}) for $\partial G_kQ_k^*\delta_y$, $|(\partial G_kQ_k^*)(b,y)-(\partial G_kQ_k^*)(b',y)|\le c_0|b-b'|^{\alpha_0}e^{-\delta_0\min(\dist(b,y),\dist(b',y))}$; sum over $y$ as in (b). For $\tilde B^{(k)}$: by \cite[(2.89)]{Ba2} each of its three terms is $\theta_{k+1}$ times a kernel sum of the same type ($\zeta^{(k)}G_kQ_k^*$, its tail $(1-\zeta^{(k)})G_kQ_k^*$, or $\zeta^{(k)}G_kQ_k^*$ composed with the covariance difference of \cref{lem:a06A:CLdiff}) applied to a block field bounded by $Cp/(\lmu t)$ on $\Lambda^{(k)'}_0$ (the restrictions \cite[(2.17), (2.55)]{Ba2}); the same computation bounds its H\"older quotient, and the separation $\ge r-2c_b-M$ between the evaluation points in $\supp\theta_{k+1}$ and the support of the field, through $E(1,\delta/2)$, gives $C^H_\kappa t_k^\kappa$ in place of $p$, as in (R2).
\end{proof}

\begin{lemma}[the gate]\label{lem:a06g:BAD}
Let $\mathrm{Bad}_k$ be the set of scale-$k$ unit cubes meeting $\supp((1-g_k)a)$. Then $\mathrm{Bad}_k\subset N_M(B^k(\Lambda^{(k)}_2))\setminus B^k(\Lambda^{(k)}_8)\subset\Omega_1$, every cube of $\mathrm{Bad}_k$ lies within $M$ of a cube of $\Lambda^{(k)}_2\cap\Lambda^{(k)c}_8$, and
\[
|\mathrm{Bad}_k|\le(2M+1)^3|\Lambda^{(k)}_2\cap\Lambda^{(k)c}_8|\le(2M+1)^3\big(|\Lambda^{(k)}_2\cap\Lambda^{(k)c}_7|+|\Lambda^{(k)}_7\setminus\Lambda^{(k)}_8|\big),\qquad\Lambda^{(k)}_2\cap\Lambda^{(k)c}_7\subset\mathrm{Col}_k.
\]
\end{lemma}

\begin{proof}
$(1-g_k)a$ vanishes outside $\supp\theta_{k+1}\subset N_M(B^k(\Lambda^{(k)}_2))$ and on $B^k(\Lambda^{(k)}_8)$, where $g_k=1$. Let $y\in\mathrm{Bad}_k$ and let $x$ be a cube of $B^k(\Lambda^{(k)}_2)$ within $M$ of $y$. If $y\in\Lambda^{(k)}_2$, then $y\in\Lambda^{(k)}_2\cap\Lambda^{(k)c}_8$ itself. If $y\notin\Lambda^{(k)}_2$, then $\dist(x,\Lambda^{(k)c}_2)\le M+1<6(r-1)\le\dist(\Lambda^{(k)}_8,\Lambda^{(k)c}_2)$, so $x\notin\Lambda^{(k)}_8$ and $x\in\Lambda^{(k)}_2\cap\Lambda^{(k)c}_8$. A cube of $\Lambda^{(k)}_2\cap\Lambda^{(k)c}_8$ has at most $(2M+1)^3$ cubes within $M$ of it. The inclusions: $\Lambda^{(k)}_8\subset\Lambda^{(k)}_7$, and $\Lambda^{(k)}_2\subset\Lambda^{(k)}_0\subset\Lambda^{(k-1)'}_7$ \cite[pp.~566, 570]{Ba2}, so $\Lambda^{(k)}_2\cap\Lambda^{(k)c}_7\subset\Lambda^{(k-1)'}_7\cap\Lambda^{(k)c}_7=\mathrm{Col}_k$ and $N_M(B^k(\Lambda^{(k)}_2))\subset B^k(\Lambda^{(k)}_1)\subset\Omega_1$.
\end{proof}

% a06_pgate_b: Step 3(F), fifth part, continued (input from sec/a06_BALtorus_ct8.tex): the terms of P^(k) whose vector leg is
% gated by g_k outside Lambda_8^(k) (the P^(k)-half of the second sentence of Ba2 p. 579), and their place in the ledger.

In the Setting above, expand the step-$k$ polynomial at the background $\theta_kA^{(k)}=\tilde B+a$ in powers of the new field $a$ \cite[(3.15), (3.44)]{Ba1}, \cite[p.~579]{Ba2}:
\[
P^{(k)}(\Omega'_1,\tilde B+a,\varphi)=\sum_{n_a=0}^{\bar n}P^{(k)}_{n_a}[a^{n_a}]+R^{(k)}_a,
\]
where $P^{(k)}_{n_a}[a^{n_a}]$ is the sum of the localised renormalised families of $P^{(k)}(\Omega'_1,\tilde B,\varphi)$ carrying $n_a$ external vector legs $a$: the vertices \cite[(1.8)--(1.11)]{Ba3} with $n'\ge1$, and the $a$-insertions on the scalar lines $G_k(\Omega_A,\tilde B+a)$ and on the external fields $\Phi(\tilde B+a)$ of \cite[(2.56)]{Ba2}; $R^{(k)}_a$ collects the $R$-vertices and the orders $>\bar n$ in $a$. On p.~579 the new legs are localised by $g_k$. Splitting every leg as $a=g_ka+(1-g_k)a$ (multilinearity), put
\[
T^{(k)}_{\mathrm{gate}}\defeq\sum_{n_a=1}^{\bar n}\big(P^{(k)}_{n_a}[a^{n_a}]-P^{(k)}_{n_a}[(g_ka)^{n_a}]\big),
\]
the sum of the terms with at least one leg carrying $(1-g_k)a$, i.e.\ with a vector leg localised where $g_k\ne1$: in $B^k(\Lambda^{(k)c}_7)$ or in $B^k(\Lambda^{(k)}_7\setminus\Lambda^{(k)}_8)$. If $P^{(k)}(\Lambda^{(k)}_7,\theta_kA^{(k)},\varphi)$ in \cite[(2.57)]{Ba2} is understood as the polynomial already built with the gated field $g_ka$, there is no such term; otherwise the gating is a move, whose cost is the following proposition.

\begin{proposition}[{the terms of $P^{(k)}$ removed by $g_k$; the $P^{(k)}$-half of the second sentence of \cite[p.~579]{Ba2}}]\label{prop:a06g:HP}
For $1\le k\le K-1$, under \cite[(2.55)]{Ba2} and the caps of Step~3(D$_2$), for every $\beta>0$,
\[
|T^{(k)}_{\mathrm{gate}}|\le C_{\mathrm{HP}}t_k^{1/2-\beta}|\mathrm{Bad}_k|\le C_{\mathrm{HP}}(2M+1)^3t_k^{1/2-\beta}\big(|\mathrm{Col}_k|+|\Lambda^{(k)}_7\setminus\Lambda^{(k)}_8|\big),
\]
with $C_{\mathrm{HP}}$ depending on the data of \cref{thm:a06f:C2} and on $c_\theta,c_{\theta2},c_g$, and free of $\eps$, $k$, $\eps_0$, the torus and the history. The remainder $R^{(k)}_a$ (the $R$-vertices and the orders $>\bar n$ in $a$) is $O(t_k^\kappa)|\Lambda^{(k)}_2|$ for every $\kappa$ (\cref{prop:a06e:B12}(iii)), the first half of the sentence.
\end{proposition}

\begin{proof}
\emph{Step 1: the weight.} Count every $a$-leg in $d_v$: the $a$-dependence of each parallel transporter is $U(\tilde B+a)=U(\tilde B)U(a)$ (the group is abelian), so each $a$-insertion carries its own factor $e_k$, and $W\defeq2d_v+4d_s-n_{\mathrm{ext}}$ is defined for a family with $a$-legs as in \cref{lem:a06f:C24}, the $a$-legs included in $d_v$. Inserting one $a$-leg into a family raises $W$ by at least $2$: on an existing gauge vertex, $d_v(v)$ rises by $1$ and $X(v)$ by $2$; as a new vertex of the type \cite[(1.8)]{Ba3} with $n=1$ (one vector leg, two scalar legs) splitting a scalar line, whether an internal line or the line $G_k(\Omega_A,\tilde B+a)$ of an external field $\Phi$, $X(\text{new})=2-2=0$ and $l_s$ rises by $1$, while $n_{\mathrm{ext}}$ is unchanged (the part of the split line beyond the new vertex is the external field $\Phi(\tilde B)$); as an insertion on the block average $Q_k^*(\tilde B+a)$ of an external field, $d_v$ rises by $1$ and nothing else changes: the external field becomes $a_kG_k(\Omega_A,\tilde B)f_a$ with the source $f_a\defeq[Q_k^*(\tilde B+a)-Q_k^*(\tilde B)]\Lambda^{(k-1)'}_6\varphi$, $\norm{f_a}_\infty\le e_k(L-1)d\sup|a|\sup|\varphi|$ (the transporter of $Q_k^*$ runs along a contour of $\eta$-length at most $(L-1)d$ in scale-$k$ units, \cite[p.~559]{Ba2}, and $q$ is antisymmetric), which the proof of \cref{lem:a06f:C23} bounds by $e_k(L-1)d\,2C_{2.91}p\cdot C_\Phi p\lambda_k^{-1/4}$, the factor $e_kC_ap$ of Step~3 up to $O(1)$. By \cref{lem:a06f:C24} every family of $P^{(k)}$ has $W\ge0$, so every term of $T^{(k)}_{\mathrm{gate}}$ has $W\ge2$ ($W$ stays even). In particular the $W=0$ objects of \cref{lem:a06f:C26} (the lone quartic and the single-line terms) occur in $T^{(k)}_{\mathrm{gate}}$ only with an $a$-leg attached, hence with $W\ge2$, and the cancellation of \cref{lem:a06f:C26}(b) is not needed.

\emph{Step 2: the fields.} The external vector field $a$ obeys \cref{lem:a06g:A}: $|a|\le2C_{2.91}p$ on $\supp\theta_{k+1}$ \cite[(2.90)--(2.91)]{Ba2}, $|\partial^\eta a|\le(C'_{2.91}+1)p$, H\"older quotient $\le C^H_{2.91}p$; so it is an admissible external vector field of \cite[p.~412]{Ba3} with norm $\le C_ap$. The external scalar fields are those of \cref{lem:a06f:C23}. The background $\tilde B$ is regular by \cite[(2.94)--(2.98)]{Ba2}, as in \cref{prop:a06e:B12}(iii).

\emph{Step 3: the bound per family.} Proposition~1 with the external vector field $a$ (the legs of \cite[(1.10)]{Ba3} with $n'\ge1$; its $O(1)$ is independent of $\eps$, $k$, the domains and the background, \cite[p.~421]{Ba3}) bounds each localised family by
\[
O(1)\,e^{d_v}\lambda^{d_s-n_{\mathrm{ext}}/4}(C_\Phi p)^{n_{\mathrm{ext}}}(C_ap)^{n_a}\,t^{W/4}e^{-\delta_0d(\{\square\})},
\]
by the computation of \cref{lem:a06f:C25} in $d=3$ ($e_k^{d_v}\lambda_k^{d_s}(p\lambda_k^{-1/4})^{n_{\mathrm{ext}}}=e^{d_v}\lambda^{d_s-n_{\mathrm{ext}}/4}p^{n_{\mathrm{ext}}}t^{W/4}$, the $e_k$ of the $a$-legs included in $e^{d_v}t^{d_v/2}$); the \cite[(1.7)]{Ba3} vertices with the finite parts $\dm_{\mathrm{fin}}$, $\dm_{K,k}$ enter through \cref{thm:a06e:T0} and \cref{lem:a06e:T1}. By Step~1, $t^{W/4}\le t^{1/2}$; $n_{\mathrm{ext}}+n_a\le C(\bar n)$, so $p^{n_{\mathrm{ext}}+n_a}\le C_\beta t^{-\beta}$ ($P(\beta,n)<\infty$); hence each family is $\le C_\beta t^{1/2-\beta}e^{-\delta_0d(\{\square\})}$.

\emph{Step 4: the sum.} Root the tree of each family at the cube of a leg carrying $(1-g_k)a$: at most $|\mathrm{Bad}_k|$ choices; sum the remaining cubes by \cref{lem:a06d:V11}; there are finitely many families of order $\le\bar n$ with $n_a\le\bar n$. This gives the first bound, and \cref{lem:a06g:BAD} the second.

\emph{Step 5: what is not gated.} The $R$-vertices and the orders $>\bar n$ in $a$ are not terms of $T^{(k)}_{\mathrm{gate}}$; they are bounded in \cref{prop:a06e:B12}(iii) by $O(t_k^\kappa)|\Lambda^{(k)}_2|$, Proposition~1 with the field restrictions and $\bar n+1$ legs $a$.
\end{proof}

\begin{corollary}[the two parts in the ledger]\label{cor:a06g:ledger}
The part of $T^{(k)}_{\mathrm{gate}}$ with a gated leg in $B^k(\Lambda^{(k)c}_7)$ is $\le C_{\mathrm{HP}}(2M+1)^3t_k^{1/2-\beta}|\mathrm{Col}_k|$: a collar term of the form of \cite[(2.57)]{Ba2}'s remainder, which enters \cite[(2.117)]{Ba2} with the collar terms (it is the $P^{(k)}$-part of the printed $O((L^k\eps)^{\kappa_0})|\Lambda^{(k)}_2\cap\Lambda^{(k)c}_7|$). The part with a gated leg in $B^k(\Lambda^{(k)}_7\setminus\Lambda^{(k)}_8)$ is $\le C_{\mathrm{HP}}(2M+1)^3|\Lambda^{(k)c}_8|$ ($t_k^{1/2-\beta}\le1$): a term of the form of \cref{prop:a06f:D2C}'s moves and of \cref{prop:a06f:H}'s bound, which enters the augmented exponent of \cref{thm:a06f:R7} with the same sign and in the same place. \Cref{thm:a06f:R7} holds for every $k$-free $c_8\ge0$; the three terms of this kind are carried together by taking
\[
c_8\ \ge\ Q_1(C_{\mathrm{D22}}+C_{\mathrm{D23}})+(2M+1)^3\big(\sup_{0<t\le1}F_H(t)+C_{\mathrm{HP}}\big),
\]
a constant free of $k$, $\eps$, $K$, the torus and the history, fixed before $\eps_0$, so that (E.ent) of \cref{thm:a06f:R7} is one cap on $\eps_0$ as before.
\end{corollary}

\begin{proof}
\Cref{lem:a06g:BAD} splits $\mathrm{Bad}_k$ into the cubes within $M$ of $\Lambda^{(k)}_2\cap\Lambda^{(k)c}_7\subset\mathrm{Col}_k$ and those within $M$ of $\Lambda^{(k)}_7\setminus\Lambda^{(k)}_8\subset\Lambda^{(k)c}_8$; the rooting of Step~4 is done in the respective set. The bookkeeping is that of the proof of \cref{prop:a06f:D2C}, and \cref{lem:a06f:C141} carries the terms through \cite[Sect.~3A]{Ba2}; $F_H$ is bounded on $(0,1]$ (\cref{prop:a06f:H}).
\end{proof}

% a06_pgate_c: Step 3(F), fifth part, end (input from sec/a06_BALtorus_ct8.tex): the final step, the field-dependent part of
% Ba2 (2.116), as Theorem thm:a06f:C2 at k = K with an empty inner domain.

\emph{The final step.} Let $K$ be the last step, $l_K=t_K=L^K\eps\in(\eps_0/L,\eps_0]$ (\cref{thm:a06:F}), in scale-$K$ units, with the Setting of \cref{thm:a06f:C2} at $k=K$: $\Omega_A=B^K(\Lambda^{(K-1)'}_2)$, $\Omega_1=B^K(\Lambda^{(K-1)'}_7)$, $\tilde B=A^{(K)}=\theta_KA^{(K)}$ on $\Omega_A$, $\Phi=\phi^{(K)}$ of \cite[(2.56)]{Ba2} at $k=K$, under \cite[(2.55)]{Ba2} at step $K$. No set $\Lambda^{(K)}_i$ is defined at the final step: \cite[(2.116)]{Ba2} bounds $P^{(K)}$ on its whole domain $\Omega_1$, and its right side is the volume $|\Lambda^{(K-1)'}_7|$ of that domain, as in the last term of \cite[(2.117)]{Ba2}. The chain is run with $E_0+E_K$ (\cref{lem:a06:CF}), so no step polynomial contains $\Efin$, which is the global factor bounded by \cref{thm:a06:F}.

\begin{proposition}[{the final step, \cite[(2.116)]{Ba2}}]\label{prop:a06g:FS}
For every $\beta>0$,
\[
P^{(K)}(\Omega_1,\tilde B,\varphi)=-\lambda_K\sum_{x\in\Omega_1}\eta^3|\Phi(x)|^4+R_{\mathrm{field}}+c_K(\Omega_1),\qquad|R_{\mathrm{field}}|\le C_\beta t_K^{1/2-\beta}|\Lambda^{(K-1)'}_7|,
\]
with $c_K(\Omega_1)$ the field-independent part, $|c_K(\Omega_1)|\le Ct_K(1+\ln t_K^{-1})^5|\Lambda^{(K-1)'}_7|\le C_\beta t_K^{1-\beta}|\Lambda^{(K-1)'}_7|$. Hence $P^{(K)}(\Omega_1,\tilde B,\varphi)\le C_\beta t_K^{1/2-\beta}|\Lambda^{(K-1)'}_7|$, which is \cite[(2.116)]{Ba2} with $\kappa_0=\frac12-\beta$. The constants are those of \cref{thm:a06f:C2}, free of $\eps$, $K$, $\eps_0$, the torus and the history.
\end{proposition}

\begin{proof}
Run the proof of \cref{thm:a06f:C2} at $k=K$ with the inner domain $\Omega'_1\defeq\emptyset$.

\emph{Step 1: the empty domain.} By \cref{lem:a06f:C21} the vertex domain enters only through $\mathbf 1_{\Omega'_1}$ on the vertices \cite[(1.6)--(1.7)]{Ba3}, through the gate on the $A'$-legs ($\supp g\subset\Omega'_1$) and through the localised counterterms. With $\Omega'_1=\emptyset$ every vertex vanishes and every $A'$-leg with it, so $P^{(K)}(\emptyset,\tilde B,\varphi)$ is field-independent and equals $c_K(\emptyset)$, which is $0$: on the chain run with $E_0+E_K$, $c_k(\Omega')=-E_{K,k}+a_k(\Omega')$ (\cref{lem:a06e:VK0}), $E_{K,K}=E_K-E_K=0$, and $a_K(\emptyset)$ has no pair vertex to place. In \cref{lem:a06f:C22}, $D_{\mathrm{good}}=\emptyset$ and $\widetilde{\mathrm{Col}}=\Omega_1$: every localised family is marked, and $P^{(K)}(\Omega_1)=P^{(K)}(\Omega_1)-P^{(K)}(\emptyset)$ is the sum over all localised families, rooted at a cube of $\Omega_1$ ($|\Omega_1|=|\Lambda^{(K-1)'}_7|$ choices).

\emph{Step 2: the three kinds of family.} The families with $W\ge2$ are bounded by \cref{lem:a06f:C25} at $k=K$: Proposition~1 on the lines $G_K(\Omega_A,\tilde B)$ with the external scalar fields of \cref{lem:a06f:C23} at $k=K$ (which uses \cite[(2.55)]{Ba2} at step $K$ and \cite[(2.24)--(2.25)]{Ba1} at points of $\Omega_1$, at depth $\ge r(t_K)-2\ge R_0$ in $\Omega_A$, $\tilde B$ regular by \cite[(2.60)]{Ba2}), the \cite[(1.7)]{Ba3} vertices through \cref{thm:a06e:T0} and \cref{lem:a06e:T1}, the sum by \cref{lem:a06d:V11} with the root in $\Omega_1$: in all $\le C_\beta t_K^{1/2-\beta}|\Lambda^{(K-1)'}_7|$. The families with $W=0$ are \cref{lem:a06f:C26}: the lone \cite[(1.6)]{Ba3} vertex, whose sum over the vertex domain is sharp, gives exactly $-\lambda_K\sum_{x\in\Omega_1}\eta^3|\Phi(x)|^4$; the single-line terms give $\le C_{\mathrm{SL}}e^2\lambda^{-1/2}p^4t_K^{1/2}$ per cube of $\Omega_1$ by \cite[Lemma~2.4, (2.65)--(2.66)]{Ba2} at every $y\in\Lambda^{(K-1)'}_7$ (\cref{thm:a06:BALA}(c)) and the slice kernels of \cref{lem:a06e:SL}, in all $\le Ct_K^{1/2-\beta}|\Lambda^{(K-1)'}_7|$. The field-independent families are \cref{lem:a06f:C27}(ii) with $\Omega'_1=\emptyset$: the vacuum families with a cube meeting $\Omega_1$ have $W\ge4$, and $|c_K(\Omega_1)|=|c_K(\Omega_1)-c_K(\emptyset)|\le Ct_K(1+\ln t_K^{-1})^5|\Lambda^{(K-1)'}_7|$; at zero background this is \cref{cor:a06d:V13}, and on the chain of \cref{lem:a06:CF} it is the bound of \cref{thm:a06e:VK}(a) on $a_K(\Omega_1)$, whose proof does not use $k<K$.

\emph{Step 3: the inputs at $k=K$.} Every input of \cref{thm:a06f:C2} is available at the final step: \cite[(2.55)]{Ba2} holds at every step, $K$ included; Proposition~1's $O(1)$ does not depend on the step; the layer distances used are those of the step-$(K-1)$ sets, $\dist(\Omega_1,\Omega_A^c)\ge r(t_K)-2$; and no set $\Lambda^{(K)}_i$ enters, since the inner domain is empty. The one-sided bound follows by dropping the quartic term, which is $\le0$.
\end{proof}

The same bound, with $|\Lambda^{(K-1)'}_7|$ replaced by $L^3|T^{(K)}|\ge|\Lambda^{(K-1)'}_7|$, is the $k=K$ term of the exponents of \cite[(3.42)]{Ba2} as \cref{prop:a06f:RS} carries it. This completes the field-dependent part of the second sentence of \cite[p.~579]{Ba2} (\cref{prop:a06g:HP}) and of the final step (\cref{prop:a06g:FS}); with \cref{thm:a06f:C2} (the collar), the three transfers of the mechanism of \cite[Prop.~2.1]{Ba2} that \cite[p.~570]{Ba2} announces (``similar conclusions hold for other expressions which will be included in the action in the later stages of the procedure'') are written.


% a06_BALtorus_A1: part of Step 3 of the proof of thm:a06:T1 (input from sec/a06_BALtorus_c.tex, after the part files d1-ct8, before Step 4).
% Group (A), part 1 of 4: setting, cube coercivity, coercivity with mixed pins, the Combes-Thomas bound.
% Continues in sec/a06_BALtorus_A2.tex (block inversion, resolvent differences, (1.18)), A3 (depth bounds, (5.5), (1.20), Ba1 Prop. 2.2, the summary theorem), A4 (the torus).

\subsubsection*{Step 3(A): the propagator estimates of group (A)}

This part proves the estimates of group (A) that Steps~3(B)--3(D) consume: the coercivity \cite[(1.8)]{Ba4} on a general region and for the mixed-pin operators \cite[(3.5)]{Ba4}, the Combes--Thomas form of \cite[Cor.~2.3]{Ba4} for every pin structure, the two-sided bounds \cite[(1.15)]{Ba4}, the decay \cite[(1.16)]{Ba4}, the conditioning-region difference \cite[(1.18)]{Ba4}, the domain differences \cite[(5.5), (1.20)]{Ba4}, \cite[Prop.~2.2, (2.27)--(2.29)]{Ba1}, the two-sided form of the depth factor of \cite[Thm., (1.12)]{Ba4}, and their validity on the torus with torus-free constants. The notation is that of Steps~3(B)--3(D): $\eta=L^{-k}$, $e_k=et_k^{1/2}$, $a_k\in[a_*,a]$ with $a_*=a(1-L^{-2})$, $Q_k(A)$, $P_k(A)=Q_k^*Q_k$, $G_k(\Omega,A)$, $\Delta^{(k)}(\Omega,A)$, $C^{(k)}_\Lambda=C_\Lambda$, and $M$ the large-block side. Distances are the sup-distances of \cite[(1.3)]{Ba1} (on a torus, the quotient distances); $|x-x'|$ is the sup-distance of unit-lattice points, and kernels of $N$-component operators are $N\times N$ matrices, $|K(x,x')|$ their operator norm.

\paragraph{Setting.} Functions on the fine lattice $\eta\Z^d$ carry the $\eta^d$-weighted $\ell^2$ product $\langle\phi,\psi\rangle=\sum_x\eta^d\phi(x)\cdot\psi(x)$, functions on the unit lattice the unweighted one; bond functions are weighted at the initial point $b_-$ of the bond. $U(A_b)\defeq\exp(qe_k\eta A_b)$ with $q$ a fixed real antisymmetric $N\times N$ matrix \cite[(1.2)]{Ba4}; all $U(A_b)$ commute and are orthogonal, and $U(A(\Gamma))\defeq\prod_{b\in\Gamma}U(A_b)$ for a contour $\Gamma$. For a union $\Omega$ of unit blocks, $(D_A\phi)(b)\defeq\eta^{-1}(U(A_b)\phi(b_+)-\phi(b_-))$ on the bonds $b\subset\Omega$, so that $\langle\phi,-\Delta^{\eta,N}_{A,\Omega}\phi\rangle=\norm{D_A\phi}^2$ is the Neumann covariant Dirichlet form \cite[(1.3)]{Ba4}. $Q_k(A)$ averages over the unit blocks $B^k(y)$, $y\in\Omega^{(k)}$, along the contours $\Gamma^{(k)}_{y,x}\subset B^k(y)$ of \cite[(2.1)--(2.2)]{Ba1}; $q_k(y)\defeq Q_k(A)^*\delta_y$ is supported in $B^k(y)$, $|q_k(y)(x)|=1$ there, $\norm{q_k(y)}=1$; $P_k(A)=Q_k^*Q_k$ is the orthogonal projection onto the transported constants $x\mapsto U(A(\Gamma^{(k)}_{y,x}))v$ of the unit blocks. $P(A)\defeq Q(A)^*Q(A)$, with $Q(A)$ the covariant average over the $L$-blocks of the unit lattice, is the orthogonal projection onto the transported constants of the $L$-blocks, with $|P(A)(y,y')|\le1$ and $P(A)(y,y')=0$ unless $y,y'$ lie in one $L$-block. The operators of this part are
\[
G_k(\Omega,\Lambda,A)\defeq T_\Lambda^{-1},\quad T_\Lambda\defeq-\Delta^{\eta,N}_{A,\Omega}+m_k^2+a_kP_k(A)(\Omega\setminus B^k\Lambda)+a_{k+1}L^{-2}P_{k+1}(A)(B^k\Lambda)
\]
\cite[(3.5)]{Ba4}, for $\Lambda\subset\Omega^{(k)}$ a union of $L$-blocks ($P_k(A)(Y)$ is the sum of the block projections of the unit blocks in $Y$, $P_{k+1}(A)(Y)$ that of the $L$-blocks), so that $T_\emptyset=G_k(\Omega,A)^{-1}$ and $T_{\Omega^{(k)}}$ is the operator of $G^\eta_{k+1}(\Omega,A)$; and
\[
X_\Lambda\defeq(\Delta^{(k)}(\Omega,A)+aL^{-2}P(A))|_\Lambda,\qquad C_\Lambda=X_\Lambda^{-1}
\]
\cite[(1.13)]{Ba4}, with $X|_\Lambda\defeq\Lambda X\Lambda$ on $\ell^2(\Lambda)$. Three facts about Ba{\l}aban's blocks and contours are used. (H1) $\Omega^{(k)}$ and $\Lambda$ are unions of large blocks, and every large block is a union of $L$-blocks (\cite[p.~607]{Ba1}: the large-block partition is the block partition with $L$ replaced by $ML$). (H2) The contours are composed, $\Gamma^{(k+1)}_{z,x}=\Gamma_{z,y}\cup\Gamma^{(k)}_{y,x}$ for $x\in B^k(y)$, $y\in B(z)$ \cite[(2.2)]{Ba1}; hence $Q_{k+1}(A)=Q(A)Q_k(A)$ and $P_{k+1}(A)=Q_k^*P(A)Q_k$ is the orthogonal projection onto the transported constants of the $L$-blocks $B^{k+1}(z)$. (H3) Each contour is a union of axis-ordered monotone lattice paths inside its block \cite[(2.1)]{Ba1}; the $\eta$-length of $\Gamma^{(k)}_{y,x}$ is at most $d(L-1)L^{-1}(1+L^{-1}+\dots)\le d$, that of $\Gamma^{(k+1)}_{z,x}$ at most $d(L-1)+d\le d(L+1)$. A vector field $A$ is \emph{regular} on $\Omega$ when $|\partial^\eta_\mu A_\nu(x)|\le ce_k^{\beta-1}$ for $x\in\Omega$, some $\beta>0$ and a universal $c$ \cite[(1.7)]{Ba4}; its curvature $F_{\nu\mu}\defeq\partial^\eta_\nu A_\mu-\partial^\eta_\mu A_\nu$ \cite[(1.6)]{Ba1} is gauge invariant and then obeys $|F_{\nu\mu}|\le2ce_k^{\beta-1}$. Every smallness condition on $e_k=et_k^{1/2}\le e\eps_0^{1/2}$ below is a cap on $\eps_0$ depending on $(d,N,\norm q,c,\beta,a,L)$ only. For $\theta>0$ put $S(\theta)\defeq\sum_{n\in\Z^d}e^{-\theta|n|}<\infty$.

\begin{lemma}[cube coercivity]\label{lem:a06A:cube}
Let $\Delta\subset\eta\Z^d$ be a cube of side $\ell\in\{1,L\}$ that is a union of unit blocks, $A$ a bond field on $\Delta$ with $|F_{\nu\mu}(A)|\le f$ on the plaquettes of $\Delta$, $\Pi^A$ the orthogonal projection onto $\{x\mapsto U(A(\Gamma_x))v:v\in\R^N\}$ for contours $\Gamma_x\subset\Delta$ from a root, each of $\eta$-length at most $\lambda_\Gamma$, and $c_p>0$. Put $\alpha\defeq(d-1)\ell f$ and assume $\norm qe_k\max\{\lambda_\Gamma,1\}\alpha\le\ln2$. Then, with $\kappa\defeq2\norm qe_k\alpha$ and $\eps_g\defeq2\norm qe_k\lambda_\Gamma\alpha$,
\[
-\Delta^{\eta,N}_{A,\Delta}+c_p\Pi^A\ \ge\ \min\{2\ell^{-2},c_p\}-2c_p\eps_g-d\kappa^2 .
\]
If $A$ is regular on $\Delta$, then $f\le2ce_k^{\beta-1}$ and $\kappa,\eps_g=O(e_k^\beta)$, so there is $e_1(d,N,\norm q,c,\beta,L,a)>0$ such that for $e_k\le e_1$ the right side is at least $\frac12\min\{2\ell^{-2},c_p\}$ for every $c_p\le a$ and every $\lambda_\Gamma\le d(L+1)$.
\end{lemma}

\begin{proof}
\emph{Gauge covariance.} For $\lambda\colon\Delta\to\R$ the gauge transform $A^\lambda_b\defeq A_b+\eta^{-1}(\lambda(b_+)-\lambda(b_-))$ gives $U(A^\lambda_b)=U(A_b)e^{qe_k(\lambda(b_+)-\lambda(b_-))}$. With the orthogonal multiplication operator $G\defeq e^{-qe_k\lambda}$ (acting on bond functions through $\lambda(b_-)$), $D_{A^\lambda}=GD_AG^{-1}$, since all the matrices commute; and $U(A^\lambda(\Gamma_x))=e^{qe_k\lambda(x)}U(A(\Gamma_x))e^{-qe_k\lambda(x_0)}$ for the root $x_0$, so the range of $\Pi^{A^\lambda}$ is $G^{-1}$ of the range of $\Pi^A$ and $\Pi^{A^\lambda}=G\Pi^AG^{-1}$. Hence $-\Delta^{\eta,N}_{A^\lambda,\Delta}+c_p\Pi^{A^\lambda}=G(-\Delta^{\eta,N}_{A,\Delta}+c_p\Pi^A)G^{-1}$, and the lower bound is gauge invariant; $F$ is unchanged.

\emph{Axial gauge.} Choose $\lambda$ so that $A'\defeq A^\lambda$ has $A'_1\equiv0$, $A'_2=0$ on the bonds of one face $\{x_1=x_1^0\}$, $A'_3=0$ on the bonds of one line $\{x_1=x_1^0,x_2=x_2^0\}$, and so on in $d$ dimensions (the bonds on which a component is prescribed form a tree, so $\lambda$ exists). Then $A'_\mu(x)$ is reached from a bond where it vanishes by at most $d-1$ axis-parallel segments of length at most $\ell$, along each of which it changes by $\eta$ times the curvature: on a segment in direction $\nu<\mu$, $\partial^\eta_\nu A'_\mu=F_{\nu\mu}+\partial^\eta_\mu A'_\nu$, and $A'_\nu$ vanishes identically on the set over which the segment runs. Hence $|A'_\mu|\le(d-1)\ell f=\alpha$ on $\Delta$.

\emph{Bond and pin defects.} $|U(A'_b)-1|\le\norm qe_k\eta|A'_b|e^{\norm qe_k\eta|A'_b|}\le2\norm qe_k\alpha\eta$, so $\kappa\ge\sup_b|U(A'_b)-1|/\eta$. The transported constants are $g(x)\defeq U(A'(\Gamma_x))=\exp(qe_k\sum_{b\in\Gamma_x}\eta A'_b)$, and $\sum_{b\in\Gamma_x}\eta|A'_b|\le\lambda_\Gamma\alpha$ gives $|g(x)-1|\le2\norm qe_k\lambda_\Gamma\alpha=\eps_g$.

\emph{Error operator.} Write $D_{A'}=D_1+E$ with $(D_1\phi)(b)=\eta^{-1}(\phi(b_+)-\phi(b_-))$ and $(E\phi)(b)=\eta^{-1}(U(A'_b)-1)\phi(b_+)$. Each site is the final point of at most $d$ bonds of $\Delta$, so $\norm{E\phi}^2\le\kappa^2\sum_b\eta^d|\phi(b_+)|^2\le d\kappa^2\norm\phi^2$, and $\norm{D_{A'}\phi}^2\ge\frac12\norm{D_1\phi}^2-d\kappa^2\norm\phi^2$.

\emph{Neumann gap.} $\Delta$ is a product of $d$ paths with $\ell/\eta$ sites and spacing $\eta$; the Neumann Laplacian of a path has eigenvalues $4\eta^{-2}\sin^2(\pi j\eta/(2\ell))$, $0\le j<\ell/\eta$, so $\norm{D_1\phi}^2\ge\lambda_\ell\norm{(1-\Pi_1)\phi}^2$ with $\Pi_1$ the projection onto the constants and $\lambda_\ell=4\eta^{-2}\sin^2(\pi\eta/(2\ell))\ge4\ell^{-2}$, by $\sin t\ge2t/\pi$ on $[0,\pi/2]$.

\emph{Pin comparison.} $\Pi^{A'}=V_gV_g^*$ with $(V_gv)(x)\defeq g(x)v/\ell^{d/2}$, an isometry $\R^N\to\ell^2(\Delta)$, and $\Pi_1=V_1V_1^*$; $\norm{V_g-V_1}\le\sup_x|g(x)-1|\le\eps_g$, so $\norm{\Pi^{A'}-\Pi_1}\le2\eps_g$.

\emph{Combination.} $-\Delta^{\eta,N}_{A',\Delta}+c_p\Pi^{A'}\ge\frac12\lambda_\ell(1-\Pi_1)+c_p\Pi_1-2c_p\eps_g-d\kappa^2\ge\min\{2\ell^{-2},c_p\}-2c_p\eps_g-d\kappa^2$. The last sentence of the statement follows since $\alpha\le2(d-1)Lce_k^{\beta-1}$ and $\kappa,\eps_g\le4d(L+1)\norm q(d-1)Lc\,e_k^\beta$.
\end{proof}

\begin{lemma}[{coercivity with mixed pins; \cite[(1.8)]{Ba4} on a general region}]\label{lem:a06A:coer}
Assume (H1)--(H3), $A$ regular on $\Omega$ and $e_k\le e_1$. Then for every $\Lambda\subset\Omega^{(k)}$ that is a union of $L$-blocks ($\Lambda=\emptyset$ and $\Lambda=\Omega^{(k)}$ included)
\[
T_\Lambda\ \ge\ \gamma_L\defeq\tfrac14\min\{2L^{-2},a_*L^{-2}\},
\]
and $T_\emptyset=G_k(\Omega,A)^{-1}\ge\gamma_\emptyset\defeq\frac14\min\{2,a_*\}$, which is \cite[(1.8)]{Ba4} on every union $\Omega$ of unit blocks, with $\gamma_0=\gamma_\emptyset$ independent of $L$, $\eta$, $k$, $\Omega$, $A$ and the torus.
\end{lemma}

\begin{proof}
Partition $\Omega$ into the unit blocks of $\Omega\setminus B^k\Lambda$ and the $L$-blocks of $B^k\Lambda$ (cells). Dropping the bonds between cells removes non-negative terms from the Dirichlet form, so $-\Delta^{\eta,N}_{A,\Omega}\ge\bigoplus_{\text{cells}}(-\Delta^{\eta,N}_{A,\text{cell}})$. The pins are block-diagonal on this partition: $P_k(A)(\Omega\setminus B^k\Lambda)$ is the sum of the projections $\Pi^A$ of the unit blocks, with contours of $\eta$-length at most $d$ (H3), and $P_{k+1}(A)(B^k\Lambda)$ is the sum of the projections of the $L$-blocks (H2), with contours of $\eta$-length at most $d(L+1)$ (H3). On each cell apply \cref{lem:a06A:cube} with $(\ell,c_p)=(1,a_k)$, resp. $(L,a_{k+1}L^{-2})$, and $c_p\le a$: the cell operator is at least $\frac12\min\{2,a_*\}\ge\gamma_L$, resp. $\frac12\min\{2L^{-2},a_*L^{-2}\}\ge\gamma_L$. Drop $m_k^2\ge0$. For $\Lambda=\emptyset$ all cells are unit blocks and the bound is $\frac12\min\{2,a_*\}\ge\gamma_\emptyset$.
\end{proof}

\begin{lemma}[{Combes--Thomas bound; the $L^2$ form of \cite[Cor.~2.3]{Ba4} for every pin structure}]\label{lem:a06A:CT}
Let $T\defeq D_A^*D_A+V+\sum_Bc_B\Pi_B$ on $\ell^2(\Omega,\eta^d)$, where $V\ge0$ is a multiplication operator, the $\Pi_B$ are orthogonal projections whose ranges are supported in pairwise disjoint sets $B$ of diameter at most $\ell_*$, $0\le c_B\le a$, and $T\ge\gamma>0$. Put
\[
\theta_0(\gamma,\ell_*)\defeq\min\Big\{\frac1{\ell_*},\ \frac{\gamma^{1/2}}{36\sqrt d},\ \frac{\gamma}{12\mathrm e^2a\ell_*}\Big\},
\]
$\mathrm e$ Euler's number. Then for $f,f'\in\ell^2(\Omega)$ and bond functions $g,g'$, with $\delta$ the distance of the supports (bonds located at $b_-$) and $\theta\defeq\theta_0(\gamma,\ell_*)$:
\begin{enumerate}[label=\textup{(\roman*)},nosep,leftmargin=*]
\item $|\langle f,T^{-1}f'\rangle|\le2\gamma^{-1}e^{-\theta\delta}\norm f\norm{f'}$;
\item $|\langle g,D_AT^{-1}f'\rangle|$ and $|\langle f,T^{-1}D_A^*g'\rangle|$ are at most $(2\gamma^{-1/2}+6\sqrt d\,\theta\gamma^{-1})e^{\theta}e^{-\theta\delta}$ times the product of the norms;
\item $|\langle g,D_AT^{-1}D_A^*g'\rangle|\le4e^{\theta}e^{-\theta\delta}\norm g\norm{g'}$.
\end{enumerate}
The constants do not depend on $\eta$, $\Omega$, $A$ or the torus. $\theta_0(\gamma,\ell_*)$ is increasing in $\gamma$ and decreasing in $\ell_*$; in particular $\theta_0(\gamma_\emptyset,1)\ge\theta_0(\gamma_L,L)$.
\end{lemma}

\begin{proof}
Let $\rho\defeq\dist(\cdot,\supp f')$ (for bond functions evaluated at $b_-$), a $1$-Lipschitz function along bonds, and $T_\theta\defeq e^{\theta\rho}Te^{-\theta\rho}$ for $0<\theta\le\theta_0(\gamma,\ell_*)$; then $T^{-1}=e^{-\theta\rho}T_\theta^{-1}e^{\theta\rho}$ and $e^{\theta\rho}f'=f'$.

\emph{Hopping.} $(e^{\theta\rho}D_Ae^{-\theta\rho}\phi)(b)=(D_A\phi)(b)+\eta^{-1}(e^{-\theta(\rho(b_+)-\rho(b_-))}-1)U(A_b)\phi(b_+)$. Since $|\rho(b_+)-\rho(b_-)|\le\eta$ and $\theta\eta\le1$, $|\eta^{-1}(e^{-\theta s}-1)|\le3\theta$ for $|s|\le\eta$, so $e^{\theta\rho}D_Ae^{-\theta\rho}=D_A+\theta B_1$ with $\norm{B_1}\le3\sqrt d$ (each site is the final point of at most $d$ bonds). Likewise $e^{\theta\rho}D_A^*e^{-\theta\rho}=(e^{-\theta\rho}D_Ae^{\theta\rho})^*=(D_A-\theta B_1')^*$ with $\norm{B_1'}\le3\sqrt d$. Hence $e^{\theta\rho}D_A^*D_Ae^{-\theta\rho}-D_A^*D_A=\theta(D_A^*B_1-B_1'^*D_A)-\theta^2B_1'^*B_1$, and
\[
|\langle\phi,(e^{\theta\rho}D_A^*D_Ae^{-\theta\rho}-D_A^*D_A)\psi\rangle|\le3\sqrt d\,\theta(\norm{D_A\phi}\norm\psi+\norm\phi\norm{D_A\psi})+9d\theta^2\norm\phi\norm\psi .
\]
$V$ commutes with $e^{\theta\rho}$.

\emph{Pins.} Write $\Pi_B=VV^*$ with $V$ an isometry onto the range, and let $\rho_B$ be the value of $\rho$ at a point of $B$. Then $e^{\theta\rho}\Pi_Be^{-\theta\rho}=(e^{\theta(\rho-\rho_B)}V)(e^{-\theta(\rho-\rho_B)}V)^*$, and $|\rho-\rho_B|\le\ell_*$ on $B$, so $\norm{e^{\pm\theta(\rho-\rho_B)}V-V}\le\theta\ell_*e^{\theta\ell_*}$ and $\norm{e^{\theta\rho}\Pi_Be^{-\theta\rho}-\Pi_B}\le2\theta\ell_*e^{\theta\ell_*}+\theta^2\ell_*^2e^{2\theta\ell_*}\le2\theta\ell_*e^{2\theta\ell_*}$, using $te^t\le2(e^t-1)$ for $t=\theta\ell_*\le1$. The sets $B$ are disjoint, so the pin changes are block-diagonal and their sum has norm at most $2a\theta\ell_*e^{2\theta\ell_*}\le2\mathrm e^2a\theta\ell_*$.

\emph{Perturbation.} Put $K\defeq T^{-1/2}(T_\theta-T)T^{-1/2}$. Since $D_A^*D_A\le T$, $\norm{D_AT^{-1/2}}\le1$, and $\norm{T^{-1/2}}\le\gamma^{-1/2}$, so
\[
\norm K\le6\sqrt d\,\theta\gamma^{-1/2}+9d\theta^2\gamma^{-1}+2\mathrm e^2a\theta\ell_*\gamma^{-1}\le\tfrac16+\tfrac1{144}+\tfrac16<\tfrac12
\]
for $\theta\le\theta_0(\gamma,\ell_*)$. Hence $T_\theta=T^{1/2}(1+K)T^{1/2}$ is invertible, $T_\theta^{-1}=T^{-1/2}(1+K)^{-1}T^{-1/2}$, and $\norm{T_\theta^{-1}}\le2\gamma^{-1}$, $\norm{D_AT_\theta^{-1}}\le2\gamma^{-1/2}$, $\norm{D_AT_\theta^{-1}D_A^*}\le2$.

\emph{Bounds.} (i) $\langle f,T^{-1}f'\rangle=\langle e^{-\theta\rho}f,T_\theta^{-1}f'\rangle$ and $\norm{e^{-\theta\rho}f}\le e^{-\theta\delta}\norm f$. (ii) $\langle g,D_AT^{-1}f'\rangle=\langle e^{-\theta\rho}g,(D_A+\theta B_1)T_\theta^{-1}f'\rangle$, bounded by $(2\gamma^{-1/2}+3\sqrt d\,\theta\cdot2\gamma^{-1})e^{-\theta\delta}\norm g\norm{f'}$; the adjoint form follows from the symmetry of $T$ with $\rho$ the distance to $\supp f$. (iii) With $\rho$ the distance to $\supp g'$, $e^{\theta\rho}g'=g'$ and $\langle g,D_AT^{-1}D_A^*g'\rangle=\langle e^{-\theta\rho}g,(D_A+\theta B_1)T_\theta^{-1}(D_A-\theta B_1')^*g'\rangle$, bounded by $(2+12\sqrt d\,\theta\gamma^{-1/2}+18d\theta^2\gamma^{-1})e^{-\theta\delta}\norm g\norm{g'}\le4e^{-\theta\delta}\norm g\norm{g'}$. The factors $e^\theta$ in (ii)--(iii) are $\ge1$. The monotonicity of $\theta_0$ is read off its definition, and $\gamma_\emptyset\ge\gamma_L$, $1\le L$.
\end{proof}

With \cref{lem:a06A:coer} ($\ell_*=L$, $\gamma=\gamma_L$: the pins of $T_\Lambda$ are the unit-block and $L$-block projections, of diameters $\le1$ and $\le L$, with $c_B\in\{a_k,a_{k+1}L^{-2}\}\subset[0,a]$), \cref{lem:a06A:CT} gives \cite[Cor.~2.3]{Ba4} for $G_k(\Omega,\Lambda,A)$ and $G^\eta_{k+1}(\Omega,A)$, with constants depending on $L$ as \cite[p.~588]{Ba4} says; with $\Lambda=\emptyset$ ($\ell_*=1$, $\gamma=\gamma_\emptyset$) it is a proof of the first clause of \cite[Cor.~2.3]{Ba4}, the bounds for $G_k(\Omega,A)$ (its $\delta G_k$ clause is not used in this appendix). Throughout Steps~3(B)--3(D), $\theta_0\defeq\theta_0(\gamma_L,L)$; by the last sentence of \cref{lem:a06A:CT} every Combes--Thomas rate of this part is at least $\theta_0$.

% a06_BALtorus_A2: continuation of a06_BALtorus_A1 (input from it). Group (A), part 2 of 4:
% the block inversion (3.8) as an identity, (1.15), (1.16); the resolvent difference (3.9)-(3.10); the conditioning-region difference (1.18).

\begin{lemma}[{block inversion; \cite[(3.8), (1.15), (1.16)]{Ba4}}]\label{lem:a06A:blockinv}
Assume (H1)--(H3), $A$ regular on $\Omega$, $e_k\le e_1$, and let $\Lambda\subset\Omega^{(k)}$ be a union of $L$-blocks. Put $P_\Lambda\defeq\Lambda P(A)\Lambda$ (an orthogonal projection on $\ell^2(\Lambda)$, since $P(A)$ is block-diagonal over the $L$-blocks) and $W\defeq Q_k^*\Lambda(1-(a_{k+1}/a_k)L^{-2}P_\Lambda)$. Then:
\begin{enumerate}[label=\textup{(\alph*)},nosep,leftmargin=*]
\item $X_\Lambda$ is invertible and $X_\Lambda^{-1}=W^*G_k(\Omega,\Lambda,A)W+a_k^{-1}-(a_{k+1}L^{-2}/a_k^2)P_\Lambda$; that is, $W\delta_y=q_k(y)-(a_{k+1}/a_k)L^{-2}q_{k+1}(y)$ with $q_{k+1}(y)\defeq Q_k^*P(A)\delta_y$, which is \cite[(3.8)]{Ba4}.
\item $\gamma_X\le X_\Lambda\le a_k+aL^{-2}\le\gamma_1\defeq a+aL^{-2}$, with $\gamma_X\defeq(a_*^{-1}+\gamma_L^{-1})^{-1}$: this is \cite[(1.15)]{Ba4}.
\item $|C_\Lambda(y,y')|\le c_2e^{-\theta_0|y-y'|}$ for $y,y'\in\Lambda$, with $c_2\defeq(2\gamma_L^{-1}+2a_*^{-1})e^{2\theta_0L}$ and $\theta_0=\theta_0(\gamma_L,L)$: this is \cite[(1.16)]{Ba4}, also for $C^{(k)}(\Omega,A)=C_{\Omega^{(k)}}$.
\end{enumerate}
The constants depend on $(d,N,\norm q,a,L)$ only.
\end{lemma}

\begin{proof}
On $\ell^2(\Lambda)\oplus\ell^2(\Omega)$ let $\mathcal M\defeq\begin{pmatrix}A_\psi&-B\\-B^*&D_\phi\end{pmatrix}$ with $A_\psi\defeq a_k+aL^{-2}P_\Lambda$, $B\defeq a_k\Lambda Q_k$, $D_\phi\defeq H+a_kP_k(A)$, $H\defeq-\Delta^{\eta,N}_{A,\Omega}+m_k^2$. Its quadratic form is
\[
a_k\norm{\psi-\Lambda Q_k\phi}^2+a_k\norm{(1-\Lambda)Q_k\phi}^2+aL^{-2}\norm{P_\Lambda\psi}^2+\langle\phi,H\phi\rangle .
\]
\emph{Schur complement in $\phi$.} $D_\phi=G_k(\Omega,A)^{-1}\ge\gamma_\emptyset$ by \cref{lem:a06A:coer}, and $A_\psi-BD_\phi^{-1}B^*=\Lambda(a_k-a_k^2Q_kG_k(\Omega,A)Q_k^*+aL^{-2}P(A))\Lambda=X_\Lambda$ by \cite[(1.14)]{Ba4}.

\emph{Schur complement in $\psi$.} $A_\psi\ge a_k>0$, and with $b\defeq aL^{-2}$, $(a_k+bP_\Lambda)^{-1}=a_k^{-1}(1-P_\Lambda)+(a_k+b)^{-1}P_\Lambda$, so $a_k-a_k^2(a_k+bP_\Lambda)^{-1}=\frac{a_kb}{a_k+b}P_\Lambda=a_{k+1}L^{-2}P_\Lambda$ by the recursion $a_{k+1}=aa_k/(aL^{-2}+a_k)$. Hence
\[
D_\phi-B^*A_\psi^{-1}B=H+a_kQ_k^*(1-\Lambda)Q_k+a_{k+1}L^{-2}Q_k^*P_\Lambda Q_k=T_\Lambda,
\]
since $Q_k^*(1-\Lambda)Q_k=P_k(A)(\Omega\setminus B^k\Lambda)$ and $Q_k^*P_\Lambda Q_k=P_{k+1}(A)(B^k\Lambda)$ by (H2). This is the computation \cite[(3.4)--(3.5)]{Ba4} done as algebra.

\emph{Block inverse.} $A_\psi>0$ and $T_\Lambda\ge\gamma_L>0$ (\cref{lem:a06A:coer}), so $\mathcal M$ is positive definite, and its inverse's $\psi\psi$ block is both $(A_\psi-BD_\phi^{-1}B^*)^{-1}=X_\Lambda^{-1}$ and $A_\psi^{-1}+A_\psi^{-1}BT_\Lambda^{-1}B^*A_\psi^{-1}$. Now $A_\psi^{-1}B=(1-\frac{b}{a_k+b}P_\Lambda)\Lambda Q_k=W^*$, because $\frac{b}{a_k+b}=\frac{a_{k+1}}{a_k}L^{-2}$, and $A_\psi^{-1}=a_k^{-1}-\frac{b}{a_k(a_k+b)}P_\Lambda=a_k^{-1}-\frac{a_{k+1}L^{-2}}{a_k^2}P_\Lambda$. This is (a); $P(A)\delta_y$ is supported in the $L$-block of $y$, which lies in $\Lambda$ by (H1), so $\Lambda P_\Lambda\delta_y=P(A)\delta_y$ and $W\delta_y$ is as stated.

\emph{(b).} $\norm{Q_k}\le1$ (Cauchy--Schwarz on blocks of volume $1$) and $\norm{1-tP_\Lambda}\le1$ for $t=(a_{k+1}/a_k)L^{-2}\in[0,1]$ (as $a_{k+1}\le a_kL^2$), so $\norm W\le1$ and $\norm{X_\Lambda^{-1}}\le a_k^{-1}+\gamma_L^{-1}\le a_*^{-1}+\gamma_L^{-1}$. The upper bound holds because $a_k^2Q_kG_kQ_k^*\ge0$ and $\norm{P_\Lambda}\le1$.

\emph{(c).} $W\delta_y$ is supported in the $L$-block $B^{k+1}$ containing $y$, a cube of side $L$ containing $y$, and $\norm{W\delta_y}\le1$. \Cref{lem:a06A:CT}(i) for $T_\Lambda$ (pins of diameter $\le L$, $\gamma=\gamma_L$) gives $|\langle W\delta_y,T_\Lambda^{-1}W\delta_{y'}\rangle|\le2\gamma_L^{-1}e^{-\theta_0(|y-y'|-2L)}$. The term $a_k^{-1}\delta_{yy'}$ is at most $a_*^{-1}$ and vanishes unless $y=y'$; the term $(a_{k+1}L^{-2}/a_k^2)P_\Lambda(y,y')$ has $a_{k+1}L^{-2}/a_k^2=\frac{b}{a_k(a_k+b)}\le a_*^{-1}$, $|P_\Lambda(y,y')|\le1$, and vanishes unless $|y-y'|\le L$. Summing, $|C_\Lambda(y,y')|\le(2\gamma_L^{-1}e^{2\theta_0L}+a_*^{-1}+a_*^{-1}e^{\theta_0L})e^{-\theta_0|y-y'|}\le c_2e^{-\theta_0|y-y'|}$. $\Lambda=\Omega^{(k)}$ is allowed.
\end{proof}

\begin{lemma}[{the resolvent difference; \cite[(3.9), (3.10)]{Ba4}}]\label{lem:a06A:resdiff}
Under the hypotheses of \cref{lem:a06A:blockinv}:
\begin{enumerate}[label=\textup{(\alph*)},nosep,leftmargin=*]
\item $G_k(\Omega,\Lambda,A)-G^\eta_{k+1}(\Omega,A)=G_k(\Omega,\Lambda,A)\,M'\,G^\eta_{k+1}(\Omega,A)$ with $M'\defeq[a_{k+1}L^{-2}P_{k+1}(A)-a_kP_k(A)](\Omega\setminus B^k\Lambda)$.
\item For $f,f'\in\ell^2(B^k\Lambda)$, with $d_f\defeq\dist(\supp f,\Omega\setminus B^k\Lambda)$,
\[
|\langle f,(G_k(\Omega,\Lambda,A)-G^\eta_{k+1}(\Omega,A))f'\rangle|\le c_3\norm f\norm{f'}\exp\big[-\tfrac{\theta_0}6\big(\dist(\supp f,\supp f')+d_f+d_{f'}\big)\big],
\]
with $c_3\defeq8a\gamma_L^{-2}e^{\theta_0L/3}S_L(\theta_0/3)S_L(\theta_0/6)$, $S_L(\theta)\defeq\sum_{n\in\Z^d}e^{-\theta L(|n|-1)_+}$. This is stronger than \cite[(3.10)]{Ba4}, whose $B^k(\Lambda^c)$ contains $\Omega\setminus B^k\Lambda$.
\end{enumerate}
\end{lemma}

\begin{proof}
(a) is the second resolvent identity, since $T_{\Omega^{(k)}}-T_\Lambda=M'$. (The bracket of \cite[(3.9)]{Ba4} is printed with the opposite sign; only norms of it are used there.)

(b) $M'$ is block-diagonal over the $L$-blocks $\Delta\subset\Omega\setminus B^k\Lambda$ (by (H2) each $P_{k+1}$ summand and the $P_k$ summands it contains live on one $L$-block), with $\norm{M'_\Delta}\le a_{k+1}L^{-2}+a_k\le2a$. Let first $f$, $f'$ be supported in single $L$-blocks $B,B'\subset B^k\Lambda$, and write $u_\Delta\defeq\dist(\Delta,\supp f)$, $v_\Delta\defeq\dist(\Delta,\supp f')$. Both $T_\Lambda$ and $T_{\Omega^{(k)}}$ are at least $\gamma_L$ (\cref{lem:a06A:coer}), so \cref{lem:a06A:CT}(i) gives $\norm{\mathbf 1_\Delta G_k(\Omega,\Lambda,A)f}\le2\gamma_L^{-1}e^{-\theta_0u_\Delta}\norm f$ and the same for $G^\eta_{k+1}(\Omega,A)f'$ with $v_\Delta$; hence
\[
|\langle f,(G_k(\Omega,\Lambda,A)-G^\eta_{k+1}(\Omega,A))f'\rangle|\le2a(2\gamma_L^{-1})^2\sum_\Delta e^{-\theta_0(u_\Delta+v_\Delta)}\norm f\norm{f'} .
\]
Since $\Delta\subset\Omega\setminus B^k\Lambda$, $u_\Delta\ge d_f$ and $v_\Delta\ge d_{f'}$; since $\Delta$ has diameter $L$, $\dist(\supp f,\supp f')\le u_\Delta+L+v_\Delta$. Therefore $u_\Delta+v_\Delta\ge\frac13(\dist(\supp f,\supp f')-L)+\frac13(d_f+d_{f'})+\frac13(u_\Delta+v_\Delta)$, and $\sum_\Delta e^{-(\theta_0/3)u_\Delta}\le S_L(\theta_0/3)$, a sum over the $L$-block lattice, uniform in the torus (on a torus, each block has a nearest lift, and distinct blocks have distinct lifts). So for block-supported $f,f'$ the bound holds with the rate $\theta_0/3$ and the constant $8a\gamma_L^{-2}e^{\theta_0L/3}S_L(\theta_0/3)$. For general $f,f'$ decompose $f=\sum_Bf_B$, $f'=\sum_{B'}f'_{B'}$ over the $L$-blocks of $B^k\Lambda$; restricting supports increases all three distances, so $\dist(\supp f_B,\supp f'_{B'})+d_{f_B}+d_{f'_{B'}}\ge\dist(\supp f,\supp f')+d_f+d_{f'}$, and splitting the rate $\theta_0/3$ into two halves, one half gives the exponential of the statement and the other leaves the kernel $e^{-(\theta_0/6)\dist(\supp f_B,\supp f'_{B'})}\le e^{-(\theta_0/6)\dist(B,B')}$, whose Schur norm on $\ell^2$ of the block indices is at most $S_L(\theta_0/6)$; since $\sum_B\norm{f_B}^2=\norm f^2$, the Schur test gives (b).
\end{proof}

\begin{lemma}[{the conditioning-region difference; \cite[(1.18)]{Ba4}}]\label{lem:a06A:CLdiff}
Under the hypotheses of \cref{lem:a06A:blockinv}, for $y,y'\in\Lambda$,
\[
|C_\Lambda(y,y')-C^{(k)}(\Omega,A;y,y')|\le c_4\exp\big[-\tfrac{\theta_0}6\big(|y-y'|+\dist(y,\Omega^{(k)}\setminus\Lambda)+\dist(y',\Omega^{(k)}\setminus\Lambda)\big)\big],
\]
with $c_4\defeq c_3e^{(\theta_0/6)(4L+2)}$. The bound decays a fortiori in the distances to the whole complement $\Lambda^c$.
\end{lemma}

\begin{proof}
$C^{(k)}(\Omega,A)=C_{\Omega^{(k)}}$. For $y\in\Lambda$, $W_{\Omega^{(k)}}\delta_y=W_\Lambda\delta_y=Q_k^*(\delta_y-(a_{k+1}/a_k)L^{-2}P(A)\delta_y)$, because $P(A)\delta_y$ lives on the $L$-block of $y$, which lies in $\Lambda$ (H1); and $P_\Lambda(y,y')=P(A)(y,y')$ for $y,y'\in\Lambda$. By \cref{lem:a06A:blockinv}(a) for both $\Lambda$ and $\Omega^{(k)}$ the constant terms cancel and
\[
C_\Lambda(y,y')-C_{\Omega^{(k)}}(y,y')=\langle W\delta_y,(G_k(\Omega,\Lambda,A)-G^\eta_{k+1}(\Omega,A))W\delta_{y'}\rangle .
\]
Apply \cref{lem:a06A:resdiff}(b) with $f=W\delta_y$, $f'=W\delta_{y'}$, both of norm $\le1$ and supported in $B^k$ of the $L$-blocks of $y$, $y'$, which lie in $B^k\Lambda$. Every point of $\supp f$ is within $L$ of $y$, and $\Omega\setminus B^k\Lambda=B^k(\Omega^{(k)}\setminus\Lambda)$ lies within $1$ of $\Omega^{(k)}\setminus\Lambda$; so $\dist(\supp f,\supp f')\ge|y-y'|-2L$ and $d_f\ge\dist(y,\Omega^{(k)}\setminus\Lambda)-L-1$, and the shifts total at most $4L+2$.
\end{proof}

% a06_BALtorus_A3: continuation of a06_BALtorus_A2 (input from it). Group (A), part 3 of 4:
% the two-sided depth factor of Ba4's Theorem, the domain differences (5.5) and (1.20), Ba1 Prop. 2.2, the summary theorem, the polylog bound.

The remaining estimates use the printed theorem of \cite{Ba4}, which is \cite[Prop.~2.1]{Ba1}: for $\alpha<1$ there are $\delta_0,c_0,R_0>0$, independent of $A$, $k$, $\Omega$, depending on $d$, $M$ (and $c_0$ on $\alpha$), such that for $e_k$ small and $f\colon\Omega\to\R^N$: the $C^{1,\alpha}$ bound \cite[(1.9)]{Ba4} for $x,x'\in\Omega$ with $\dist(\{x,x'\},\Omega^c)\ge R_0$; $|(D^\eta_{A,\mu}G_k(\Omega,A)f)(x)|,\,|(G_k(\Omega,A)f)(x)|\le c_0e^{-\delta_0\dist(x,\supp f)}\norm f_\infty$ for $x\in\Omega$ with $\dist(x,\Omega^c)\ge R_0$ \cite[(1.10)]{Ba4}; and, for $\Omega\subset\Omega_0$ and $\delta G_k(\Omega,\Omega_0,A)\defeq G_k(\Omega,A)-G_k(\Omega_0,A)$ \cite[(1.11)]{Ba4}, the same bounds with the same restrictions and an additional factor \cite[(1.12)]{Ba4}, printed in \cite[(2.26)]{Ba1} as $\exp(-\delta_0\dist(\supp f,\Omega^c)-\delta_0\dist(\{x,x'\},\Omega^c))$. When $\Omega$ is the whole torus the depth condition is void \cite[p.~611]{Ba1}. Only this restricted form (points at depth at least $R_0$, or $\Omega$ the whole torus) is used anywhere in this appendix; the clause of \cite[p.~611]{Ba1} and \cite[p.~573]{Ba4} that drops the depth condition on rectangular parallelepipeds is not used.

\begin{lemma}[the depth factor, two-sided]\label{lem:a06A:twosided}
Let $\Omega\subset\Omega_0$ be unions of large blocks, $A$ regular on $\Omega_0$, $e_k$ small, $\supp f$ of diameter at most $s$, and $x\in\Omega$ with $\dist(x,\Omega^c)\ge R_0$. Then
\[
|(\delta G_k(\Omega,\Omega_0,A)f)(x)|\le c_0e^{\delta_0s/2}\exp\big[-\tfrac{\delta_0}2\big(\dist(x,\supp f)+\dist(x,\Omega^c)+\dist(\supp f,\Omega^c)\big)\big]\norm f_\infty,
\]
and the same for $D^\eta_{A,\mu}\delta G_k(\Omega,\Omega_0,A)f$, and for the H\"older quotient of \cite[(1.9)]{Ba4} with $|x-x'|\le1$ and $\dist(\{x,x'\},\Omega^c)\ge R_0$, with $\dist(\{x,x'\},\cdot)$ in place of $\dist(x,\cdot)$. For general $f$ (no bound on the diameter) the bounds hold at the rate $\delta_0/4$ with the constant $c_0e^{3\delta_0/4}S(\delta_0/4)$ ($c_0e^{\delta_0}S(\delta_0/4)$ for the H\"older quotient).
\end{lemma}

\begin{proof}
$\dist(x,\Omega^c)\le\dist(x,\supp f)+s+\dist(\supp f,\Omega^c)$, so $\dist(x,\supp f)+\dist(x,\Omega^c)+\dist(\supp f,\Omega^c)\le2\dist(x,\supp f)+2\dist(\supp f,\Omega^c)+s$. The printed factor, read with both of its terms as $\dist(\supp f,\Omega^c)$, gives $c_0e^{-\delta_0\dist(x,\supp f)}e^{-2\delta_0\dist(\supp f,\Omega^c)}\norm f_\infty$, which is at most the right side; read as printed in \cite[(2.26)]{Ba1}, it gives the right side without the factor $e^{\delta_0s/2}$ and at the rate $\delta_0$. So the bound holds under either reading. For the H\"older quotient the printed factor carries $\dist(\{x,x'\},\cdot)$ and the same arithmetic applies, because $\dist(\{x,x'\},\Omega^c)\le\dist(\{x,x'\},\supp f)+s+\dist(\supp f,\Omega^c)$ (the inequality for $x$ and for $x'$ separately, then the minimum over the two points). For general $f$ decompose $\supp f$ into the unit blocks $\square$ it meets, apply the bound to $f\mathbf 1_\square$ ($s\le1$; $\dist(x,\square\cap\supp f)\ge\dist(x,\supp f)$ and $\dist(\square\cap\supp f,\Omega^c)\ge\dist(\supp f,\Omega^c)$), and sum: with $\square_n$ the unit block of $n\in\Z^d$ and $x\in\square_{n_x}$, $\dist(x,\square_n)\ge(|n-n_x|-1)_+$, so $\sum_\square e^{-(\delta_0/2)\dist(x,\square\cap\supp f)}\le e^{-(\delta_0/4)\dist(x,\supp f)}\sum_{n\in\Z^d}e^{-(\delta_0/4)(|n-n_x|-1)_+}\le e^{-(\delta_0/4)\dist(x,\supp f)}e^{\delta_0/4}S(\delta_0/4)$; for the H\"older quotient $\dist(\{x,x'\},\square_n)\ge\dist(x,\square_n)-1$ costs one more factor $e^{\delta_0/4}$.
\end{proof}

\begin{lemma}[{domain differences; \cite[(5.5), (1.20)]{Ba4}}]\label{lem:a06A:domdiff}
Let $\Omega\subset\Omega_0$ be unions of large blocks, (H1)--(H3), $A$ regular on $\Omega_0$, $e_k\le e_1$ and small enough for the printed theorem. Put
\[
\delta_7\defeq\min\{\delta_0/2,\,\theta_0(\gamma_\emptyset,1)\},\qquad\delta^S_7\defeq\min\{\theta_0,\,\delta_7\}=\min\{\theta_0(\gamma_L,L),\,\delta_0/2\}.
\]
\begin{enumerate}[label=\textup{(\alph*)},nosep,leftmargin=*]
\item For $y,y'\in\Omega^{(k)}$, with $d(y)\defeq\dist(B^k(y),\Omega^c)$,
\[
|(\Delta^{(k)}(\Omega,A)-\Delta^{(k)}(\Omega_0,A))(y,y')|\le c_6\exp\big[-\delta_7\big(|y-y'|+d(y)+d(y')\big)\big],
\]
$c_6\defeq\max\{a^2c_0e^{\delta_0/2},\,4a^2\gamma_\emptyset^{-1}e^{2\theta_0(\gamma_\emptyset,1)R_0}\}e^{\delta_7}$: this is \cite[(5.5)]{Ba4}.
\item For every $\Lambda\subset\Omega^{(k)}$ that is a union of $L$-blocks and $x,x'\in\Lambda$, with $\tilde d(x)\defeq\dist(x,\Omega^{(k)c})$,
\[
|C_\Lambda(\Omega,A;x,x')-C_\Lambda(\Omega_0,A;x,x')|\le c_{20}\exp\big[-\tfrac{\delta^S_7}4\big(|x-x'|+\tilde d(x)+\tilde d(x')\big)\big],
\]
$c_{20}\defeq c_2^2c_6e^{\delta^S_7/2}S(\delta^S_7/2)^2$: this is \cite[(1.20)]{Ba4}.
\end{enumerate}
\end{lemma}

\begin{proof}
(a) $Q_k(A)$ on $\Omega$ and on $\Omega_0$ agree on the blocks $B^k(y)$, $y\in\Omega^{(k)}$, so by \cite[(1.14)]{Ba4} the kernel is $\delta\Delta(y,y')=-a_k^2\langle q_k(y),\delta G_k(\Omega,\Omega_0,A)q_k(y')\rangle$, with $q_k(y),q_k(y')$ supported in $\Omega$. Three cases. If $d(y)\ge R_0$, every $x\in B^k(y)$ has $\dist(x,\Omega^c)\ge R_0$, and \cref{lem:a06A:twosided} with $f=q_k(y')$ ($\norm f_\infty=1$, $s\le1$) gives $|(\delta G_kq_k(y'))(x)|\le c_0e^{\delta_0/2}e^{-(\delta_0/2)(\dist(B^k(y),B^k(y'))+d(y)+d(y'))}$ for $x\in B^k(y)$; pairing with $q_k(y)$, whose entries have modulus $1$ on a block of volume $1$, gives $|\delta\Delta(y,y')|\le a_k^2c_0e^{\delta_0/2}e^{-(\delta_0/2)(\dist(B^k(y),B^k(y'))+d(y)+d(y'))}$. If $d(y')\ge R_0$, use $\delta\Delta(y,y')=\delta\Delta(y',y)^{\mathsf T}$. If $d(y),d(y')<R_0$, apply \cref{lem:a06A:CT}(i) to $G_k(\Omega,A)$ and to $G_k(\Omega_0,A)$ separately (\cref{lem:a06A:coer} with $\Lambda=\emptyset$ on $\Omega$ and on $\Omega_0$; $\ell_*=1$, $\gamma=\gamma_\emptyset$): $|\delta\Delta(y,y')|\le4a_k^2\gamma_\emptyset^{-1}e^{-\theta_0(\gamma_\emptyset,1)\dist(B^k(y),B^k(y'))}\le4a_k^2\gamma_\emptyset^{-1}e^{2\theta_0(\gamma_\emptyset,1)R_0}e^{-\theta_0(\gamma_\emptyset,1)(\dist(B^k(y),B^k(y'))+d(y)+d(y'))}$. Finally $\dist(B^k(y),B^k(y'))\ge|y-y'|-1$.

(b) The inverses are bounded by \cref{lem:a06A:blockinv}(b) on both domains, so $C_\Lambda(\Omega,A)-C_\Lambda(\Omega_0,A)=X_\Lambda(\Omega,A)^{-1}[X_\Lambda(\Omega_0,A)-X_\Lambda(\Omega,A)]X_\Lambda(\Omega_0,A)^{-1}$, and $X_\Lambda(\Omega_0,A)-X_\Lambda(\Omega,A)=-\delta\Delta|_\Lambda$, because $\Lambda\subset\Omega^{(k)}$ and $P(A)$ is the same on $\Lambda$ for both domains. With \cref{lem:a06A:blockinv}(c) for both domains and (a),
\[
|\delta C(x,x')|\le c_2^2c_6\sum_{y,y'\in\Lambda}e^{-\delta^S_7\Sigma},\qquad\Sigma\defeq|x-y|+|y-y'|+d(y)+d(y')+|y'-x'| .
\]
Since $\Omega^c=B^k(\Omega^{(k)c})$ and $\dist(B^k(y),B^k(z))\ge|y-z|-1$, $\dist(y,\Omega^{(k)c})\le d(y)+1$; so $\tilde d(x)\le|x-y|+d(y)+1$, $\tilde d(x')\le|x'-y'|+d(y')+1$, and $|x-x'|+\tilde d(x)+\tilde d(x')\le2\Sigma+2$. Hence $e^{-\delta^S_7\Sigma}\le e^{\delta^S_7/2}e^{-(\delta^S_7/4)(|x-x'|+\tilde d(x)+\tilde d(x'))}e^{-(\delta^S_7/2)(|x-y|+|y'-x'|)}$, and the sums over $y$ and $y'$ give $S(\delta^S_7/2)^2$.
\end{proof}

\begin{proposition}[{\cite[Prop.~2.2, (2.27)--(2.29)]{Ba1}}]\label{prop:a06A:Ba1P22}
Assume (H1)--(H3), $A$ regular on $\Omega_0\supset\Omega$, $e_k$ as in \cref{lem:a06A:domdiff}. Then for $y,y'\in\Omega^{(k)}$
\[
|\Delta^{(k)}(\Omega,A;y,y')|\le c_7e^{-\delta_7|y-y'|},\qquad|\Delta^{(k)}(\Omega,A;y,y')-\Delta^{(k)}(\Omega_0,A;y,y')|\le c_7\exp\big[-\delta_7\big(|y-y'|+\dist(y,\Omega^{(k)c})+\dist(y',\Omega^{(k)c})\big)\big],
\]
with $c_7\defeq\max\{a+2a^2\gamma_\emptyset^{-1}e^{\theta_0(\gamma_\emptyset,1)},\,c_6e^{2\delta_7}\}$. The first bound holds on every union $\Omega$ of unit blocks and uses only \cref{lem:a06A:coer,lem:a06A:CT}; its rate is in fact $\theta_0(\gamma_\emptyset,1)\ge\delta_7$.
\end{proposition}

\begin{proof}
$\Delta^{(k)}(y,y')=a_k\delta_{yy'}-a_k^2\langle q_k(y),G_k(\Omega,A)q_k(y')\rangle$ \cite[(1.14)]{Ba4}. \Cref{lem:a06A:coer} with $\Lambda=\emptyset$ and \cref{lem:a06A:CT}(i) with $\ell_*=1$ give $|\langle q_k(y),G_kq_k(y')\rangle|\le2\gamma_\emptyset^{-1}e^{-\theta_0(\gamma_\emptyset,1)\dist(B^k(y),B^k(y'))}$, and $\dist(B^k(y),B^k(y'))\ge|y-y'|-1$; this is the first bound. The second is \cref{lem:a06A:domdiff}(a) with $d(y)\ge\dist(y,\Omega^{(k)c})-1$.
\end{proof}

\begin{theorem}[the propagator estimates of group (A)]\label{thm:a06A:groupA}
Let $d=3$, $L=5$, (H1)--(H3), $\Omega\subset\Omega_0$ unions of large blocks of $T^{(k)}_1$ (the whole torus included), $\Lambda\subset\Omega^{(k)}$ a union of $L$-blocks, $A$ regular on $\Omega_0$, and $\eps_0$ below the caps of \cref{lem:a06A:cube,lem:a06A:domdiff}. Then \cite[(1.8)]{Ba4} on $\Omega$ (\cref{lem:a06A:coer}); the $G_k$ clause of \cite[Cor.~2.3]{Ba4} for $G_k(\Omega,A)$, $G_k(\Omega,\Lambda,A)$ and $G^\eta_{k+1}(\Omega,A)$ (\cref{lem:a06A:CT}); \cite[Prop.~2.3]{Ba1}, i.e.\ \cite[(1.15), (1.16), (1.18), (1.20)]{Ba4} (\cref{lem:a06A:blockinv}(b), (c), \cref{lem:a06A:CLdiff}, \cref{lem:a06A:domdiff}(b)); \cite[(5.5)]{Ba4} (\cref{lem:a06A:domdiff}(a)); and \cite[Prop.~2.2]{Ba1} (\cref{prop:a06A:Ba1P22}) hold with the constants
\[
\gamma_\emptyset,\ \gamma_L,\ \gamma_X,\ \gamma_1,\ \theta_0=\theta_0(\gamma_L,L),\ \theta_0(\gamma_\emptyset,1),\ c_2,\ c_3,\ c_4,\ c_6,\ c_7,\ c_{20},\ \delta_7,\ \delta^S_7
\]
defined above. They depend on $(d,N,\norm q,c,\beta,a,L)$, on $M$ through the constants $\delta_0,c_0,R_0$ of \cite[Thm.]{Ba4} (which enter \cref{lem:a06A:twosided,lem:a06A:domdiff}, \cref{prop:a06A:Ba1P22} and nothing else), and on the cap on $\eps_0$; they do not depend on $\eta$, $k$, $\Omega$, $\Omega_0$, $\Lambda$, $A$, $\eps$ or the torus.
\end{theorem}

\begin{proof}
Each estimate is the lemma named after it; the dependences are read off the definitions of the constants. The order of the arguments is not circular: \cref{lem:a06A:cube,lem:a06A:coer,lem:a06A:CT,lem:a06A:blockinv,lem:a06A:resdiff,lem:a06A:CLdiff} and the first bound of \cref{prop:a06A:Ba1P22} use no estimate of \cite{Ba4} beyond its definitions, while \cref{lem:a06A:twosided,lem:a06A:domdiff} and the second bound of \cref{prop:a06A:Ba1P22} use the printed theorem of \cite{Ba4} in its restricted form (supplied on a torus, with the constants of \cite{Ba4}, by \cref{lem:a06A:torus} below), whose proof \cite[Sect.~2]{Ba4} rests on \cite[Lemmas~2.1, 2.2, 2.4]{Ba4}; the $j$-uniform form of \cite[Lemma~2.4]{Ba4} is \cref{thm:a06:A5}, whose clause \cite[(2.37)]{Ba4} is \cref{lem:a06A:blockinv}(c) at $A=0$, $\Lambda=\Omega^{(k)}$.
\end{proof}

\begin{remark}[dependence on $L$; the depth margin]\label{rem:a06A:L}
(i) $\gamma_\emptyset=\frac14\min\{2,a_*\}$ and the constants of \cref{thm:a06:A5} do not depend on $L$, with $a_*\ge3a/4$ for $L\ge2$; $\gamma_L$, $\gamma_X$, $\theta_0$ and $c_2$ do, through $\gamma_L\propto L^{-2}$, as \cite[p.~588]{Ba4} says of the constants of \cite[(2.37)]{Ba4} and \cite[Prop.~2.3]{Ba1}. At the fixed value $L=5$ all are numbers. (ii) The depth condition of the printed theorem is met wherever it is invoked in Steps~3(B)--3(D): with $r(t)=R(1+\ln t^{-1})^\varrho$ and $R>R_0$ \cite[(2.7)]{Ba2}, every point within distance $1$ of a set at distance more than $2r(t)$ from a boundary has depth more than $2R-1\ge R_0$ (as $R_0\ge1$ may be assumed, enlarging $R_0$ weakening nothing in the theorem); pairs $x,x'$ with $|x-x'|>1$ in \cite[(1.9)]{Ba4} reduce to the sup bounds \cite[p.~578]{Ba4}.
\end{remark}

\begin{lemma}[the polylogarithmic layer]\label{lem:a06A:polylog}
Let $r(t)=R(1+\ln t^{-1})^\varrho$ with $\varrho>1$, and $\delta,R,\kappa>0$. Then $e^{-\delta r(t)}\le t^\kappa$ for all $t\in(0,1]$ if $\delta R\ge\kappa$, and in general for $t\le t_0\defeq\exp(-(\kappa/(\delta R))^{1/(\varrho-1)})$. Consequently $e^{-\delta r(t)}\le C_\kappa t^\kappa$ on $(0,1]$ with $C_\kappa\defeq t_0^{-\kappa}=\exp(\kappa(\kappa/(\delta R))^{1/(\varrho-1)})$.
\end{lemma}

\begin{proof}
With $u\defeq\ln t^{-1}\ge0$, $\delta R(1+u)^\varrho\ge\delta R\max\{1+u,u^\varrho\}$. This is at least $\kappa u$ if $\delta R\ge\kappa$, or if $u^{\varrho-1}\ge\kappa/(\delta R)$, i.e.\ $t\le t_0$. For $t\in[t_0,1]$, $e^{-\delta r(t)}\le1\le t_0^{-\kappa}t^\kappa$.
\end{proof}

\input{sec/a06_BALtorus_A4}




\paragraph{Step 4: the mass.} The scalar mass is a parameter of the theorems: \cite[(1.14)]{Ba1} and \cite[(1.3)]{Ba2} are stated for every $M^2>0$, with constants depending on $M$, and nothing in Steps 1--3 distinguishes $m$ from $M_1$. The one printed clause on the mass, \cite[Prop.~3.1]{Ba2} ``we assume $m^2\le O(1)$ also'', is the cap $(L^k\eps)^2M^2\le c_m(d,a)$ on the dimensionless mass at the scales at which the proposition is applied (\cref{rem:s14:BALm}), one more one-sided cap on $\eps_0$; taking $\eps_0$ below the caps for both masses gives \cref{thm:a06:T1} at $m$ and at $M_1$ with one stopping scale. This proves \cref{thm:a06:T1}, and \cref{hyp:s03:T1} follows as stated there.

